% Generated human-review packet template.  Source records, Lean statements,
% and raw-source-to-Spec screening fields are substituted by
% scripts/review_dashboard_packet.py.
% Do not hand-edit generated packets; annotate the PDF or reconcile notes into
% the paper-local human-review log.
\documentclass[10pt]{article}
\usepackage[margin=0.43in,footskip=0.2in]{geometry}
\usepackage{fontspec}
% Use a Unicode-complete main face as well as a Unicode-complete code face.
% Reviewer reasons and source labels can contain exact Greek symbols outside
% verbatim blocks; silently dropping one would corrupt the review packet.
\setmainfont{DejaVu Serif}
% The broad DejaVu Sans Unicode coverage preserves Lean's mathematical glyphs
% (including U+2216 set difference and superscript identifiers) in verbatim
% review blocks.  Its proportional metrics are preferable to silently missing
% exact semantic text from a narrower monospace font.
\setmonofont{DejaVu Sans}
% DejaVu lacks the indexed infimum/supremum glyphs used in expanded Lean.
% Switch just those characters, including inside verbatim, without changing
% the exact review text or the surrounding font.
\newfontfamily{\reviewmathfont}{Noto Sans Math}
\newXeTeXintercharclass\reviewmathclass
\XeTeXcharclass"2A05=\reviewmathclass
\XeTeXcharclass"2A06=\reviewmathclass
\XeTeXinterchartoks 0 \reviewmathclass={\begingroup\reviewmathfont}
\XeTeXinterchartoks \reviewmathclass 0={\endgroup}
\XeTeXinterchartoks 4095 \reviewmathclass={\begingroup\reviewmathfont}
\XeTeXinterchartoks \reviewmathclass 4095={\endgroup}
\XeTeXinterchartokenstate=1
\usepackage{hyperref}
\usepackage{amssymb}
\usepackage{fvextra}
\usepackage{enumitem}
\setlength{\parindent}{0pt}
\setlength{\parskip}{0.15em}
\linespread{0.92}
\hypersetup{colorlinks=true,linkcolor=blue,urlcolor=blue,breaklinks=true}
\DefineVerbatimEnvironment{ReviewVerbatim}{Verbatim}{breaklines=true,breakanywhere=true,fontsize=\scriptsize}
\newcommand{\reviewerbox}{%
  \par\noindent\fbox{\parbox{0.96\linewidth}{\rule{0pt}{0.38in}}}\par
}
\newcommand{\reviewmatch}[1]{%
  \par\noindent\textbf{Reviewer check:} $\square$\; #1\par
  \noindent\emph{If left unchecked, explain the discrepancy or uncertainty below.}\par
}

\begin{document}
\fontsize{9}{10}\selectfont

\begin{center}
{\LARGE Human review packet: GolzHaghtalabYang2025Distortion}\\[0.35em]
\small Generated from byte-pinned source excerpts, the expanded PaperInterface
specifications, and hash-bound raw-source-to-Spec screening records.
\end{center}

\noindent\textbf{Paper:} GolzHaghtalabYang2025Distortion\\
\textbf{Paper id:} \texttt{GolzHaghtalabYang2025Distortion}\\
\textbf{Source version:} \parbox[t]{0.76\linewidth}{\raggedright NeurIPS 2025 published conference PDF}\\
\textbf{Source record:} \texttt{audit/paper\_statement\_map.json}\\
\textbf{Rows:} 16\\
\textbf{Generated:} 2026-09-27\\
\textbf{Source URL:} \href{https://arxiv.org/abs/2505.23749}{official source}





\section*{How to use this packet}
For each source claim, compare the exact source excerpts with the expanded
PaperInterface specification. Mark up this PDF or fill the blank reviewer box in the TeX source;
return the annotated file for reconciliation into the paper-local human-review
log.

\section*{Open the interactive dashboard}
The interactive dashboard is an optional alternative to using this PDF; it
presents the same review surface rather than an additional review requirement.
After forking and cloning (or pulling) AppliedModelingLib, run the following from the
repository root. The cache command is needed only the first time in a new clone.
\begin{ReviewVerbatim}
cd <your-repository-checkout>
lake exe cache get
python3 scripts/review_dashboard.py --paper GolzHaghtalabYang2025Distortion --serve
\end{ReviewVerbatim}
Open \url{http://127.0.0.1:8765/} in a browser and keep that terminal running
while reviewing. The dashboard records saved annotations in this paper's local
reviewer-owned review record.

\section*{Contents}
\begin{itemize}[leftmargin=1.2em]
\item \hyperlink{governing-declaration-section-5ef5ef0364b6939c}{Governing Lean declarations (143; supporting context)}
\item \hyperlink{library-section-e44db2071aa7a1c7}{Semantic library prerequisites (11; not paper claims)}
\begin{itemize}[leftmargin=1.2em]
\item \hyperlink{library-prerequisite-56c0734f8877ddbe}{Applied\allowbreak{}Modeling\allowbreak{}Lib.\allowbreak{}Alignment.\allowbreak{}Welfare.\allowbreak{}Finite\allowbreak{}Utility\allowbreak{}Profile}
\item \hyperlink{library-prerequisite-d3573881ae45d4c6}{Applied\allowbreak{}Modeling\allowbreak{}Lib.\allowbreak{}finite\allowbreak{}KL\allowbreak{}Divergence}
\item \hyperlink{library-prerequisite-1d625de57bf58a6a}{Applied\allowbreak{}Modeling\allowbreak{}Lib.\allowbreak{}Alignment.\allowbreak{}Welfare.\allowbreak{}In\allowbreak{}Finite\allowbreak{}KL\allowbreak{}Ball}
\item \hyperlink{library-prerequisite-e35e7dd3bb4adb9e}{Applied\allowbreak{}Modeling\allowbreak{}Lib.\allowbreak{}Alignment.\allowbreak{}Welfare.\allowbreak{}Reference\allowbreak{}Alternative\allowbreak{}Policy}
\item \hyperlink{library-prerequisite-cbe8afa0d6d1ba78}{Applied\allowbreak{}Modeling\allowbreak{}Lib.\allowbreak{}Alignment.\allowbreak{}Welfare.\allowbreak{}Unit\allowbreak{}Interval\allowbreak{}Utility\allowbreak{}Profile}
\item \hyperlink{library-prerequisite-15f4b442749deab2}{Applied\allowbreak{}Modeling\allowbreak{}Lib.\allowbreak{}Alignment.\allowbreak{}Welfare.\allowbreak{}population\allowbreak{}Average\allowbreak{}Utility}
\item \hyperlink{library-prerequisite-184c74a4db4311c2}{Applied\allowbreak{}Modeling\allowbreak{}Lib.\allowbreak{}Alignment.\allowbreak{}Welfare.\allowbreak{}policy\allowbreak{}Average\allowbreak{}Utility}
\item \hyperlink{library-prerequisite-7e2db20cf087d9ba}{Applied\allowbreak{}Modeling\allowbreak{}Lib.\allowbreak{}Learning.\allowbreak{}Human\allowbreak{}Feedback.\allowbreak{}Pairwise\allowbreak{}Preference}
\item \hyperlink{library-prerequisite-cf06d15d7f8b7c31}{Applied\allowbreak{}Modeling\allowbreak{}Lib.\allowbreak{}Alignment.\allowbreak{}Welfare.\allowbreak{}population\allowbreak{}Bradley\allowbreak{}Terry\allowbreak{}Preference}
\item \hyperlink{library-prerequisite-268c959ee7d2d148}{Applied\allowbreak{}Modeling\allowbreak{}Lib.\allowbreak{}PMF\allowbreak{}Full\allowbreak{}Support}
\item \hyperlink{library-prerequisite-de8e2a8bfe221a38}{Applied\allowbreak{}Modeling\allowbreak{}Lib.\allowbreak{}pmf\allowbreak{}Prod}
\end{itemize}
\item \hyperlink{paper-prerequisite-section-5ef5ef0364b6939c}{Paper-specific semantic prerequisites (17; not paper claims)}
\begin{itemize}[leftmargin=1.2em]
\item \hyperlink{paper-prerequisite-53af1666fc341b0f}{Golz\allowbreak{}Haghtalab\allowbreak{}Yang2025\allowbreak{}Distortion.\allowbreak{}Is\allowbreak{}Empirical\allowbreak{}Maximal\allowbreak{}Lottery}
\item \hyperlink{paper-prerequisite-54fbc7c16362dbdd}{Golz\allowbreak{}Haghtalab\allowbreak{}Yang2025\allowbreak{}Distortion.\allowbreak{}Is\allowbreak{}Source\allowbreak{}Alignment\allowbreak{}Benchmark}
\item \hyperlink{paper-prerequisite-dbd778917da393ac}{Golz\allowbreak{}Haghtalab\allowbreak{}Yang2025\allowbreak{}Distortion.\allowbreak{}source\allowbreak{}Maximal\allowbreak{}Lottery\allowbreak{}Margin}
\item \hyperlink{paper-prerequisite-bbe7079711c4e442}{Golz\allowbreak{}Haghtalab\allowbreak{}Yang2025\allowbreak{}Distortion.\allowbreak{}source\allowbreak{}Maximal\allowbreak{}Lottery\allowbreak{}Value}
\item \hyperlink{paper-prerequisite-91b5f077c58b43c1}{Golz\allowbreak{}Haghtalab\allowbreak{}Yang2025\allowbreak{}Distortion.\allowbreak{}Is\allowbreak{}Source\allowbreak{}Maximal\allowbreak{}Lottery}
\item \hyperlink{paper-prerequisite-ceb3f23b84b56503}{Golz\allowbreak{}Haghtalab\allowbreak{}Yang2025\allowbreak{}Distortion.\allowbreak{}Is\allowbreak{}Source\allowbreak{}Social\allowbreak{}Choice\allowbreak{}Optimal\allowbreak{}Average\allowbreak{}Utility}
\item \hyperlink{paper-prerequisite-54cac251e94c8e83}{Golz\allowbreak{}Haghtalab\allowbreak{}Yang2025\allowbreak{}Distortion.\allowbreak{}corollary4\allowbreak{}Iid\allowbreak{}User\allowbreak{}Empirical\allowbreak{}Maximal\allowbreak{}Lottery}
\item \hyperlink{paper-prerequisite-674d1bc00409609f}{Golz\allowbreak{}Haghtalab\allowbreak{}Yang2025\allowbreak{}Distortion.\allowbreak{}empirical\allowbreak{}Borda\allowbreak{}Definition\allowbreak{}Spec}
\item \hyperlink{paper-prerequisite-acf649baada45fd3}{Golz\allowbreak{}Haghtalab\allowbreak{}Yang2025\allowbreak{}Distortion.\allowbreak{}nlhf\allowbreak{}Definition\allowbreak{}Spec}
\item \hyperlink{paper-prerequisite-8df126bd40665caf}{Golz\allowbreak{}Haghtalab\allowbreak{}Yang2025\allowbreak{}Distortion.\allowbreak{}population\allowbreak{}Borda\allowbreak{}Definition\allowbreak{}Spec}
\item \hyperlink{paper-prerequisite-58ee6f1d89c4f942}{Golz\allowbreak{}Haghtalab\allowbreak{}Yang2025\allowbreak{}Distortion.\allowbreak{}rlhf\allowbreak{}Definition\allowbreak{}Spec}
\item \hyperlink{paper-prerequisite-e2a06898e9a1c2db}{Golz\allowbreak{}Haghtalab\allowbreak{}Yang2025\allowbreak{}Distortion.\allowbreak{}source\allowbreak{}Alignment\allowbreak{}Average\allowbreak{}Utility}
\item \hyperlink{paper-prerequisite-44ab9adb1cd36c3c}{Golz\allowbreak{}Haghtalab\allowbreak{}Yang2025\allowbreak{}Distortion.\allowbreak{}source\allowbreak{}Alignment\allowbreak{}Distortion}
\item \hyperlink{paper-prerequisite-d7683af95771516b}{Golz\allowbreak{}Haghtalab\allowbreak{}Yang2025\allowbreak{}Distortion.\allowbreak{}source\allowbreak{}Social\allowbreak{}Choice\allowbreak{}Average\allowbreak{}Utility}
\item \hyperlink{paper-prerequisite-e6860339dbf51649}{Golz\allowbreak{}Haghtalab\allowbreak{}Yang2025\allowbreak{}Distortion.\allowbreak{}source\allowbreak{}Social\allowbreak{}Choice\allowbreak{}Distortion}
\item \hyperlink{paper-prerequisite-a9bfbfb57b04c106}{Golz\allowbreak{}Haghtalab\allowbreak{}Yang2025\allowbreak{}Distortion.\allowbreak{}source\allowbreak{}Worst\allowbreak{}Case\allowbreak{}Alignment\allowbreak{}Distortion}
\item \hyperlink{paper-prerequisite-43538553aea8c3ab}{Golz\allowbreak{}Haghtalab\allowbreak{}Yang2025\allowbreak{}Distortion.\allowbreak{}source\allowbreak{}Worst\allowbreak{}Case\allowbreak{}Social\allowbreak{}Choice\allowbreak{}Distortion}
\end{itemize}
\item \hyperlink{source-section-f10b5f68e4acf3dc}{Source claims (16)}
\begin{itemize}[leftmargin=1.2em]
\item \hyperlink{source-claim-979ce61abfc9e53f}{lemma1\allowbreak{}Spec}
\item \hyperlink{source-claim-43099ab484dcdf43}{lemma10\allowbreak{}Spec}
\item \hyperlink{source-claim-069675611a64cc9c}{lemma11\allowbreak{}Spec}
\item \hyperlink{source-claim-99f78f35feb19dfc}{lemma15\allowbreak{}Spec}
\item \hyperlink{source-claim-c85b480e847e3345}{theorem2\allowbreak{}Spec}
\item \hyperlink{source-claim-c933c69486621c74}{theorem3\allowbreak{}Spec}
\item \hyperlink{source-claim-27fd47f63d6ea396}{theorem5\allowbreak{}Spec}
\item \hyperlink{source-claim-8c95df698cbd00a0}{theorem12\allowbreak{}Spec}
\item \hyperlink{source-claim-b1ccdc13d898700c}{theorem6\allowbreak{}Spec}
\item \hyperlink{source-claim-5067ef3ec8c43462}{theorem7\allowbreak{}Spec}
\item \hyperlink{source-claim-067d55f193fc0110}{corollary4\allowbreak{}Spec}
\item \hyperlink{source-claim-9824dd64b1ddd3ad}{proposition13\allowbreak{}Spec}
\item \hyperlink{source-claim-fd234c18942cf2cc}{corollary8\allowbreak{}Spec}
\item \hyperlink{source-claim-ddcbcef3944ac5d9}{proposition14\allowbreak{}Spec}
\item \hyperlink{source-claim-f1c80c98e2c74036}{theorem9\allowbreak{}Spec}
\item \hyperlink{source-claim-cc5e88fff44863d2}{appendix\allowbreak{}F3\allowbreak{}Dpo\allowbreak{}Rlhf\allowbreak{}Equivalence\allowbreak{}Spec}
\end{itemize}
\end{itemize}

\clearpage
\hypertarget{governing-declaration-section-5ef5ef0364b6939c}{}
\section*{Governing Lean declarations}
These exact graph-bound declaration bodies are retained by the displayed specifications or reviewed prerequisites. They are supporting context and do not add source-result rows or semantic judgments.
\hypertarget{governing-declaration-b7cf9e04e23161e5}{}
\subsection*{\small\ttfamily\raggedright Applied\allowbreak{}Modeling\allowbreak{}Lib.\allowbreak{}Alignment.\allowbreak{}Welfare.\allowbreak{}context\allowbreak{}Free\allowbreak{}Alternative\allowbreak{}Policy}
\textbf{Declaration location:} reusable library
\paragraph{Lean declaration body}
\begin{ReviewVerbatim}
type:
{Alternative : Type u_1} → PMF Alternative → PUnit.{u_2 + 1} → PMF Alternative

value:
fun {Alternative : Type u_1} (policy : PMF Alternative) (a : PUnit.{u_2 + 1}) => policy
\end{ReviewVerbatim}


\hypertarget{governing-declaration-a4d623b8c608b566}{}
\subsection*{\small\ttfamily\raggedright Applied\allowbreak{}Modeling\allowbreak{}Lib.\allowbreak{}Alignment.\allowbreak{}Welfare.\allowbreak{}sigmoid\allowbreak{}Chord\allowbreak{}Slope}
\textbf{Declaration location:} reusable library
\paragraph{Lean declaration body}
\begin{ReviewVerbatim}
type:
ℝ → ℝ

value:
fun (beta : ℝ) => (beta.sigmoid - 1 / 2) / beta
\end{ReviewVerbatim}


\hypertarget{governing-declaration-bc3d6ca8e7e0ac9d}{}
\subsection*{\small\ttfamily\raggedright Applied\allowbreak{}Modeling\allowbreak{}Lib.\allowbreak{}Learning.\allowbreak{}Human\allowbreak{}Feedback.\allowbreak{}Binary\allowbreak{}Pairwise\allowbreak{}Report}
\textbf{Declaration location:} reusable library
\paragraph{Lean declaration body}
\begin{ReviewVerbatim}
type:
Type u → Type u

value:
fun (Alternative : Type u) => (Alternative × Alternative) × Bool
\end{ReviewVerbatim}


\hypertarget{governing-declaration-588529a2add6b6a9}{}
\subsection*{\small\ttfamily\raggedright Applied\allowbreak{}Modeling\allowbreak{}Lib.\allowbreak{}Learning.\allowbreak{}Human\allowbreak{}Feedback.\allowbreak{}Correlated\allowbreak{}Pair\allowbreak{}Sampling}
\textbf{Declaration location:} reusable library
\paragraph{Lean declaration body}
\begin{ReviewVerbatim}
type:
Type u_1 → Type u_1

value:
fun (Response : Type u_1) => PMF (Response × Response)
\end{ReviewVerbatim}


\hypertarget{governing-declaration-4ded7e185e3a56c5}{}
\subsection*{\small\ttfamily\raggedright Applied\allowbreak{}Modeling\allowbreak{}Lib.\allowbreak{}Learning.\allowbreak{}Human\allowbreak{}Feedback.\allowbreak{}Correlated\allowbreak{}Pair\allowbreak{}Sampling.\allowbreak{}Has\allowbreak{}Full\allowbreak{}Off\allowbreak{}Diagonal\allowbreak{}Support}
\textbf{Declaration location:} reusable library
\paragraph{Lean declaration body}
\begin{ReviewVerbatim}
type:
{Response : Type u_1} → AppliedModelingLib.Learning.HumanFeedback.CorrelatedPairSampling Response → Prop

value:
fun {Response : Type u_1} (sampling : AppliedModelingLib.Learning.HumanFeedback.CorrelatedPairSampling Response) =>
  ∀ (first second : Response), first ≠ second → 0 < (sampling (first, second)).toReal
\end{ReviewVerbatim}

\paragraph{Direct retained dependencies}
\begin{itemize}\raggedright
\item \hyperlink{governing-declaration-588529a2add6b6a9}{Applied\allowbreak{}Modeling\allowbreak{}Lib.\allowbreak{}Learning.\allowbreak{}Human\allowbreak{}Feedback.\allowbreak{}Correlated\allowbreak{}Pair\allowbreak{}Sampling}
\end{itemize}
\hypertarget{governing-declaration-28f7d7214bf02489}{}
\subsection*{\small\ttfamily\raggedright Applied\allowbreak{}Modeling\allowbreak{}Lib.\allowbreak{}Learning.\allowbreak{}Human\allowbreak{}Feedback.\allowbreak{}Correlated\allowbreak{}Pair\allowbreak{}Sampling.\allowbreak{}Has\allowbreak{}Zero\allowbreak{}Diagonal}
\textbf{Declaration location:} reusable library
\paragraph{Lean declaration body}
\begin{ReviewVerbatim}
type:
{Response : Type u_1} → AppliedModelingLib.Learning.HumanFeedback.CorrelatedPairSampling Response → Prop

value:
fun {Response : Type u_1} (sampling : AppliedModelingLib.Learning.HumanFeedback.CorrelatedPairSampling Response) =>
  ∀ (response : Response), (sampling (response, response)).toReal = 0
\end{ReviewVerbatim}

\paragraph{Direct retained dependencies}
\begin{itemize}\raggedright
\item \hyperlink{governing-declaration-588529a2add6b6a9}{Applied\allowbreak{}Modeling\allowbreak{}Lib.\allowbreak{}Learning.\allowbreak{}Human\allowbreak{}Feedback.\allowbreak{}Correlated\allowbreak{}Pair\allowbreak{}Sampling}
\end{itemize}
\hypertarget{governing-declaration-11432d593d43fd7d}{}
\subsection*{\small\ttfamily\raggedright Applied\allowbreak{}Modeling\allowbreak{}Lib.\allowbreak{}Learning.\allowbreak{}Human\allowbreak{}Feedback.\allowbreak{}Correlated\allowbreak{}Pair\allowbreak{}Sampling.\allowbreak{}Is\allowbreak{}Symmetric}
\textbf{Declaration location:} reusable library
\paragraph{Lean declaration body}
\begin{ReviewVerbatim}
type:
{Response : Type u_1} → AppliedModelingLib.Learning.HumanFeedback.CorrelatedPairSampling Response → Prop

value:
fun {Response : Type u_1} (sampling : AppliedModelingLib.Learning.HumanFeedback.CorrelatedPairSampling Response) =>
  ∀ (first second : Response), (sampling (first, second)).toReal = (sampling (second, first)).toReal
\end{ReviewVerbatim}

\paragraph{Direct retained dependencies}
\begin{itemize}\raggedright
\item \hyperlink{governing-declaration-588529a2add6b6a9}{Applied\allowbreak{}Modeling\allowbreak{}Lib.\allowbreak{}Learning.\allowbreak{}Human\allowbreak{}Feedback.\allowbreak{}Correlated\allowbreak{}Pair\allowbreak{}Sampling}
\end{itemize}
\hypertarget{governing-declaration-2ba23f394a8d08a2}{}
\subsection*{\small\ttfamily\raggedright Applied\allowbreak{}Modeling\allowbreak{}Lib.\allowbreak{}Learning.\allowbreak{}Human\allowbreak{}Feedback.\allowbreak{}Pairwise\allowbreak{}Count\allowbreak{}Dataset}
\textbf{Declaration location:} reusable library
\paragraph{Lean declaration body}
\begin{ReviewVerbatim}
field count:
{Alternative : Type u_1} →
  AppliedModelingLib.Learning.HumanFeedback.PairwiseCountDataset Alternative → Alternative → Alternative → ℕ

field diagonal_zero:
∀ {Alternative : Type u_1} (self : AppliedModelingLib.Learning.HumanFeedback.PairwiseCountDataset Alternative)
  (alternative : Alternative), self.count alternative alternative = 0
\end{ReviewVerbatim}


\hypertarget{governing-declaration-3449f3559185d0fd}{}
\subsection*{\small\ttfamily\raggedright Applied\allowbreak{}Modeling\allowbreak{}Lib.\allowbreak{}Learning.\allowbreak{}Human\allowbreak{}Feedback.\allowbreak{}Pairwise\allowbreak{}Count\allowbreak{}Dataset.\allowbreak{}CDF\allowbreak{}Like\allowbreak{}Pairwise\allowbreak{}Link}
\textbf{Declaration location:} reusable library
\paragraph{Lean declaration body}
\begin{ReviewVerbatim}
field toFun:
AppliedModelingLib.Learning.HumanFeedback.PairwiseCountDataset.CDFLikePairwiseLink → ℝ → ℝ

field nonneg:
∀ (self : AppliedModelingLib.Learning.HumanFeedback.PairwiseCountDataset.CDFLikePairwiseLink) (gap : ℝ),
  0 ≤ self.toFun gap

field le_one:
∀ (self : AppliedModelingLib.Learning.HumanFeedback.PairwiseCountDataset.CDFLikePairwiseLink) (gap : ℝ),
  self.toFun gap ≤ 1

field monotone:
∀ (self : AppliedModelingLib.Learning.HumanFeedback.PairwiseCountDataset.CDFLikePairwiseLink), Monotone self.toFun

field complementary:
∀ (self : AppliedModelingLib.Learning.HumanFeedback.PairwiseCountDataset.CDFLikePairwiseLink) (gap : ℝ),
  self.toFun (-gap) + self.toFun gap = 1

field tendsto_atBot_zero:
∀ (self : AppliedModelingLib.Learning.HumanFeedback.PairwiseCountDataset.CDFLikePairwiseLink),
  Filter.Tendsto self.toFun Filter.atBot (nhds 0)

field tendsto_atTop_one:
∀ (self : AppliedModelingLib.Learning.HumanFeedback.PairwiseCountDataset.CDFLikePairwiseLink),
  Filter.Tendsto self.toFun Filter.atTop (nhds 1)
\end{ReviewVerbatim}


\hypertarget{governing-declaration-87eda978eb1b0c9d}{}
\subsection*{\small\ttfamily\raggedright Applied\allowbreak{}Modeling\allowbreak{}Lib.\allowbreak{}Learning.\allowbreak{}Human\allowbreak{}Feedback.\allowbreak{}Pairwise\allowbreak{}Count\allowbreak{}Dataset.\allowbreak{}Score\allowbreak{}Vector}
\textbf{Declaration location:} reusable library
\paragraph{Lean declaration body}
\begin{ReviewVerbatim}
type:
Type u_1 → Type u_1

value:
fun (Alternative : Type u_1) => Alternative → ℝ
\end{ReviewVerbatim}


\hypertarget{governing-declaration-183c779b24cebcca}{}
\subsection*{\small\ttfamily\raggedright Applied\allowbreak{}Modeling\allowbreak{}Lib.\allowbreak{}Learning.\allowbreak{}Human\allowbreak{}Feedback.\allowbreak{}Pairwise\allowbreak{}Count\allowbreak{}Dataset.\allowbreak{}is\allowbreak{}Reference\allowbreak{}Normalized}
\textbf{Declaration location:} reusable library
\paragraph{Lean declaration body}
\begin{ReviewVerbatim}
type:
{Alternative : Type u_1} →
  Alternative → AppliedModelingLib.Learning.HumanFeedback.PairwiseCountDataset.ScoreVector Alternative → Prop

value:
fun {Alternative : Type u_1} (reference : Alternative)
    (score : AppliedModelingLib.Learning.HumanFeedback.PairwiseCountDataset.ScoreVector Alternative) =>
  score reference = 0
\end{ReviewVerbatim}

\paragraph{Direct retained dependencies}
\begin{itemize}\raggedright
\item \hyperlink{governing-declaration-87eda978eb1b0c9d}{Applied\allowbreak{}Modeling\allowbreak{}Lib.\allowbreak{}Learning.\allowbreak{}Human\allowbreak{}Feedback.\allowbreak{}Pairwise\allowbreak{}Count\allowbreak{}Dataset.\allowbreak{}Score\allowbreak{}Vector}
\end{itemize}
\hypertarget{governing-declaration-4f800616f542ca3c}{}
\subsection*{\small\ttfamily\raggedright Applied\allowbreak{}Modeling\allowbreak{}Lib.\allowbreak{}Learning.\allowbreak{}Human\allowbreak{}Feedback.\allowbreak{}Pairwise\allowbreak{}Count\allowbreak{}Dataset.\allowbreak{}random\allowbreak{}Utility\allowbreak{}Win\allowbreak{}Probability}
\textbf{Declaration location:} reusable library
\paragraph{Lean declaration body}
\begin{ReviewVerbatim}
type:
{Alternative : Type u_1} →
  (ℝ → ℝ) →
    AppliedModelingLib.Learning.HumanFeedback.PairwiseCountDataset.ScoreVector Alternative →
      Alternative → Alternative → ℝ

value:
fun {Alternative : Type u_1} (link : ℝ → ℝ)
    (score : AppliedModelingLib.Learning.HumanFeedback.PairwiseCountDataset.ScoreVector Alternative)
    (first second : Alternative) =>
  link (score first - score second)
\end{ReviewVerbatim}

\paragraph{Direct retained dependencies}
\begin{itemize}\raggedright
\item \hyperlink{governing-declaration-87eda978eb1b0c9d}{Applied\allowbreak{}Modeling\allowbreak{}Lib.\allowbreak{}Learning.\allowbreak{}Human\allowbreak{}Feedback.\allowbreak{}Pairwise\allowbreak{}Count\allowbreak{}Dataset.\allowbreak{}Score\allowbreak{}Vector}
\end{itemize}
\hypertarget{governing-declaration-10adb9e7c1ff9f4c}{}
\subsection*{\small\ttfamily\raggedright Applied\allowbreak{}Modeling\allowbreak{}Lib.\allowbreak{}Learning.\allowbreak{}Human\allowbreak{}Feedback.\allowbreak{}Pairwise\allowbreak{}Count\allowbreak{}Dataset.\allowbreak{}pairwise\allowbreak{}Log\allowbreak{}Likelihood}
\textbf{Declaration location:} reusable library
\paragraph{Lean declaration body}
\begin{ReviewVerbatim}
type:
{Alternative : Type u_1} →
  [Fintype Alternative] →
    [DecidableEq Alternative] →
      AppliedModelingLib.Learning.HumanFeedback.PairwiseCountDataset Alternative →
        (ℝ → ℝ) → AppliedModelingLib.Learning.HumanFeedback.PairwiseCountDataset.ScoreVector Alternative → ℝ

value:
fun {Alternative : Type u_1} [Fintype Alternative] [DecidableEq Alternative]
    (dataset : AppliedModelingLib.Learning.HumanFeedback.PairwiseCountDataset Alternative) (link : ℝ → ℝ)
    (score : AppliedModelingLib.Learning.HumanFeedback.PairwiseCountDataset.ScoreVector Alternative) =>
  ∑ pair ∈ Finset.univ.offDiag,
    ↑(dataset.count pair.1 pair.2) *
      Real.log
        (AppliedModelingLib.Learning.HumanFeedback.PairwiseCountDataset.randomUtilityWinProbability link score pair.1
          pair.2)
\end{ReviewVerbatim}

\paragraph{Direct retained dependencies}
\begin{itemize}\raggedright
\item \hyperlink{governing-declaration-2ba23f394a8d08a2}{Applied\allowbreak{}Modeling\allowbreak{}Lib.\allowbreak{}Learning.\allowbreak{}Human\allowbreak{}Feedback.\allowbreak{}Pairwise\allowbreak{}Count\allowbreak{}Dataset}
\item \hyperlink{governing-declaration-87eda978eb1b0c9d}{Applied\allowbreak{}Modeling\allowbreak{}Lib.\allowbreak{}Learning.\allowbreak{}Human\allowbreak{}Feedback.\allowbreak{}Pairwise\allowbreak{}Count\allowbreak{}Dataset.\allowbreak{}Score\allowbreak{}Vector}
\item \hyperlink{governing-declaration-4f800616f542ca3c}{Applied\allowbreak{}Modeling\allowbreak{}Lib.\allowbreak{}Learning.\allowbreak{}Human\allowbreak{}Feedback.\allowbreak{}Pairwise\allowbreak{}Count\allowbreak{}Dataset.\allowbreak{}random\allowbreak{}Utility\allowbreak{}Win\allowbreak{}Probability}
\end{itemize}
\hypertarget{governing-declaration-cba21d23e0e5aa10}{}
\subsection*{\small\ttfamily\raggedright Applied\allowbreak{}Modeling\allowbreak{}Lib.\allowbreak{}Learning.\allowbreak{}Human\allowbreak{}Feedback.\allowbreak{}Pairwise\allowbreak{}Count\allowbreak{}Dataset.\allowbreak{}is\allowbreak{}Pairwise\allowbreak{}MLE}
\textbf{Declaration location:} reusable library
\paragraph{Lean declaration body}
\begin{ReviewVerbatim}
type:
{Alternative : Type u_1} →
  [Fintype Alternative] →
    [DecidableEq Alternative] →
      AppliedModelingLib.Learning.HumanFeedback.PairwiseCountDataset Alternative →
        (ℝ → ℝ) →
          Alternative → AppliedModelingLib.Learning.HumanFeedback.PairwiseCountDataset.ScoreVector Alternative → Prop

value:
fun {Alternative : Type u_1} [Fintype Alternative] [DecidableEq Alternative]
    (dataset : AppliedModelingLib.Learning.HumanFeedback.PairwiseCountDataset Alternative) (link : ℝ → ℝ)
    (reference : Alternative)
    (score : AppliedModelingLib.Learning.HumanFeedback.PairwiseCountDataset.ScoreVector Alternative) =>
  AppliedModelingLib.Learning.HumanFeedback.PairwiseCountDataset.isReferenceNormalized reference score ∧
    IsMaxOn (dataset.pairwiseLogLikelihood link)
      {candidate : AppliedModelingLib.Learning.HumanFeedback.PairwiseCountDataset.ScoreVector Alternative |
        AppliedModelingLib.Learning.HumanFeedback.PairwiseCountDataset.isReferenceNormalized reference candidate}
      score
\end{ReviewVerbatim}

\paragraph{Direct retained dependencies}
\begin{itemize}\raggedright
\item \hyperlink{governing-declaration-2ba23f394a8d08a2}{Applied\allowbreak{}Modeling\allowbreak{}Lib.\allowbreak{}Learning.\allowbreak{}Human\allowbreak{}Feedback.\allowbreak{}Pairwise\allowbreak{}Count\allowbreak{}Dataset}
\item \hyperlink{governing-declaration-87eda978eb1b0c9d}{Applied\allowbreak{}Modeling\allowbreak{}Lib.\allowbreak{}Learning.\allowbreak{}Human\allowbreak{}Feedback.\allowbreak{}Pairwise\allowbreak{}Count\allowbreak{}Dataset.\allowbreak{}Score\allowbreak{}Vector}
\item \hyperlink{governing-declaration-183c779b24cebcca}{Applied\allowbreak{}Modeling\allowbreak{}Lib.\allowbreak{}Learning.\allowbreak{}Human\allowbreak{}Feedback.\allowbreak{}Pairwise\allowbreak{}Count\allowbreak{}Dataset.\allowbreak{}is\allowbreak{}Reference\allowbreak{}Normalized}
\item \hyperlink{governing-declaration-10adb9e7c1ff9f4c}{Applied\allowbreak{}Modeling\allowbreak{}Lib.\allowbreak{}Learning.\allowbreak{}Human\allowbreak{}Feedback.\allowbreak{}Pairwise\allowbreak{}Count\allowbreak{}Dataset.\allowbreak{}pairwise\allowbreak{}Log\allowbreak{}Likelihood}
\end{itemize}
\hypertarget{governing-declaration-ed6af38b0210f5ec}{}
\subsection*{\small\ttfamily\raggedright Applied\allowbreak{}Modeling\allowbreak{}Lib.\allowbreak{}Learning.\allowbreak{}Human\allowbreak{}Feedback.\allowbreak{}Pairwise\allowbreak{}Count\allowbreak{}Dataset.\allowbreak{}sigmoid\allowbreak{}CDF\allowbreak{}Like\allowbreak{}Pairwise\allowbreak{}Link}
\textbf{Declaration location:} reusable library
\paragraph{Lean declaration body}
\begin{ReviewVerbatim}
type:
AppliedModelingLib.Learning.HumanFeedback.PairwiseCountDataset.CDFLikePairwiseLink

value:
{ toFun := Real.sigmoid, nonneg := Real.sigmoid_nonneg, le_one := Real.sigmoid_le_one,
  monotone := Real.sigmoid_monotone,
  complementary := AppliedModelingLib.Learning.HumanFeedback.PairwiseCountDataset.sigmoidCDFLikePairwiseLink._proof_1,
  tendsto_atBot_zero := Real.tendsto_sigmoid_atBot, tendsto_atTop_one := Real.tendsto_sigmoid_atTop }
\end{ReviewVerbatim}

\paragraph{Direct retained dependencies}
\begin{itemize}\raggedright
\item \hyperlink{governing-declaration-3449f3559185d0fd}{Applied\allowbreak{}Modeling\allowbreak{}Lib.\allowbreak{}Learning.\allowbreak{}Human\allowbreak{}Feedback.\allowbreak{}Pairwise\allowbreak{}Count\allowbreak{}Dataset.\allowbreak{}CDF\allowbreak{}Like\allowbreak{}Pairwise\allowbreak{}Link}
\end{itemize}
\hypertarget{governing-declaration-16b50286b4cd31a4}{}
\subsection*{\small\ttfamily\raggedright Applied\allowbreak{}Modeling\allowbreak{}Lib.\allowbreak{}Learning.\allowbreak{}Human\allowbreak{}Feedback.\allowbreak{}Reward\allowbreak{}Equivalent}
\textbf{Declaration location:} reusable library
\paragraph{Lean declaration body}
\begin{ReviewVerbatim}
type:
{Context : Type u_1} → {Response : Type u_2} → (Context → Response → ℝ) → (Context → Response → ℝ) → Prop

value:
fun {Context : Type u_1} {Response : Type u_2} (first second : Context → Response → ℝ) =>
  ∃ (shift : Context → ℝ),
    ∀ (context : Context) (response : Response), second context response = first context response + shift context
\end{ReviewVerbatim}


\hypertarget{governing-declaration-8f11e069c88cc540}{}
\subsection*{\small\ttfamily\raggedright Applied\allowbreak{}Modeling\allowbreak{}Lib.\allowbreak{}Learning.\allowbreak{}Human\allowbreak{}Feedback.\allowbreak{}binary\allowbreak{}Report\allowbreak{}Win}
\textbf{Declaration location:} reusable library
\paragraph{Lean declaration body}
\begin{ReviewVerbatim}
type:
{Alternative : Type u_1} →
  [DecidableEq Alternative] →
    Alternative → Alternative → AppliedModelingLib.Learning.HumanFeedback.BinaryPairwiseReport Alternative → Prop

value:
fun {Alternative : Type u_1} [DecidableEq Alternative] (winner loser : Alternative)
    (report : AppliedModelingLib.Learning.HumanFeedback.BinaryPairwiseReport Alternative) =>
  winner ≠ loser ∧
    (report.1.1 = winner ∧ report.1.2 = loser ∧ report.2 = true ∨
      report.1.1 = loser ∧ report.1.2 = winner ∧ report.2 = false)
\end{ReviewVerbatim}

\paragraph{Direct retained dependencies}
\begin{itemize}\raggedright
\item \hyperlink{governing-declaration-bc3d6ca8e7e0ac9d}{Applied\allowbreak{}Modeling\allowbreak{}Lib.\allowbreak{}Learning.\allowbreak{}Human\allowbreak{}Feedback.\allowbreak{}Binary\allowbreak{}Pairwise\allowbreak{}Report}
\end{itemize}
\hypertarget{governing-declaration-449f87e91bf37ce0}{}
\subsection*{\small\ttfamily\raggedright Applied\allowbreak{}Modeling\allowbreak{}Lib.\allowbreak{}Learning.\allowbreak{}Human\allowbreak{}Feedback.\allowbreak{}Pairwise\allowbreak{}Count\allowbreak{}Dataset.\allowbreak{}of\allowbreak{}Binary\allowbreak{}Reports}
\textbf{Declaration location:} reusable library
\paragraph{Lean declaration body}
\begin{ReviewVerbatim}
type:
{Alternative : Type u_1} →
  [DecidableEq Alternative] →
    {horizon : ℕ} →
      (Fin horizon → AppliedModelingLib.Learning.HumanFeedback.BinaryPairwiseReport Alternative) →
        AppliedModelingLib.Learning.HumanFeedback.PairwiseCountDataset Alternative

value:
fun {Alternative : Type u_1} [DecidableEq Alternative] {horizon : ℕ}
    (sample : Fin horizon → AppliedModelingLib.Learning.HumanFeedback.BinaryPairwiseReport Alternative) =>
  {
    count := fun (winner loser : Alternative) =>
      {index : Fin horizon |
          AppliedModelingLib.Learning.HumanFeedback.binaryReportWin winner loser (sample index)}.card,
    diagonal_zero := AppliedModelingLib.Learning.HumanFeedback.PairwiseCountDataset.ofBinaryReports._proof_1 sample }
\end{ReviewVerbatim}

\paragraph{Direct retained dependencies}
\begin{itemize}\raggedright
\item \hyperlink{governing-declaration-bc3d6ca8e7e0ac9d}{Applied\allowbreak{}Modeling\allowbreak{}Lib.\allowbreak{}Learning.\allowbreak{}Human\allowbreak{}Feedback.\allowbreak{}Binary\allowbreak{}Pairwise\allowbreak{}Report}
\item \hyperlink{governing-declaration-2ba23f394a8d08a2}{Applied\allowbreak{}Modeling\allowbreak{}Lib.\allowbreak{}Learning.\allowbreak{}Human\allowbreak{}Feedback.\allowbreak{}Pairwise\allowbreak{}Count\allowbreak{}Dataset}
\item \hyperlink{governing-declaration-8f11e069c88cc540}{Applied\allowbreak{}Modeling\allowbreak{}Lib.\allowbreak{}Learning.\allowbreak{}Human\allowbreak{}Feedback.\allowbreak{}binary\allowbreak{}Report\allowbreak{}Win}
\end{itemize}
\hypertarget{governing-declaration-acbe2789b3d93e71}{}
\subsection*{\small\ttfamily\raggedright Applied\allowbreak{}Modeling\allowbreak{}Lib.\allowbreak{}Learning.\allowbreak{}Human\allowbreak{}Feedback.\allowbreak{}correlated\allowbreak{}Bernoulli\allowbreak{}Report}
\textbf{Declaration location:} reusable library
\paragraph{Lean declaration body}
\begin{ReviewVerbatim}
type:
(probability : ℝ) → 0 ≤ probability → probability ≤ 1 → PMF Bool

value:
fun (probability : ℝ) (hnonneg : 0 ≤ probability) (hle_one : probability ≤ 1) =>
  PMF.bernoulli ⟨probability, hnonneg⟩
    (AppliedModelingLib.Learning.HumanFeedback.correlatedBernoulliReport._proof_1 probability hnonneg hle_one)
\end{ReviewVerbatim}


\hypertarget{governing-declaration-33b603cd884f09c4}{}
\subsection*{\small\ttfamily\raggedright Applied\allowbreak{}Modeling\allowbreak{}Lib.\allowbreak{}Learning.\allowbreak{}Human\allowbreak{}Feedback.\allowbreak{}correlated\allowbreak{}One\allowbreak{}Comparison\allowbreak{}Report\allowbreak{}Law}
\textbf{Declaration location:} reusable library
\paragraph{Lean declaration body}
\begin{ReviewVerbatim}
type:
{Response : Type u_1} →
  [Fintype Response] →
    [DecidableEq Response] →
      AppliedModelingLib.Learning.HumanFeedback.CorrelatedPairSampling Response →
        AppliedModelingLib.Learning.HumanFeedback.PairwisePreference PUnit.{u_2 + 1} Response →
          PMF (AppliedModelingLib.Learning.HumanFeedback.BinaryPairwiseReport Response)

value:
fun {Response : Type u_1} [Fintype Response] [DecidableEq Response]
    (sampling : AppliedModelingLib.Learning.HumanFeedback.CorrelatedPairSampling Response)
    (preference : AppliedModelingLib.Learning.HumanFeedback.PairwisePreference PUnit.{u_2 + 1} Response) =>
  PMF.bind sampling fun (pair : Response × Response) =>
    PMF.map (fun (outcome : Bool) => (pair, outcome))
      (AppliedModelingLib.Learning.HumanFeedback.correlatedBernoulliReport (preference.prob PUnit.unit pair.1 pair.2)
        (AppliedModelingLib.Learning.HumanFeedback.correlatedOneComparisonReportLaw._proof_1 preference pair)
        (AppliedModelingLib.Learning.HumanFeedback.correlatedOneComparisonReportLaw._proof_2 preference pair))
\end{ReviewVerbatim}

\paragraph{Direct retained dependencies}
\begin{itemize}\raggedright
\item \hyperlink{governing-declaration-bc3d6ca8e7e0ac9d}{Applied\allowbreak{}Modeling\allowbreak{}Lib.\allowbreak{}Learning.\allowbreak{}Human\allowbreak{}Feedback.\allowbreak{}Binary\allowbreak{}Pairwise\allowbreak{}Report}
\item \hyperlink{governing-declaration-588529a2add6b6a9}{Applied\allowbreak{}Modeling\allowbreak{}Lib.\allowbreak{}Learning.\allowbreak{}Human\allowbreak{}Feedback.\allowbreak{}Correlated\allowbreak{}Pair\allowbreak{}Sampling}
\item \hyperlink{library-prerequisite-7e2db20cf087d9ba}{Applied\allowbreak{}Modeling\allowbreak{}Lib.\allowbreak{}Learning.\allowbreak{}Human\allowbreak{}Feedback.\allowbreak{}Pairwise\allowbreak{}Preference}
\item \hyperlink{governing-declaration-acbe2789b3d93e71}{Applied\allowbreak{}Modeling\allowbreak{}Lib.\allowbreak{}Learning.\allowbreak{}Human\allowbreak{}Feedback.\allowbreak{}correlated\allowbreak{}Bernoulli\allowbreak{}Report}
\end{itemize}
\hypertarget{governing-declaration-959259c0698492d3}{}
\subsection*{\small\ttfamily\raggedright Applied\allowbreak{}Modeling\allowbreak{}Lib.\allowbreak{}PMF\allowbreak{}Kernel}
\textbf{Declaration location:} reusable library
\paragraph{Lean declaration body}
\begin{ReviewVerbatim}
type:
Type u_1 → Type u_2 → Type (max u_1 u_2)

value:
fun (Input : Type u_1) (Output : Type u_2) => Input → PMF Output
\end{ReviewVerbatim}


\hypertarget{governing-declaration-577316c60f45ef11}{}
\subsection*{\small\ttfamily\raggedright Applied\allowbreak{}Modeling\allowbreak{}Lib.\allowbreak{}Learning.\allowbreak{}Human\allowbreak{}Feedback.\allowbreak{}Finite\allowbreak{}Policy}
\textbf{Declaration location:} reusable library
\paragraph{Lean declaration body}
\begin{ReviewVerbatim}
type:
Type u_1 → Type u_2 → Type (max u_1 u_2)

value:
fun (Context : Type u_1) (Response : Type u_2) => AppliedModelingLib.PMFKernel Context Response
\end{ReviewVerbatim}

\paragraph{Direct retained dependencies}
\begin{itemize}\raggedright
\item \hyperlink{governing-declaration-959259c0698492d3}{Applied\allowbreak{}Modeling\allowbreak{}Lib.\allowbreak{}PMF\allowbreak{}Kernel}
\end{itemize}
\hypertarget{governing-declaration-0ade06a0ad9064a9}{}
\subsection*{\small\ttfamily\raggedright Applied\allowbreak{}Modeling\allowbreak{}Lib.\allowbreak{}Learning.\allowbreak{}Human\allowbreak{}Feedback.\allowbreak{}policy\allowbreak{}Log\allowbreak{}Ratio}
\textbf{Declaration location:} reusable library
\paragraph{Lean declaration body}
\begin{ReviewVerbatim}
type:
{Context : Type u_1} →
  {Response : Type u_2} →
    AppliedModelingLib.Learning.HumanFeedback.FinitePolicy Context Response →
      AppliedModelingLib.Learning.HumanFeedback.FinitePolicy Context Response → Context → Response → ℝ

value:
fun {Context : Type u_1} {Response : Type u_2}
    (candidate reference : AppliedModelingLib.Learning.HumanFeedback.FinitePolicy Context Response) (context : Context)
    (response : Response) =>
  Real.log ((candidate context) response).toReal - Real.log ((reference context) response).toReal
\end{ReviewVerbatim}

\paragraph{Direct retained dependencies}
\begin{itemize}\raggedright
\item \hyperlink{governing-declaration-577316c60f45ef11}{Applied\allowbreak{}Modeling\allowbreak{}Lib.\allowbreak{}Learning.\allowbreak{}Human\allowbreak{}Feedback.\allowbreak{}Finite\allowbreak{}Policy}
\end{itemize}
\hypertarget{governing-declaration-95571fb7643367ad}{}
\subsection*{\small\ttfamily\raggedright Applied\allowbreak{}Modeling\allowbreak{}Lib.\allowbreak{}Learning.\allowbreak{}Human\allowbreak{}Feedback.\allowbreak{}dpo\allowbreak{}Implicit\allowbreak{}Reward}
\textbf{Declaration location:} reusable library
\paragraph{Lean declaration body}
\begin{ReviewVerbatim}
type:
{Context : Type u_1} →
  {Response : Type u_2} →
    AppliedModelingLib.Learning.HumanFeedback.FinitePolicy Context Response →
      AppliedModelingLib.Learning.HumanFeedback.FinitePolicy Context Response → ℝ → Context → Response → ℝ

value:
fun {Context : Type u_1} {Response : Type u_2}
    (candidate reference : AppliedModelingLib.Learning.HumanFeedback.FinitePolicy Context Response) (klWeight : ℝ)
    (a : Context) (a_1 : Response) =>
  klWeight * AppliedModelingLib.Learning.HumanFeedback.policyLogRatio candidate reference a a_1
\end{ReviewVerbatim}

\paragraph{Direct retained dependencies}
\begin{itemize}\raggedright
\item \hyperlink{governing-declaration-577316c60f45ef11}{Applied\allowbreak{}Modeling\allowbreak{}Lib.\allowbreak{}Learning.\allowbreak{}Human\allowbreak{}Feedback.\allowbreak{}Finite\allowbreak{}Policy}
\item \hyperlink{governing-declaration-0ade06a0ad9064a9}{Applied\allowbreak{}Modeling\allowbreak{}Lib.\allowbreak{}Learning.\allowbreak{}Human\allowbreak{}Feedback.\allowbreak{}policy\allowbreak{}Log\allowbreak{}Ratio}
\end{itemize}
\hypertarget{governing-declaration-65283b6bc1ebd8c0}{}
\subsection*{\small\ttfamily\raggedright Applied\allowbreak{}Modeling\allowbreak{}Lib.\allowbreak{}Probability.\allowbreak{}finite\allowbreak{}Iid\allowbreak{}Score\allowbreak{}Sum}
\textbf{Declaration location:} reusable library
\paragraph{Lean declaration body}
\begin{ReviewVerbatim}
type:
{α : Type u_1} → {ι : Type u_2} → [Fintype ι] → (α → ℝ) → (ι → α) → ℝ

value:
fun {α : Type u_1} {ι : Type u_2} [Fintype ι] (score : α → ℝ) (sample : ι → α) => ∑ i : ι, score (sample i)
\end{ReviewVerbatim}


\hypertarget{governing-declaration-0294a37f5d0d89a5}{}
\subsection*{\small\ttfamily\raggedright Applied\allowbreak{}Modeling\allowbreak{}Lib.\allowbreak{}Probability.\allowbreak{}finite\allowbreak{}MGF}
\textbf{Declaration location:} reusable library
\paragraph{Lean declaration body}
\begin{ReviewVerbatim}
type:
{α : Type u_1} → [Fintype α] → PMF α → (α → ℝ) → ℝ → ℝ

value:
fun {α : Type u_1} [Fintype α] (μ : PMF α) (score : α → ℝ) (z : ℝ) => ∑ a : α, (μ a).toReal * Real.exp (z * score a)
\end{ReviewVerbatim}


\hypertarget{governing-declaration-8fd2dabd8bfaa1ab}{}
\subsection*{\small\ttfamily\raggedright Applied\allowbreak{}Modeling\allowbreak{}Lib.\allowbreak{}binary\allowbreak{}Mixture\allowbreak{}PMF}
\textbf{Declaration location:} reusable library
\paragraph{Lean declaration body}
\begin{ReviewVerbatim}
type:
{α : Type u_1} → (p : NNReal) → p ≤ 1 → PMF α → PMF α → PMF α

value:
fun {α : Type u_1} (p : NNReal) (hp : p ≤ 1) (selected unselected : PMF α) =>
  (PMF.bernoulli p hp).bind fun (takeSelected : Bool) => if takeSelected = true then selected else unselected
\end{ReviewVerbatim}


\hypertarget{governing-declaration-a7c19260714e6ad9}{}
\subsection*{\small\ttfamily\raggedright Applied\allowbreak{}Modeling\allowbreak{}Lib.\allowbreak{}exponential\allowbreak{}Tilt}
\textbf{Declaration location:} reusable library
\paragraph{Lean declaration body}
\begin{ReviewVerbatim}
type:
{Outcome : Type u_1} → [Fintype Outcome] → [DecidableEq Outcome] → PMF Outcome → (Outcome → ℝ) → ℝ → PMF Outcome

value:
fun {Outcome : Type u_1} [Fintype Outcome] [DecidableEq Outcome] (reference : PMF Outcome) (utility : Outcome → ℝ)
    (inverseTemperature : ℝ) =>
  PMF.ofFintype
    (fun (outcome : Outcome) =>
      ENNReal.ofReal
        ((reference outcome).toReal * Real.exp (inverseTemperature * utility outcome) /
          AppliedModelingLib.Probability.finiteMGF reference utility inverseTemperature))
    (AppliedModelingLib.exponentialTilt._proof_1 reference utility inverseTemperature)
\end{ReviewVerbatim}

\paragraph{Direct retained dependencies}
\begin{itemize}\raggedright
\item \hyperlink{governing-declaration-0294a37f5d0d89a5}{Applied\allowbreak{}Modeling\allowbreak{}Lib.\allowbreak{}Probability.\allowbreak{}finite\allowbreak{}MGF}
\end{itemize}
\hypertarget{governing-declaration-a868460489b0ccee}{}
\subsection*{\small\ttfamily\raggedright Applied\allowbreak{}Modeling\allowbreak{}Lib.\allowbreak{}Learning.\allowbreak{}Human\allowbreak{}Feedback.\allowbreak{}context\allowbreak{}Exponential\allowbreak{}Tilt\allowbreak{}Policy}
\textbf{Declaration location:} reusable library
\paragraph{Lean declaration body}
\begin{ReviewVerbatim}
type:
{Context : Type u_1} →
  {Response : Type u_2} →
    [Fintype Response] →
      [DecidableEq Response] →
        AppliedModelingLib.Learning.HumanFeedback.FinitePolicy Context Response →
          (Context → Response → ℝ) → ℝ → Context → PMF Response

value:
fun {Context : Type u_1} {Response : Type u_2} [Fintype Response] [DecidableEq Response]
    (reference : AppliedModelingLib.Learning.HumanFeedback.FinitePolicy Context Response)
    (utility : Context → Response → ℝ) (inverseTemperature : ℝ) (a : Context) =>
  AppliedModelingLib.exponentialTilt (reference a) (utility a) inverseTemperature
\end{ReviewVerbatim}

\paragraph{Direct retained dependencies}
\begin{itemize}\raggedright
\item \hyperlink{governing-declaration-577316c60f45ef11}{Applied\allowbreak{}Modeling\allowbreak{}Lib.\allowbreak{}Learning.\allowbreak{}Human\allowbreak{}Feedback.\allowbreak{}Finite\allowbreak{}Policy}
\item \hyperlink{governing-declaration-a7c19260714e6ad9}{Applied\allowbreak{}Modeling\allowbreak{}Lib.\allowbreak{}exponential\allowbreak{}Tilt}
\end{itemize}
\hypertarget{governing-declaration-eeb142baa1b200d1}{}
\subsection*{\small\ttfamily\raggedright Applied\allowbreak{}Modeling\allowbreak{}Lib.\allowbreak{}Learning.\allowbreak{}Human\allowbreak{}Feedback.\allowbreak{}dpo\allowbreak{}Optimal\allowbreak{}Policy}
\textbf{Declaration location:} reusable library
\paragraph{Lean declaration body}
\begin{ReviewVerbatim}
type:
{Context : Type u_1} →
  {Response : Type u_2} →
    [Fintype Response] →
      [DecidableEq Response] →
        AppliedModelingLib.Learning.HumanFeedback.FinitePolicy Context Response →
          (Context → Response → ℝ) → ℝ → Context → PMF Response

value:
fun {Context : Type u_1} {Response : Type u_2} [Fintype Response] [DecidableEq Response]
    (reference : AppliedModelingLib.Learning.HumanFeedback.FinitePolicy Context Response)
    (reward : Context → Response → ℝ) (klWeight : ℝ) (a : Context) =>
  AppliedModelingLib.Learning.HumanFeedback.contextExponentialTiltPolicy reference reward klWeight⁻¹ a
\end{ReviewVerbatim}

\paragraph{Direct retained dependencies}
\begin{itemize}\raggedright
\item \hyperlink{governing-declaration-577316c60f45ef11}{Applied\allowbreak{}Modeling\allowbreak{}Lib.\allowbreak{}Learning.\allowbreak{}Human\allowbreak{}Feedback.\allowbreak{}Finite\allowbreak{}Policy}
\item \hyperlink{governing-declaration-a868460489b0ccee}{Applied\allowbreak{}Modeling\allowbreak{}Lib.\allowbreak{}Learning.\allowbreak{}Human\allowbreak{}Feedback.\allowbreak{}context\allowbreak{}Exponential\allowbreak{}Tilt\allowbreak{}Policy}
\end{itemize}
\hypertarget{governing-declaration-a1c0a8e99d75c298}{}
\subsection*{\small\ttfamily\raggedright Applied\allowbreak{}Modeling\allowbreak{}Lib.\allowbreak{}Learning.\allowbreak{}Human\allowbreak{}Feedback.\allowbreak{}pointwise\allowbreak{}Policy\allowbreak{}KL\allowbreak{}Divergence}
\textbf{Declaration location:} reusable library
\paragraph{Lean declaration body}
\begin{ReviewVerbatim}
type:
{Context : Type u_1} →
  {Response : Type u_2} →
    [Fintype Response] →
      [DecidableEq Response] →
        AppliedModelingLib.Learning.HumanFeedback.FinitePolicy Context Response →
          AppliedModelingLib.Learning.HumanFeedback.FinitePolicy Context Response → Context → ℝ

value:
fun {Context : Type u_1} {Response : Type u_2} [Fintype Response] [DecidableEq Response]
    (first second : AppliedModelingLib.Learning.HumanFeedback.FinitePolicy Context Response) (context : Context) =>
  AppliedModelingLib.finiteKLDivergence (first context) (second context)
\end{ReviewVerbatim}

\paragraph{Direct retained dependencies}
\begin{itemize}\raggedright
\item \hyperlink{governing-declaration-577316c60f45ef11}{Applied\allowbreak{}Modeling\allowbreak{}Lib.\allowbreak{}Learning.\allowbreak{}Human\allowbreak{}Feedback.\allowbreak{}Finite\allowbreak{}Policy}
\item \hyperlink{library-prerequisite-d3573881ae45d4c6}{Applied\allowbreak{}Modeling\allowbreak{}Lib.\allowbreak{}finite\allowbreak{}KL\allowbreak{}Divergence}
\end{itemize}
\hypertarget{governing-declaration-779fe1555fbbcea5}{}
\subsection*{\small\ttfamily\raggedright Applied\allowbreak{}Modeling\allowbreak{}Lib.\allowbreak{}finite\allowbreak{}PMF\allowbreak{}Iid\allowbreak{}Path\allowbreak{}Measure}
\textbf{Declaration location:} reusable library
\paragraph{Lean declaration body}
\begin{ReviewVerbatim}
type:
{Signal : Type u_1} →
  [inst : MeasurableSpace Signal] →
    [Fintype Signal] → [DecidableEq Signal] → PMF Signal → MeasureTheory.Measure (ℕ → Signal)

value:
fun {Signal : Type u_1} [MeasurableSpace Signal] [Fintype Signal] [DecidableEq Signal] (law : PMF Signal) =>
  MeasureTheory.Measure.infinitePi fun (x : ℕ) => law.toMeasure
\end{ReviewVerbatim}


\hypertarget{governing-declaration-89219e02ea5145c7}{}
\subsection*{\small\ttfamily\raggedright Applied\allowbreak{}Modeling\allowbreak{}Lib.\allowbreak{}finite\allowbreak{}PMF\allowbreak{}Iid\allowbreak{}Path\allowbreak{}Prefix}
\textbf{Declaration location:} reusable library
\paragraph{Lean declaration body}
\begin{ReviewVerbatim}
type:
{Signal : Type u_1} → (ℕ → Signal) → (N : ℕ) → Fin N → Signal

value:
fun {Signal : Type u_1} (path : ℕ → Signal) (N : ℕ) (a : Fin N) => path ↑a
\end{ReviewVerbatim}


\hypertarget{governing-declaration-0153ce2f02e7dc12}{}
\subsection*{\small\ttfamily\raggedright Applied\allowbreak{}Modeling\allowbreak{}Lib.\allowbreak{}finite\allowbreak{}Weighted\allowbreak{}PMF}
\textbf{Declaration location:} reusable library
\paragraph{Lean declaration body}
\begin{ReviewVerbatim}
type:
{α : Type u_1} → [inst : Fintype α] → (weight : α → ℝ) → (∀ (a : α), 0 ≤ weight a) → 0 < ∑ a : α, weight a → PMF α

value:
fun {α : Type u_1} [Fintype α] (weight : α → ℝ) (hweight_nonneg : ∀ (a : α), 0 ≤ weight a)
    (htotal_pos : 0 < ∑ a : α, weight a) =>
  PMF.ofFintype (fun (a : α) => ENNReal.ofReal (weight a / ∑ b : α, weight b))
    (AppliedModelingLib.finiteWeightedPMF._proof_1 weight hweight_nonneg htotal_pos)
\end{ReviewVerbatim}


\hypertarget{governing-declaration-8e4b4389a4f6454b}{}
\subsection*{\small\ttfamily\raggedright Applied\allowbreak{}Modeling\allowbreak{}Lib.\allowbreak{}pmf\allowbreak{}Exp}
\textbf{Declaration location:} reusable library
\paragraph{Lean declaration body}
\begin{ReviewVerbatim}
type:
{α : Type u_1} → [Fintype α] → PMF α → (α → ℝ) → ℝ

value:
fun {α : Type u_1} [Fintype α] (μ : PMF α) (f : α → ℝ) => ∑ a : α, (μ a).toReal * f a
\end{ReviewVerbatim}


\hypertarget{governing-declaration-3808cf59ea39b435}{}
\subsection*{\small\ttfamily\raggedright Applied\allowbreak{}Modeling\allowbreak{}Lib.\allowbreak{}Alignment.\allowbreak{}Welfare.\allowbreak{}pairwise\allowbreak{}Borda\allowbreak{}Score}
\textbf{Declaration location:} reusable library
\paragraph{Lean declaration body}
\begin{ReviewVerbatim}
type:
{Alternative : Type u_1} →
  [Fintype Alternative] →
    [DecidableEq Alternative] →
      PMF Alternative →
        AppliedModelingLib.Learning.HumanFeedback.PairwisePreference PUnit.{u_2 + 1} Alternative → Alternative → ℝ

value:
fun {Alternative : Type u_1} [Fintype Alternative] [DecidableEq Alternative] (sampling : PMF Alternative)
    (preference : AppliedModelingLib.Learning.HumanFeedback.PairwisePreference PUnit.{u_2 + 1} Alternative)
    (alternative : Alternative) =>
  AppliedModelingLib.pmfExp sampling fun (opponent : Alternative) => preference.prob PUnit.unit alternative opponent
\end{ReviewVerbatim}

\paragraph{Direct retained dependencies}
\begin{itemize}\raggedright
\item \hyperlink{library-prerequisite-7e2db20cf087d9ba}{Applied\allowbreak{}Modeling\allowbreak{}Lib.\allowbreak{}Learning.\allowbreak{}Human\allowbreak{}Feedback.\allowbreak{}Pairwise\allowbreak{}Preference}
\item \hyperlink{governing-declaration-8e4b4389a4f6454b}{Applied\allowbreak{}Modeling\allowbreak{}Lib.\allowbreak{}pmf\allowbreak{}Exp}
\end{itemize}
\hypertarget{governing-declaration-603c21f4ad34025c}{}
\subsection*{\small\ttfamily\raggedright Applied\allowbreak{}Modeling\allowbreak{}Lib.\allowbreak{}Is\allowbreak{}Finite\allowbreak{}KL\allowbreak{}Extended\allowbreak{}Regularized\allowbreak{}Utility\allowbreak{}Max}
\textbf{Declaration location:} reusable library
\paragraph{Lean declaration body}
\begin{ReviewVerbatim}
type:
{Outcome : Type u_1} →
  [Fintype Outcome] → [DecidableEq Outcome] → PMF Outcome → PMF Outcome → (Outcome → ℝ) → WithTop ℝ → Prop

value:
fun {Outcome : Type u_1} [Fintype Outcome] [DecidableEq Outcome] (reference candidate : PMF Outcome)
    (utility : Outcome → ℝ) (klWeight : WithTop ℝ) =>
  (∃ (finiteWeight : ℝ),
      klWeight = ↑finiteWeight ∧
        0 ≤ finiteWeight ∧
          ∀ (other : PMF Outcome),
            AppliedModelingLib.pmfExp other utility -
                finiteWeight * AppliedModelingLib.finiteKLDivergence other reference ≤
              AppliedModelingLib.pmfExp candidate utility -
                finiteWeight * AppliedModelingLib.finiteKLDivergence candidate reference) ∨
    klWeight = ⊤ ∧ candidate = reference
\end{ReviewVerbatim}

\paragraph{Direct retained dependencies}
\begin{itemize}\raggedright
\item \hyperlink{library-prerequisite-d3573881ae45d4c6}{Applied\allowbreak{}Modeling\allowbreak{}Lib.\allowbreak{}finite\allowbreak{}KL\allowbreak{}Divergence}
\item \hyperlink{governing-declaration-8e4b4389a4f6454b}{Applied\allowbreak{}Modeling\allowbreak{}Lib.\allowbreak{}pmf\allowbreak{}Exp}
\end{itemize}
\hypertarget{governing-declaration-c528a6568d75fffd}{}
\subsection*{\small\ttfamily\raggedright Applied\allowbreak{}Modeling\allowbreak{}Lib.\allowbreak{}Learning.\allowbreak{}Human\allowbreak{}Feedback.\allowbreak{}context\allowbreak{}Averaged\allowbreak{}Policy\allowbreak{}KL\allowbreak{}Divergence}
\textbf{Declaration location:} reusable library
\paragraph{Lean declaration body}
\begin{ReviewVerbatim}
type:
{Context : Type u_1} →
  {Response : Type u_2} →
    [Fintype Context] →
      [DecidableEq Context] →
        [Fintype Response] →
          [DecidableEq Response] →
            PMF Context →
              AppliedModelingLib.Learning.HumanFeedback.FinitePolicy Context Response →
                AppliedModelingLib.Learning.HumanFeedback.FinitePolicy Context Response → ℝ

value:
fun {Context : Type u_1} {Response : Type u_2} [Fintype Context] [DecidableEq Context] [Fintype Response]
    [DecidableEq Response] (contextLaw : PMF Context)
    (first second : AppliedModelingLib.Learning.HumanFeedback.FinitePolicy Context Response) =>
  AppliedModelingLib.pmfExp contextLaw
    (AppliedModelingLib.Learning.HumanFeedback.pointwisePolicyKLDivergence first second)
\end{ReviewVerbatim}

\paragraph{Direct retained dependencies}
\begin{itemize}\raggedright
\item \hyperlink{governing-declaration-577316c60f45ef11}{Applied\allowbreak{}Modeling\allowbreak{}Lib.\allowbreak{}Learning.\allowbreak{}Human\allowbreak{}Feedback.\allowbreak{}Finite\allowbreak{}Policy}
\item \hyperlink{governing-declaration-a1c0a8e99d75c298}{Applied\allowbreak{}Modeling\allowbreak{}Lib.\allowbreak{}Learning.\allowbreak{}Human\allowbreak{}Feedback.\allowbreak{}pointwise\allowbreak{}Policy\allowbreak{}KL\allowbreak{}Divergence}
\item \hyperlink{governing-declaration-8e4b4389a4f6454b}{Applied\allowbreak{}Modeling\allowbreak{}Lib.\allowbreak{}pmf\allowbreak{}Exp}
\end{itemize}
\hypertarget{governing-declaration-4daa53fb741a5ea7}{}
\subsection*{\small\ttfamily\raggedright Applied\allowbreak{}Modeling\allowbreak{}Lib.\allowbreak{}Learning.\allowbreak{}Human\allowbreak{}Feedback.\allowbreak{}policy\allowbreak{}Expected\allowbreak{}Score}
\textbf{Declaration location:} reusable library
\paragraph{Lean declaration body}
\begin{ReviewVerbatim}
type:
{Context : Type u_1} →
  {Response : Type u_2} →
    [Fintype Context] →
      [DecidableEq Context] →
        [Fintype Response] →
          [DecidableEq Response] →
            PMF Context →
              AppliedModelingLib.Learning.HumanFeedback.FinitePolicy Context Response → (Context → Response → ℝ) → ℝ

value:
fun {Context : Type u_1} {Response : Type u_2} [Fintype Context] [DecidableEq Context] [Fintype Response]
    [DecidableEq Response] (contextLaw : PMF Context)
    (policy : AppliedModelingLib.Learning.HumanFeedback.FinitePolicy Context Response)
    (score : Context → Response → ℝ) =>
  AppliedModelingLib.pmfExp contextLaw fun (context : Context) =>
    AppliedModelingLib.pmfExp (policy context) (score context)
\end{ReviewVerbatim}

\paragraph{Direct retained dependencies}
\begin{itemize}\raggedright
\item \hyperlink{governing-declaration-577316c60f45ef11}{Applied\allowbreak{}Modeling\allowbreak{}Lib.\allowbreak{}Learning.\allowbreak{}Human\allowbreak{}Feedback.\allowbreak{}Finite\allowbreak{}Policy}
\item \hyperlink{governing-declaration-8e4b4389a4f6454b}{Applied\allowbreak{}Modeling\allowbreak{}Lib.\allowbreak{}pmf\allowbreak{}Exp}
\end{itemize}
\hypertarget{governing-declaration-09677001d465ec9d}{}
\subsection*{\small\ttfamily\raggedright Applied\allowbreak{}Modeling\allowbreak{}Lib.\allowbreak{}pmf\allowbreak{}Pair\allowbreak{}Exp}
\textbf{Declaration location:} reusable library
\paragraph{Lean declaration body}
\begin{ReviewVerbatim}
type:
{α : Type u_1} →
  {β : Type u_2} → [Fintype α] → [DecidableEq α] → [Fintype β] → [DecidableEq β] → PMF α → PMF β → (α → β → ℝ) → ℝ

value:
fun {α : Type u_1} {β : Type u_2} [Fintype α] [DecidableEq α] [Fintype β] [DecidableEq β] (μ : PMF α) (ν : PMF β)
    (f : α → β → ℝ) =>
  AppliedModelingLib.pmfExp μ fun (a : α) => AppliedModelingLib.pmfExp ν fun (b : β) => f a b
\end{ReviewVerbatim}

\paragraph{Direct retained dependencies}
\begin{itemize}\raggedright
\item \hyperlink{governing-declaration-8e4b4389a4f6454b}{Applied\allowbreak{}Modeling\allowbreak{}Lib.\allowbreak{}pmf\allowbreak{}Exp}
\end{itemize}
\hypertarget{governing-declaration-5c3611f5a76c4e77}{}
\subsection*{\small\ttfamily\raggedright Applied\allowbreak{}Modeling\allowbreak{}Lib.\allowbreak{}Alignment.\allowbreak{}Welfare.\allowbreak{}population\allowbreak{}Bradley\allowbreak{}Terry\allowbreak{}Policy\allowbreak{}Margin}
\textbf{Declaration location:} reusable library
\paragraph{Lean declaration body}
\begin{ReviewVerbatim}
type:
{User : Type u_1} →
  {Alternative : Type u_2} →
    [Fintype User] →
      [DecidableEq User] →
        [Fintype Alternative] →
          [DecidableEq Alternative] →
            PMF User →
              AppliedModelingLib.Alignment.Welfare.FiniteUtilityProfile User Alternative →
                ℝ → PMF Alternative → PMF Alternative → ℝ

value:
fun {User : Type u_1} {Alternative : Type u_2} [Fintype User] [DecidableEq User] [Fintype Alternative]
    [DecidableEq Alternative] (population : PMF User)
    (utility : AppliedModelingLib.Alignment.Welfare.FiniteUtilityProfile User Alternative) (btScale : ℝ)
    (firstPolicy secondPolicy : PMF Alternative) =>
  AppliedModelingLib.pmfPairExp firstPolicy secondPolicy fun (first second : Alternative) =>
    (AppliedModelingLib.Alignment.Welfare.populationBradleyTerryPreference population utility btScale).prob PUnit.unit
        first second -
      1 / 2
\end{ReviewVerbatim}

\paragraph{Direct retained dependencies}
\begin{itemize}\raggedright
\item \hyperlink{library-prerequisite-56c0734f8877ddbe}{Applied\allowbreak{}Modeling\allowbreak{}Lib.\allowbreak{}Alignment.\allowbreak{}Welfare.\allowbreak{}Finite\allowbreak{}Utility\allowbreak{}Profile}
\item \hyperlink{library-prerequisite-cf06d15d7f8b7c31}{Applied\allowbreak{}Modeling\allowbreak{}Lib.\allowbreak{}Alignment.\allowbreak{}Welfare.\allowbreak{}population\allowbreak{}Bradley\allowbreak{}Terry\allowbreak{}Preference}
\item \hyperlink{library-prerequisite-7e2db20cf087d9ba}{Applied\allowbreak{}Modeling\allowbreak{}Lib.\allowbreak{}Learning.\allowbreak{}Human\allowbreak{}Feedback.\allowbreak{}Pairwise\allowbreak{}Preference}
\item \hyperlink{governing-declaration-09677001d465ec9d}{Applied\allowbreak{}Modeling\allowbreak{}Lib.\allowbreak{}pmf\allowbreak{}Pair\allowbreak{}Exp}
\end{itemize}
\hypertarget{governing-declaration-904cae1c822ccc01}{}
\subsection*{\small\ttfamily\raggedright Applied\allowbreak{}Modeling\allowbreak{}Lib.\allowbreak{}Learning.\allowbreak{}Human\allowbreak{}Feedback.\allowbreak{}policy\allowbreak{}Preference}
\textbf{Declaration location:} reusable library
\paragraph{Lean declaration body}
\begin{ReviewVerbatim}
type:
{Context : Type u_1} →
  {Response : Type u_2} →
    [Fintype Context] →
      [DecidableEq Context] →
        [Fintype Response] →
          [DecidableEq Response] →
            PMF Context →
              AppliedModelingLib.Learning.HumanFeedback.PairwisePreference Context Response →
                AppliedModelingLib.Learning.HumanFeedback.FinitePolicy Context Response →
                  AppliedModelingLib.Learning.HumanFeedback.FinitePolicy Context Response → ℝ

value:
fun {Context : Type u_1} {Response : Type u_2} [Fintype Context] [DecidableEq Context] [Fintype Response]
    [DecidableEq Response] (contextLaw : PMF Context)
    (preference : AppliedModelingLib.Learning.HumanFeedback.PairwisePreference Context Response)
    (first second : AppliedModelingLib.Learning.HumanFeedback.FinitePolicy Context Response) =>
  AppliedModelingLib.pmfExp contextLaw fun (context : Context) =>
    AppliedModelingLib.pmfPairExp (first context) (second context) (preference.prob context)
\end{ReviewVerbatim}

\paragraph{Direct retained dependencies}
\begin{itemize}\raggedright
\item \hyperlink{governing-declaration-577316c60f45ef11}{Applied\allowbreak{}Modeling\allowbreak{}Lib.\allowbreak{}Learning.\allowbreak{}Human\allowbreak{}Feedback.\allowbreak{}Finite\allowbreak{}Policy}
\item \hyperlink{library-prerequisite-7e2db20cf087d9ba}{Applied\allowbreak{}Modeling\allowbreak{}Lib.\allowbreak{}Learning.\allowbreak{}Human\allowbreak{}Feedback.\allowbreak{}Pairwise\allowbreak{}Preference}
\item \hyperlink{governing-declaration-8e4b4389a4f6454b}{Applied\allowbreak{}Modeling\allowbreak{}Lib.\allowbreak{}pmf\allowbreak{}Exp}
\item \hyperlink{governing-declaration-09677001d465ec9d}{Applied\allowbreak{}Modeling\allowbreak{}Lib.\allowbreak{}pmf\allowbreak{}Pair\allowbreak{}Exp}
\end{itemize}
\hypertarget{governing-declaration-8e5c390864970969}{}
\subsection*{\small\ttfamily\raggedright Applied\allowbreak{}Modeling\allowbreak{}Lib.\allowbreak{}Game\allowbreak{}Theory.\allowbreak{}Preference\allowbreak{}Game.\allowbreak{}preference\allowbreak{}Game\allowbreak{}Payoff}
\textbf{Declaration location:} reusable library
\paragraph{Lean declaration body}
\begin{ReviewVerbatim}
type:
{Context : Type u_1} →
  {Response : Type u_2} →
    [Fintype Context] →
      [DecidableEq Context] →
        [Fintype Response] →
          [DecidableEq Response] →
            PMF Context →
              AppliedModelingLib.Learning.HumanFeedback.PairwisePreference Context Response →
                AppliedModelingLib.Learning.HumanFeedback.FinitePolicy Context Response →
                  AppliedModelingLib.Learning.HumanFeedback.FinitePolicy Context Response → ℝ

value:
fun {Context : Type u_1} {Response : Type u_2} [Fintype Context] [DecidableEq Context] [Fintype Response]
    [DecidableEq Response] (contextLaw : PMF Context)
    (preference : AppliedModelingLib.Learning.HumanFeedback.PairwisePreference Context Response)
    (first second : AppliedModelingLib.Learning.HumanFeedback.FinitePolicy Context Response) =>
  AppliedModelingLib.Learning.HumanFeedback.policyPreference contextLaw preference first second
\end{ReviewVerbatim}

\paragraph{Direct retained dependencies}
\begin{itemize}\raggedright
\item \hyperlink{governing-declaration-577316c60f45ef11}{Applied\allowbreak{}Modeling\allowbreak{}Lib.\allowbreak{}Learning.\allowbreak{}Human\allowbreak{}Feedback.\allowbreak{}Finite\allowbreak{}Policy}
\item \hyperlink{library-prerequisite-7e2db20cf087d9ba}{Applied\allowbreak{}Modeling\allowbreak{}Lib.\allowbreak{}Learning.\allowbreak{}Human\allowbreak{}Feedback.\allowbreak{}Pairwise\allowbreak{}Preference}
\item \hyperlink{governing-declaration-904cae1c822ccc01}{Applied\allowbreak{}Modeling\allowbreak{}Lib.\allowbreak{}Learning.\allowbreak{}Human\allowbreak{}Feedback.\allowbreak{}policy\allowbreak{}Preference}
\end{itemize}
\hypertarget{governing-declaration-ccdbd7fdd43ef252}{}
\subsection*{\small\ttfamily\raggedright Applied\allowbreak{}Modeling\allowbreak{}Lib.\allowbreak{}Game\allowbreak{}Theory.\allowbreak{}Preference\allowbreak{}Game.\allowbreak{}centered\allowbreak{}Preference\allowbreak{}Game\allowbreak{}Payoff}
\textbf{Declaration location:} reusable library
\paragraph{Lean declaration body}
\begin{ReviewVerbatim}
type:
{Context : Type u_1} →
  {Response : Type u_2} →
    [Fintype Context] →
      [DecidableEq Context] →
        [Fintype Response] →
          [DecidableEq Response] →
            PMF Context →
              AppliedModelingLib.Learning.HumanFeedback.PairwisePreference Context Response →
                AppliedModelingLib.Learning.HumanFeedback.FinitePolicy Context Response →
                  AppliedModelingLib.Learning.HumanFeedback.FinitePolicy Context Response → ℝ

value:
fun {Context : Type u_1} {Response : Type u_2} [Fintype Context] [DecidableEq Context] [Fintype Response]
    [DecidableEq Response] (contextLaw : PMF Context)
    (preference : AppliedModelingLib.Learning.HumanFeedback.PairwisePreference Context Response)
    (first second : AppliedModelingLib.Learning.HumanFeedback.FinitePolicy Context Response) =>
  AppliedModelingLib.GameTheory.PreferenceGame.preferenceGamePayoff contextLaw preference first second - 1 / 2
\end{ReviewVerbatim}

\paragraph{Direct retained dependencies}
\begin{itemize}\raggedright
\item \hyperlink{governing-declaration-8e5c390864970969}{Applied\allowbreak{}Modeling\allowbreak{}Lib.\allowbreak{}Game\allowbreak{}Theory.\allowbreak{}Preference\allowbreak{}Game.\allowbreak{}preference\allowbreak{}Game\allowbreak{}Payoff}
\item \hyperlink{governing-declaration-577316c60f45ef11}{Applied\allowbreak{}Modeling\allowbreak{}Lib.\allowbreak{}Learning.\allowbreak{}Human\allowbreak{}Feedback.\allowbreak{}Finite\allowbreak{}Policy}
\item \hyperlink{library-prerequisite-7e2db20cf087d9ba}{Applied\allowbreak{}Modeling\allowbreak{}Lib.\allowbreak{}Learning.\allowbreak{}Human\allowbreak{}Feedback.\allowbreak{}Pairwise\allowbreak{}Preference}
\end{itemize}
\hypertarget{governing-declaration-44fb06319bbd7683}{}
\subsection*{\small\ttfamily\raggedright Applied\allowbreak{}Modeling\allowbreak{}Lib.\allowbreak{}Game\allowbreak{}Theory.\allowbreak{}Preference\allowbreak{}Game.\allowbreak{}context\allowbreak{}Free\allowbreak{}Centered\allowbreak{}Preference\allowbreak{}Payoff}
\textbf{Declaration location:} reusable library
\paragraph{Lean declaration body}
\begin{ReviewVerbatim}
type:
{Response : Type u_1} →
  [Fintype Response] →
    [DecidableEq Response] →
      AppliedModelingLib.Learning.HumanFeedback.PairwisePreference PUnit.{u_2 + 1} Response →
        PMF Response → PMF Response → ℝ

value:
fun {Response : Type u_1} [Fintype Response] [DecidableEq Response]
    (preference : AppliedModelingLib.Learning.HumanFeedback.PairwisePreference PUnit.{u_2 + 1} Response)
    (first second : PMF Response) =>
  AppliedModelingLib.GameTheory.PreferenceGame.centeredPreferenceGamePayoff (PMF.pure PUnit.unit) preference
    (fun (x : PUnit.{u_2 + 1}) => first) fun (x : PUnit.{u_2 + 1}) => second
\end{ReviewVerbatim}

\paragraph{Direct retained dependencies}
\begin{itemize}\raggedright
\item \hyperlink{governing-declaration-ccdbd7fdd43ef252}{Applied\allowbreak{}Modeling\allowbreak{}Lib.\allowbreak{}Game\allowbreak{}Theory.\allowbreak{}Preference\allowbreak{}Game.\allowbreak{}centered\allowbreak{}Preference\allowbreak{}Game\allowbreak{}Payoff}
\item \hyperlink{library-prerequisite-7e2db20cf087d9ba}{Applied\allowbreak{}Modeling\allowbreak{}Lib.\allowbreak{}Learning.\allowbreak{}Human\allowbreak{}Feedback.\allowbreak{}Pairwise\allowbreak{}Preference}
\end{itemize}
\hypertarget{governing-declaration-578b1c4cb8565a0d}{}
\subsection*{\small\ttfamily\raggedright Applied\allowbreak{}Modeling\allowbreak{}Lib.\allowbreak{}Game\allowbreak{}Theory.\allowbreak{}Preference\allowbreak{}Game.\allowbreak{}Is\allowbreak{}Context\allowbreak{}Free\allowbreak{}Extended\allowbreak{}Regularized\allowbreak{}Preference\allowbreak{}Game\allowbreak{}Equilibrium}
\textbf{Declaration location:} reusable library
\paragraph{Lean declaration body}
\begin{ReviewVerbatim}
type:
{Response : Type u_1} →
  [Fintype Response] →
    [DecidableEq Response] →
      AppliedModelingLib.Learning.HumanFeedback.PairwisePreference PUnit.{u_2 + 1} Response →
        PMF Response → PMF Response → WithTop ℝ → Prop

value:
fun {Response : Type u_1} [Fintype Response] [DecidableEq Response]
    (preference : AppliedModelingLib.Learning.HumanFeedback.PairwisePreference PUnit.{u_2 + 1} Response)
    (reference policy : PMF Response) (klWeight : WithTop ℝ) =>
  (∃ (finiteWeight : ℝ),
      klWeight = ↑finiteWeight ∧
        0 ≤ finiteWeight ∧
          ∀ (opponent : PMF Response),
            0 ≤
              AppliedModelingLib.GameTheory.PreferenceGame.contextFreeCenteredPreferencePayoff preference policy
                  opponent +
                finiteWeight *
                  (AppliedModelingLib.finiteKLDivergence opponent reference -
                    AppliedModelingLib.finiteKLDivergence policy reference)) ∨
    klWeight = ⊤ ∧ policy = reference
\end{ReviewVerbatim}

\paragraph{Direct retained dependencies}
\begin{itemize}\raggedright
\item \hyperlink{governing-declaration-44fb06319bbd7683}{Applied\allowbreak{}Modeling\allowbreak{}Lib.\allowbreak{}Game\allowbreak{}Theory.\allowbreak{}Preference\allowbreak{}Game.\allowbreak{}context\allowbreak{}Free\allowbreak{}Centered\allowbreak{}Preference\allowbreak{}Payoff}
\item \hyperlink{library-prerequisite-7e2db20cf087d9ba}{Applied\allowbreak{}Modeling\allowbreak{}Lib.\allowbreak{}Learning.\allowbreak{}Human\allowbreak{}Feedback.\allowbreak{}Pairwise\allowbreak{}Preference}
\item \hyperlink{library-prerequisite-d3573881ae45d4c6}{Applied\allowbreak{}Modeling\allowbreak{}Lib.\allowbreak{}finite\allowbreak{}KL\allowbreak{}Divergence}
\end{itemize}
\hypertarget{governing-declaration-d8827c2845a356ad}{}
\subsection*{\small\ttfamily\raggedright Applied\allowbreak{}Modeling\allowbreak{}Lib.\allowbreak{}Game\allowbreak{}Theory.\allowbreak{}Preference\allowbreak{}Game.\allowbreak{}Is\allowbreak{}Finite\allowbreak{}KL\allowbreak{}Preference\allowbreak{}Maximin}
\textbf{Declaration location:} reusable library
\paragraph{Lean declaration body}
\begin{ReviewVerbatim}
type:
{Response : Type u_1} →
  [Fintype Response] →
    [DecidableEq Response] →
      AppliedModelingLib.Learning.HumanFeedback.PairwisePreference PUnit.{u_2 + 1} Response →
        PMF Response → ℝ → PMF Response → Prop

value:
fun {Response : Type u_1} [Fintype Response] [DecidableEq Response]
    (preference : AppliedModelingLib.Learning.HumanFeedback.PairwisePreference PUnit.{u_2 + 1} Response)
    (reference : PMF Response) (klBudget : ℝ) (policy : PMF Response) =>
  AppliedModelingLib.finiteKLDivergence policy reference ≤ klBudget ∧
    ∃ (worst : PMF Response),
      AppliedModelingLib.finiteKLDivergence worst reference ≤ klBudget ∧
        (∀ (opponent : PMF Response),
            AppliedModelingLib.finiteKLDivergence opponent reference ≤ klBudget →
              AppliedModelingLib.GameTheory.PreferenceGame.contextFreeCenteredPreferencePayoff preference policy worst ≤
                AppliedModelingLib.GameTheory.PreferenceGame.contextFreeCenteredPreferencePayoff preference policy
                  opponent) ∧
          ∀ (candidate : PMF Response),
            AppliedModelingLib.finiteKLDivergence candidate reference ≤ klBudget →
              AppliedModelingLib.GameTheory.PreferenceGame.contextFreeCenteredPreferencePayoff preference candidate
                  worst ≤
                AppliedModelingLib.GameTheory.PreferenceGame.contextFreeCenteredPreferencePayoff preference policy worst
\end{ReviewVerbatim}

\paragraph{Direct retained dependencies}
\begin{itemize}\raggedright
\item \hyperlink{governing-declaration-44fb06319bbd7683}{Applied\allowbreak{}Modeling\allowbreak{}Lib.\allowbreak{}Game\allowbreak{}Theory.\allowbreak{}Preference\allowbreak{}Game.\allowbreak{}context\allowbreak{}Free\allowbreak{}Centered\allowbreak{}Preference\allowbreak{}Payoff}
\item \hyperlink{library-prerequisite-7e2db20cf087d9ba}{Applied\allowbreak{}Modeling\allowbreak{}Lib.\allowbreak{}Learning.\allowbreak{}Human\allowbreak{}Feedback.\allowbreak{}Pairwise\allowbreak{}Preference}
\item \hyperlink{library-prerequisite-d3573881ae45d4c6}{Applied\allowbreak{}Modeling\allowbreak{}Lib.\allowbreak{}finite\allowbreak{}KL\allowbreak{}Divergence}
\end{itemize}
\hypertarget{governing-declaration-2b862fbe69d5aa04}{}
\subsection*{\small\ttfamily\raggedright Applied\allowbreak{}Modeling\allowbreak{}Lib.\allowbreak{}Alignment.\allowbreak{}Welfare.\allowbreak{}Is\allowbreak{}Constrained\allowbreak{}Population\allowbreak{}Bradley\allowbreak{}Terry\allowbreak{}Maximin}
\textbf{Declaration location:} reusable library
\paragraph{Lean declaration body}
\begin{ReviewVerbatim}
type:
{User : Type u_1} →
  {Alternative : Type u_2} →
    [Fintype User] →
      [DecidableEq User] →
        [Fintype Alternative] →
          [DecidableEq Alternative] →
            PMF User →
              AppliedModelingLib.Alignment.Welfare.FiniteUtilityProfile User Alternative →
                ℝ → PMF Alternative → ℝ → PMF Alternative → Prop

value:
fun {User : Type u_1} {Alternative : Type u_2} [Fintype User] [DecidableEq User] [Fintype Alternative]
    [DecidableEq Alternative] (population : PMF User)
    (utility : AppliedModelingLib.Alignment.Welfare.FiniteUtilityProfile User Alternative) (btScale : ℝ)
    (reference : PMF Alternative) (klBudget : ℝ) (policy : PMF Alternative) =>
  AppliedModelingLib.GameTheory.PreferenceGame.IsFiniteKLPreferenceMaximin
    (have this : AppliedModelingLib.Learning.HumanFeedback.PairwisePreference PUnit.{1} Alternative :=
      AppliedModelingLib.Alignment.Welfare.populationBradleyTerryPreference population utility btScale;
    this)
    reference klBudget policy
\end{ReviewVerbatim}

\paragraph{Direct retained dependencies}
\begin{itemize}\raggedright
\item \hyperlink{library-prerequisite-56c0734f8877ddbe}{Applied\allowbreak{}Modeling\allowbreak{}Lib.\allowbreak{}Alignment.\allowbreak{}Welfare.\allowbreak{}Finite\allowbreak{}Utility\allowbreak{}Profile}
\item \hyperlink{library-prerequisite-cf06d15d7f8b7c31}{Applied\allowbreak{}Modeling\allowbreak{}Lib.\allowbreak{}Alignment.\allowbreak{}Welfare.\allowbreak{}population\allowbreak{}Bradley\allowbreak{}Terry\allowbreak{}Preference}
\item \hyperlink{governing-declaration-d8827c2845a356ad}{Applied\allowbreak{}Modeling\allowbreak{}Lib.\allowbreak{}Game\allowbreak{}Theory.\allowbreak{}Preference\allowbreak{}Game.\allowbreak{}Is\allowbreak{}Finite\allowbreak{}KL\allowbreak{}Preference\allowbreak{}Maximin}
\item \hyperlink{library-prerequisite-7e2db20cf087d9ba}{Applied\allowbreak{}Modeling\allowbreak{}Lib.\allowbreak{}Learning.\allowbreak{}Human\allowbreak{}Feedback.\allowbreak{}Pairwise\allowbreak{}Preference}
\end{itemize}
\hypertarget{governing-declaration-957ca8dbdbec4af9}{}
\subsection*{\small\ttfamily\raggedright Applied\allowbreak{}Modeling\allowbreak{}Lib.\allowbreak{}Game\allowbreak{}Theory.\allowbreak{}Preference\allowbreak{}Game.\allowbreak{}regularized\allowbreak{}Preference\allowbreak{}Game\allowbreak{}Payoff}
\textbf{Declaration location:} reusable library
\paragraph{Lean declaration body}
\begin{ReviewVerbatim}
type:
{Context : Type u_1} →
  {Response : Type u_2} →
    [Fintype Context] →
      [DecidableEq Context] →
        [Fintype Response] →
          [DecidableEq Response] →
            PMF Context →
              AppliedModelingLib.Learning.HumanFeedback.PairwisePreference Context Response →
                AppliedModelingLib.Learning.HumanFeedback.FinitePolicy Context Response →
                  AppliedModelingLib.Learning.HumanFeedback.FinitePolicy Context Response →
                    AppliedModelingLib.Learning.HumanFeedback.FinitePolicy Context Response → ℝ → ℝ

value:
fun {Context : Type u_1} {Response : Type u_2} [Fintype Context] [DecidableEq Context] [Fintype Response]
    [DecidableEq Response] (contextLaw : PMF Context)
    (preference : AppliedModelingLib.Learning.HumanFeedback.PairwisePreference Context Response)
    (reference first second : AppliedModelingLib.Learning.HumanFeedback.FinitePolicy Context Response)
    (klRegularization : ℝ) =>
  AppliedModelingLib.GameTheory.PreferenceGame.preferenceGamePayoff contextLaw preference first second -
      klRegularization *
        AppliedModelingLib.Learning.HumanFeedback.contextAveragedPolicyKLDivergence contextLaw first reference +
    klRegularization *
      AppliedModelingLib.Learning.HumanFeedback.contextAveragedPolicyKLDivergence contextLaw second reference
\end{ReviewVerbatim}

\paragraph{Direct retained dependencies}
\begin{itemize}\raggedright
\item \hyperlink{governing-declaration-8e5c390864970969}{Applied\allowbreak{}Modeling\allowbreak{}Lib.\allowbreak{}Game\allowbreak{}Theory.\allowbreak{}Preference\allowbreak{}Game.\allowbreak{}preference\allowbreak{}Game\allowbreak{}Payoff}
\item \hyperlink{governing-declaration-577316c60f45ef11}{Applied\allowbreak{}Modeling\allowbreak{}Lib.\allowbreak{}Learning.\allowbreak{}Human\allowbreak{}Feedback.\allowbreak{}Finite\allowbreak{}Policy}
\item \hyperlink{library-prerequisite-7e2db20cf087d9ba}{Applied\allowbreak{}Modeling\allowbreak{}Lib.\allowbreak{}Learning.\allowbreak{}Human\allowbreak{}Feedback.\allowbreak{}Pairwise\allowbreak{}Preference}
\item \hyperlink{governing-declaration-c528a6568d75fffd}{Applied\allowbreak{}Modeling\allowbreak{}Lib.\allowbreak{}Learning.\allowbreak{}Human\allowbreak{}Feedback.\allowbreak{}context\allowbreak{}Averaged\allowbreak{}Policy\allowbreak{}KL\allowbreak{}Divergence}
\end{itemize}
\hypertarget{governing-declaration-49e6e1fe2a2cc849}{}
\subsection*{\small\ttfamily\raggedright Applied\allowbreak{}Modeling\allowbreak{}Lib.\allowbreak{}Game\allowbreak{}Theory.\allowbreak{}Preference\allowbreak{}Game.\allowbreak{}Is\allowbreak{}Regularized\allowbreak{}Preference\allowbreak{}Game\allowbreak{}Equilibrium}
\textbf{Declaration location:} reusable library
\paragraph{Lean declaration body}
\begin{ReviewVerbatim}
type:
{Context : Type u_1} →
  {Response : Type u_2} →
    [Fintype Context] →
      [DecidableEq Context] →
        [Fintype Response] →
          [DecidableEq Response] →
            PMF Context →
              AppliedModelingLib.Learning.HumanFeedback.PairwisePreference Context Response →
                AppliedModelingLib.Learning.HumanFeedback.FinitePolicy Context Response →
                  AppliedModelingLib.Learning.HumanFeedback.FinitePolicy Context Response → ℝ → Prop

value:
fun {Context : Type u_1} {Response : Type u_2} [Fintype Context] [DecidableEq Context] [Fintype Response]
    [DecidableEq Response] (contextLaw : PMF Context)
    (preference : AppliedModelingLib.Learning.HumanFeedback.PairwisePreference Context Response)
    (reference policy : AppliedModelingLib.Learning.HumanFeedback.FinitePolicy Context Response)
    (klRegularization : ℝ) =>
  ∀ (alternative : AppliedModelingLib.Learning.HumanFeedback.FinitePolicy Context Response),
    1 / 2 ≤
      AppliedModelingLib.GameTheory.PreferenceGame.regularizedPreferenceGamePayoff contextLaw preference reference
        policy alternative klRegularization
\end{ReviewVerbatim}

\paragraph{Direct retained dependencies}
\begin{itemize}\raggedright
\item \hyperlink{governing-declaration-957ca8dbdbec4af9}{Applied\allowbreak{}Modeling\allowbreak{}Lib.\allowbreak{}Game\allowbreak{}Theory.\allowbreak{}Preference\allowbreak{}Game.\allowbreak{}regularized\allowbreak{}Preference\allowbreak{}Game\allowbreak{}Payoff}
\item \hyperlink{governing-declaration-577316c60f45ef11}{Applied\allowbreak{}Modeling\allowbreak{}Lib.\allowbreak{}Learning.\allowbreak{}Human\allowbreak{}Feedback.\allowbreak{}Finite\allowbreak{}Policy}
\item \hyperlink{library-prerequisite-7e2db20cf087d9ba}{Applied\allowbreak{}Modeling\allowbreak{}Lib.\allowbreak{}Learning.\allowbreak{}Human\allowbreak{}Feedback.\allowbreak{}Pairwise\allowbreak{}Preference}
\end{itemize}
\hypertarget{governing-declaration-9c38d53b9d664221}{}
\subsection*{\small\ttfamily\raggedright Applied\allowbreak{}Modeling\allowbreak{}Lib.\allowbreak{}pmf\allowbreak{}Prob}
\textbf{Declaration location:} reusable library
\paragraph{Lean declaration body}
\begin{ReviewVerbatim}
type:
{α : Type u_1} → [Fintype α] → PMF α → (p : α → Prop) → [DecidablePred p] → ℝ

value:
fun {α : Type u_1} [Fintype α] (μ : PMF α) (p : α → Prop) [DecidablePred p] =>
  AppliedModelingLib.pmfExp μ fun (a : α) => if p a then 1 else 0
\end{ReviewVerbatim}

\paragraph{Direct retained dependencies}
\begin{itemize}\raggedright
\item \hyperlink{governing-declaration-8e4b4389a4f6454b}{Applied\allowbreak{}Modeling\allowbreak{}Lib.\allowbreak{}pmf\allowbreak{}Exp}
\end{itemize}
\hypertarget{governing-declaration-5ac12d0a52b69068}{}
\subsection*{\small\ttfamily\raggedright Applied\allowbreak{}Modeling\allowbreak{}Lib.\allowbreak{}pmf\allowbreak{}Product}
\textbf{Declaration location:} reusable library
\paragraph{Lean declaration body}
\begin{ReviewVerbatim}
type:
(ι : Type u_1) → (α : Type u_2) → [Fintype ι] → [DecidableEq ι] → [Fintype α] → PMF α → PMF (ι → α)

value:
fun (ι : Type u_1) (α : Type u_2) [Fintype ι] [DecidableEq ι] [Fintype α] (μ : PMF α) =>
  PMF.ofFintype (fun (f : ι → α) => ∏ i : ι, μ (f i)) (AppliedModelingLib.pmfProduct._proof_1 ι α μ)
\end{ReviewVerbatim}


\hypertarget{governing-declaration-48f10f2715c04df3}{}
\subsection*{\small\ttfamily\raggedright Applied\allowbreak{}Modeling\allowbreak{}Lib.\allowbreak{}uniform\allowbreak{}PMF}
\textbf{Declaration location:} reusable library
\paragraph{Lean declaration body}
\begin{ReviewVerbatim}
type:
(α : Type u_1) → [Fintype α] → [Nonempty α] → PMF α

value:
fun (α : Type u_1) [Fintype α] [Nonempty α] =>
  PMF.ofFintype (fun (x : α) => (↑(Fintype.card α))⁻¹) (AppliedModelingLib.uniformPMF._proof_1 α)
\end{ReviewVerbatim}


\hypertarget{governing-declaration-d99b0ff87003286e}{}
\subsection*{\small\ttfamily\raggedright Golz\allowbreak{}Haghtalab\allowbreak{}Yang2025\allowbreak{}Distortion.\allowbreak{}lemma15\allowbreak{}Drop\allowbreak{}Add}
\textbf{Declaration location:} paper-local
\paragraph{Lean declaration body}
\begin{ReviewVerbatim}
type:
(Fin 3 → ℝ) → Fin 3 → ℝ → Fin 3 → ℝ

value:
fun (utility : Fin 3 → ℝ) (peak : Fin 3) (increment : ℝ) (a : Fin 3) => if a = peak then 0 else utility a + increment
\end{ReviewVerbatim}


\hypertarget{governing-declaration-019c4b857af92023}{}
\subsection*{\small\ttfamily\raggedright Golz\allowbreak{}Haghtalab\allowbreak{}Yang2025\allowbreak{}Distortion.\allowbreak{}lemma15\allowbreak{}Increment}
\textbf{Declaration location:} paper-local
\paragraph{Lean declaration body}
\begin{ReviewVerbatim}
type:
ℝ → ℝ

value:
fun (delta : ℝ) => Real.log ((Real.exp (delta / 2) + 1) ^ 3 / (2 * (Real.exp delta + 3)))
\end{ReviewVerbatim}


\hypertarget{governing-declaration-a3540c93d773dd6c}{}
\subsection*{\small\ttfamily\raggedright Golz\allowbreak{}Haghtalab\allowbreak{}Yang2025\allowbreak{}Distortion.\allowbreak{}lemma15\allowbreak{}Peak}
\textbf{Declaration location:} paper-local
\paragraph{Lean declaration body}
\begin{ReviewVerbatim}
type:
(Fin 3 → ℝ) → Fin 3

value:
fun (utility : Fin 3 → ℝ) => Classical.choose (GolzHaghtalabYang2025Distortion.lemma15Peak._proof_1 utility)
\end{ReviewVerbatim}


\hypertarget{governing-declaration-877274851b50572c}{}
\subsection*{\small\ttfamily\raggedright Golz\allowbreak{}Haghtalab\allowbreak{}Yang2025\allowbreak{}Distortion.\allowbreak{}lemma15\allowbreak{}Population}
\textbf{Declaration location:} paper-local
\paragraph{Lean declaration body}
\begin{ReviewVerbatim}
type:
PMF (Fin 3)

value:
AppliedModelingLib.uniformPMF (Fin 3)
\end{ReviewVerbatim}

\paragraph{Direct retained dependencies}
\begin{itemize}\raggedright
\item \hyperlink{governing-declaration-48f10f2715c04df3}{Applied\allowbreak{}Modeling\allowbreak{}Lib.\allowbreak{}uniform\allowbreak{}PMF}
\end{itemize}
\hypertarget{governing-declaration-218b6dfa710b7a0d}{}
\subsection*{\small\ttfamily\raggedright Golz\allowbreak{}Haghtalab\allowbreak{}Yang2025\allowbreak{}Distortion.\allowbreak{}lemma15\allowbreak{}Step}
\textbf{Declaration location:} paper-local
\paragraph{Lean declaration body}
\begin{ReviewVerbatim}
type:
ℝ → (Fin 3 → ℝ) → Fin 3 → ℝ

value:
fun (beta : ℝ) (utility : Fin 3 → ℝ) (a : Fin 3) =>
  (have peak : Fin 3 := GolzHaghtalabYang2025Distortion.lemma15Peak utility;
    GolzHaghtalabYang2025Distortion.lemma15DropAdd utility peak
      (GolzHaghtalabYang2025Distortion.lemma15Increment (beta * utility peak) / beta))
    a
\end{ReviewVerbatim}

\paragraph{Direct retained dependencies}
\begin{itemize}\raggedright
\item \hyperlink{governing-declaration-d99b0ff87003286e}{Golz\allowbreak{}Haghtalab\allowbreak{}Yang2025\allowbreak{}Distortion.\allowbreak{}lemma15\allowbreak{}Drop\allowbreak{}Add}
\item \hyperlink{governing-declaration-019c4b857af92023}{Golz\allowbreak{}Haghtalab\allowbreak{}Yang2025\allowbreak{}Distortion.\allowbreak{}lemma15\allowbreak{}Increment}
\item \hyperlink{governing-declaration-a3540c93d773dd6c}{Golz\allowbreak{}Haghtalab\allowbreak{}Yang2025\allowbreak{}Distortion.\allowbreak{}lemma15\allowbreak{}Peak}
\end{itemize}
\hypertarget{governing-declaration-f1a69dba69c50eae}{}
\subsection*{\small\ttfamily\raggedright Golz\allowbreak{}Haghtalab\allowbreak{}Yang2025\allowbreak{}Distortion.\allowbreak{}lemma15\allowbreak{}Utility}
\textbf{Declaration location:} paper-local
\paragraph{Lean declaration body}
\begin{ReviewVerbatim}
type:
ℝ → ℕ → Fin 3 → ℝ

value:
fun (beta : ℝ) (a : ℕ) (a_1 : Fin 3) =>
  Nat.brecOn (motive := fun (x : ℕ) => Fin 3 → ℝ) a (GolzHaghtalabYang2025Distortion.lemma15Utility._f beta) a_1
\end{ReviewVerbatim}

\paragraph{Direct retained dependencies}
\begin{itemize}\raggedright
\item \hyperlink{governing-declaration-218b6dfa710b7a0d}{Golz\allowbreak{}Haghtalab\allowbreak{}Yang2025\allowbreak{}Distortion.\allowbreak{}lemma15\allowbreak{}Step}
\end{itemize}
\hypertarget{governing-declaration-0f2c5b70a888d5b3}{}
\subsection*{\small\ttfamily\raggedright Golz\allowbreak{}Haghtalab\allowbreak{}Yang2025\allowbreak{}Distortion.\allowbreak{}source\allowbreak{}Dpo\allowbreak{}Binary\allowbreak{}Report\allowbreak{}Loss}
\textbf{Declaration location:} paper-local
\paragraph{Lean declaration body}
\begin{ReviewVerbatim}
type:
{Alternative : Type u_1} →
  [Fintype Alternative] →
    [DecidableEq Alternative] →
      {horizon : ℕ} →
        (Fin horizon → AppliedModelingLib.Learning.HumanFeedback.BinaryPairwiseReport Alternative) →
          PMF Alternative → PMF Alternative → ℝ → ℝ

value:
fun {Alternative : Type u_1} [Fintype Alternative] [DecidableEq Alternative] {horizon : ℕ}
    (sample : Fin horizon → AppliedModelingLib.Learning.HumanFeedback.BinaryPairwiseReport Alternative)
    (reference policy : PMF Alternative) (klWeight : ℝ) =>
  -(AppliedModelingLib.Learning.HumanFeedback.PairwiseCountDataset.ofBinaryReports sample).pairwiseLogLikelihood
      Real.sigmoid fun (alternative : Alternative) =>
      AppliedModelingLib.Learning.HumanFeedback.dpoImplicitReward
        (AppliedModelingLib.Alignment.Welfare.contextFreeAlternativePolicy policy)
        (AppliedModelingLib.Alignment.Welfare.contextFreeAlternativePolicy reference) klWeight PUnit.unit alternative
\end{ReviewVerbatim}

\paragraph{Direct retained dependencies}
\begin{itemize}\raggedright
\item \hyperlink{governing-declaration-b7cf9e04e23161e5}{Applied\allowbreak{}Modeling\allowbreak{}Lib.\allowbreak{}Alignment.\allowbreak{}Welfare.\allowbreak{}context\allowbreak{}Free\allowbreak{}Alternative\allowbreak{}Policy}
\item \hyperlink{governing-declaration-bc3d6ca8e7e0ac9d}{Applied\allowbreak{}Modeling\allowbreak{}Lib.\allowbreak{}Learning.\allowbreak{}Human\allowbreak{}Feedback.\allowbreak{}Binary\allowbreak{}Pairwise\allowbreak{}Report}
\item \hyperlink{governing-declaration-449f87e91bf37ce0}{Applied\allowbreak{}Modeling\allowbreak{}Lib.\allowbreak{}Learning.\allowbreak{}Human\allowbreak{}Feedback.\allowbreak{}Pairwise\allowbreak{}Count\allowbreak{}Dataset.\allowbreak{}of\allowbreak{}Binary\allowbreak{}Reports}
\item \hyperlink{governing-declaration-10adb9e7c1ff9f4c}{Applied\allowbreak{}Modeling\allowbreak{}Lib.\allowbreak{}Learning.\allowbreak{}Human\allowbreak{}Feedback.\allowbreak{}Pairwise\allowbreak{}Count\allowbreak{}Dataset.\allowbreak{}pairwise\allowbreak{}Log\allowbreak{}Likelihood}
\item \hyperlink{governing-declaration-95571fb7643367ad}{Applied\allowbreak{}Modeling\allowbreak{}Lib.\allowbreak{}Learning.\allowbreak{}Human\allowbreak{}Feedback.\allowbreak{}dpo\allowbreak{}Implicit\allowbreak{}Reward}
\end{itemize}
\hypertarget{governing-declaration-4cdae0c3705e3d00}{}
\subsection*{\small\ttfamily\raggedright Golz\allowbreak{}Haghtalab\allowbreak{}Yang2025\allowbreak{}Distortion.\allowbreak{}Is\allowbreak{}Source\allowbreak{}Dpo\allowbreak{}Loss\allowbreak{}Minimizer}
\textbf{Declaration location:} paper-local
\paragraph{Lean declaration body}
\begin{ReviewVerbatim}
type:
{Alternative : Type u_1} →
  [Fintype Alternative] →
    [DecidableEq Alternative] →
      {horizon : ℕ} →
        (Fin horizon → AppliedModelingLib.Learning.HumanFeedback.BinaryPairwiseReport Alternative) →
          PMF Alternative → ℝ → PMF Alternative → Prop

value:
fun {Alternative : Type u_1} [Fintype Alternative] [DecidableEq Alternative] {horizon : ℕ}
    (sample : Fin horizon → AppliedModelingLib.Learning.HumanFeedback.BinaryPairwiseReport Alternative)
    (reference : PMF Alternative) (klWeight : ℝ) (policy : PMF Alternative) =>
  AppliedModelingLib.PMFFullSupport policy ∧
    ∀ (candidate : PMF Alternative),
      AppliedModelingLib.PMFFullSupport candidate →
        GolzHaghtalabYang2025Distortion.sourceDpoBinaryReportLoss sample reference policy klWeight ≤
          GolzHaghtalabYang2025Distortion.sourceDpoBinaryReportLoss sample reference candidate klWeight
\end{ReviewVerbatim}

\paragraph{Direct retained dependencies}
\begin{itemize}\raggedright
\item \hyperlink{governing-declaration-bc3d6ca8e7e0ac9d}{Applied\allowbreak{}Modeling\allowbreak{}Lib.\allowbreak{}Learning.\allowbreak{}Human\allowbreak{}Feedback.\allowbreak{}Binary\allowbreak{}Pairwise\allowbreak{}Report}
\item \hyperlink{library-prerequisite-268c959ee7d2d148}{Applied\allowbreak{}Modeling\allowbreak{}Lib.\allowbreak{}PMF\allowbreak{}Full\allowbreak{}Support}
\item \hyperlink{governing-declaration-0f2c5b70a888d5b3}{Golz\allowbreak{}Haghtalab\allowbreak{}Yang2025\allowbreak{}Distortion.\allowbreak{}source\allowbreak{}Dpo\allowbreak{}Binary\allowbreak{}Report\allowbreak{}Loss}
\end{itemize}
\hypertarget{governing-declaration-d547690922524bc2}{}
\subsection*{\small\ttfamily\raggedright Golz\allowbreak{}Haghtalab\allowbreak{}Yang2025\allowbreak{}Distortion.\allowbreak{}source\allowbreak{}Dpo\allowbreak{}Normalized\allowbreak{}MLE\allowbreak{}Reward}
\textbf{Declaration location:} paper-local
\paragraph{Lean declaration body}
\begin{ReviewVerbatim}
type:
{Alternative : Type u_1} →
  [Fintype Alternative] →
    [DecidableEq Alternative] → Alternative → PMF Alternative → PMF Alternative → ℝ → Alternative → ℝ

value:
fun {Alternative : Type u_1} [Fintype Alternative] [DecidableEq Alternative] (scoreReference : Alternative)
    (reference policy : PMF Alternative) (klWeight : ℝ) (a : Alternative) =>
  AppliedModelingLib.Learning.HumanFeedback.dpoImplicitReward
      (AppliedModelingLib.Alignment.Welfare.contextFreeAlternativePolicy policy)
      (AppliedModelingLib.Alignment.Welfare.contextFreeAlternativePolicy reference) klWeight PUnit.unit a -
    AppliedModelingLib.Learning.HumanFeedback.dpoImplicitReward
      (AppliedModelingLib.Alignment.Welfare.contextFreeAlternativePolicy policy)
      (AppliedModelingLib.Alignment.Welfare.contextFreeAlternativePolicy reference) klWeight PUnit.unit scoreReference
\end{ReviewVerbatim}

\paragraph{Direct retained dependencies}
\begin{itemize}\raggedright
\item \hyperlink{governing-declaration-b7cf9e04e23161e5}{Applied\allowbreak{}Modeling\allowbreak{}Lib.\allowbreak{}Alignment.\allowbreak{}Welfare.\allowbreak{}context\allowbreak{}Free\allowbreak{}Alternative\allowbreak{}Policy}
\item \hyperlink{governing-declaration-95571fb7643367ad}{Applied\allowbreak{}Modeling\allowbreak{}Lib.\allowbreak{}Learning.\allowbreak{}Human\allowbreak{}Feedback.\allowbreak{}dpo\allowbreak{}Implicit\allowbreak{}Reward}
\end{itemize}
\hypertarget{governing-declaration-7f6992add20edc0b}{}
\subsection*{\small\ttfamily\raggedright Golz\allowbreak{}Haghtalab\allowbreak{}Yang2025\allowbreak{}Distortion.\allowbreak{}source\allowbreak{}Dpo\allowbreak{}Policy}
\textbf{Declaration location:} paper-local
\paragraph{Lean declaration body}
\begin{ReviewVerbatim}
type:
{Alternative : Type u_1} →
  [Fintype Alternative] →
    [DecidableEq Alternative] →
      PMF Alternative →
        AppliedModelingLib.Learning.HumanFeedback.PairwiseCountDataset.ScoreVector Alternative → ℝ → PMF Alternative

value:
fun {Alternative : Type u_1} [Fintype Alternative] [DecidableEq Alternative] (reference : PMF Alternative)
    (reward : AppliedModelingLib.Learning.HumanFeedback.PairwiseCountDataset.ScoreVector Alternative) (klWeight : ℝ) =>
  AppliedModelingLib.Learning.HumanFeedback.dpoOptimalPolicy
    (AppliedModelingLib.Alignment.Welfare.contextFreeAlternativePolicy reference) (fun (x : PUnit.{1}) => reward)
    klWeight PUnit.unit
\end{ReviewVerbatim}

\paragraph{Direct retained dependencies}
\begin{itemize}\raggedright
\item \hyperlink{governing-declaration-b7cf9e04e23161e5}{Applied\allowbreak{}Modeling\allowbreak{}Lib.\allowbreak{}Alignment.\allowbreak{}Welfare.\allowbreak{}context\allowbreak{}Free\allowbreak{}Alternative\allowbreak{}Policy}
\item \hyperlink{governing-declaration-87eda978eb1b0c9d}{Applied\allowbreak{}Modeling\allowbreak{}Lib.\allowbreak{}Learning.\allowbreak{}Human\allowbreak{}Feedback.\allowbreak{}Pairwise\allowbreak{}Count\allowbreak{}Dataset.\allowbreak{}Score\allowbreak{}Vector}
\item \hyperlink{governing-declaration-eeb142baa1b200d1}{Applied\allowbreak{}Modeling\allowbreak{}Lib.\allowbreak{}Learning.\allowbreak{}Human\allowbreak{}Feedback.\allowbreak{}dpo\allowbreak{}Optimal\allowbreak{}Policy}
\end{itemize}
\hypertarget{governing-declaration-9dc701ec2b5767ff}{}
\subsection*{\small\ttfamily\raggedright Golz\allowbreak{}Haghtalab\allowbreak{}Yang2025\allowbreak{}Distortion.\allowbreak{}source\allowbreak{}Rlhf\allowbreak{}Binary\allowbreak{}Report\allowbreak{}Loss}
\textbf{Declaration location:} paper-local
\paragraph{Lean declaration body}
\begin{ReviewVerbatim}
type:
{Alternative : Type u_1} →
  [Fintype Alternative] →
    [DecidableEq Alternative] →
      {horizon : ℕ} →
        (Fin horizon → AppliedModelingLib.Learning.HumanFeedback.BinaryPairwiseReport Alternative) →
          AppliedModelingLib.Learning.HumanFeedback.PairwiseCountDataset.ScoreVector Alternative → ℝ

value:
fun {Alternative : Type u_1} [Fintype Alternative] [DecidableEq Alternative] {horizon : ℕ}
    (sample : Fin horizon → AppliedModelingLib.Learning.HumanFeedback.BinaryPairwiseReport Alternative)
    (reward : AppliedModelingLib.Learning.HumanFeedback.PairwiseCountDataset.ScoreVector Alternative) =>
  -(AppliedModelingLib.Learning.HumanFeedback.PairwiseCountDataset.ofBinaryReports sample).pairwiseLogLikelihood
      Real.sigmoid reward
\end{ReviewVerbatim}

\paragraph{Direct retained dependencies}
\begin{itemize}\raggedright
\item \hyperlink{governing-declaration-bc3d6ca8e7e0ac9d}{Applied\allowbreak{}Modeling\allowbreak{}Lib.\allowbreak{}Learning.\allowbreak{}Human\allowbreak{}Feedback.\allowbreak{}Binary\allowbreak{}Pairwise\allowbreak{}Report}
\item \hyperlink{governing-declaration-87eda978eb1b0c9d}{Applied\allowbreak{}Modeling\allowbreak{}Lib.\allowbreak{}Learning.\allowbreak{}Human\allowbreak{}Feedback.\allowbreak{}Pairwise\allowbreak{}Count\allowbreak{}Dataset.\allowbreak{}Score\allowbreak{}Vector}
\item \hyperlink{governing-declaration-449f87e91bf37ce0}{Applied\allowbreak{}Modeling\allowbreak{}Lib.\allowbreak{}Learning.\allowbreak{}Human\allowbreak{}Feedback.\allowbreak{}Pairwise\allowbreak{}Count\allowbreak{}Dataset.\allowbreak{}of\allowbreak{}Binary\allowbreak{}Reports}
\item \hyperlink{governing-declaration-10adb9e7c1ff9f4c}{Applied\allowbreak{}Modeling\allowbreak{}Lib.\allowbreak{}Learning.\allowbreak{}Human\allowbreak{}Feedback.\allowbreak{}Pairwise\allowbreak{}Count\allowbreak{}Dataset.\allowbreak{}pairwise\allowbreak{}Log\allowbreak{}Likelihood}
\end{itemize}
\hypertarget{governing-declaration-55b65ae526406849}{}
\subsection*{\small\ttfamily\raggedright Golz\allowbreak{}Haghtalab\allowbreak{}Yang2025\allowbreak{}Distortion.\allowbreak{}source\allowbreak{}Rlhf\allowbreak{}Policy}
\textbf{Declaration location:} paper-local
\paragraph{Lean declaration body}
\begin{ReviewVerbatim}
type:
{Alternative : Type u_1} →
  [Fintype Alternative] →
    [DecidableEq Alternative] →
      PMF Alternative →
        AppliedModelingLib.Learning.HumanFeedback.PairwiseCountDataset.ScoreVector Alternative → ℝ → PMF Alternative

value:
fun {Alternative : Type u_1} [Fintype Alternative] [DecidableEq Alternative] (reference : PMF Alternative)
    (reward : AppliedModelingLib.Learning.HumanFeedback.PairwiseCountDataset.ScoreVector Alternative) (klWeight : ℝ) =>
  AppliedModelingLib.Learning.HumanFeedback.contextExponentialTiltPolicy
    (AppliedModelingLib.Alignment.Welfare.contextFreeAlternativePolicy reference) (fun (x : PUnit.{1}) => reward)
    klWeight⁻¹ PUnit.unit
\end{ReviewVerbatim}

\paragraph{Direct retained dependencies}
\begin{itemize}\raggedright
\item \hyperlink{governing-declaration-b7cf9e04e23161e5}{Applied\allowbreak{}Modeling\allowbreak{}Lib.\allowbreak{}Alignment.\allowbreak{}Welfare.\allowbreak{}context\allowbreak{}Free\allowbreak{}Alternative\allowbreak{}Policy}
\item \hyperlink{governing-declaration-87eda978eb1b0c9d}{Applied\allowbreak{}Modeling\allowbreak{}Lib.\allowbreak{}Learning.\allowbreak{}Human\allowbreak{}Feedback.\allowbreak{}Pairwise\allowbreak{}Count\allowbreak{}Dataset.\allowbreak{}Score\allowbreak{}Vector}
\item \hyperlink{governing-declaration-a868460489b0ccee}{Applied\allowbreak{}Modeling\allowbreak{}Lib.\allowbreak{}Learning.\allowbreak{}Human\allowbreak{}Feedback.\allowbreak{}context\allowbreak{}Exponential\allowbreak{}Tilt\allowbreak{}Policy}
\end{itemize}
\hypertarget{governing-declaration-aaf6dc2063f9888f}{}
\subsection*{\small\ttfamily\raggedright Golz\allowbreak{}Haghtalab\allowbreak{}Yang2025\allowbreak{}Distortion.\allowbreak{}theorem10\allowbreak{}Source\allowbreak{}Rate}
\textbf{Declaration location:} paper-local
\paragraph{Lean declaration body}
\begin{ReviewVerbatim}
type:
ℕ → ℕ → ℝ → ℝ → ℝ

value:
fun (users comparisonsPerUser : ℕ) (minimumMass failureBudget : ℝ) =>
  8 * √(Real.log (4 / failureBudget) / (↑users * min 1 (↑comparisonsPerUser * minimumMass ^ 2))) +
    16 * Real.log (4 / failureBudget) / (↑users * minimumMass ^ 2)
\end{ReviewVerbatim}


\hypertarget{governing-declaration-e63568da8eac1b3c}{}
\subsection*{\small\ttfamily\raggedright Golz\allowbreak{}Haghtalab\allowbreak{}Yang2025\allowbreak{}Distortion.\allowbreak{}theorem10\allowbreak{}Unordered\allowbreak{}Pair\allowbreak{}Indicator}
\textbf{Declaration location:} paper-local
\paragraph{Lean declaration body}
\begin{ReviewVerbatim}
type:
{Alternative : Type u_1} → [DecidableEq Alternative] → Alternative → Alternative → Alternative × Alternative → ℝ

value:
fun {Alternative : Type u_1} [DecidableEq Alternative] (first second : Alternative) (a : Alternative × Alternative) =>
  match a with
  | (left, right) => if left = first ∧ right = second ∨ left = second ∧ right = first then 1 else 0
\end{ReviewVerbatim}


\hypertarget{governing-declaration-2fe7d1ffcbbbb71a}{}
\subsection*{\small\ttfamily\raggedright Golz\allowbreak{}Haghtalab\allowbreak{}Yang2025\allowbreak{}Distortion.\allowbreak{}theorem10\allowbreak{}User\allowbreak{}Unordered\allowbreak{}Pair\allowbreak{}Count}
\textbf{Declaration location:} paper-local
\paragraph{Lean declaration body}
\begin{ReviewVerbatim}
type:
{Alternative : Type u_1} →
  [DecidableEq Alternative] →
    {comparisonsPerUser : ℕ} → Alternative → Alternative → (Fin comparisonsPerUser → Alternative × Alternative) → ℝ

value:
fun {Alternative : Type u_1} [DecidableEq Alternative] {comparisonsPerUser : ℕ} (first second : Alternative)
    (labels : Fin comparisonsPerUser → Alternative × Alternative) =>
  ∑ position : Fin comparisonsPerUser,
    GolzHaghtalabYang2025Distortion.theorem10UnorderedPairIndicator first second (labels position)
\end{ReviewVerbatim}

\paragraph{Direct retained dependencies}
\begin{itemize}\raggedright
\item \hyperlink{governing-declaration-e63568da8eac1b3c}{Golz\allowbreak{}Haghtalab\allowbreak{}Yang2025\allowbreak{}Distortion.\allowbreak{}theorem10\allowbreak{}Unordered\allowbreak{}Pair\allowbreak{}Indicator}
\end{itemize}
\hypertarget{governing-declaration-96705d9adbc097bb}{}
\subsection*{\small\ttfamily\raggedright Golz\allowbreak{}Haghtalab\allowbreak{}Yang2025\allowbreak{}Distortion.\allowbreak{}theorem12\allowbreak{}Bernoulli\allowbreak{}Report}
\textbf{Declaration location:} paper-local
\paragraph{Lean declaration body}
\begin{ReviewVerbatim}
type:
(probability : ℝ) → 0 ≤ probability → probability ≤ 1 → PMF Bool

value:
fun (probability : ℝ) (hnonneg : 0 ≤ probability) (hle_one : probability ≤ 1) =>
  PMF.bernoulli ⟨probability, hnonneg⟩
    (GolzHaghtalabYang2025Distortion.theorem12BernoulliReport._proof_1 probability hnonneg hle_one)
\end{ReviewVerbatim}


\hypertarget{governing-declaration-ec899679d7cb6165}{}
\subsection*{\small\ttfamily\raggedright Golz\allowbreak{}Haghtalab\allowbreak{}Yang2025\allowbreak{}Distortion.\allowbreak{}theorem12\allowbreak{}Clone\allowbreak{}Collapse}
\textbf{Declaration location:} paper-local
\paragraph{Lean declaration body}
\begin{ReviewVerbatim}
type:
{Alternative : Type v} → Alternative → (clones : ℕ) → Alternative ⊕ Fin clones → Alternative

value:
fun {Alternative : Type v} (split : Alternative) (clones : ℕ) (a : Alternative ⊕ Fin clones) =>
  match a with
  | Sum.inl alternative => alternative
  | Sum.inr val => split
\end{ReviewVerbatim}


\hypertarget{governing-declaration-88811107533a14f3}{}
\subsection*{\small\ttfamily\raggedright Golz\allowbreak{}Haghtalab\allowbreak{}Yang2025\allowbreak{}Distortion.\allowbreak{}theorem12\allowbreak{}Empirical\allowbreak{}Borda\allowbreak{}Incidences}
\textbf{Declaration location:} paper-local
\paragraph{Lean declaration body}
\begin{ReviewVerbatim}
type:
{Report : Type u_1} →
  {Alternative : Type u_2} →
    [Fintype Alternative] → (Alternative → Report → ℝ) → Alternative → {horizon : ℕ} → (Fin horizon → Report) → ℝ

value:
fun {Report : Type u_1} {Alternative : Type u_2} [Fintype Alternative] (incidences : Alternative → Report → ℝ)
    (alternative : Alternative) {horizon : ℕ} (sample : Fin horizon → Report) =>
  AppliedModelingLib.Probability.finiteIidScoreSum (incidences alternative) sample
\end{ReviewVerbatim}

\paragraph{Direct retained dependencies}
\begin{itemize}\raggedright
\item \hyperlink{governing-declaration-65283b6bc1ebd8c0}{Applied\allowbreak{}Modeling\allowbreak{}Lib.\allowbreak{}Probability.\allowbreak{}finite\allowbreak{}Iid\allowbreak{}Score\allowbreak{}Sum}
\end{itemize}
\hypertarget{governing-declaration-76ef72f9364a0835}{}
\subsection*{\small\ttfamily\raggedright Golz\allowbreak{}Haghtalab\allowbreak{}Yang2025\allowbreak{}Distortion.\allowbreak{}theorem12\allowbreak{}Empirical\allowbreak{}Borda\allowbreak{}Wins}
\textbf{Declaration location:} paper-local
\paragraph{Lean declaration body}
\begin{ReviewVerbatim}
type:
{Report : Type u_1} →
  {Alternative : Type u_2} →
    [Fintype Alternative] → (Alternative → Report → ℝ) → Alternative → {horizon : ℕ} → (Fin horizon → Report) → ℝ

value:
fun {Report : Type u_1} {Alternative : Type u_2} [Fintype Alternative] (wins : Alternative → Report → ℝ)
    (alternative : Alternative) {horizon : ℕ} (sample : Fin horizon → Report) =>
  AppliedModelingLib.Probability.finiteIidScoreSum (wins alternative) sample
\end{ReviewVerbatim}

\paragraph{Direct retained dependencies}
\begin{itemize}\raggedright
\item \hyperlink{governing-declaration-65283b6bc1ebd8c0}{Applied\allowbreak{}Modeling\allowbreak{}Lib.\allowbreak{}Probability.\allowbreak{}finite\allowbreak{}Iid\allowbreak{}Score\allowbreak{}Sum}
\end{itemize}
\hypertarget{governing-declaration-b2a8d113fc3130c2}{}
\subsection*{\small\ttfamily\raggedright Golz\allowbreak{}Haghtalab\allowbreak{}Yang2025\allowbreak{}Distortion.\allowbreak{}theorem12\allowbreak{}Empirical\allowbreak{}Borda\allowbreak{}Score}
\textbf{Declaration location:} paper-local
\paragraph{Lean declaration body}
\begin{ReviewVerbatim}
type:
{Report : Type u_1} →
  {Alternative : Type u_2} →
    [Fintype Alternative] →
      (Alternative → Report → ℝ) → (Alternative → Report → ℝ) → Alternative → {horizon : ℕ} → (Fin horizon → Report) → ℝ

value:
fun {Report : Type u_1} {Alternative : Type u_2} [Fintype Alternative] (wins incidences : Alternative → Report → ℝ)
    (alternative : Alternative) {horizon : ℕ} (sample : Fin horizon → Report) =>
  if GolzHaghtalabYang2025Distortion.theorem12EmpiricalBordaIncidences incidences alternative sample = 0 then 0
  else
    GolzHaghtalabYang2025Distortion.theorem12EmpiricalBordaWins wins alternative sample /
      GolzHaghtalabYang2025Distortion.theorem12EmpiricalBordaIncidences incidences alternative sample
\end{ReviewVerbatim}

\paragraph{Direct retained dependencies}
\begin{itemize}\raggedright
\item \hyperlink{governing-declaration-88811107533a14f3}{Golz\allowbreak{}Haghtalab\allowbreak{}Yang2025\allowbreak{}Distortion.\allowbreak{}theorem12\allowbreak{}Empirical\allowbreak{}Borda\allowbreak{}Incidences}
\item \hyperlink{governing-declaration-76ef72f9364a0835}{Golz\allowbreak{}Haghtalab\allowbreak{}Yang2025\allowbreak{}Distortion.\allowbreak{}theorem12\allowbreak{}Empirical\allowbreak{}Borda\allowbreak{}Wins}
\end{itemize}
\hypertarget{governing-declaration-aec5a9f50c9e967a}{}
\subsection*{\small\ttfamily\raggedright Golz\allowbreak{}Haghtalab\allowbreak{}Yang2025\allowbreak{}Distortion.\allowbreak{}theorem12\allowbreak{}Eq10\allowbreak{}Coefficient}
\textbf{Declaration location:} paper-local
\paragraph{Lean declaration body}
\begin{ReviewVerbatim}
type:
ℝ → ℝ → ℝ

value:
fun (beta gamma : ℝ) =>
  beta / 4 / (beta.sigmoid - 1 / 2) * (1 - gamma + ((beta * gamma).sigmoid - 1 / 2) / (beta.sigmoid - 1 / 2))
\end{ReviewVerbatim}


\hypertarget{governing-declaration-e5cae00c98b32d81}{}
\subsection*{\small\ttfamily\raggedright Golz\allowbreak{}Haghtalab\allowbreak{}Yang2025\allowbreak{}Distortion.\allowbreak{}theorem12\allowbreak{}Iid\allowbreak{}Borda\allowbreak{}Report\allowbreak{}Batch\allowbreak{}Law}
\textbf{Declaration location:} paper-local
\paragraph{Lean declaration body}
\begin{ReviewVerbatim}
type:
{Report : Type u_1} → [Fintype Report] → [DecidableEq Report] → PMF Report → (horizon : ℕ) → PMF (Fin horizon → Report)

value:
fun {Report : Type u_1} [Fintype Report] [DecidableEq Report] (reportLaw : PMF Report) (horizon : ℕ) =>
  AppliedModelingLib.pmfProduct (Fin horizon) Report reportLaw
\end{ReviewVerbatim}

\paragraph{Direct retained dependencies}
\begin{itemize}\raggedright
\item \hyperlink{governing-declaration-5ac12d0a52b69068}{Applied\allowbreak{}Modeling\allowbreak{}Lib.\allowbreak{}pmf\allowbreak{}Product}
\end{itemize}
\hypertarget{governing-declaration-3940e2eb509393ea}{}
\subsection*{\small\ttfamily\raggedright Golz\allowbreak{}Haghtalab\allowbreak{}Yang2025\allowbreak{}Distortion.\allowbreak{}theorem12\allowbreak{}Iid\allowbreak{}Empirical\allowbreak{}Borda\allowbreak{}Failure}
\textbf{Declaration location:} paper-local
\paragraph{Lean declaration body}
\begin{ReviewVerbatim}
type:
{Report : Type u_1} →
  {Alternative : Type u_2} →
    [Fintype Alternative] →
      [DecidableEq Alternative] →
        (Alternative → Report → ℝ) →
          (Alternative → Report → ℝ) → Alternative → {horizon : ℕ} → (Fin horizon → Report) → Prop

value:
fun {Report : Type u_1} {Alternative : Type u_2} [Fintype Alternative] [DecidableEq Alternative]
    (wins incidences : Alternative → Report → ℝ) (winner : Alternative) {horizon : ℕ} (sample : Fin horizon → Report) =>
  ∃ (ordinary : { ordinary : Alternative // ordinary ≠ winner }),
    GolzHaghtalabYang2025Distortion.theorem12EmpiricalBordaScore wins incidences winner sample ≤
      GolzHaghtalabYang2025Distortion.theorem12EmpiricalBordaScore wins incidences (↑ordinary) sample
\end{ReviewVerbatim}

\paragraph{Direct retained dependencies}
\begin{itemize}\raggedright
\item \hyperlink{governing-declaration-b2a8d113fc3130c2}{Golz\allowbreak{}Haghtalab\allowbreak{}Yang2025\allowbreak{}Distortion.\allowbreak{}theorem12\allowbreak{}Empirical\allowbreak{}Borda\allowbreak{}Score}
\end{itemize}
\hypertarget{governing-declaration-df5d365b55936500}{}
\subsection*{\small\ttfamily\raggedright Golz\allowbreak{}Haghtalab\allowbreak{}Yang2025\allowbreak{}Distortion.\allowbreak{}theorem12\allowbreak{}Log\allowbreak{}Choice}
\textbf{Declaration location:} paper-local
\paragraph{Lean declaration body}
\begin{ReviewVerbatim}
type:
ℝ → ℝ

value:
fun (beta : ℝ) => Real.log (beta + 1) / beta
\end{ReviewVerbatim}


\hypertarget{governing-declaration-baa615edc2baf254}{}
\subsection*{\small\ttfamily\raggedright Golz\allowbreak{}Haghtalab\allowbreak{}Yang2025\allowbreak{}Distortion.\allowbreak{}theorem12\allowbreak{}One\allowbreak{}Comparison\allowbreak{}Incidences}
\textbf{Declaration location:} paper-local
\paragraph{Lean declaration body}
\begin{ReviewVerbatim}
type:
{Alternative : Type u_1} → [DecidableEq Alternative] → Alternative → (Alternative × Alternative) × Bool → ℝ

value:
fun {Alternative : Type u_1} [DecidableEq Alternative] (alternative : Alternative)
    (a : (Alternative × Alternative) × Bool) =>
  match a with
  | ((first, second), snd) => (if first = alternative then 1 else 0) + if second = alternative then 1 else 0
\end{ReviewVerbatim}


\hypertarget{governing-declaration-9f5f3a71633fcfda}{}
\subsection*{\small\ttfamily\raggedright Golz\allowbreak{}Haghtalab\allowbreak{}Yang2025\allowbreak{}Distortion.\allowbreak{}theorem12\allowbreak{}One\allowbreak{}Comparison\allowbreak{}Report\allowbreak{}Law}
\textbf{Declaration location:} paper-local
\paragraph{Lean declaration body}
\begin{ReviewVerbatim}
type:
{Alternative : Type u_1} →
  [Fintype Alternative] →
    [DecidableEq Alternative] →
      PMF Alternative →
        AppliedModelingLib.Learning.HumanFeedback.PairwisePreference PUnit.{u_2 + 1} Alternative →
          PMF ((Alternative × Alternative) × Bool)

value:
fun {Alternative : Type u_1} [Fintype Alternative] [DecidableEq Alternative] (sampling : PMF Alternative)
    (preference : AppliedModelingLib.Learning.HumanFeedback.PairwisePreference PUnit.{u_2 + 1} Alternative) =>
  sampling.bind fun (first : Alternative) =>
    sampling.bind fun (second : Alternative) =>
      PMF.map (fun (outcome : Bool) => ((first, second), outcome))
        (GolzHaghtalabYang2025Distortion.theorem12BernoulliReport (preference.prob PUnit.unit first second)
          (preference.nonneg PUnit.unit first second) (preference.le_one PUnit.unit first second))
\end{ReviewVerbatim}

\paragraph{Direct retained dependencies}
\begin{itemize}\raggedright
\item \hyperlink{library-prerequisite-7e2db20cf087d9ba}{Applied\allowbreak{}Modeling\allowbreak{}Lib.\allowbreak{}Learning.\allowbreak{}Human\allowbreak{}Feedback.\allowbreak{}Pairwise\allowbreak{}Preference}
\item \hyperlink{governing-declaration-96705d9adbc097bb}{Golz\allowbreak{}Haghtalab\allowbreak{}Yang2025\allowbreak{}Distortion.\allowbreak{}theorem12\allowbreak{}Bernoulli\allowbreak{}Report}
\end{itemize}
\hypertarget{governing-declaration-0d61e98c23f20a22}{}
\subsection*{\small\ttfamily\raggedright Golz\allowbreak{}Haghtalab\allowbreak{}Yang2025\allowbreak{}Distortion.\allowbreak{}theorem12\allowbreak{}One\allowbreak{}Comparison\allowbreak{}Wins}
\textbf{Declaration location:} paper-local
\paragraph{Lean declaration body}
\begin{ReviewVerbatim}
type:
{Alternative : Type u_1} → [DecidableEq Alternative] → Alternative → (Alternative × Alternative) × Bool → ℝ

value:
fun {Alternative : Type u_1} [DecidableEq Alternative] (alternative : Alternative)
    (a : (Alternative × Alternative) × Bool) =>
  match a with
  | ((first, second), outcome) =>
    (if first = alternative ∧ outcome = true then 1 else 0) + if second = alternative ∧ ¬outcome = true then 1 else 0
\end{ReviewVerbatim}


\hypertarget{governing-declaration-35d01f4df2faa71d}{}
\subsection*{\small\ttfamily\raggedright Golz\allowbreak{}Haghtalab\allowbreak{}Yang2025\allowbreak{}Distortion.\allowbreak{}theorem12\allowbreak{}Split\allowbreak{}Utility}
\textbf{Declaration location:} paper-local
\paragraph{Lean declaration body}
\begin{ReviewVerbatim}
type:
{User : Type u} →
  {OldAlternative : Type v} →
    {NewAlternative : Type w} →
      AppliedModelingLib.Alignment.Welfare.FiniteUtilityProfile User OldAlternative →
        (NewAlternative → OldAlternative) → User → NewAlternative → ℝ

value:
fun {User : Type u} {OldAlternative : Type v} {NewAlternative : Type w}
    (utility : AppliedModelingLib.Alignment.Welfare.FiniteUtilityProfile User OldAlternative)
    (collapse : NewAlternative → OldAlternative) (a : User) (a_1 : NewAlternative) =>
  utility a (collapse a_1)
\end{ReviewVerbatim}

\paragraph{Direct retained dependencies}
\begin{itemize}\raggedright
\item \hyperlink{library-prerequisite-56c0734f8877ddbe}{Applied\allowbreak{}Modeling\allowbreak{}Lib.\allowbreak{}Alignment.\allowbreak{}Welfare.\allowbreak{}Finite\allowbreak{}Utility\allowbreak{}Profile}
\end{itemize}
\hypertarget{governing-declaration-a73ce8ecb43a8abd}{}
\subsection*{\small\ttfamily\raggedright Golz\allowbreak{}Haghtalab\allowbreak{}Yang2025\allowbreak{}Distortion.\allowbreak{}theorem12\allowbreak{}Utility}
\textbf{Declaration location:} paper-local
\paragraph{Lean declaration body}
\begin{ReviewVerbatim}
type:
ℝ → ℝ → ℝ → Fin 3 → Fin 3 → ℝ

value:
fun (epsilon epsilonZero gamma : ℝ) (a a_1 : Fin 3) =>
  match a, a_1 with
  | 0, 0 => 1 - gamma
  | 0, 1 => 1
  | 0, 2 => 0
  | 1, 0 => 1
  | 1, 1 => 0
  | 1, 2 => epsilon
  | 2, 0 => 0
  | 2, 1 => 0
  | 2, 2 => epsilon + epsilonZero
\end{ReviewVerbatim}


\hypertarget{governing-declaration-5d4c92feb59536c0}{}
\subsection*{\small\ttfamily\raggedright Golz\allowbreak{}Haghtalab\allowbreak{}Yang2025\allowbreak{}Distortion.\allowbreak{}theorem2\allowbreak{}Source\allowbreak{}Rate}
\textbf{Declaration location:} paper-local
\paragraph{Lean declaration body}
\begin{ReviewVerbatim}
type:
ℕ → ℕ → ℕ → ℝ → ℝ → ℝ

value:
fun (alternativeCount users comparisonsPerUser : ℕ) (minimumMass delta : ℝ) =>
  have logLevel : ℝ := Real.log (8 * ↑alternativeCount ^ 2 / delta);
  20 * √(↑users * logLevel / min 1 (↑comparisonsPerUser * minimumMass ^ 2)) / ↑users +
    8 * (↑alternativeCount + 1) * logLevel / (↑users * minimumMass)
\end{ReviewVerbatim}


\hypertarget{governing-declaration-2302e6a831f3044f}{}
\subsection*{\small\ttfamily\raggedright Golz\allowbreak{}Haghtalab\allowbreak{}Yang2025\allowbreak{}Distortion.\allowbreak{}theorem2\allowbreak{}User\allowbreak{}Batch\allowbreak{}Report}
\textbf{Declaration location:} paper-local
\paragraph{Lean declaration body}
\begin{ReviewVerbatim}
type:
Type u_1 → ℕ → Type (max 0 u_1)

value:
fun (Alternative : Type u_1) (comparisonsPerUser : ℕ) => Fin comparisonsPerUser → (Alternative × Alternative) × Bool
\end{ReviewVerbatim}


\hypertarget{governing-declaration-0d0f32b50ad2d6cd}{}
\subsection*{\small\ttfamily\raggedright Golz\allowbreak{}Haghtalab\allowbreak{}Yang2025\allowbreak{}Distortion.\allowbreak{}theorem2\allowbreak{}User\allowbreak{}Batch\allowbreak{}Borda\allowbreak{}Incidences}
\textbf{Declaration location:} paper-local
\paragraph{Lean declaration body}
\begin{ReviewVerbatim}
type:
{Alternative : Type u_1} →
  [DecidableEq Alternative] →
    {comparisonsPerUser : ℕ} →
      Alternative → GolzHaghtalabYang2025Distortion.theorem2UserBatchReport Alternative comparisonsPerUser → ℝ

value:
fun {Alternative : Type u_1} [DecidableEq Alternative] {comparisonsPerUser : ℕ} (alternative : Alternative)
    (report : GolzHaghtalabYang2025Distortion.theorem2UserBatchReport Alternative comparisonsPerUser) =>
  ∑ position : Fin comparisonsPerUser,
    GolzHaghtalabYang2025Distortion.theorem12OneComparisonIncidences alternative (report position)
\end{ReviewVerbatim}

\paragraph{Direct retained dependencies}
\begin{itemize}\raggedright
\item \hyperlink{governing-declaration-baa615edc2baf254}{Golz\allowbreak{}Haghtalab\allowbreak{}Yang2025\allowbreak{}Distortion.\allowbreak{}theorem12\allowbreak{}One\allowbreak{}Comparison\allowbreak{}Incidences}
\item \hyperlink{governing-declaration-2302e6a831f3044f}{Golz\allowbreak{}Haghtalab\allowbreak{}Yang2025\allowbreak{}Distortion.\allowbreak{}theorem2\allowbreak{}User\allowbreak{}Batch\allowbreak{}Report}
\end{itemize}
\hypertarget{governing-declaration-2888dfae9507fa11}{}
\subsection*{\small\ttfamily\raggedright Golz\allowbreak{}Haghtalab\allowbreak{}Yang2025\allowbreak{}Distortion.\allowbreak{}theorem2\allowbreak{}User\allowbreak{}Batch\allowbreak{}Borda\allowbreak{}Wins}
\textbf{Declaration location:} paper-local
\paragraph{Lean declaration body}
\begin{ReviewVerbatim}
type:
{Alternative : Type u_1} →
  [DecidableEq Alternative] →
    {comparisonsPerUser : ℕ} →
      Alternative → GolzHaghtalabYang2025Distortion.theorem2UserBatchReport Alternative comparisonsPerUser → ℝ

value:
fun {Alternative : Type u_1} [DecidableEq Alternative] {comparisonsPerUser : ℕ} (alternative : Alternative)
    (report : GolzHaghtalabYang2025Distortion.theorem2UserBatchReport Alternative comparisonsPerUser) =>
  ∑ position : Fin comparisonsPerUser,
    GolzHaghtalabYang2025Distortion.theorem12OneComparisonWins alternative (report position)
\end{ReviewVerbatim}

\paragraph{Direct retained dependencies}
\begin{itemize}\raggedright
\item \hyperlink{governing-declaration-0d61e98c23f20a22}{Golz\allowbreak{}Haghtalab\allowbreak{}Yang2025\allowbreak{}Distortion.\allowbreak{}theorem12\allowbreak{}One\allowbreak{}Comparison\allowbreak{}Wins}
\item \hyperlink{governing-declaration-2302e6a831f3044f}{Golz\allowbreak{}Haghtalab\allowbreak{}Yang2025\allowbreak{}Distortion.\allowbreak{}theorem2\allowbreak{}User\allowbreak{}Batch\allowbreak{}Report}
\end{itemize}
\hypertarget{governing-declaration-98c3ee9a11a59ef9}{}
\subsection*{\small\ttfamily\raggedright Golz\allowbreak{}Haghtalab\allowbreak{}Yang2025\allowbreak{}Distortion.\allowbreak{}theorem2\allowbreak{}User\allowbreak{}Pair\allowbreak{}Label\allowbreak{}Law}
\textbf{Declaration location:} paper-local
\paragraph{Lean declaration body}
\begin{ReviewVerbatim}
type:
{Alternative : Type u_1} → [Fintype Alternative] → PMF Alternative → PMF (Alternative × Alternative)

value:
fun {Alternative : Type u_1} [Fintype Alternative] (sampling : PMF Alternative) =>
  AppliedModelingLib.pmfProd sampling sampling
\end{ReviewVerbatim}

\paragraph{Direct retained dependencies}
\begin{itemize}\raggedright
\item \hyperlink{library-prerequisite-de8e2a8bfe221a38}{Applied\allowbreak{}Modeling\allowbreak{}Lib.\allowbreak{}pmf\allowbreak{}Prod}
\end{itemize}
\hypertarget{governing-declaration-c5d0273cbd7ddcc9}{}
\subsection*{\small\ttfamily\raggedright Golz\allowbreak{}Haghtalab\allowbreak{}Yang2025\allowbreak{}Distortion.\allowbreak{}theorem2\allowbreak{}User\allowbreak{}Response\allowbreak{}Table}
\textbf{Declaration location:} paper-local
\paragraph{Lean declaration body}
\begin{ReviewVerbatim}
type:
Type u_1 → Type u_1

value:
fun (Alternative : Type u_1) => Alternative → Alternative → Bool
\end{ReviewVerbatim}


\hypertarget{governing-declaration-c9fd7bce4d99fa56}{}
\subsection*{\small\ttfamily\raggedright Golz\allowbreak{}Haghtalab\allowbreak{}Yang2025\allowbreak{}Distortion.\allowbreak{}theorem10\allowbreak{}Iid\allowbreak{}User\allowbreak{}Unordered\allowbreak{}Pair\allowbreak{}Count}
\textbf{Declaration location:} paper-local
\paragraph{Lean declaration body}
\begin{ReviewVerbatim}
type:
{Alternative : Type u_1} →
  [DecidableEq Alternative] →
    {users comparisonsPerUser : ℕ} →
      Alternative →
        Alternative →
          (Fin users →
              GolzHaghtalabYang2025Distortion.theorem2UserResponseTable Alternative ×
                (Fin comparisonsPerUser → Alternative × Alternative)) →
            ℝ

value:
fun {Alternative : Type u_1} [DecidableEq Alternative] {users comparisonsPerUser : ℕ} (first second : Alternative)
    (sample :
      Fin users →
        GolzHaghtalabYang2025Distortion.theorem2UserResponseTable Alternative ×
          (Fin comparisonsPerUser → Alternative × Alternative)) =>
  ∑ user : Fin users, GolzHaghtalabYang2025Distortion.theorem10UserUnorderedPairCount first second (sample user).2
\end{ReviewVerbatim}

\paragraph{Direct retained dependencies}
\begin{itemize}\raggedright
\item \hyperlink{governing-declaration-2fe7d1ffcbbbb71a}{Golz\allowbreak{}Haghtalab\allowbreak{}Yang2025\allowbreak{}Distortion.\allowbreak{}theorem10\allowbreak{}User\allowbreak{}Unordered\allowbreak{}Pair\allowbreak{}Count}
\item \hyperlink{governing-declaration-c5d0273cbd7ddcc9}{Golz\allowbreak{}Haghtalab\allowbreak{}Yang2025\allowbreak{}Distortion.\allowbreak{}theorem2\allowbreak{}User\allowbreak{}Response\allowbreak{}Table}
\end{itemize}
\hypertarget{governing-declaration-b49ea9dfb34ee95a}{}
\subsection*{\small\ttfamily\raggedright Golz\allowbreak{}Haghtalab\allowbreak{}Yang2025\allowbreak{}Distortion.\allowbreak{}theorem10\allowbreak{}User\allowbreak{}Unordered\allowbreak{}Pair\allowbreak{}Wins}
\textbf{Declaration location:} paper-local
\paragraph{Lean declaration body}
\begin{ReviewVerbatim}
type:
{Alternative : Type u_1} →
  [DecidableEq Alternative] →
    {comparisonsPerUser : ℕ} →
      Alternative →
        Alternative →
          GolzHaghtalabYang2025Distortion.theorem2UserResponseTable Alternative →
            (Fin comparisonsPerUser → Alternative × Alternative) → ℝ

value:
fun {Alternative : Type u_1} [DecidableEq Alternative] {comparisonsPerUser : ℕ} (first second : Alternative)
    (response : GolzHaghtalabYang2025Distortion.theorem2UserResponseTable Alternative)
    (labels : Fin comparisonsPerUser → Alternative × Alternative) =>
  ∑ position : Fin comparisonsPerUser,
    GolzHaghtalabYang2025Distortion.theorem10UnorderedPairIndicator first second (labels position) *
      GolzHaghtalabYang2025Distortion.theorem12OneComparisonWins first
        (labels position, response (labels position).1 (labels position).2)
\end{ReviewVerbatim}

\paragraph{Direct retained dependencies}
\begin{itemize}\raggedright
\item \hyperlink{governing-declaration-e63568da8eac1b3c}{Golz\allowbreak{}Haghtalab\allowbreak{}Yang2025\allowbreak{}Distortion.\allowbreak{}theorem10\allowbreak{}Unordered\allowbreak{}Pair\allowbreak{}Indicator}
\item \hyperlink{governing-declaration-0d61e98c23f20a22}{Golz\allowbreak{}Haghtalab\allowbreak{}Yang2025\allowbreak{}Distortion.\allowbreak{}theorem12\allowbreak{}One\allowbreak{}Comparison\allowbreak{}Wins}
\item \hyperlink{governing-declaration-c5d0273cbd7ddcc9}{Golz\allowbreak{}Haghtalab\allowbreak{}Yang2025\allowbreak{}Distortion.\allowbreak{}theorem2\allowbreak{}User\allowbreak{}Response\allowbreak{}Table}
\end{itemize}
\hypertarget{governing-declaration-3febd1c4a8d5d851}{}
\subsection*{\small\ttfamily\raggedright Golz\allowbreak{}Haghtalab\allowbreak{}Yang2025\allowbreak{}Distortion.\allowbreak{}theorem10\allowbreak{}Iid\allowbreak{}User\allowbreak{}Unordered\allowbreak{}Pair\allowbreak{}Wins}
\textbf{Declaration location:} paper-local
\paragraph{Lean declaration body}
\begin{ReviewVerbatim}
type:
{Alternative : Type u_1} →
  [DecidableEq Alternative] →
    {users comparisonsPerUser : ℕ} →
      Alternative →
        Alternative →
          (Fin users →
              GolzHaghtalabYang2025Distortion.theorem2UserResponseTable Alternative ×
                (Fin comparisonsPerUser → Alternative × Alternative)) →
            ℝ

value:
fun {Alternative : Type u_1} [DecidableEq Alternative] {users comparisonsPerUser : ℕ} (first second : Alternative)
    (sample :
      Fin users →
        GolzHaghtalabYang2025Distortion.theorem2UserResponseTable Alternative ×
          (Fin comparisonsPerUser → Alternative × Alternative)) =>
  ∑ user : Fin users,
    GolzHaghtalabYang2025Distortion.theorem10UserUnorderedPairWins first second (sample user).1 (sample user).2
\end{ReviewVerbatim}

\paragraph{Direct retained dependencies}
\begin{itemize}\raggedright
\item \hyperlink{governing-declaration-b49ea9dfb34ee95a}{Golz\allowbreak{}Haghtalab\allowbreak{}Yang2025\allowbreak{}Distortion.\allowbreak{}theorem10\allowbreak{}User\allowbreak{}Unordered\allowbreak{}Pair\allowbreak{}Wins}
\item \hyperlink{governing-declaration-c5d0273cbd7ddcc9}{Golz\allowbreak{}Haghtalab\allowbreak{}Yang2025\allowbreak{}Distortion.\allowbreak{}theorem2\allowbreak{}User\allowbreak{}Response\allowbreak{}Table}
\end{itemize}
\hypertarget{governing-declaration-5548b54799d0baf0}{}
\subsection*{\small\ttfamily\raggedright Golz\allowbreak{}Haghtalab\allowbreak{}Yang2025\allowbreak{}Distortion.\allowbreak{}theorem10\allowbreak{}Iid\allowbreak{}User\allowbreak{}Empirical\allowbreak{}Win\allowbreak{}Rate}
\textbf{Declaration location:} paper-local
\paragraph{Lean declaration body}
\begin{ReviewVerbatim}
type:
{Alternative : Type u_1} →
  [DecidableEq Alternative] →
    {users comparisonsPerUser : ℕ} →
      Alternative →
        Alternative →
          (Fin users →
              GolzHaghtalabYang2025Distortion.theorem2UserResponseTable Alternative ×
                (Fin comparisonsPerUser → Alternative × Alternative)) →
            ℝ

value:
fun {Alternative : Type u_1} [DecidableEq Alternative] {users comparisonsPerUser : ℕ} (first second : Alternative)
    (sample :
      Fin users →
        GolzHaghtalabYang2025Distortion.theorem2UserResponseTable Alternative ×
          (Fin comparisonsPerUser → Alternative × Alternative)) =>
  if GolzHaghtalabYang2025Distortion.theorem10IidUserUnorderedPairCount first second sample = 0 then 0
  else
    GolzHaghtalabYang2025Distortion.theorem10IidUserUnorderedPairWins first second sample /
      GolzHaghtalabYang2025Distortion.theorem10IidUserUnorderedPairCount first second sample
\end{ReviewVerbatim}

\paragraph{Direct retained dependencies}
\begin{itemize}\raggedright
\item \hyperlink{governing-declaration-c9fd7bce4d99fa56}{Golz\allowbreak{}Haghtalab\allowbreak{}Yang2025\allowbreak{}Distortion.\allowbreak{}theorem10\allowbreak{}Iid\allowbreak{}User\allowbreak{}Unordered\allowbreak{}Pair\allowbreak{}Count}
\item \hyperlink{governing-declaration-3febd1c4a8d5d851}{Golz\allowbreak{}Haghtalab\allowbreak{}Yang2025\allowbreak{}Distortion.\allowbreak{}theorem10\allowbreak{}Iid\allowbreak{}User\allowbreak{}Unordered\allowbreak{}Pair\allowbreak{}Wins}
\item \hyperlink{governing-declaration-c5d0273cbd7ddcc9}{Golz\allowbreak{}Haghtalab\allowbreak{}Yang2025\allowbreak{}Distortion.\allowbreak{}theorem2\allowbreak{}User\allowbreak{}Response\allowbreak{}Table}
\end{itemize}
\hypertarget{governing-declaration-a1441da0ef6b2fa2}{}
\subsection*{\small\ttfamily\raggedright Golz\allowbreak{}Haghtalab\allowbreak{}Yang2025\allowbreak{}Distortion.\allowbreak{}theorem10\allowbreak{}Iid\allowbreak{}User\allowbreak{}Paper\allowbreak{}Win\allowbreak{}Rate}
\textbf{Declaration location:} paper-local
\paragraph{Lean declaration body}
\begin{ReviewVerbatim}
type:
{Alternative : Type u_1} →
  [DecidableEq Alternative] →
    {users comparisonsPerUser : ℕ} →
      Alternative →
        Alternative →
          (Fin users →
              GolzHaghtalabYang2025Distortion.theorem2UserResponseTable Alternative ×
                (Fin comparisonsPerUser → Alternative × Alternative)) →
            ℝ

value:
fun {Alternative : Type u_1} [DecidableEq Alternative] {users comparisonsPerUser : ℕ} (first second : Alternative)
    (sample :
      Fin users →
        GolzHaghtalabYang2025Distortion.theorem2UserResponseTable Alternative ×
          (Fin comparisonsPerUser → Alternative × Alternative)) =>
  if first = second then 1 / 2 else GolzHaghtalabYang2025Distortion.theorem10IidUserEmpiricalWinRate first second sample
\end{ReviewVerbatim}

\paragraph{Direct retained dependencies}
\begin{itemize}\raggedright
\item \hyperlink{governing-declaration-5548b54799d0baf0}{Golz\allowbreak{}Haghtalab\allowbreak{}Yang2025\allowbreak{}Distortion.\allowbreak{}theorem10\allowbreak{}Iid\allowbreak{}User\allowbreak{}Empirical\allowbreak{}Win\allowbreak{}Rate}
\item \hyperlink{governing-declaration-c5d0273cbd7ddcc9}{Golz\allowbreak{}Haghtalab\allowbreak{}Yang2025\allowbreak{}Distortion.\allowbreak{}theorem2\allowbreak{}User\allowbreak{}Response\allowbreak{}Table}
\end{itemize}
\hypertarget{governing-declaration-1241febbc6cbfb66}{}
\subsection*{\small\ttfamily\raggedright Golz\allowbreak{}Haghtalab\allowbreak{}Yang2025\allowbreak{}Distortion.\allowbreak{}theorem11\allowbreak{}Iid\allowbreak{}User\allowbreak{}Latent\allowbreak{}Borda\allowbreak{}Incidences}
\textbf{Declaration location:} paper-local
\paragraph{Lean declaration body}
\begin{ReviewVerbatim}
type:
{Alternative : Type u_1} →
  [DecidableEq Alternative] →
    {users comparisonsPerUser : ℕ} →
      Alternative →
        (Fin users →
            GolzHaghtalabYang2025Distortion.theorem2UserResponseTable Alternative ×
              (Fin comparisonsPerUser → Alternative × Alternative)) →
          ℝ

value:
fun {Alternative : Type u_1} [DecidableEq Alternative] {users comparisonsPerUser : ℕ} (alternative : Alternative)
    (sample :
      Fin users →
        GolzHaghtalabYang2025Distortion.theorem2UserResponseTable Alternative ×
          (Fin comparisonsPerUser → Alternative × Alternative)) =>
  ∑ user : Fin users,
    GolzHaghtalabYang2025Distortion.theorem2UserBatchBordaIncidences alternative
      fun (position : Fin comparisonsPerUser) =>
      ((sample user).2 position, (sample user).1 ((sample user).2 position).1 ((sample user).2 position).2)
\end{ReviewVerbatim}

\paragraph{Direct retained dependencies}
\begin{itemize}\raggedright
\item \hyperlink{governing-declaration-0d0f32b50ad2d6cd}{Golz\allowbreak{}Haghtalab\allowbreak{}Yang2025\allowbreak{}Distortion.\allowbreak{}theorem2\allowbreak{}User\allowbreak{}Batch\allowbreak{}Borda\allowbreak{}Incidences}
\item \hyperlink{governing-declaration-c5d0273cbd7ddcc9}{Golz\allowbreak{}Haghtalab\allowbreak{}Yang2025\allowbreak{}Distortion.\allowbreak{}theorem2\allowbreak{}User\allowbreak{}Response\allowbreak{}Table}
\end{itemize}
\hypertarget{governing-declaration-243fdcdf030ecaae}{}
\subsection*{\small\ttfamily\raggedright Golz\allowbreak{}Haghtalab\allowbreak{}Yang2025\allowbreak{}Distortion.\allowbreak{}theorem11\allowbreak{}Iid\allowbreak{}User\allowbreak{}Latent\allowbreak{}Borda\allowbreak{}Wins}
\textbf{Declaration location:} paper-local
\paragraph{Lean declaration body}
\begin{ReviewVerbatim}
type:
{Alternative : Type u_1} →
  [DecidableEq Alternative] →
    {users comparisonsPerUser : ℕ} →
      Alternative →
        (Fin users →
            GolzHaghtalabYang2025Distortion.theorem2UserResponseTable Alternative ×
              (Fin comparisonsPerUser → Alternative × Alternative)) →
          ℝ

value:
fun {Alternative : Type u_1} [DecidableEq Alternative] {users comparisonsPerUser : ℕ} (alternative : Alternative)
    (sample :
      Fin users →
        GolzHaghtalabYang2025Distortion.theorem2UserResponseTable Alternative ×
          (Fin comparisonsPerUser → Alternative × Alternative)) =>
  ∑ user : Fin users,
    GolzHaghtalabYang2025Distortion.theorem2UserBatchBordaWins alternative fun (position : Fin comparisonsPerUser) =>
      ((sample user).2 position, (sample user).1 ((sample user).2 position).1 ((sample user).2 position).2)
\end{ReviewVerbatim}

\paragraph{Direct retained dependencies}
\begin{itemize}\raggedright
\item \hyperlink{governing-declaration-2888dfae9507fa11}{Golz\allowbreak{}Haghtalab\allowbreak{}Yang2025\allowbreak{}Distortion.\allowbreak{}theorem2\allowbreak{}User\allowbreak{}Batch\allowbreak{}Borda\allowbreak{}Wins}
\item \hyperlink{governing-declaration-c5d0273cbd7ddcc9}{Golz\allowbreak{}Haghtalab\allowbreak{}Yang2025\allowbreak{}Distortion.\allowbreak{}theorem2\allowbreak{}User\allowbreak{}Response\allowbreak{}Table}
\end{itemize}
\hypertarget{governing-declaration-7506044331f97cc9}{}
\subsection*{\small\ttfamily\raggedright Golz\allowbreak{}Haghtalab\allowbreak{}Yang2025\allowbreak{}Distortion.\allowbreak{}theorem11\allowbreak{}Iid\allowbreak{}User\allowbreak{}Latent\allowbreak{}Borda\allowbreak{}Score}
\textbf{Declaration location:} paper-local
\paragraph{Lean declaration body}
\begin{ReviewVerbatim}
type:
{Alternative : Type u_1} →
  [DecidableEq Alternative] →
    {users comparisonsPerUser : ℕ} →
      Alternative →
        (Fin users →
            GolzHaghtalabYang2025Distortion.theorem2UserResponseTable Alternative ×
              (Fin comparisonsPerUser → Alternative × Alternative)) →
          ℝ

value:
fun {Alternative : Type u_1} [DecidableEq Alternative] {users comparisonsPerUser : ℕ} (alternative : Alternative)
    (sample :
      Fin users →
        GolzHaghtalabYang2025Distortion.theorem2UserResponseTable Alternative ×
          (Fin comparisonsPerUser → Alternative × Alternative)) =>
  if GolzHaghtalabYang2025Distortion.theorem11IidUserLatentBordaIncidences alternative sample = 0 then 0
  else
    GolzHaghtalabYang2025Distortion.theorem11IidUserLatentBordaWins alternative sample /
      GolzHaghtalabYang2025Distortion.theorem11IidUserLatentBordaIncidences alternative sample
\end{ReviewVerbatim}

\paragraph{Direct retained dependencies}
\begin{itemize}\raggedright
\item \hyperlink{governing-declaration-1241febbc6cbfb66}{Golz\allowbreak{}Haghtalab\allowbreak{}Yang2025\allowbreak{}Distortion.\allowbreak{}theorem11\allowbreak{}Iid\allowbreak{}User\allowbreak{}Latent\allowbreak{}Borda\allowbreak{}Incidences}
\item \hyperlink{governing-declaration-243fdcdf030ecaae}{Golz\allowbreak{}Haghtalab\allowbreak{}Yang2025\allowbreak{}Distortion.\allowbreak{}theorem11\allowbreak{}Iid\allowbreak{}User\allowbreak{}Latent\allowbreak{}Borda\allowbreak{}Wins}
\item \hyperlink{governing-declaration-c5d0273cbd7ddcc9}{Golz\allowbreak{}Haghtalab\allowbreak{}Yang2025\allowbreak{}Distortion.\allowbreak{}theorem2\allowbreak{}User\allowbreak{}Response\allowbreak{}Table}
\end{itemize}
\hypertarget{governing-declaration-080b9388e51a6018}{}
\subsection*{\small\ttfamily\raggedright Golz\allowbreak{}Haghtalab\allowbreak{}Yang2025\allowbreak{}Distortion.\allowbreak{}theorem2\allowbreak{}User\allowbreak{}Response\allowbreak{}Calibrated}
\textbf{Declaration location:} paper-local
\paragraph{Lean declaration body}
\begin{ReviewVerbatim}
type:
{Alternative : Type u_1} →
  [Fintype Alternative] →
    [DecidableEq Alternative] →
      PMF (GolzHaghtalabYang2025Distortion.theorem2UserResponseTable Alternative) →
        AppliedModelingLib.Learning.HumanFeedback.PairwisePreference PUnit.{u_2 + 1} Alternative → Prop

value:
fun {Alternative : Type u_1} [Fintype Alternative] [DecidableEq Alternative]
    (responseLaw : PMF (GolzHaghtalabYang2025Distortion.theorem2UserResponseTable Alternative))
    (preference : AppliedModelingLib.Learning.HumanFeedback.PairwisePreference PUnit.{u_2 + 1} Alternative) =>
  ∀ (first second : Alternative),
    (AppliedModelingLib.pmfExp responseLaw
        fun (response : GolzHaghtalabYang2025Distortion.theorem2UserResponseTable Alternative) =>
        if response first second = true then 1 else 0) =
      preference.prob PUnit.unit first second
\end{ReviewVerbatim}

\paragraph{Direct retained dependencies}
\begin{itemize}\raggedright
\item \hyperlink{library-prerequisite-7e2db20cf087d9ba}{Applied\allowbreak{}Modeling\allowbreak{}Lib.\allowbreak{}Learning.\allowbreak{}Human\allowbreak{}Feedback.\allowbreak{}Pairwise\allowbreak{}Preference}
\item \hyperlink{governing-declaration-8e4b4389a4f6454b}{Applied\allowbreak{}Modeling\allowbreak{}Lib.\allowbreak{}pmf\allowbreak{}Exp}
\item \hyperlink{governing-declaration-c5d0273cbd7ddcc9}{Golz\allowbreak{}Haghtalab\allowbreak{}Yang2025\allowbreak{}Distortion.\allowbreak{}theorem2\allowbreak{}User\allowbreak{}Response\allowbreak{}Table}
\end{itemize}
\hypertarget{governing-declaration-1898dea7a2cffadf}{}
\subsection*{\small\ttfamily\raggedright Golz\allowbreak{}Haghtalab\allowbreak{}Yang2025\allowbreak{}Distortion.\allowbreak{}theorem3\allowbreak{}Balance\allowbreak{}Denominator}
\textbf{Declaration location:} paper-local
\paragraph{Lean declaration body}
\begin{ReviewVerbatim}
type:
ℝ → ℝ → ℝ

value:
fun (beta epsilon : ℝ) => beta.sigmoid + (beta * epsilon).sigmoid - 1
\end{ReviewVerbatim}


\hypertarget{governing-declaration-0f03655f38b9e33e}{}
\subsection*{\small\ttfamily\raggedright Golz\allowbreak{}Haghtalab\allowbreak{}Yang2025\allowbreak{}Distortion.\allowbreak{}theorem3\allowbreak{}D2\allowbreak{}Diagonal\allowbreak{}Epsilon}
\textbf{Declaration location:} paper-local
\paragraph{Lean declaration body}
\begin{ReviewVerbatim}
type:
ℕ → ℝ

value:
fun (n : ℕ) => (↑(n + 3))⁻¹
\end{ReviewVerbatim}


\hypertarget{governing-declaration-783b316e1c8a0ee2}{}
\subsection*{\small\ttfamily\raggedright Golz\allowbreak{}Haghtalab\allowbreak{}Yang2025\allowbreak{}Distortion.\allowbreak{}theorem3\allowbreak{}D2\allowbreak{}Diagonal\allowbreak{}Xi}
\textbf{Declaration location:} paper-local
\paragraph{Lean declaration body}
\begin{ReviewVerbatim}
type:
ℕ → ℝ

value:
fun (n : ℕ) => 1 + GolzHaghtalabYang2025Distortion.theorem3D2DiagonalEpsilon n
\end{ReviewVerbatim}

\paragraph{Direct retained dependencies}
\begin{itemize}\raggedright
\item \hyperlink{governing-declaration-0f03655f38b9e33e}{Golz\allowbreak{}Haghtalab\allowbreak{}Yang2025\allowbreak{}Distortion.\allowbreak{}theorem3\allowbreak{}D2\allowbreak{}Diagonal\allowbreak{}Epsilon}
\end{itemize}
\hypertarget{governing-declaration-a7437eb8cad7f1a3}{}
\subsection*{\small\ttfamily\raggedright Golz\allowbreak{}Haghtalab\allowbreak{}Yang2025\allowbreak{}Distortion.\allowbreak{}theorem3\allowbreak{}Diagonal\allowbreak{}Epsilon}
\textbf{Declaration location:} paper-local
\paragraph{Lean declaration body}
\begin{ReviewVerbatim}
type:
ℕ → ℝ

value:
fun (n : ℕ) => (↑(n + 2))⁻¹
\end{ReviewVerbatim}


\hypertarget{governing-declaration-8d2e05cae4023172}{}
\subsection*{\small\ttfamily\raggedright Golz\allowbreak{}Haghtalab\allowbreak{}Yang2025\allowbreak{}Distortion.\allowbreak{}theorem3\allowbreak{}Iid\allowbreak{}Empirical\allowbreak{}Strict\allowbreak{}Condorcet\allowbreak{}Loser}
\textbf{Declaration location:} paper-local
\paragraph{Lean declaration body}
\begin{ReviewVerbatim}
type:
{Report : Type u_1} →
  {Alternative : Type u_2} →
    [Fintype Alternative] →
      [DecidableEq Alternative] →
        (Alternative → Report → ℝ) → Alternative → {horizon : ℕ} → (Fin horizon → Report) → Prop

value:
fun {Report : Type u_1} {Alternative : Type u_2} [Fintype Alternative] [DecidableEq Alternative]
    (margin : Alternative → Report → ℝ) (special : Alternative) {horizon : ℕ} (sample : Fin horizon → Report) =>
  ∀ (ordinary : Alternative),
    ordinary ≠ special → 0 < AppliedModelingLib.Probability.finiteIidScoreSum (margin ordinary) sample
\end{ReviewVerbatim}

\paragraph{Direct retained dependencies}
\begin{itemize}\raggedright
\item \hyperlink{governing-declaration-65283b6bc1ebd8c0}{Applied\allowbreak{}Modeling\allowbreak{}Lib.\allowbreak{}Probability.\allowbreak{}finite\allowbreak{}Iid\allowbreak{}Score\allowbreak{}Sum}
\end{itemize}
\hypertarget{governing-declaration-7fc5db381e1a8f71}{}
\subsection*{\small\ttfamily\raggedright Golz\allowbreak{}Haghtalab\allowbreak{}Yang2025\allowbreak{}Distortion.\allowbreak{}theorem3\allowbreak{}Iid\allowbreak{}User\allowbreak{}Report\allowbreak{}Batch\allowbreak{}Law}
\textbf{Declaration location:} paper-local
\paragraph{Lean declaration body}
\begin{ReviewVerbatim}
type:
{Report : Type u_1} → [Fintype Report] → [DecidableEq Report] → PMF Report → (horizon : ℕ) → PMF (Fin horizon → Report)

value:
fun {Report : Type u_1} [Fintype Report] [DecidableEq Report] (reportLaw : PMF Report) (horizon : ℕ) =>
  AppliedModelingLib.pmfProduct (Fin horizon) Report reportLaw
\end{ReviewVerbatim}

\paragraph{Direct retained dependencies}
\begin{itemize}\raggedright
\item \hyperlink{governing-declaration-5ac12d0a52b69068}{Applied\allowbreak{}Modeling\allowbreak{}Lib.\allowbreak{}pmf\allowbreak{}Product}
\end{itemize}
\hypertarget{governing-declaration-bfef2d422eb25497}{}
\subsection*{\small\ttfamily\raggedright Golz\allowbreak{}Haghtalab\allowbreak{}Yang2025\allowbreak{}Distortion.\allowbreak{}theorem3\allowbreak{}Many\allowbreak{}Alternative\allowbreak{}Utility}
\textbf{Declaration location:} paper-local
\paragraph{Lean declaration body}
\begin{ReviewVerbatim}
type:
{Alternative : Type u_1} → [DecidableEq Alternative] → ℝ → ℝ → Alternative → Bool → Alternative → ℝ

value:
fun {Alternative : Type u_1} [DecidableEq Alternative] (epsilon xi : ℝ) (special : Alternative) (a : Bool)
    (a_1 : Alternative) =>
  match a, a_1 with
  | true, alternative => if alternative = special then 1 else 0
  | false, alternative => if alternative = special then 0 else xi * epsilon
\end{ReviewVerbatim}


\hypertarget{governing-declaration-f48897680b4251cd}{}
\subsection*{\small\ttfamily\raggedright Golz\allowbreak{}Haghtalab\allowbreak{}Yang2025\allowbreak{}Distortion.\allowbreak{}theorem3\allowbreak{}Observation\allowbreak{}Selection\allowbreak{}Law}
\textbf{Declaration location:} paper-local
\paragraph{Lean declaration body}
\begin{ReviewVerbatim}
type:
{Observation : Type u_1} →
  {Alternative : Type u_2} → PMF Observation → (Observation → PMF Alternative) → PMF Alternative

value:
fun {Observation : Type u_1} {Alternative : Type u_2} (observationLaw : PMF Observation)
    (rule : Observation → PMF Alternative) =>
  observationLaw.bind rule
\end{ReviewVerbatim}


\hypertarget{governing-declaration-3f23f0dd0e01aa44}{}
\subsection*{\small\ttfamily\raggedright Golz\allowbreak{}Haghtalab\allowbreak{}Yang2025\allowbreak{}Distortion.\allowbreak{}theorem3\allowbreak{}Ordinary\allowbreak{}Special\allowbreak{}Comparison\allowbreak{}Margin}
\textbf{Declaration location:} paper-local
\paragraph{Lean declaration body}
\begin{ReviewVerbatim}
type:
{Alternative : Type u_1} →
  [DecidableEq Alternative] → Alternative → Alternative → (Alternative × Alternative) × Bool → ℝ

value:
fun {Alternative : Type u_1} [DecidableEq Alternative] (special ordinary : Alternative)
    (a : (Alternative × Alternative) × Bool) =>
  match a with
  | ((first, second), outcome) =>
    if first = ordinary ∧ second = special then if outcome = true then 1 else -1
    else if first = special ∧ second = ordinary then if outcome = true then -1 else 1 else 0
\end{ReviewVerbatim}


\hypertarget{governing-declaration-fa9a73ecc3173034}{}
\subsection*{\small\ttfamily\raggedright Golz\allowbreak{}Haghtalab\allowbreak{}Yang2025\allowbreak{}Distortion.\allowbreak{}theorem3\allowbreak{}Probabilistic\allowbreak{}Condorcet\allowbreak{}Loser\allowbreak{}Criterion}
\textbf{Declaration location:} paper-local
\paragraph{Lean declaration body}
\begin{ReviewVerbatim}
type:
{Observation : Type u_1} →
  {Alternative : Type u_2} →
    [Fintype Alternative] →
      [DecidableEq Alternative] → (Observation → PMF Alternative) → (Observation → Alternative → Prop) → Prop

value:
fun {Observation : Type u_1} {Alternative : Type u_2} [Fintype Alternative] [DecidableEq Alternative]
    (rule : Observation → PMF Alternative) (isCondorcetLoser : Observation → Alternative → Prop) =>
  ∀ (observation : Observation) (alternative : Alternative),
    isCondorcetLoser observation alternative → ((rule observation) alternative).toReal ≤ (↑(Fintype.card Alternative))⁻¹
\end{ReviewVerbatim}


\hypertarget{governing-declaration-5e11e0b8b649abae}{}
\subsection*{\small\ttfamily\raggedright Golz\allowbreak{}Haghtalab\allowbreak{}Yang2025\allowbreak{}Distortion.\allowbreak{}theorem3\allowbreak{}Reported\allowbreak{}Ordinary\allowbreak{}Special\allowbreak{}Margin}
\textbf{Declaration location:} paper-local
\paragraph{Lean declaration body}
\begin{ReviewVerbatim}
type:
{Alternative : Type u_1} →
  [DecidableEq Alternative] → {d : ℕ} → Alternative → Alternative → (Fin d → (Alternative × Alternative) × Bool) → ℝ

value:
fun {Alternative : Type u_1} [DecidableEq Alternative] {d : ℕ} (special ordinary : Alternative)
    (report : Fin d → (Alternative × Alternative) × Bool) =>
  ∑ comparisonIndex : Fin d,
    GolzHaghtalabYang2025Distortion.theorem3OrdinarySpecialComparisonMargin special ordinary (report comparisonIndex)
\end{ReviewVerbatim}

\paragraph{Direct retained dependencies}
\begin{itemize}\raggedright
\item \hyperlink{governing-declaration-3f23f0dd0e01aa44}{Golz\allowbreak{}Haghtalab\allowbreak{}Yang2025\allowbreak{}Distortion.\allowbreak{}theorem3\allowbreak{}Ordinary\allowbreak{}Special\allowbreak{}Comparison\allowbreak{}Margin}
\end{itemize}
\hypertarget{governing-declaration-d300ff02fbceb879}{}
\subsection*{\small\ttfamily\raggedright Golz\allowbreak{}Haghtalab\allowbreak{}Yang2025\allowbreak{}Distortion.\allowbreak{}theorem3\allowbreak{}Special\allowbreak{}Type\allowbreak{}Mass}
\textbf{Declaration location:} paper-local
\paragraph{Lean declaration body}
\begin{ReviewVerbatim}
type:
ℝ → ℝ → ℝ

value:
fun (beta epsilon : ℝ) =>
  ((beta * epsilon).sigmoid - 1 / 2) / GolzHaghtalabYang2025Distortion.theorem3BalanceDenominator beta epsilon
\end{ReviewVerbatim}

\paragraph{Direct retained dependencies}
\begin{itemize}\raggedright
\item \hyperlink{governing-declaration-1898dea7a2cffadf}{Golz\allowbreak{}Haghtalab\allowbreak{}Yang2025\allowbreak{}Distortion.\allowbreak{}theorem3\allowbreak{}Balance\allowbreak{}Denominator}
\end{itemize}
\hypertarget{governing-declaration-07a43ced0cebd3e9}{}
\subsection*{\small\ttfamily\raggedright Golz\allowbreak{}Haghtalab\allowbreak{}Yang2025\allowbreak{}Distortion.\allowbreak{}theorem12\allowbreak{}Type\allowbreak{}A\allowbreak{}Mass}
\textbf{Declaration location:} paper-local
\paragraph{Lean declaration body}
\begin{ReviewVerbatim}
type:
ℝ → ℝ → ℝ

value:
fun (beta epsilon : ℝ) => GolzHaghtalabYang2025Distortion.theorem3SpecialTypeMass beta epsilon
\end{ReviewVerbatim}

\paragraph{Direct retained dependencies}
\begin{itemize}\raggedright
\item \hyperlink{governing-declaration-d300ff02fbceb879}{Golz\allowbreak{}Haghtalab\allowbreak{}Yang2025\allowbreak{}Distortion.\allowbreak{}theorem3\allowbreak{}Special\allowbreak{}Type\allowbreak{}Mass}
\end{itemize}
\hypertarget{governing-declaration-2346ad50a709c645}{}
\subsection*{\small\ttfamily\raggedright Golz\allowbreak{}Haghtalab\allowbreak{}Yang2025\allowbreak{}Distortion.\allowbreak{}theorem12\allowbreak{}Type\allowbreak{}B\allowbreak{}Mass}
\textbf{Declaration location:} paper-local
\paragraph{Lean declaration body}
\begin{ReviewVerbatim}
type:
ℝ → ℝ → ℝ → ℝ

value:
fun (beta epsilon gamma : ℝ) =>
  GolzHaghtalabYang2025Distortion.theorem12TypeAMass beta epsilon * ((beta * gamma).sigmoid - 1 / 2) /
    (beta.sigmoid - 1 / 2)
\end{ReviewVerbatim}

\paragraph{Direct retained dependencies}
\begin{itemize}\raggedright
\item \hyperlink{governing-declaration-07a43ced0cebd3e9}{Golz\allowbreak{}Haghtalab\allowbreak{}Yang2025\allowbreak{}Distortion.\allowbreak{}theorem12\allowbreak{}Type\allowbreak{}A\allowbreak{}Mass}
\end{itemize}
\hypertarget{governing-declaration-a909169d83d1cd68}{}
\subsection*{\small\ttfamily\raggedright Golz\allowbreak{}Haghtalab\allowbreak{}Yang2025\allowbreak{}Distortion.\allowbreak{}theorem12\allowbreak{}Type\allowbreak{}Weight}
\textbf{Declaration location:} paper-local
\paragraph{Lean declaration body}
\begin{ReviewVerbatim}
type:
ℝ → ℝ → ℝ → Fin 3 → ℝ

value:
fun (beta epsilon gamma : ℝ) (a : Fin 3) =>
  match a with
  | 0 => GolzHaghtalabYang2025Distortion.theorem12TypeAMass beta epsilon
  | 1 => GolzHaghtalabYang2025Distortion.theorem12TypeBMass beta epsilon gamma
  | 2 =>
    1 - GolzHaghtalabYang2025Distortion.theorem12TypeAMass beta epsilon -
      GolzHaghtalabYang2025Distortion.theorem12TypeBMass beta epsilon gamma
\end{ReviewVerbatim}

\paragraph{Direct retained dependencies}
\begin{itemize}\raggedright
\item \hyperlink{governing-declaration-07a43ced0cebd3e9}{Golz\allowbreak{}Haghtalab\allowbreak{}Yang2025\allowbreak{}Distortion.\allowbreak{}theorem12\allowbreak{}Type\allowbreak{}A\allowbreak{}Mass}
\item \hyperlink{governing-declaration-2346ad50a709c645}{Golz\allowbreak{}Haghtalab\allowbreak{}Yang2025\allowbreak{}Distortion.\allowbreak{}theorem12\allowbreak{}Type\allowbreak{}B\allowbreak{}Mass}
\end{itemize}
\hypertarget{governing-declaration-787bb3a4dec78d2f}{}
\subsection*{\small\ttfamily\raggedright Golz\allowbreak{}Haghtalab\allowbreak{}Yang2025\allowbreak{}Distortion.\allowbreak{}theorem12\allowbreak{}Population}
\textbf{Declaration location:} paper-local
\paragraph{Lean declaration body}
\begin{ReviewVerbatim}
type:
(beta epsilon gamma : ℝ) → 0 < beta → 0 < epsilon → epsilon < 1 → 0 < gamma → gamma < 1 → PMF (Fin 3)

value:
fun (beta epsilon gamma : ℝ) (hbeta : 0 < beta) (hepsilon : 0 < epsilon) (hepsilon_lt_one : epsilon < 1)
    (hgamma : 0 < gamma) (hgamma_lt_one : gamma < 1) =>
  AppliedModelingLib.finiteWeightedPMF (GolzHaghtalabYang2025Distortion.theorem12TypeWeight beta epsilon gamma)
    (GolzHaghtalabYang2025Distortion.theorem12TypeWeight_nonneg hbeta hepsilon hepsilon_lt_one hgamma hgamma_lt_one)
    (GolzHaghtalabYang2025Distortion.theorem12Population._proof_1 beta epsilon gamma)
\end{ReviewVerbatim}

\paragraph{Direct retained dependencies}
\begin{itemize}\raggedright
\item \hyperlink{governing-declaration-0153ce2f02e7dc12}{Applied\allowbreak{}Modeling\allowbreak{}Lib.\allowbreak{}finite\allowbreak{}Weighted\allowbreak{}PMF}
\item \hyperlink{governing-declaration-a909169d83d1cd68}{Golz\allowbreak{}Haghtalab\allowbreak{}Yang2025\allowbreak{}Distortion.\allowbreak{}theorem12\allowbreak{}Type\allowbreak{}Weight}
\end{itemize}
\hypertarget{governing-declaration-5256b014ca2f6f1b}{}
\subsection*{\small\ttfamily\raggedright Golz\allowbreak{}Haghtalab\allowbreak{}Yang2025\allowbreak{}Distortion.\allowbreak{}theorem3\allowbreak{}Two\allowbreak{}Type\allowbreak{}Population}
\textbf{Declaration location:} paper-local
\paragraph{Lean declaration body}
\begin{ReviewVerbatim}
type:
(beta epsilon : ℝ) → 0 < beta → 0 < epsilon → PMF Bool

value:
fun (beta epsilon : ℝ) (hbeta : 0 < beta) (hepsilon : 0 < epsilon) =>
  PMF.bernoulli
    ⟨GolzHaghtalabYang2025Distortion.theorem3SpecialTypeMass beta epsilon,
      GolzHaghtalabYang2025Distortion.theorem3TwoTypePopulation._proof_1 beta epsilon hbeta hepsilon⟩
    (GolzHaghtalabYang2025Distortion.theorem3TwoTypePopulation._proof_2 beta epsilon hbeta hepsilon)
\end{ReviewVerbatim}

\paragraph{Direct retained dependencies}
\begin{itemize}\raggedright
\item \hyperlink{governing-declaration-d300ff02fbceb879}{Golz\allowbreak{}Haghtalab\allowbreak{}Yang2025\allowbreak{}Distortion.\allowbreak{}theorem3\allowbreak{}Special\allowbreak{}Type\allowbreak{}Mass}
\end{itemize}
\hypertarget{governing-declaration-2c2aa314cf0c1edd}{}
\subsection*{\small\ttfamily\raggedright Golz\allowbreak{}Haghtalab\allowbreak{}Yang2025\allowbreak{}Distortion.\allowbreak{}theorem3\allowbreak{}Many\allowbreak{}Alternative\allowbreak{}One\allowbreak{}Comparison\allowbreak{}Law}
\textbf{Declaration location:} paper-local
\paragraph{Lean declaration body}
\begin{ReviewVerbatim}
type:
{Alternative : Type u_1} →
  [Fintype Alternative] →
    [DecidableEq Alternative] →
      PMF (Alternative × Alternative) →
        (beta epsilon : ℝ) → 0 < beta → 0 < epsilon → Alternative → PMF ((Alternative × Alternative) × Bool)

value:
fun {Alternative : Type u_1} [Fintype Alternative] [DecidableEq Alternative]
    (pairSampling : PMF (Alternative × Alternative)) (beta epsilon : ℝ) (hbeta : 0 < beta) (hepsilon : 0 < epsilon)
    (special : Alternative) =>
  pairSampling.bind fun (pair : Alternative × Alternative) =>
    let preference : AppliedModelingLib.Learning.HumanFeedback.PairwisePreference PUnit.{1} Alternative :=
      AppliedModelingLib.Alignment.Welfare.populationBradleyTerryPreference
        (GolzHaghtalabYang2025Distortion.theorem3TwoTypePopulation beta epsilon hbeta hepsilon)
        (GolzHaghtalabYang2025Distortion.theorem3ManyAlternativeUtility epsilon 1 special) beta;
    PMF.map (fun (outcome : Bool) => (pair, outcome))
      (PMF.bernoulli
        ⟨preference.prob PUnit.unit pair.1 pair.2,
          GolzHaghtalabYang2025Distortion.theorem3ManyAlternativeOneComparisonLaw._proof_1 beta epsilon hbeta hepsilon
            special pair⟩
        (GolzHaghtalabYang2025Distortion.theorem3ManyAlternativeOneComparisonLaw._proof_2 beta epsilon hbeta hepsilon
          special pair))
\end{ReviewVerbatim}

\paragraph{Direct retained dependencies}
\begin{itemize}\raggedright
\item \hyperlink{library-prerequisite-cf06d15d7f8b7c31}{Applied\allowbreak{}Modeling\allowbreak{}Lib.\allowbreak{}Alignment.\allowbreak{}Welfare.\allowbreak{}population\allowbreak{}Bradley\allowbreak{}Terry\allowbreak{}Preference}
\item \hyperlink{library-prerequisite-7e2db20cf087d9ba}{Applied\allowbreak{}Modeling\allowbreak{}Lib.\allowbreak{}Learning.\allowbreak{}Human\allowbreak{}Feedback.\allowbreak{}Pairwise\allowbreak{}Preference}
\item \hyperlink{governing-declaration-bfef2d422eb25497}{Golz\allowbreak{}Haghtalab\allowbreak{}Yang2025\allowbreak{}Distortion.\allowbreak{}theorem3\allowbreak{}Many\allowbreak{}Alternative\allowbreak{}Utility}
\item \hyperlink{governing-declaration-5256b014ca2f6f1b}{Golz\allowbreak{}Haghtalab\allowbreak{}Yang2025\allowbreak{}Distortion.\allowbreak{}theorem3\allowbreak{}Two\allowbreak{}Type\allowbreak{}Population}
\end{itemize}
\hypertarget{governing-declaration-8cd7a7c5a45b332d}{}
\subsection*{\small\ttfamily\raggedright Golz\allowbreak{}Haghtalab\allowbreak{}Yang2025\allowbreak{}Distortion.\allowbreak{}theorem3\allowbreak{}Many\allowbreak{}Alternative\allowbreak{}Observation\allowbreak{}Law}
\textbf{Declaration location:} paper-local
\paragraph{Lean declaration body}
\begin{ReviewVerbatim}
type:
{Alternative : Type u_1} →
  [Fintype Alternative] →
    [DecidableEq Alternative] →
      PMF (Alternative × Alternative) →
        (beta epsilon : ℝ) →
          0 < beta → 0 < epsilon → Alternative → (horizon : ℕ) → PMF (Fin horizon → (Alternative × Alternative) × Bool)

value:
fun {Alternative : Type u_1} [Fintype Alternative] [DecidableEq Alternative]
    (pairSampling : PMF (Alternative × Alternative)) (beta epsilon : ℝ) (hbeta : 0 < beta) (hepsilon : 0 < epsilon)
    (special : Alternative) (horizon : ℕ) =>
  AppliedModelingLib.pmfProduct (Fin horizon) ((Alternative × Alternative) × Bool)
    (GolzHaghtalabYang2025Distortion.theorem3ManyAlternativeOneComparisonLaw pairSampling beta epsilon hbeta hepsilon
      special)
\end{ReviewVerbatim}

\paragraph{Direct retained dependencies}
\begin{itemize}\raggedright
\item \hyperlink{governing-declaration-5ac12d0a52b69068}{Applied\allowbreak{}Modeling\allowbreak{}Lib.\allowbreak{}pmf\allowbreak{}Product}
\item \hyperlink{governing-declaration-2c2aa314cf0c1edd}{Golz\allowbreak{}Haghtalab\allowbreak{}Yang2025\allowbreak{}Distortion.\allowbreak{}theorem3\allowbreak{}Many\allowbreak{}Alternative\allowbreak{}One\allowbreak{}Comparison\allowbreak{}Law}
\end{itemize}
\hypertarget{governing-declaration-f3bdeb2b79da5cb7}{}
\subsection*{\small\ttfamily\raggedright Golz\allowbreak{}Haghtalab\allowbreak{}Yang2025\allowbreak{}Distortion.\allowbreak{}theorem3\allowbreak{}Population\allowbreak{}Comparison\allowbreak{}Outcome\allowbreak{}Law}
\textbf{Declaration location:} paper-local
\paragraph{Lean declaration body}
\begin{ReviewVerbatim}
type:
{Alternative : Type u_1} →
  [Fintype Alternative] →
    [DecidableEq Alternative] →
      (beta epsilon : ℝ) → ℝ → 0 < beta → 0 < epsilon → Alternative → Alternative × Alternative → PMF Bool

value:
fun {Alternative : Type u_1} [Fintype Alternative] [DecidableEq Alternative] (beta epsilon xi : ℝ) (hbeta : 0 < beta)
    (hepsilon : 0 < epsilon) (special : Alternative) (pair : Alternative × Alternative) =>
  let preference : AppliedModelingLib.Learning.HumanFeedback.PairwisePreference PUnit.{1} Alternative :=
    AppliedModelingLib.Alignment.Welfare.populationBradleyTerryPreference
      (GolzHaghtalabYang2025Distortion.theorem3TwoTypePopulation beta epsilon hbeta hepsilon)
      (GolzHaghtalabYang2025Distortion.theorem3ManyAlternativeUtility epsilon xi special) beta;
  PMF.bernoulli
    ⟨preference.prob PUnit.unit pair.1 pair.2,
      GolzHaghtalabYang2025Distortion.theorem3PopulationComparisonOutcomeLaw._proof_1 beta epsilon xi hbeta hepsilon
        special pair⟩
    (GolzHaghtalabYang2025Distortion.theorem3PopulationComparisonOutcomeLaw._proof_2 beta epsilon xi hbeta hepsilon
      special pair)
\end{ReviewVerbatim}

\paragraph{Direct retained dependencies}
\begin{itemize}\raggedright
\item \hyperlink{library-prerequisite-cf06d15d7f8b7c31}{Applied\allowbreak{}Modeling\allowbreak{}Lib.\allowbreak{}Alignment.\allowbreak{}Welfare.\allowbreak{}population\allowbreak{}Bradley\allowbreak{}Terry\allowbreak{}Preference}
\item \hyperlink{library-prerequisite-7e2db20cf087d9ba}{Applied\allowbreak{}Modeling\allowbreak{}Lib.\allowbreak{}Learning.\allowbreak{}Human\allowbreak{}Feedback.\allowbreak{}Pairwise\allowbreak{}Preference}
\item \hyperlink{governing-declaration-bfef2d422eb25497}{Golz\allowbreak{}Haghtalab\allowbreak{}Yang2025\allowbreak{}Distortion.\allowbreak{}theorem3\allowbreak{}Many\allowbreak{}Alternative\allowbreak{}Utility}
\item \hyperlink{governing-declaration-5256b014ca2f6f1b}{Golz\allowbreak{}Haghtalab\allowbreak{}Yang2025\allowbreak{}Distortion.\allowbreak{}theorem3\allowbreak{}Two\allowbreak{}Type\allowbreak{}Population}
\end{itemize}
\hypertarget{governing-declaration-94724c60280865f4}{}
\subsection*{\small\ttfamily\raggedright Golz\allowbreak{}Haghtalab\allowbreak{}Yang2025\allowbreak{}Distortion.\allowbreak{}theorem3\allowbreak{}Many\allowbreak{}Alternative\allowbreak{}One\allowbreak{}Comparison\allowbreak{}Law\allowbreak{}Xi}
\textbf{Declaration location:} paper-local
\paragraph{Lean declaration body}
\begin{ReviewVerbatim}
type:
{Alternative : Type u_1} →
  [Fintype Alternative] →
    [DecidableEq Alternative] →
      PMF (Alternative × Alternative) →
        (beta epsilon : ℝ) → ℝ → 0 < beta → 0 < epsilon → Alternative → PMF ((Alternative × Alternative) × Bool)

value:
fun {Alternative : Type u_1} [Fintype Alternative] [DecidableEq Alternative]
    (pairSampling : PMF (Alternative × Alternative)) (beta epsilon xi : ℝ) (hbeta : 0 < beta) (hepsilon : 0 < epsilon)
    (special : Alternative) =>
  pairSampling.bind fun (pair : Alternative × Alternative) =>
    PMF.map (fun (outcome : Bool) => (pair, outcome))
      (GolzHaghtalabYang2025Distortion.theorem3PopulationComparisonOutcomeLaw beta epsilon xi hbeta hepsilon special
        pair)
\end{ReviewVerbatim}

\paragraph{Direct retained dependencies}
\begin{itemize}\raggedright
\item \hyperlink{governing-declaration-f3bdeb2b79da5cb7}{Golz\allowbreak{}Haghtalab\allowbreak{}Yang2025\allowbreak{}Distortion.\allowbreak{}theorem3\allowbreak{}Population\allowbreak{}Comparison\allowbreak{}Outcome\allowbreak{}Law}
\end{itemize}
\hypertarget{governing-declaration-3451d41a5ff413f7}{}
\subsection*{\small\ttfamily\raggedright Golz\allowbreak{}Haghtalab\allowbreak{}Yang2025\allowbreak{}Distortion.\allowbreak{}theorem6\allowbreak{}Alternative}
\textbf{Declaration location:} paper-local
\paragraph{Lean declaration body}
\begin{ReviewVerbatim}
type:
ℕ → Type

value:
fun (copies : ℕ) => Fin 2 ⊕ Fin copies
\end{ReviewVerbatim}


\hypertarget{governing-declaration-7b9a203f54b7d337}{}
\subsection*{\small\ttfamily\raggedright Golz\allowbreak{}Haghtalab\allowbreak{}Yang2025\allowbreak{}Distortion.\allowbreak{}theorem6\allowbreak{}A}
\textbf{Declaration location:} paper-local
\paragraph{Lean declaration body}
\begin{ReviewVerbatim}
type:
{copies : ℕ} → GolzHaghtalabYang2025Distortion.theorem6Alternative copies

value:
fun {copies : ℕ} => Sum.inl 0
\end{ReviewVerbatim}

\paragraph{Direct retained dependencies}
\begin{itemize}\raggedright
\item \hyperlink{governing-declaration-3451d41a5ff413f7}{Golz\allowbreak{}Haghtalab\allowbreak{}Yang2025\allowbreak{}Distortion.\allowbreak{}theorem6\allowbreak{}Alternative}
\end{itemize}
\hypertarget{governing-declaration-2108be1d9586e812}{}
\subsection*{\small\ttfamily\raggedright Golz\allowbreak{}Haghtalab\allowbreak{}Yang2025\allowbreak{}Distortion.\allowbreak{}theorem6\allowbreak{}Epsilon}
\textbf{Declaration location:} paper-local
\paragraph{Lean declaration body}
\begin{ReviewVerbatim}
type:
ℝ → ℝ

value:
fun (eta : ℝ) => Real.exp (-2 / eta)
\end{ReviewVerbatim}


\hypertarget{governing-declaration-1b9e5a272efa4c90}{}
\subsection*{\small\ttfamily\raggedright Golz\allowbreak{}Haghtalab\allowbreak{}Yang2025\allowbreak{}Distortion.\allowbreak{}theorem6\allowbreak{}Minority\allowbreak{}Mass}
\textbf{Declaration location:} paper-local
\paragraph{Lean declaration body}
\begin{ReviewVerbatim}
type:
ℝ → ℝ

value:
fun (beta : ℝ) => 10 / (10 + Real.exp beta)
\end{ReviewVerbatim}


\hypertarget{governing-declaration-a14b7af0e1436262}{}
\subsection*{\small\ttfamily\raggedright Golz\allowbreak{}Haghtalab\allowbreak{}Yang2025\allowbreak{}Distortion.\allowbreak{}theorem6\allowbreak{}Copy\allowbreak{}Count}
\textbf{Declaration location:} paper-local
\paragraph{Lean declaration body}
\begin{ReviewVerbatim}
type:
ℝ → ℕ

value:
fun (beta : ℝ) => ⌈2 / (GolzHaghtalabYang2025Distortion.theorem6MinorityMass beta * (1 / 3 - Real.exp 1 / 10))⌉₊
\end{ReviewVerbatim}

\paragraph{Direct retained dependencies}
\begin{itemize}\raggedright
\item \hyperlink{governing-declaration-1b9e5a272efa4c90}{Golz\allowbreak{}Haghtalab\allowbreak{}Yang2025\allowbreak{}Distortion.\allowbreak{}theorem6\allowbreak{}Minority\allowbreak{}Mass}
\end{itemize}
\hypertarget{governing-declaration-21a83f84d7ec54d5}{}
\subsection*{\small\ttfamily\raggedright Golz\allowbreak{}Haghtalab\allowbreak{}Yang2025\allowbreak{}Distortion.\allowbreak{}theorem6\allowbreak{}Reference\allowbreak{}Weight}
\textbf{Declaration location:} paper-local
\paragraph{Lean declaration body}
\begin{ReviewVerbatim}
type:
ℝ → {copies : ℕ} → GolzHaghtalabYang2025Distortion.theorem6Alternative copies → ℝ

value:
fun (epsilon : ℝ) {copies : ℕ} (a : GolzHaghtalabYang2025Distortion.theorem6Alternative copies) =>
  match a with
  | Sum.inl val => (1 - epsilon) / 2
  | Sum.inr val => epsilon / ↑copies
\end{ReviewVerbatim}

\paragraph{Direct retained dependencies}
\begin{itemize}\raggedright
\item \hyperlink{governing-declaration-3451d41a5ff413f7}{Golz\allowbreak{}Haghtalab\allowbreak{}Yang2025\allowbreak{}Distortion.\allowbreak{}theorem6\allowbreak{}Alternative}
\end{itemize}
\hypertarget{governing-declaration-0527d6356fcaa0f2}{}
\subsection*{\small\ttfamily\raggedright Golz\allowbreak{}Haghtalab\allowbreak{}Yang2025\allowbreak{}Distortion.\allowbreak{}theorem6\allowbreak{}Reference}
\textbf{Declaration location:} paper-local
\paragraph{Lean declaration body}
\begin{ReviewVerbatim}
type:
(epsilon : ℝ) →
  {copies : ℕ} →
    0 < epsilon → epsilon < 1 → 0 < copies → PMF (GolzHaghtalabYang2025Distortion.theorem6Alternative copies)

value:
fun (epsilon : ℝ) {copies : ℕ} (hepsilon_pos : 0 < epsilon) (hepsilon_lt_one : epsilon < 1)
    (hcopies_pos : 0 < copies) =>
  AppliedModelingLib.finiteWeightedPMF (GolzHaghtalabYang2025Distortion.theorem6ReferenceWeight epsilon)
    (GolzHaghtalabYang2025Distortion.theorem6Reference._proof_1 epsilon hepsilon_pos hepsilon_lt_one hcopies_pos)
    (GolzHaghtalabYang2025Distortion.theorem6Reference._proof_2 epsilon hcopies_pos)
\end{ReviewVerbatim}

\paragraph{Direct retained dependencies}
\begin{itemize}\raggedright
\item \hyperlink{governing-declaration-0153ce2f02e7dc12}{Applied\allowbreak{}Modeling\allowbreak{}Lib.\allowbreak{}finite\allowbreak{}Weighted\allowbreak{}PMF}
\item \hyperlink{governing-declaration-3451d41a5ff413f7}{Golz\allowbreak{}Haghtalab\allowbreak{}Yang2025\allowbreak{}Distortion.\allowbreak{}theorem6\allowbreak{}Alternative}
\item \hyperlink{governing-declaration-21a83f84d7ec54d5}{Golz\allowbreak{}Haghtalab\allowbreak{}Yang2025\allowbreak{}Distortion.\allowbreak{}theorem6\allowbreak{}Reference\allowbreak{}Weight}
\end{itemize}
\hypertarget{governing-declaration-c37e8478335061af}{}
\subsection*{\small\ttfamily\raggedright Golz\allowbreak{}Haghtalab\allowbreak{}Yang2025\allowbreak{}Distortion.\allowbreak{}theorem6\allowbreak{}Type\allowbreak{}Weight}
\textbf{Declaration location:} paper-local
\paragraph{Lean declaration body}
\begin{ReviewVerbatim}
type:
ℝ → Fin 2 → ℝ

value:
fun (beta : ℝ) (a : Fin 2) =>
  match a with
  | 0 => GolzHaghtalabYang2025Distortion.theorem6MinorityMass beta
  | 1 => 1 - GolzHaghtalabYang2025Distortion.theorem6MinorityMass beta
\end{ReviewVerbatim}

\paragraph{Direct retained dependencies}
\begin{itemize}\raggedright
\item \hyperlink{governing-declaration-1b9e5a272efa4c90}{Golz\allowbreak{}Haghtalab\allowbreak{}Yang2025\allowbreak{}Distortion.\allowbreak{}theorem6\allowbreak{}Minority\allowbreak{}Mass}
\end{itemize}
\hypertarget{governing-declaration-75eccaf6a07b78d3}{}
\subsection*{\small\ttfamily\raggedright Golz\allowbreak{}Haghtalab\allowbreak{}Yang2025\allowbreak{}Distortion.\allowbreak{}theorem6\allowbreak{}Population}
\textbf{Declaration location:} paper-local
\paragraph{Lean declaration body}
\begin{ReviewVerbatim}
type:
ℝ → PMF (Fin 2)

value:
fun (beta : ℝ) =>
  AppliedModelingLib.finiteWeightedPMF (GolzHaghtalabYang2025Distortion.theorem6TypeWeight beta)
    (GolzHaghtalabYang2025Distortion.theorem6Population._proof_1 beta)
    (GolzHaghtalabYang2025Distortion.theorem6Population._proof_2 beta)
\end{ReviewVerbatim}

\paragraph{Direct retained dependencies}
\begin{itemize}\raggedright
\item \hyperlink{governing-declaration-0153ce2f02e7dc12}{Applied\allowbreak{}Modeling\allowbreak{}Lib.\allowbreak{}finite\allowbreak{}Weighted\allowbreak{}PMF}
\item \hyperlink{governing-declaration-c37e8478335061af}{Golz\allowbreak{}Haghtalab\allowbreak{}Yang2025\allowbreak{}Distortion.\allowbreak{}theorem6\allowbreak{}Type\allowbreak{}Weight}
\end{itemize}
\hypertarget{governing-declaration-3f17ddfde71d8d00}{}
\subsection*{\small\ttfamily\raggedright Golz\allowbreak{}Haghtalab\allowbreak{}Yang2025\allowbreak{}Distortion.\allowbreak{}theorem6\allowbreak{}Utility}
\textbf{Declaration location:} paper-local
\paragraph{Lean declaration body}
\begin{ReviewVerbatim}
type:
ℝ → {copies : ℕ} → Fin 2 → GolzHaghtalabYang2025Distortion.theorem6Alternative copies → ℝ

value:
fun (beta : ℝ) {copies : ℕ} (a : Fin 2) (a_1 : GolzHaghtalabYang2025Distortion.theorem6Alternative copies) =>
  match a, a_1 with
  | 0, Sum.inl 0 => 0
  | 0, Sum.inl 1 => 1
  | 0, Sum.inr val => 0
  | 1, Sum.inl 0 => 1 / beta
  | 1, Sum.inl 1 => 0
  | 1, Sum.inr val => 1
\end{ReviewVerbatim}

\paragraph{Direct retained dependencies}
\begin{itemize}\raggedright
\item \hyperlink{governing-declaration-3451d41a5ff413f7}{Golz\allowbreak{}Haghtalab\allowbreak{}Yang2025\allowbreak{}Distortion.\allowbreak{}theorem6\allowbreak{}Alternative}
\end{itemize}
\hypertarget{governing-declaration-b14cd971fad8513c}{}
\subsection*{\small\ttfamily\raggedright Golz\allowbreak{}Haghtalab\allowbreak{}Yang2025\allowbreak{}Distortion.\allowbreak{}theorem6\allowbreak{}Preference}
\textbf{Declaration location:} paper-local
\paragraph{Lean declaration body}
\begin{ReviewVerbatim}
type:
ℝ →
  {copies : ℕ} →
    AppliedModelingLib.Learning.HumanFeedback.PairwisePreference PUnit.{1}
      (GolzHaghtalabYang2025Distortion.theorem6Alternative copies)

value:
fun (beta : ℝ) {copies : ℕ} =>
  AppliedModelingLib.Alignment.Welfare.populationBradleyTerryPreference
    (GolzHaghtalabYang2025Distortion.theorem6Population beta) (GolzHaghtalabYang2025Distortion.theorem6Utility beta)
    beta
\end{ReviewVerbatim}

\paragraph{Direct retained dependencies}
\begin{itemize}\raggedright
\item \hyperlink{library-prerequisite-cf06d15d7f8b7c31}{Applied\allowbreak{}Modeling\allowbreak{}Lib.\allowbreak{}Alignment.\allowbreak{}Welfare.\allowbreak{}population\allowbreak{}Bradley\allowbreak{}Terry\allowbreak{}Preference}
\item \hyperlink{library-prerequisite-7e2db20cf087d9ba}{Applied\allowbreak{}Modeling\allowbreak{}Lib.\allowbreak{}Learning.\allowbreak{}Human\allowbreak{}Feedback.\allowbreak{}Pairwise\allowbreak{}Preference}
\item \hyperlink{governing-declaration-3451d41a5ff413f7}{Golz\allowbreak{}Haghtalab\allowbreak{}Yang2025\allowbreak{}Distortion.\allowbreak{}theorem6\allowbreak{}Alternative}
\item \hyperlink{governing-declaration-75eccaf6a07b78d3}{Golz\allowbreak{}Haghtalab\allowbreak{}Yang2025\allowbreak{}Distortion.\allowbreak{}theorem6\allowbreak{}Population}
\item \hyperlink{governing-declaration-3f17ddfde71d8d00}{Golz\allowbreak{}Haghtalab\allowbreak{}Yang2025\allowbreak{}Distortion.\allowbreak{}theorem6\allowbreak{}Utility}
\end{itemize}
\hypertarget{governing-declaration-ea43bde8cd383b24}{}
\subsection*{\small\ttfamily\raggedright Golz\allowbreak{}Haghtalab\allowbreak{}Yang2025\allowbreak{}Distortion.\allowbreak{}theorem6\allowbreak{}Welfare}
\textbf{Declaration location:} paper-local
\paragraph{Lean declaration body}
\begin{ReviewVerbatim}
type:
ℝ → {copies : ℕ} → GolzHaghtalabYang2025Distortion.theorem6Alternative copies → ℝ

value:
fun (beta : ℝ) {copies : ℕ} (a : GolzHaghtalabYang2025Distortion.theorem6Alternative copies) =>
  AppliedModelingLib.Alignment.Welfare.populationAverageUtility
    (GolzHaghtalabYang2025Distortion.theorem6Population beta) (GolzHaghtalabYang2025Distortion.theorem6Utility beta) a
\end{ReviewVerbatim}

\paragraph{Direct retained dependencies}
\begin{itemize}\raggedright
\item \hyperlink{library-prerequisite-15f4b442749deab2}{Applied\allowbreak{}Modeling\allowbreak{}Lib.\allowbreak{}Alignment.\allowbreak{}Welfare.\allowbreak{}population\allowbreak{}Average\allowbreak{}Utility}
\item \hyperlink{governing-declaration-3451d41a5ff413f7}{Golz\allowbreak{}Haghtalab\allowbreak{}Yang2025\allowbreak{}Distortion.\allowbreak{}theorem6\allowbreak{}Alternative}
\item \hyperlink{governing-declaration-75eccaf6a07b78d3}{Golz\allowbreak{}Haghtalab\allowbreak{}Yang2025\allowbreak{}Distortion.\allowbreak{}theorem6\allowbreak{}Population}
\item \hyperlink{governing-declaration-3f17ddfde71d8d00}{Golz\allowbreak{}Haghtalab\allowbreak{}Yang2025\allowbreak{}Distortion.\allowbreak{}theorem6\allowbreak{}Utility}
\end{itemize}
\hypertarget{governing-declaration-2ccaa1524be736e8}{}
\subsection*{\small\ttfamily\raggedright Golz\allowbreak{}Haghtalab\allowbreak{}Yang2025\allowbreak{}Distortion.\allowbreak{}theorem7\allowbreak{}Empirical\allowbreak{}Preference}
\textbf{Declaration location:} paper-local
\paragraph{Lean declaration body}
\begin{ReviewVerbatim}
type:
{Alternative : Type u_1} →
  (empiricalRate : Alternative → Alternative → ℝ) →
    (∀ (first second : Alternative), 0 ≤ empiricalRate first second) →
      (∀ (first second : Alternative), empiricalRate first second ≤ 1) →
        AppliedModelingLib.Learning.HumanFeedback.PairwisePreference PUnit.{1} Alternative

value:
fun {Alternative : Type u_1} (empiricalRate : Alternative → Alternative → ℝ)
    (hempirical_nonneg : ∀ (first second : Alternative), 0 ≤ empiricalRate first second)
    (hempirical_le_one : ∀ (first second : Alternative), empiricalRate first second ≤ 1) =>
  {
    prob := fun (x : PUnit.{1}) (first second : Alternative) =>
      1 / 2 + (empiricalRate first second - empiricalRate second first) / 2,
    nonneg :=
      GolzHaghtalabYang2025Distortion.theorem7EmpiricalPreference._proof_2 empiricalRate hempirical_nonneg
        hempirical_le_one,
    le_one :=
      GolzHaghtalabYang2025Distortion.theorem7EmpiricalPreference._proof_3 empiricalRate hempirical_nonneg
        hempirical_le_one,
    complementary := GolzHaghtalabYang2025Distortion.theorem7EmpiricalPreference._proof_4 empiricalRate }
\end{ReviewVerbatim}

\paragraph{Direct retained dependencies}
\begin{itemize}\raggedright
\item \hyperlink{library-prerequisite-7e2db20cf087d9ba}{Applied\allowbreak{}Modeling\allowbreak{}Lib.\allowbreak{}Learning.\allowbreak{}Human\allowbreak{}Feedback.\allowbreak{}Pairwise\allowbreak{}Preference}
\end{itemize}
\hypertarget{governing-declaration-485a7e8af83d2b11}{}
\subsection*{\small\ttfamily\raggedright Golz\allowbreak{}Haghtalab\allowbreak{}Yang2025\allowbreak{}Distortion.\allowbreak{}theorem7\allowbreak{}Iid\allowbreak{}User\allowbreak{}Batch\allowbreak{}Law}
\textbf{Declaration location:} paper-local
\paragraph{Lean declaration body}
\begin{ReviewVerbatim}
type:
{Alternative : Type u_1} →
  [Fintype Alternative] →
    [DecidableEq Alternative] →
      PMF (GolzHaghtalabYang2025Distortion.theorem2UserResponseTable Alternative) →
        PMF Alternative →
          (users comparisonsPerUser : ℕ) →
            PMF
              (Fin users →
                GolzHaghtalabYang2025Distortion.theorem2UserResponseTable Alternative ×
                  (Fin comparisonsPerUser → Alternative × Alternative))

value:
fun {Alternative : Type u_1} [Fintype Alternative] [DecidableEq Alternative]
    (responseLaw : PMF (GolzHaghtalabYang2025Distortion.theorem2UserResponseTable Alternative))
    (sampling : PMF Alternative) (users comparisonsPerUser : ℕ) =>
  AppliedModelingLib.pmfProduct (Fin users)
    (GolzHaghtalabYang2025Distortion.theorem2UserResponseTable Alternative ×
      (Fin comparisonsPerUser → Alternative × Alternative))
    (AppliedModelingLib.pmfProd responseLaw
      (AppliedModelingLib.pmfProduct (Fin comparisonsPerUser) (Alternative × Alternative)
        (GolzHaghtalabYang2025Distortion.theorem2UserPairLabelLaw sampling)))
\end{ReviewVerbatim}

\paragraph{Direct retained dependencies}
\begin{itemize}\raggedright
\item \hyperlink{library-prerequisite-de8e2a8bfe221a38}{Applied\allowbreak{}Modeling\allowbreak{}Lib.\allowbreak{}pmf\allowbreak{}Prod}
\item \hyperlink{governing-declaration-5ac12d0a52b69068}{Applied\allowbreak{}Modeling\allowbreak{}Lib.\allowbreak{}pmf\allowbreak{}Product}
\item \hyperlink{governing-declaration-98c3ee9a11a59ef9}{Golz\allowbreak{}Haghtalab\allowbreak{}Yang2025\allowbreak{}Distortion.\allowbreak{}theorem2\allowbreak{}User\allowbreak{}Pair\allowbreak{}Label\allowbreak{}Law}
\item \hyperlink{governing-declaration-c5d0273cbd7ddcc9}{Golz\allowbreak{}Haghtalab\allowbreak{}Yang2025\allowbreak{}Distortion.\allowbreak{}theorem2\allowbreak{}User\allowbreak{}Response\allowbreak{}Table}
\end{itemize}
\hypertarget{governing-declaration-71b9e9aed968af43}{}
\subsection*{\small\ttfamily\raggedright Golz\allowbreak{}Haghtalab\allowbreak{}Yang2025\allowbreak{}Distortion.\allowbreak{}theorem7\allowbreak{}Iid\allowbreak{}User\allowbreak{}Empirical\allowbreak{}Policy}
\textbf{Declaration location:} paper-local
\paragraph{Lean declaration body}
\begin{ReviewVerbatim}
type:
{Alternative : Type u_1} →
  [Fintype Alternative] →
    [DecidableEq Alternative] →
      {users comparisonsPerUser : ℕ} →
        PMF Alternative →
          (klBudget : ℝ) →
            0 ≤ klBudget →
              (Fin users →
                  GolzHaghtalabYang2025Distortion.theorem2UserResponseTable Alternative ×
                    (Fin comparisonsPerUser → Alternative × Alternative)) →
                PMF Alternative

value:
fun {Alternative : Type u_1} [Fintype Alternative] [DecidableEq Alternative] {users comparisonsPerUser : ℕ}
    (reference : PMF Alternative) (klBudget : ℝ) (hbudget : 0 ≤ klBudget)
    (sample :
      Fin users →
        GolzHaghtalabYang2025Distortion.theorem2UserResponseTable Alternative ×
          (Fin comparisonsPerUser → Alternative × Alternative)) =>
  Classical.choose
    (GolzHaghtalabYang2025Distortion.theorem7IidUserEmpiricalPolicy._proof_1 reference klBudget hbudget sample)
\end{ReviewVerbatim}

\paragraph{Direct retained dependencies}
\begin{itemize}\raggedright
\item \hyperlink{governing-declaration-d8827c2845a356ad}{Applied\allowbreak{}Modeling\allowbreak{}Lib.\allowbreak{}Game\allowbreak{}Theory.\allowbreak{}Preference\allowbreak{}Game.\allowbreak{}Is\allowbreak{}Finite\allowbreak{}KL\allowbreak{}Preference\allowbreak{}Maximin}
\item \hyperlink{governing-declaration-a1441da0ef6b2fa2}{Golz\allowbreak{}Haghtalab\allowbreak{}Yang2025\allowbreak{}Distortion.\allowbreak{}theorem10\allowbreak{}Iid\allowbreak{}User\allowbreak{}Paper\allowbreak{}Win\allowbreak{}Rate}
\item \hyperlink{governing-declaration-c5d0273cbd7ddcc9}{Golz\allowbreak{}Haghtalab\allowbreak{}Yang2025\allowbreak{}Distortion.\allowbreak{}theorem2\allowbreak{}User\allowbreak{}Response\allowbreak{}Table}
\item \hyperlink{governing-declaration-2ccaa1524be736e8}{Golz\allowbreak{}Haghtalab\allowbreak{}Yang2025\allowbreak{}Distortion.\allowbreak{}theorem7\allowbreak{}Empirical\allowbreak{}Preference}
\end{itemize}
\hypertarget{governing-declaration-019481f0e8de9cf6}{}
\subsection*{\small\ttfamily\raggedright Golz\allowbreak{}Haghtalab\allowbreak{}Yang2025\allowbreak{}Distortion.\allowbreak{}theorem9\allowbreak{}Off\allowbreak{}Path\allowbreak{}Pairs}
\textbf{Declaration location:} paper-local
\paragraph{Lean declaration body}
\begin{ReviewVerbatim}
type:
(m : ℕ) → Finset (Fin m × Fin m)

value:
fun (m : ℕ) => {pair : Fin m × Fin m | pair.1 ≠ pair.2 ∧ ↑pair.1 + 1 ≠ ↑pair.2 ∧ ↑pair.2 + 1 ≠ ↑pair.1}
\end{ReviewVerbatim}


\hypertarget{governing-declaration-a7ee8a646050a6d0}{}
\subsection*{\small\ttfamily\raggedright Golz\allowbreak{}Haghtalab\allowbreak{}Yang2025\allowbreak{}Distortion.\allowbreak{}theorem9\allowbreak{}Off\allowbreak{}Path\allowbreak{}Pair}
\textbf{Declaration location:} paper-local
\paragraph{Lean declaration body}
\begin{ReviewVerbatim}
type:
ℕ → Type

value:
fun (m : ℕ) => ↥(GolzHaghtalabYang2025Distortion.theorem9OffPathPairs m)
\end{ReviewVerbatim}

\paragraph{Direct retained dependencies}
\begin{itemize}\raggedright
\item \hyperlink{governing-declaration-019481f0e8de9cf6}{Golz\allowbreak{}Haghtalab\allowbreak{}Yang2025\allowbreak{}Distortion.\allowbreak{}theorem9\allowbreak{}Off\allowbreak{}Path\allowbreak{}Pairs}
\end{itemize}
\hypertarget{governing-declaration-c1c27046e580b1b9}{}
\subsection*{\small\ttfamily\raggedright Golz\allowbreak{}Haghtalab\allowbreak{}Yang2025\allowbreak{}Distortion.\allowbreak{}theorem9\allowbreak{}Off\allowbreak{}Path\allowbreak{}Sampling}
\textbf{Declaration location:} paper-local
\paragraph{Lean declaration body}
\begin{ReviewVerbatim}
type:
(m : ℕ) → 3 ≤ m → AppliedModelingLib.Learning.HumanFeedback.CorrelatedPairSampling (Fin m)

value:
fun (m : ℕ) (hm : 3 ≤ m) =>
  PMF.map Subtype.val (AppliedModelingLib.uniformPMF (GolzHaghtalabYang2025Distortion.theorem9OffPathPair m))
\end{ReviewVerbatim}

\paragraph{Direct retained dependencies}
\begin{itemize}\raggedright
\item \hyperlink{governing-declaration-588529a2add6b6a9}{Applied\allowbreak{}Modeling\allowbreak{}Lib.\allowbreak{}Learning.\allowbreak{}Human\allowbreak{}Feedback.\allowbreak{}Correlated\allowbreak{}Pair\allowbreak{}Sampling}
\item \hyperlink{governing-declaration-48f10f2715c04df3}{Applied\allowbreak{}Modeling\allowbreak{}Lib.\allowbreak{}uniform\allowbreak{}PMF}
\item \hyperlink{governing-declaration-a7ee8a646050a6d0}{Golz\allowbreak{}Haghtalab\allowbreak{}Yang2025\allowbreak{}Distortion.\allowbreak{}theorem9\allowbreak{}Off\allowbreak{}Path\allowbreak{}Pair}
\item \hyperlink{governing-declaration-019481f0e8de9cf6}{Golz\allowbreak{}Haghtalab\allowbreak{}Yang2025\allowbreak{}Distortion.\allowbreak{}theorem9\allowbreak{}Off\allowbreak{}Path\allowbreak{}Pairs}
\end{itemize}
\hypertarget{governing-declaration-c4719ba77c259e9a}{}
\subsection*{\small\ttfamily\raggedright Golz\allowbreak{}Haghtalab\allowbreak{}Yang2025\allowbreak{}Distortion.\allowbreak{}theorem9\allowbreak{}Path\allowbreak{}Orientation}
\textbf{Declaration location:} paper-local
\paragraph{Lean declaration body}
\begin{ReviewVerbatim}
type:
{m : ℕ} → 2 ≤ m → Fin (m - 1) × Bool → Fin m × Fin m

value:
fun {m : ℕ} (hm : 2 ≤ m) (edge : Fin (m - 1) × Bool) =>
  have lower : Fin m := ⟨↑edge.1, GolzHaghtalabYang2025Distortion.theorem9PathOrientation._proof_3 hm edge⟩;
  have upper : Fin m := ⟨↑edge.1 + 1, GolzHaghtalabYang2025Distortion.theorem9PathOrientation._proof_4 hm edge⟩;
  if edge.2 = true then (lower, upper) else (upper, lower)
\end{ReviewVerbatim}


\hypertarget{governing-declaration-5f44c8b3d4c13c80}{}
\subsection*{\small\ttfamily\raggedright Golz\allowbreak{}Haghtalab\allowbreak{}Yang2025\allowbreak{}Distortion.\allowbreak{}theorem9\allowbreak{}Path\allowbreak{}Sampling}
\textbf{Declaration location:} paper-local
\paragraph{Lean declaration body}
\begin{ReviewVerbatim}
type:
(m : ℕ) → 2 ≤ m → AppliedModelingLib.Learning.HumanFeedback.CorrelatedPairSampling (Fin m)

value:
fun (m : ℕ) (hm : 2 ≤ m) =>
  PMF.map (GolzHaghtalabYang2025Distortion.theorem9PathOrientation hm)
    (AppliedModelingLib.uniformPMF (Fin (m - 1) × Bool))
\end{ReviewVerbatim}

\paragraph{Direct retained dependencies}
\begin{itemize}\raggedright
\item \hyperlink{governing-declaration-588529a2add6b6a9}{Applied\allowbreak{}Modeling\allowbreak{}Lib.\allowbreak{}Learning.\allowbreak{}Human\allowbreak{}Feedback.\allowbreak{}Correlated\allowbreak{}Pair\allowbreak{}Sampling}
\item \hyperlink{governing-declaration-48f10f2715c04df3}{Applied\allowbreak{}Modeling\allowbreak{}Lib.\allowbreak{}uniform\allowbreak{}PMF}
\item \hyperlink{governing-declaration-c4719ba77c259e9a}{Golz\allowbreak{}Haghtalab\allowbreak{}Yang2025\allowbreak{}Distortion.\allowbreak{}theorem9\allowbreak{}Path\allowbreak{}Orientation}
\end{itemize}
\hypertarget{governing-declaration-9f5c608cbac97312}{}
\subsection*{\small\ttfamily\raggedright Golz\allowbreak{}Haghtalab\allowbreak{}Yang2025\allowbreak{}Distortion.\allowbreak{}theorem9\allowbreak{}Comparison\allowbreak{}Sampling}
\textbf{Declaration location:} paper-local
\paragraph{Lean declaration body}
\begin{ReviewVerbatim}
type:
(m : ℕ) → 3 ≤ m → NNReal → AppliedModelingLib.Learning.HumanFeedback.CorrelatedPairSampling (Fin m)

value:
fun (m : ℕ) (hm : 3 ≤ m) (epsilon : NNReal) =>
  AppliedModelingLib.binaryMixturePMF (1 - epsilon)
    (GolzHaghtalabYang2025Distortion.theorem9ComparisonSampling._proof_1 epsilon)
    (GolzHaghtalabYang2025Distortion.theorem9PathSampling m
      (GolzHaghtalabYang2025Distortion.theorem9ComparisonSampling._proof_2 m hm))
    (GolzHaghtalabYang2025Distortion.theorem9OffPathSampling m hm)
\end{ReviewVerbatim}

\paragraph{Direct retained dependencies}
\begin{itemize}\raggedright
\item \hyperlink{governing-declaration-588529a2add6b6a9}{Applied\allowbreak{}Modeling\allowbreak{}Lib.\allowbreak{}Learning.\allowbreak{}Human\allowbreak{}Feedback.\allowbreak{}Correlated\allowbreak{}Pair\allowbreak{}Sampling}
\item \hyperlink{governing-declaration-8fd2dabd8bfaa1ab}{Applied\allowbreak{}Modeling\allowbreak{}Lib.\allowbreak{}binary\allowbreak{}Mixture\allowbreak{}PMF}
\item \hyperlink{governing-declaration-c1c27046e580b1b9}{Golz\allowbreak{}Haghtalab\allowbreak{}Yang2025\allowbreak{}Distortion.\allowbreak{}theorem9\allowbreak{}Off\allowbreak{}Path\allowbreak{}Sampling}
\item \hyperlink{governing-declaration-5f44c8b3d4c13c80}{Golz\allowbreak{}Haghtalab\allowbreak{}Yang2025\allowbreak{}Distortion.\allowbreak{}theorem9\allowbreak{}Path\allowbreak{}Sampling}
\end{itemize}
\hypertarget{governing-declaration-a2115bb0ddee2a78}{}
\subsection*{\small\ttfamily\raggedright Golz\allowbreak{}Haghtalab\allowbreak{}Yang2025\allowbreak{}Distortion.\allowbreak{}theorem9\allowbreak{}Prefix\allowbreak{}Profile}
\textbf{Declaration location:} paper-local
\paragraph{Lean declaration body}
\begin{ReviewVerbatim}
type:
ℝ → (m : ℕ) → Fin 3 → Fin m → ℝ

value:
fun (beta : ℝ) (m : ℕ) (a : Fin 3) (a_1 : Fin m) => GolzHaghtalabYang2025Distortion.lemma15Utility beta (↑a_1) a
\end{ReviewVerbatim}

\paragraph{Direct retained dependencies}
\begin{itemize}\raggedright
\item \hyperlink{governing-declaration-f1a69dba69c50eae}{Golz\allowbreak{}Haghtalab\allowbreak{}Yang2025\allowbreak{}Distortion.\allowbreak{}lemma15\allowbreak{}Utility}
\end{itemize}
\hypertarget{governing-declaration-8c0c1da56a27f976}{}
\subsection*{\small\ttfamily\raggedright Golz\allowbreak{}Haghtalab\allowbreak{}Yang2025\allowbreak{}Distortion.\allowbreak{}theorem9\allowbreak{}Prefix\allowbreak{}Preference}
\textbf{Declaration location:} paper-local
\paragraph{Lean declaration body}
\begin{ReviewVerbatim}
type:
ℝ → (m : ℕ) → AppliedModelingLib.Learning.HumanFeedback.PairwisePreference PUnit.{1} (Fin m)

value:
fun (beta : ℝ) (m : ℕ) =>
  AppliedModelingLib.Alignment.Welfare.populationBradleyTerryPreference
    GolzHaghtalabYang2025Distortion.lemma15Population (GolzHaghtalabYang2025Distortion.theorem9PrefixProfile beta m)
    beta
\end{ReviewVerbatim}

\paragraph{Direct retained dependencies}
\begin{itemize}\raggedright
\item \hyperlink{library-prerequisite-cf06d15d7f8b7c31}{Applied\allowbreak{}Modeling\allowbreak{}Lib.\allowbreak{}Alignment.\allowbreak{}Welfare.\allowbreak{}population\allowbreak{}Bradley\allowbreak{}Terry\allowbreak{}Preference}
\item \hyperlink{library-prerequisite-7e2db20cf087d9ba}{Applied\allowbreak{}Modeling\allowbreak{}Lib.\allowbreak{}Learning.\allowbreak{}Human\allowbreak{}Feedback.\allowbreak{}Pairwise\allowbreak{}Preference}
\item \hyperlink{governing-declaration-877274851b50572c}{Golz\allowbreak{}Haghtalab\allowbreak{}Yang2025\allowbreak{}Distortion.\allowbreak{}lemma15\allowbreak{}Population}
\item \hyperlink{governing-declaration-a2115bb0ddee2a78}{Golz\allowbreak{}Haghtalab\allowbreak{}Yang2025\allowbreak{}Distortion.\allowbreak{}theorem9\allowbreak{}Prefix\allowbreak{}Profile}
\end{itemize}
\hypertarget{governing-declaration-e63d2c356e95b6b5}{}
\subsection*{\small\ttfamily\raggedright Golz\allowbreak{}Haghtalab\allowbreak{}Yang2025\allowbreak{}Distortion.\allowbreak{}theorem9\allowbreak{}One\allowbreak{}Comparison\allowbreak{}Report\allowbreak{}Law}
\textbf{Declaration location:} paper-local
\paragraph{Lean declaration body}
\begin{ReviewVerbatim}
type:
ℝ → (m : ℕ) → 3 ≤ m → NNReal → PMF (AppliedModelingLib.Learning.HumanFeedback.BinaryPairwiseReport (Fin m))

value:
fun (beta : ℝ) (m : ℕ) (hm : 3 ≤ m) (epsilon : NNReal) =>
  AppliedModelingLib.Learning.HumanFeedback.correlatedOneComparisonReportLaw
    (GolzHaghtalabYang2025Distortion.theorem9ComparisonSampling m hm epsilon)
    (GolzHaghtalabYang2025Distortion.theorem9PrefixPreference beta m)
\end{ReviewVerbatim}

\paragraph{Direct retained dependencies}
\begin{itemize}\raggedright
\item \hyperlink{governing-declaration-bc3d6ca8e7e0ac9d}{Applied\allowbreak{}Modeling\allowbreak{}Lib.\allowbreak{}Learning.\allowbreak{}Human\allowbreak{}Feedback.\allowbreak{}Binary\allowbreak{}Pairwise\allowbreak{}Report}
\item \hyperlink{governing-declaration-33b603cd884f09c4}{Applied\allowbreak{}Modeling\allowbreak{}Lib.\allowbreak{}Learning.\allowbreak{}Human\allowbreak{}Feedback.\allowbreak{}correlated\allowbreak{}One\allowbreak{}Comparison\allowbreak{}Report\allowbreak{}Law}
\item \hyperlink{governing-declaration-9f5c608cbac97312}{Golz\allowbreak{}Haghtalab\allowbreak{}Yang2025\allowbreak{}Distortion.\allowbreak{}theorem9\allowbreak{}Comparison\allowbreak{}Sampling}
\item \hyperlink{governing-declaration-8c0c1da56a27f976}{Golz\allowbreak{}Haghtalab\allowbreak{}Yang2025\allowbreak{}Distortion.\allowbreak{}theorem9\allowbreak{}Prefix\allowbreak{}Preference}
\end{itemize}
\hypertarget{governing-declaration-35b3f5ca150bffc2}{}
\subsection*{\small\ttfamily\raggedright Golz\allowbreak{}Haghtalab\allowbreak{}Yang2025\allowbreak{}Distortion.\allowbreak{}theorem9\allowbreak{}Uniform\allowbreak{}Reference}
\textbf{Declaration location:} paper-local
\paragraph{Lean declaration body}
\begin{ReviewVerbatim}
type:
(m : ℕ) → 1 ≤ m → PMF (Fin m)

value:
fun (m : ℕ) (hm : 1 ≤ m) => AppliedModelingLib.uniformPMF (Fin m)
\end{ReviewVerbatim}

\paragraph{Direct retained dependencies}
\begin{itemize}\raggedright
\item \hyperlink{governing-declaration-48f10f2715c04df3}{Applied\allowbreak{}Modeling\allowbreak{}Lib.\allowbreak{}uniform\allowbreak{}PMF}
\end{itemize}
\clearpage
\hypertarget{library-section-e44db2071aa7a1c7}{}
\section*{Material library prerequisites}
\hypertarget{library-prerequisite-56c0734f8877ddbe}{}
\subsection*{\small\ttfamily\raggedright Applied\allowbreak{}Modeling\allowbreak{}Lib.\allowbreak{}Alignment.\allowbreak{}Welfare.\allowbreak{}Finite\allowbreak{}Utility\allowbreak{}Profile}
\textbf{Source locator:} {\footnotesize\raggedright cited publication, lines 174-\allowbreak{}179\par}
\paragraph{Verbatim paper-source connection}
\begin{ReviewVerbatim}
Let A = {1, . . . , m} be a finite set of alternatives. The population of users is described by a
probability distribution D over utility vectors u = (u(1), . . . , u(m)), whose entries 0 ≤ u(x) ≤ 1
indicate a user’s utility for alternative x.3 The objective is to find an alternative x such that its average
utility AvgUtil(x) := Eu∼D [u(x)] across the user population (also known as the utilitarian social
welfare) is as high as possible. We extend this notation to probability distributions π over alternatives
by setting AvgUtil(π) := Ex∼π [AvgUtil(x)].
\end{ReviewVerbatim}



\paragraph{Lean-expanded library semantic target}
\begin{ReviewVerbatim}
type:
Type u_1 → Type u_2 → Type (max u_1 u_2)

value:
fun (User : Type u_1) (Alternative : Type u_2) => User → Alternative → ℝ
\end{ReviewVerbatim}

\paragraph{Exact Lean library declaration}
\begin{ReviewVerbatim}
/-- A user's cardinal utility for each finite alternative. -/
abbrev FiniteUtilityProfile (User Alternative : Type*) := User → Alternative → ℝ
\end{ReviewVerbatim}

\paragraph{Recorded source-to-library screening}
\textbf{Verdict:} \texttt{matches}
\\
{\raggedright
\textbf{Reason:} This represents a cardinal utility vector as user to alternative to real;\allowbreak{} finiteness of alternatives is imposed at source-\allowbreak{}facing use sites.\allowbreak{}
\par}
\reviewmatch{Matches source input}
\noindent\textbf{Reviewer annotation}\par
\reviewerbox
\clearpage
\hypertarget{library-prerequisite-d3573881ae45d4c6}{}
\subsection*{\small\ttfamily\raggedright Applied\allowbreak{}Modeling\allowbreak{}Lib.\allowbreak{}finite\allowbreak{}KL\allowbreak{}Divergence}
\textbf{Source locator:} {\footnotesize\raggedright cited publication, lines 237-\allowbreak{}253\par}
\paragraph{Verbatim paper-source connection}
\begin{ReviewVerbatim}
Alignment Setting. The alignment setting generalizes the social choice setting in two ways: first,
user utilities ui (y | x) may depend on a state x; second, the goal in determining a policy π is not
purely to maximize the reward AvgUtil(π | x), but a trade-off between this reward and the goal of
remaining close to a reference policy πref (· | x) ∈ ∆(A) in terms of the KL divergence DKL (· k πref ).
For theoretical tractability, we focus our analysis on a single state x, which we from here on omit
from the notation. Conceptually, this treatment of alignment on a state-by-state basis corresponds to
an assumption that our policy class is expressive enough so that it can take the optimal distribution of
actions at each state7 and abstracts from the generalization problem of estimating the population’s
preference between a pair of alternatives at the given state x based on preferences in similar states x0 .
Having set aside the dependency between states, we focus on how the regularization with respect
to a reference policy impacts the ability to optimize the average utility of three alignment methods:
RLHF, DPO, and NLHF. In addition to the pairwise comparisons, these methods take in a reference
policy πref ∈ ∆(A) and a KL bound τ ≥ 0, and map these inputs to a policy π in the KL-ball
Bτ (πref ) := {π ∈ ∆(A) | DKL (π k πref ) ≤ τ } around the reference policy. RLHF, DPO, and NLHF
are typically implemented with a KL regularization rather than our KL-constrained formulation,
which we adopt to enable a comparison on equal terms. In Section 5, we show that these perspectives
are equivalent, and that distortion upper bounds carry over to regularized alignment methods.

[Next verbatim source excerpt]

      If τ = 0, the only feasible policy is πref , and the claim clearly holds for λ = ∞. We also assume that
πref (a) > 0 for all a ∈ M . Otherwise if πref (a) = 0 for some a ∈ M , then both the regularized and constrained
versions forbid any policy to put nonzero probability on a, which leads to an effectively smaller candidate set.
\end{ReviewVerbatim}



\paragraph{Lean-expanded library semantic target}
\begin{ReviewVerbatim}
type:
{α : Type u_1} → [Fintype α] → [DecidableEq α] → PMF α → PMF α → ℝ

value:
fun {α : Type u_1} [Fintype α] [DecidableEq α] (first second : PMF α) =>
  ∑ outcome : α, (first outcome).toReal * (Real.log (first outcome).toReal - Real.log (second outcome).toReal)
\end{ReviewVerbatim}

\paragraph{Exact Lean library declaration}
\begin{ReviewVerbatim}
/-- The finite real-valued KL expression between two probability mass functions. -/
noncomputable def finiteKLDivergence {α : Type*} [Fintype α] [DecidableEq α]
    (first second : PMF α) : ℝ :=
  ∑ outcome : α,
    (first outcome).toReal *
      (Real.log (first outcome).toReal - Real.log (second outcome).toReal)
\end{ReviewVerbatim}

\paragraph{Recorded source-to-library screening}
\textbf{Verdict:} \texttt{matches}
\\
{\raggedright
\textbf{Reason:} This is the finite-\allowbreak{}sum KL expression;\allowbreak{} bundled full support makes it the ordinary finite KL divergence and zero first-\allowbreak{}policy mass contributes zero.\allowbreak{}
\par}
\reviewmatch{Matches source input}
\noindent\textbf{Reviewer annotation}\par
\reviewerbox
\clearpage
\hypertarget{library-prerequisite-1d625de57bf58a6a}{}
\subsection*{\small\ttfamily\raggedright Applied\allowbreak{}Modeling\allowbreak{}Lib.\allowbreak{}Alignment.\allowbreak{}Welfare.\allowbreak{}In\allowbreak{}Finite\allowbreak{}KL\allowbreak{}Ball}
\textbf{Source locator:} {\footnotesize\raggedright cited publication, lines 237-\allowbreak{}253\par}
\paragraph{Verbatim paper-source connection}
\begin{ReviewVerbatim}
Alignment Setting. The alignment setting generalizes the social choice setting in two ways: first,
user utilities ui (y | x) may depend on a state x; second, the goal in determining a policy π is not
purely to maximize the reward AvgUtil(π | x), but a trade-off between this reward and the goal of
remaining close to a reference policy πref (· | x) ∈ ∆(A) in terms of the KL divergence DKL (· k πref ).
For theoretical tractability, we focus our analysis on a single state x, which we from here on omit
from the notation. Conceptually, this treatment of alignment on a state-by-state basis corresponds to
an assumption that our policy class is expressive enough so that it can take the optimal distribution of
actions at each state7 and abstracts from the generalization problem of estimating the population’s
preference between a pair of alternatives at the given state x based on preferences in similar states x0 .
Having set aside the dependency between states, we focus on how the regularization with respect
to a reference policy impacts the ability to optimize the average utility of three alignment methods:
RLHF, DPO, and NLHF. In addition to the pairwise comparisons, these methods take in a reference
policy πref ∈ ∆(A) and a KL bound τ ≥ 0, and map these inputs to a policy π in the KL-ball
Bτ (πref ) := {π ∈ ∆(A) | DKL (π k πref ) ≤ τ } around the reference policy. RLHF, DPO, and NLHF
are typically implemented with a KL regularization rather than our KL-constrained formulation,
which we adopt to enable a comparison on equal terms. In Section 5, we show that these perspectives
are equivalent, and that distortion upper bounds carry over to regularized alignment methods.

[Next verbatim source excerpt]

      If τ = 0, the only feasible policy is πref , and the claim clearly holds for λ = ∞. We also assume that
πref (a) > 0 for all a ∈ M . Otherwise if πref (a) = 0 for some a ∈ M , then both the regularized and constrained
versions forbid any policy to put nonzero probability on a, which leads to an effectively smaller candidate set.
\end{ReviewVerbatim}



\paragraph{Lean-expanded library semantic target}
\begin{ReviewVerbatim}
type:
{Alternative : Type u_1} →
  [Fintype Alternative] → [DecidableEq Alternative] → PMF Alternative → ℝ → PMF Alternative → Prop

value:
fun {Alternative : Type u_1} [Fintype Alternative] [DecidableEq Alternative] (reference : PMF Alternative)
    (klBudget : ℝ) (policy : PMF Alternative) =>
  AppliedModelingLib.finiteKLDivergence policy reference ≤ klBudget
\end{ReviewVerbatim}

\paragraph{Exact Lean library declaration}
\begin{ReviewVerbatim}
/-- A finite policy belongs to the source KL ball around `reference`. -/
def InFiniteKLBall {Alternative : Type*} [Fintype Alternative] [DecidableEq Alternative]
    (reference : PMF Alternative) (klBudget : ℝ) (policy : PMF Alternative) : Prop :=
  finiteKLDivergence policy reference ≤ klBudget
\end{ReviewVerbatim}

\paragraph{Direct retained dependencies}
\begin{itemize}\raggedright
\item \hyperlink{library-prerequisite-d3573881ae45d4c6}{Applied\allowbreak{}Modeling\allowbreak{}Lib.\allowbreak{}finite\allowbreak{}KL\allowbreak{}Divergence}
\end{itemize}
\paragraph{Recorded source-to-library screening}
\textbf{Verdict:} \texttt{matches}
\\
{\raggedright
\textbf{Reason:} This exactly expresses the source normalization 0 <= u(x) <= 1 for every user and alternative.\allowbreak{}
\par}
\reviewmatch{Matches source input}
\noindent\textbf{Reviewer annotation}\par
\reviewerbox
\clearpage
\hypertarget{library-prerequisite-e35e7dd3bb4adb9e}{}
\subsection*{\small\ttfamily\raggedright Applied\allowbreak{}Modeling\allowbreak{}Lib.\allowbreak{}Alignment.\allowbreak{}Welfare.\allowbreak{}Reference\allowbreak{}Alternative\allowbreak{}Policy}
\textbf{Source locator:} {\footnotesize\raggedright cited publication, lines 237-\allowbreak{}253\par}
\paragraph{Verbatim paper-source connection}
\begin{ReviewVerbatim}
Alignment Setting. The alignment setting generalizes the social choice setting in two ways: first,
user utilities ui (y | x) may depend on a state x; second, the goal in determining a policy π is not
purely to maximize the reward AvgUtil(π | x), but a trade-off between this reward and the goal of
remaining close to a reference policy πref (· | x) ∈ ∆(A) in terms of the KL divergence DKL (· k πref ).
For theoretical tractability, we focus our analysis on a single state x, which we from here on omit
from the notation. Conceptually, this treatment of alignment on a state-by-state basis corresponds to
an assumption that our policy class is expressive enough so that it can take the optimal distribution of
actions at each state7 and abstracts from the generalization problem of estimating the population’s
preference between a pair of alternatives at the given state x based on preferences in similar states x0 .
Having set aside the dependency between states, we focus on how the regularization with respect
to a reference policy impacts the ability to optimize the average utility of three alignment methods:
RLHF, DPO, and NLHF. In addition to the pairwise comparisons, these methods take in a reference
policy πref ∈ ∆(A) and a KL bound τ ≥ 0, and map these inputs to a policy π in the KL-ball
Bτ (πref ) := {π ∈ ∆(A) | DKL (π k πref ) ≤ τ } around the reference policy. RLHF, DPO, and NLHF
are typically implemented with a KL regularization rather than our KL-constrained formulation,
which we adopt to enable a comparison on equal terms. In Section 5, we show that these perspectives
are equivalent, and that distortion upper bounds carry over to regularized alignment methods.

[Next verbatim source excerpt]

      If τ = 0, the only feasible policy is πref , and the claim clearly holds for λ = ∞. We also assume that
πref (a) > 0 for all a ∈ M . Otherwise if πref (a) = 0 for some a ∈ M , then both the regularized and constrained
versions forbid any policy to put nonzero probability on a, which leads to an effectively smaller candidate set.
\end{ReviewVerbatim}



\paragraph{Lean-expanded library semantic target}
\begin{ReviewVerbatim}
type:
Type u_1 → Type u_1

value:
fun (Alternative : Type u_1) => PMF Alternative
\end{ReviewVerbatim}

\paragraph{Exact Lean library declaration}
\begin{ReviewVerbatim}
/-- A context-free reference policy over alternatives, kept distinct from comparison sampling. -/
abbrev ReferenceAlternativePolicy (Alternative : Type*) := PMF Alternative
\end{ReviewVerbatim}

\paragraph{Recorded source-to-library screening}
\textbf{Verdict:} \texttt{matches}
\\
{\raggedright
\textbf{Reason:} The finite-\allowbreak{}population PMF specialization is precisely expectation of u(x) under the population law.\allowbreak{}
\par}
\reviewmatch{Matches source input}
\noindent\textbf{Reviewer annotation}\par
\reviewerbox
\clearpage
\hypertarget{library-prerequisite-cbe8afa0d6d1ba78}{}
\subsection*{\small\ttfamily\raggedright Applied\allowbreak{}Modeling\allowbreak{}Lib.\allowbreak{}Alignment.\allowbreak{}Welfare.\allowbreak{}Unit\allowbreak{}Interval\allowbreak{}Utility\allowbreak{}Profile}
\textbf{Source locator:} {\footnotesize\raggedright cited publication, lines 174-\allowbreak{}179\par}
\paragraph{Verbatim paper-source connection}
\begin{ReviewVerbatim}
Let A = {1, . . . , m} be a finite set of alternatives. The population of users is described by a
probability distribution D over utility vectors u = (u(1), . . . , u(m)), whose entries 0 ≤ u(x) ≤ 1
indicate a user’s utility for alternative x.3 The objective is to find an alternative x such that its average
utility AvgUtil(x) := Eu∼D [u(x)] across the user population (also known as the utilitarian social
welfare) is as high as possible. We extend this notation to probability distributions π over alternatives
by setting AvgUtil(π) := Ex∼π [AvgUtil(x)].
\end{ReviewVerbatim}



\paragraph{Lean-expanded library semantic target}
\begin{ReviewVerbatim}
type:
{User : Type u_1} →
  {Alternative : Type u_2} → AppliedModelingLib.Alignment.Welfare.FiniteUtilityProfile User Alternative → Prop

value:
fun {User : Type u_1} {Alternative : Type u_2}
    (utility : AppliedModelingLib.Alignment.Welfare.FiniteUtilityProfile User Alternative) =>
  ∀ (user : User) (alternative : Alternative), 0 ≤ utility user alternative ∧ utility user alternative ≤ 1
\end{ReviewVerbatim}

\paragraph{Exact Lean library declaration}
\begin{ReviewVerbatim}
/-- The source distortion model's utility normalization to the unit interval. -/
def UnitIntervalUtilityProfile {User Alternative : Type*}
    (utility : FiniteUtilityProfile User Alternative) : Prop :=
  ∀ user alternative, 0 ≤ utility user alternative ∧ utility user alternative ≤ 1
\end{ReviewVerbatim}

\paragraph{Direct retained dependencies}
\begin{itemize}\raggedright
\item \hyperlink{library-prerequisite-56c0734f8877ddbe}{Applied\allowbreak{}Modeling\allowbreak{}Lib.\allowbreak{}Alignment.\allowbreak{}Welfare.\allowbreak{}Finite\allowbreak{}Utility\allowbreak{}Profile}
\end{itemize}
\paragraph{Recorded source-to-library screening}
\textbf{Verdict:} \texttt{matches}
\\
{\raggedright
\textbf{Reason:} This exactly nests the alternative-\allowbreak{}policy expectation over population average utilities.\allowbreak{}
\par}
\reviewmatch{Matches source input}
\noindent\textbf{Reviewer annotation}\par
\reviewerbox
\clearpage
\hypertarget{library-prerequisite-15f4b442749deab2}{}
\subsection*{\small\ttfamily\raggedright Applied\allowbreak{}Modeling\allowbreak{}Lib.\allowbreak{}Alignment.\allowbreak{}Welfare.\allowbreak{}population\allowbreak{}Average\allowbreak{}Utility}
\textbf{Source locator:} {\footnotesize\raggedright cited publication, lines 174-\allowbreak{}179\par}
\paragraph{Verbatim paper-source connection}
\begin{ReviewVerbatim}
Let A = {1, . . . , m} be a finite set of alternatives. The population of users is described by a
probability distribution D over utility vectors u = (u(1), . . . , u(m)), whose entries 0 ≤ u(x) ≤ 1
indicate a user’s utility for alternative x.3 The objective is to find an alternative x such that its average
utility AvgUtil(x) := Eu∼D [u(x)] across the user population (also known as the utilitarian social
welfare) is as high as possible. We extend this notation to probability distributions π over alternatives
by setting AvgUtil(π) := Ex∼π [AvgUtil(x)].
\end{ReviewVerbatim}



\paragraph{Lean-expanded library semantic target}
\begin{ReviewVerbatim}
type:
{User : Type u_1} →
  {Alternative : Type u_2} →
    [Fintype User] →
      [DecidableEq User] →
        PMF User → AppliedModelingLib.Alignment.Welfare.FiniteUtilityProfile User Alternative → Alternative → ℝ

value:
fun {User : Type u_1} {Alternative : Type u_2} [Fintype User] [DecidableEq User] (population : PMF User)
    (utility : AppliedModelingLib.Alignment.Welfare.FiniteUtilityProfile User Alternative)
    (alternative : Alternative) =>
  AppliedModelingLib.pmfExp population fun (user : User) => utility user alternative
\end{ReviewVerbatim}

\paragraph{Exact Lean library declaration}
\begin{ReviewVerbatim}
/-- The utilitarian average utility of one alternative over a user population. -/
noncomputable def populationAverageUtility
    {User Alternative : Type*} [Fintype User] [DecidableEq User]
    (population : PMF User) (utility : FiniteUtilityProfile User Alternative)
    (alternative : Alternative) : ℝ :=
  pmfExp population (fun user => utility user alternative)
\end{ReviewVerbatim}

\paragraph{Direct retained dependencies}
\begin{itemize}\raggedright
\item \hyperlink{library-prerequisite-56c0734f8877ddbe}{Applied\allowbreak{}Modeling\allowbreak{}Lib.\allowbreak{}Alignment.\allowbreak{}Welfare.\allowbreak{}Finite\allowbreak{}Utility\allowbreak{}Profile}
\item \hyperlink{governing-declaration-8e4b4389a4f6454b}{Applied\allowbreak{}Modeling\allowbreak{}Lib.\allowbreak{}pmf\allowbreak{}Exp}
\end{itemize}
\paragraph{Recorded source-to-library screening}
\textbf{Verdict:} \texttt{matches}
\\
{\raggedright
\textbf{Reason:} This is the appropriate reusable interface for a valid pairwise marginal probability, with range and complementarity;\allowbreak{} the preceding constructor supplies the Bradley-\allowbreak{}Terry formula.\allowbreak{}
\par}
\reviewmatch{Matches source input}
\noindent\textbf{Reviewer annotation}\par
\reviewerbox
\clearpage
\hypertarget{library-prerequisite-184c74a4db4311c2}{}
\subsection*{\small\ttfamily\raggedright Applied\allowbreak{}Modeling\allowbreak{}Lib.\allowbreak{}Alignment.\allowbreak{}Welfare.\allowbreak{}policy\allowbreak{}Average\allowbreak{}Utility}
\textbf{Source locator:} {\footnotesize\raggedright cited publication, lines 174-\allowbreak{}179\par}
\paragraph{Verbatim paper-source connection}
\begin{ReviewVerbatim}
Let A = {1, . . . , m} be a finite set of alternatives. The population of users is described by a
probability distribution D over utility vectors u = (u(1), . . . , u(m)), whose entries 0 ≤ u(x) ≤ 1
indicate a user’s utility for alternative x.3 The objective is to find an alternative x such that its average
utility AvgUtil(x) := Eu∼D [u(x)] across the user population (also known as the utilitarian social
welfare) is as high as possible. We extend this notation to probability distributions π over alternatives
by setting AvgUtil(π) := Ex∼π [AvgUtil(x)].
\end{ReviewVerbatim}



\paragraph{Lean-expanded library semantic target}
\begin{ReviewVerbatim}
type:
{User : Type u_1} →
  {Alternative : Type u_2} →
    [Fintype User] →
      [DecidableEq User] →
        [Fintype Alternative] →
          [DecidableEq Alternative] →
            PMF User → AppliedModelingLib.Alignment.Welfare.FiniteUtilityProfile User Alternative → PMF Alternative → ℝ

value:
fun {User : Type u_1} {Alternative : Type u_2} [Fintype User] [DecidableEq User] [Fintype Alternative]
    [DecidableEq Alternative] (population : PMF User)
    (utility : AppliedModelingLib.Alignment.Welfare.FiniteUtilityProfile User Alternative) (policy : PMF Alternative) =>
  AppliedModelingLib.pmfExp policy (AppliedModelingLib.Alignment.Welfare.populationAverageUtility population utility)
\end{ReviewVerbatim}

\paragraph{Exact Lean library declaration}
\begin{ReviewVerbatim}
/-- The average utility of a randomized alternative policy. -/
noncomputable def policyAverageUtility
    {User Alternative : Type*} [Fintype User] [DecidableEq User]
    [Fintype Alternative] [DecidableEq Alternative]
    (population : PMF User) (utility : FiniteUtilityProfile User Alternative)
    (policy : PMF Alternative) : ℝ :=
  pmfExp policy (populationAverageUtility population utility)
\end{ReviewVerbatim}

\paragraph{Direct retained dependencies}
\begin{itemize}\raggedright
\item \hyperlink{library-prerequisite-56c0734f8877ddbe}{Applied\allowbreak{}Modeling\allowbreak{}Lib.\allowbreak{}Alignment.\allowbreak{}Welfare.\allowbreak{}Finite\allowbreak{}Utility\allowbreak{}Profile}
\item \hyperlink{library-prerequisite-15f4b442749deab2}{Applied\allowbreak{}Modeling\allowbreak{}Lib.\allowbreak{}Alignment.\allowbreak{}Welfare.\allowbreak{}population\allowbreak{}Average\allowbreak{}Utility}
\item \hyperlink{governing-declaration-8e4b4389a4f6454b}{Applied\allowbreak{}Modeling\allowbreak{}Lib.\allowbreak{}pmf\allowbreak{}Exp}
\end{itemize}
\paragraph{Recorded source-to-library screening}
\textbf{Verdict:} \texttt{matches}
\\
{\raggedright
\textbf{Reason:} This computes the population expectation of the Bradley-\allowbreak{}Terry sigmoid of beta times the utility gap;\allowbreak{} the source-\allowbreak{}facing theorems retain positive beta.\allowbreak{}
\par}
\reviewmatch{Matches source input}
\noindent\textbf{Reviewer annotation}\par
\reviewerbox
\clearpage
\hypertarget{library-prerequisite-7e2db20cf087d9ba}{}
\subsection*{\small\ttfamily\raggedright Applied\allowbreak{}Modeling\allowbreak{}Lib.\allowbreak{}Learning.\allowbreak{}Human\allowbreak{}Feedback.\allowbreak{}Pairwise\allowbreak{}Preference}
\textbf{Source locator:} {\footnotesize\raggedright cited publication, lines 180-\allowbreak{}200\par}
\paragraph{Verbatim paper-source connection}
\begin{ReviewVerbatim}
Both voting rules and alignment methods observe comparisons from n users. We model each
i = 1, . . . , n as a fresh user with independently drawn utility vector ui ∼ D. For exposition, we
assume that each user i provides an equal number d ≥ 1 of pairwise comparisons. For each i,
and for j = 1, . . . , d, we independently draw alternatives xji , yij from a fixed distribution µ over

    3
     Our assumption that utilities be in [0, 1] is weaker than any of the three assumptions — unit-sum, unit-range,
or approval [23] — made in classic distortion to avoid trivial infinite lower bounds.


                                                           3

[form-feed in source extraction]
alternatives, in which the minimum probability mass µmin := minx∈A µ(x) is positive.4 User i
then compares each pair {xji , yij } (for j = 1, . . . , d) through a Bradley-Terry model based on i’s
utilities: i prefers xji over yij (written “xji [U+001F control character]i yij ”) with probability σ β · (ui (xji ) − ui (yij )) , where
                                                                                                        [U+0001 control character]

σ(t) := 1/(1+e−t ) is the logistic sigmoid function and β > 0 is a temperature parameter, and prefers
yij over xji (“yij [U+001F control character]i xji ”) otherwise.5 Whereas this specifies the marginal probability of each pairwise
comparison, we make no assumption about the correlation between i’s choices. For example, i might
derive the pairwise comparisons from a Plackett-Luce ranking, ensuring that      [U+0002 control character] the user’s comparisons  [U+0001 control character][U+0003 control character]
are always consistent.6 We set p(x [U+001F control character] y) for the expected win rate Eu∼D σ β · (u(x) − u(y)) .
\end{ReviewVerbatim}



\paragraph{Lean declaration type (metadata)}
\begin{ReviewVerbatim}
field prob:
{Context : Type u_1} →
  {Response : Type u_2} →
    AppliedModelingLib.Learning.HumanFeedback.PairwisePreference Context Response → Context → Response → Response → ℝ

field nonneg:
∀ {Context : Type u_1} {Response : Type u_2}
  (self : AppliedModelingLib.Learning.HumanFeedback.PairwisePreference Context Response) (context : Context)
  (first second : Response), 0 ≤ self.prob context first second

field le_one:
∀ {Context : Type u_1} {Response : Type u_2}
  (self : AppliedModelingLib.Learning.HumanFeedback.PairwisePreference Context Response) (context : Context)
  (first second : Response), self.prob context first second ≤ 1

field complementary:
∀ {Context : Type u_1} {Response : Type u_2}
  (self : AppliedModelingLib.Learning.HumanFeedback.PairwisePreference Context Response) (context : Context)
  (first second : Response), self.prob context first second + self.prob context second first = 1
\end{ReviewVerbatim}

\paragraph{Exact Lean library declaration}
\begin{ReviewVerbatim}
/-- A valid probability that one response is preferred to another in a context. -/
structure PairwisePreference (Context Response : Type*) where
  /-- Probability that the first response is preferred to the second. -/
  prob : Context → Response → Response → ℝ
  nonneg : ∀ context first second, 0 ≤ prob context first second
  le_one : ∀ context first second, prob context first second ≤ 1
  complementary : ∀ context first second,
    prob context first second + prob context second first = 1
\end{ReviewVerbatim}

\paragraph{Recorded source-to-library screening}
\textbf{Verdict:} \texttt{matches}
\\
{\raggedright
\textbf{Reason:} This is exactly the context-\allowbreak{}free reference policy in the finite alternative simplex, distinct from comparison sampling as in the source.\allowbreak{}
\par}
\reviewmatch{Matches source input}
\noindent\textbf{Reviewer annotation}\par
\reviewerbox
\clearpage
\hypertarget{library-prerequisite-cf06d15d7f8b7c31}{}
\subsection*{\small\ttfamily\raggedright Applied\allowbreak{}Modeling\allowbreak{}Lib.\allowbreak{}Alignment.\allowbreak{}Welfare.\allowbreak{}population\allowbreak{}Bradley\allowbreak{}Terry\allowbreak{}Preference}
\textbf{Source locator:} {\footnotesize\raggedright cited publication, lines 180-\allowbreak{}200\par}
\paragraph{Verbatim paper-source connection}
\begin{ReviewVerbatim}
Both voting rules and alignment methods observe comparisons from n users. We model each
i = 1, . . . , n as a fresh user with independently drawn utility vector ui ∼ D. For exposition, we
assume that each user i provides an equal number d ≥ 1 of pairwise comparisons. For each i,
and for j = 1, . . . , d, we independently draw alternatives xji , yij from a fixed distribution µ over

    3
     Our assumption that utilities be in [0, 1] is weaker than any of the three assumptions — unit-sum, unit-range,
or approval [23] — made in classic distortion to avoid trivial infinite lower bounds.


                                                           3

[form-feed in source extraction]
alternatives, in which the minimum probability mass µmin := minx∈A µ(x) is positive.4 User i
then compares each pair {xji , yij } (for j = 1, . . . , d) through a Bradley-Terry model based on i’s
utilities: i prefers xji over yij (written “xji [U+001F control character]i yij ”) with probability σ β · (ui (xji ) − ui (yij )) , where
                                                                                                        [U+0001 control character]

σ(t) := 1/(1+e−t ) is the logistic sigmoid function and β > 0 is a temperature parameter, and prefers
yij over xji (“yij [U+001F control character]i xji ”) otherwise.5 Whereas this specifies the marginal probability of each pairwise
comparison, we make no assumption about the correlation between i’s choices. For example, i might
derive the pairwise comparisons from a Plackett-Luce ranking, ensuring that      [U+0002 control character] the user’s comparisons  [U+0001 control character][U+0003 control character]
are always consistent.6 We set p(x [U+001F control character] y) for the expected win rate Eu∼D σ β · (u(x) − u(y)) .
\end{ReviewVerbatim}



\paragraph{Lean-expanded library semantic target}
\begin{ReviewVerbatim}
type:
{User : Type u_1} →
  {Alternative : Type u_2} →
    [Fintype User] →
      [DecidableEq User] →
        PMF User →
          AppliedModelingLib.Alignment.Welfare.FiniteUtilityProfile User Alternative →
            ℝ → AppliedModelingLib.Learning.HumanFeedback.PairwisePreference PUnit.{u_3 + 1} Alternative

value:
fun {User : Type u_1} {Alternative : Type u_2} [Fintype User] [DecidableEq User] (population : PMF User)
    (utility : AppliedModelingLib.Alignment.Welfare.FiniteUtilityProfile User Alternative) (btScale : ℝ) =>
  {
    prob := fun (x : PUnit.{u_3 + 1}) (first second : Alternative) =>
      AppliedModelingLib.pmfExp population fun (user : User) =>
        (btScale * (utility user first - utility user second)).sigmoid,
    nonneg := AppliedModelingLib.Alignment.Welfare.populationBradleyTerryPreference._proof_1 population utility btScale,
    le_one := AppliedModelingLib.Alignment.Welfare.populationBradleyTerryPreference._proof_2 population utility btScale,
    complementary :=
      AppliedModelingLib.Alignment.Welfare.populationBradleyTerryPreference._proof_3 population utility btScale }
\end{ReviewVerbatim}

\paragraph{Exact Lean library declaration}
\begin{ReviewVerbatim}
/-- Aggregate independent user Bradley--Terry comparisons into one valid preference model. -/
noncomputable def populationBradleyTerryPreference
    {User Alternative : Type*} [Fintype User] [DecidableEq User]
    (population : PMF User) (utility : FiniteUtilityProfile User Alternative)
    (btScale : ℝ) : PairwisePreference PUnit Alternative where
  prob _ first second :=
    pmfExp population fun user =>
      Real.sigmoid (btScale * (utility user first - utility user second))
  nonneg _ first second :=
    pmfExp_nonneg_of_forall_nonneg population _ fun user => Real.sigmoid_nonneg _
  le_one _ first second :=
    pmfExp_le_of_forall_le population _ 1 fun user => Real.sigmoid_le_one _
  complementary := by
    intro _ first second
    rw [← pmfExp_add, ← pmfExp_const population 1]
    refine pmfExp_congr population fun user => ?_
    have hneg : btScale * (utility user second - utility user first) =
        -(btScale * (utility user first - utility user second)) := by ring
    rw [hneg, Real.sigmoid_neg]
    ring
\end{ReviewVerbatim}

\paragraph{Direct retained dependencies}
\begin{itemize}\raggedright
\item \hyperlink{library-prerequisite-56c0734f8877ddbe}{Applied\allowbreak{}Modeling\allowbreak{}Lib.\allowbreak{}Alignment.\allowbreak{}Welfare.\allowbreak{}Finite\allowbreak{}Utility\allowbreak{}Profile}
\item \hyperlink{library-prerequisite-7e2db20cf087d9ba}{Applied\allowbreak{}Modeling\allowbreak{}Lib.\allowbreak{}Learning.\allowbreak{}Human\allowbreak{}Feedback.\allowbreak{}Pairwise\allowbreak{}Preference}
\item \hyperlink{governing-declaration-8e4b4389a4f6454b}{Applied\allowbreak{}Modeling\allowbreak{}Lib.\allowbreak{}pmf\allowbreak{}Exp}
\end{itemize}
\paragraph{Recorded source-to-library screening}
\textbf{Verdict:} \texttt{matches}
\\
{\raggedright
\textbf{Reason:} This is the exact independent-\allowbreak{}product construction used for independently sampled displayed alternatives, instantiated with the source sampling law in both factors.\allowbreak{}
\par}
\reviewmatch{Matches source input}
\noindent\textbf{Reviewer annotation}\par
\reviewerbox
\clearpage
\hypertarget{library-prerequisite-268c959ee7d2d148}{}
\subsection*{\small\ttfamily\raggedright Applied\allowbreak{}Modeling\allowbreak{}Lib.\allowbreak{}PMF\allowbreak{}Full\allowbreak{}Support}
\textbf{Source locator:} {\footnotesize\raggedright cited publication, lines 237-\allowbreak{}253\par}
\paragraph{Verbatim paper-source connection}
\begin{ReviewVerbatim}
Alignment Setting. The alignment setting generalizes the social choice setting in two ways: first,
user utilities ui (y | x) may depend on a state x; second, the goal in determining a policy π is not
purely to maximize the reward AvgUtil(π | x), but a trade-off between this reward and the goal of
remaining close to a reference policy πref (· | x) ∈ ∆(A) in terms of the KL divergence DKL (· k πref ).
For theoretical tractability, we focus our analysis on a single state x, which we from here on omit
from the notation. Conceptually, this treatment of alignment on a state-by-state basis corresponds to
an assumption that our policy class is expressive enough so that it can take the optimal distribution of
actions at each state7 and abstracts from the generalization problem of estimating the population’s
preference between a pair of alternatives at the given state x based on preferences in similar states x0 .
Having set aside the dependency between states, we focus on how the regularization with respect
to a reference policy impacts the ability to optimize the average utility of three alignment methods:
RLHF, DPO, and NLHF. In addition to the pairwise comparisons, these methods take in a reference
policy πref ∈ ∆(A) and a KL bound τ ≥ 0, and map these inputs to a policy π in the KL-ball
Bτ (πref ) := {π ∈ ∆(A) | DKL (π k πref ) ≤ τ } around the reference policy. RLHF, DPO, and NLHF
are typically implemented with a KL regularization rather than our KL-constrained formulation,
which we adopt to enable a comparison on equal terms. In Section 5, we show that these perspectives
are equivalent, and that distortion upper bounds carry over to regularized alignment methods.

[Next verbatim source excerpt]

      If τ = 0, the only feasible policy is πref , and the claim clearly holds for λ = ∞. We also assume that
πref (a) > 0 for all a ∈ M . Otherwise if πref (a) = 0 for some a ∈ M , then both the regularized and constrained
versions forbid any policy to put nonzero probability on a, which leads to an effectively smaller candidate set.
\end{ReviewVerbatim}



\paragraph{Lean-expanded library semantic target}
\begin{ReviewVerbatim}
type:
{α : Type u_1} → PMF α → Prop

value:
fun {α : Type u_1} (law : PMF α) => ∀ (outcome : α), 0 < (law outcome).toReal
\end{ReviewVerbatim}

\paragraph{Exact Lean library declaration}
\begin{ReviewVerbatim}
/-- Every atom of a finite probability mass function has strictly positive real mass. -/
def PMFFullSupport {α : Type*} (law : PMF α) : Prop :=
  ∀ outcome, 0 < (law outcome).toReal
\end{ReviewVerbatim}

\paragraph{Recorded source-to-library screening}
\textbf{Verdict:} \texttt{matches}
\\
{\raggedright
\textbf{Reason:} This is the source clarification's exact positive-\allowbreak{}mass condition for every candidate under the reference policy.\allowbreak{}
\par}
\reviewmatch{Matches source input}
\noindent\textbf{Reviewer annotation}\par
\reviewerbox
\clearpage
\hypertarget{library-prerequisite-de8e2a8bfe221a38}{}
\subsection*{\small\ttfamily\raggedright Applied\allowbreak{}Modeling\allowbreak{}Lib.\allowbreak{}pmf\allowbreak{}Prod}
\textbf{Source locator:} {\footnotesize\raggedright cited publication, lines 180-\allowbreak{}200\par}
\paragraph{Verbatim paper-source connection}
\begin{ReviewVerbatim}
Both voting rules and alignment methods observe comparisons from n users. We model each
i = 1, . . . , n as a fresh user with independently drawn utility vector ui ∼ D. For exposition, we
assume that each user i provides an equal number d ≥ 1 of pairwise comparisons. For each i,
and for j = 1, . . . , d, we independently draw alternatives xji , yij from a fixed distribution µ over

    3
     Our assumption that utilities be in [0, 1] is weaker than any of the three assumptions — unit-sum, unit-range,
or approval [23] — made in classic distortion to avoid trivial infinite lower bounds.


                                                           3

[form-feed in source extraction]
alternatives, in which the minimum probability mass µmin := minx∈A µ(x) is positive.4 User i
then compares each pair {xji , yij } (for j = 1, . . . , d) through a Bradley-Terry model based on i’s
utilities: i prefers xji over yij (written “xji [U+001F control character]i yij ”) with probability σ β · (ui (xji ) − ui (yij )) , where
                                                                                                        [U+0001 control character]

σ(t) := 1/(1+e−t ) is the logistic sigmoid function and β > 0 is a temperature parameter, and prefers
yij over xji (“yij [U+001F control character]i xji ”) otherwise.5 Whereas this specifies the marginal probability of each pairwise
comparison, we make no assumption about the correlation between i’s choices. For example, i might
derive the pairwise comparisons from a Plackett-Luce ranking, ensuring that      [U+0002 control character] the user’s comparisons  [U+0001 control character][U+0003 control character]
are always consistent.6 We set p(x [U+001F control character] y) for the expected win rate Eu∼D σ β · (u(x) − u(y)) .
\end{ReviewVerbatim}



\paragraph{Lean-expanded library semantic target}
\begin{ReviewVerbatim}
type:
{α : Type u_1} → {β : Type u_2} → [Fintype α] → [Fintype β] → PMF α → PMF β → PMF (α × β)

value:
fun {α : Type u_1} {β : Type u_2} [Fintype α] [Fintype β] (μ : PMF α) (ν : PMF β) =>
  PMF.ofFintype (fun (p : α × β) => μ p.1 * ν p.2) (AppliedModelingLib.pmfProd._proof_1 μ ν)
\end{ReviewVerbatim}

\paragraph{Exact Lean library declaration}
\begin{ReviewVerbatim}
/-- Independent product PMF on `α × β`. -/
noncomputable def pmfProd {α β : Type*} [Fintype α] [Fintype β]
    (μ : PMF α) (ν : PMF β) : PMF (α × β) :=
  PMF.ofFintype (fun p : α × β => μ p.1 * ν p.2) (by
    classical
    have hμ : ∑ a : α, μ a = 1 := by
      rw [← PMF.tsum_coe μ, tsum_fintype]
    have hν : ∑ b : β, ν b = 1 := by
      rw [← PMF.tsum_coe ν, tsum_fintype]
    calc
      ∑ p : α × β, μ p.1 * ν p.2
          = ∑ a : α, ∑ b : β, μ a * ν b := by
            rw [Fintype.sum_prod_type]
      _ = ∑ a : α, μ a * (∑ b : β, ν b) := by
            simp [Finset.mul_sum]
      _ = 1 := by
            simp [hμ, hν])
\end{ReviewVerbatim}

\paragraph{Recorded source-to-library screening}
\textbf{Verdict:} \texttt{matches}
\\
{\raggedright
\textbf{Reason:} This exactly defines the KL ball around the reference policy.\allowbreak{} Nonnegative radius and full-\allowbreak{}support reference remain explicit surrounding source-\allowbreak{}model assumptions.\allowbreak{}
\par}
\reviewmatch{Matches source input}
\noindent\textbf{Reviewer annotation}\par
\reviewerbox
\clearpage
\hypertarget{paper-prerequisite-section-5ef5ef0364b6939c}{}
\section*{Paper-specific semantic prerequisites}
\hypertarget{paper-prerequisite-53af1666fc341b0f}{}
\subsection*{\small\ttfamily\raggedright Golz\allowbreak{}Haghtalab\allowbreak{}Yang2025\allowbreak{}Distortion.\allowbreak{}Is\allowbreak{}Empirical\allowbreak{}Maximal\allowbreak{}Lottery}
\noindent\textbf{Paper-local declaration:}\par
{\footnotesize\ttfamily\raggedright Golz\allowbreak{}Haghtalab\allowbreak{}Yang2025\allowbreak{}Distortion.\allowbreak{}Is\allowbreak{}Empirical\allowbreak{}Maximal\allowbreak{}Lottery\par}
\textbf{Source locator:} {\footnotesize\raggedright cited publication, lines 230-\allowbreak{}236\par}
\paragraph{Verbatim paper-source connection}
\begin{ReviewVerbatim}
the winner uniformly among all alternatives with maximum Borda score. The Maximal Lotteries vot-
ing rule first computes the margin matrix M ∈ Rm×m , where Mx,y = #(x[U+001F control character]y)−#(y[U+001F control character]x)
                                                                          #(x[U+001F control character]y)+#(y[U+001F control character]x) . It then con-
siders a symmetric two-player zero-sum game in which player 1 selects alternative x1 , player 2 selects
alternative x2 , and the payoffs are Mx1 ,x2 for player 1 and Mx2 ,x1 = −Mx1 ,x2 for player 2. The max-
imal lotteries rule returns a distribution π ∈ argmaxπ1 ∈∆(A) minπ2 ∈∆(A) Ex1 ∼π1 ,x2 ∼π2 [Mx1 ,x2 ],
i.e., a mixed strategy in Nash equilibrium. When several such π exist, our results hold for any choice.
\end{ReviewVerbatim}



\paragraph{Lean-expanded paper semantic target}
\begin{ReviewVerbatim}
type:
{Alternative : Type u_1} →
  [Fintype Alternative] → [DecidableEq Alternative] → (Alternative → Alternative → ℝ) → PMF Alternative → Prop

value:
fun {Alternative : Type u_1} [Fintype Alternative] [DecidableEq Alternative]
    (empiricalRate : Alternative → Alternative → ℝ) (lottery : PMF Alternative) =>
  ∀ (opponent : PMF Alternative),
    0 ≤
      AppliedModelingLib.pmfPairExp lottery opponent fun (first second : Alternative) =>
        empiricalRate first second - empiricalRate second first
\end{ReviewVerbatim}

\paragraph{Direct retained dependencies}
\begin{itemize}\raggedright
\item \hyperlink{governing-declaration-09677001d465ec9d}{Applied\allowbreak{}Modeling\allowbreak{}Lib.\allowbreak{}pmf\allowbreak{}Pair\allowbreak{}Exp}
\end{itemize}
\paragraph{Recorded source-to-declaration screening}
\textbf{Verdict:} \texttt{matches}
\\
{\raggedright
\textbf{Reason:} This is the exact finite zero-\allowbreak{}sum empirical-\allowbreak{}equilibrium condition against every opponent.\allowbreak{} Its unobserved-\allowbreak{}pair rule is separately disclosed.\allowbreak{}
\par}
\reviewmatch{Matches source input}
\noindent\textbf{Reviewer annotation}\par
\reviewerbox
\clearpage
\hypertarget{paper-prerequisite-54fbc7c16362dbdd}{}
\subsection*{\small\ttfamily\raggedright Golz\allowbreak{}Haghtalab\allowbreak{}Yang2025\allowbreak{}Distortion.\allowbreak{}Is\allowbreak{}Source\allowbreak{}Alignment\allowbreak{}Benchmark}
\noindent\textbf{Paper-local declaration:}\par
{\footnotesize\ttfamily\raggedright Golz\allowbreak{}Haghtalab\allowbreak{}Yang2025\allowbreak{}Distortion.\allowbreak{}Is\allowbreak{}Source\allowbreak{}Alignment\allowbreak{}Benchmark\par}
\textbf{Source locator:} {\footnotesize\raggedright cited publication, lines 281-\allowbreak{}285\par}
\paragraph{Verbatim paper-source connection}
\begin{ReviewVerbatim}
To generalize the definition of distortion to alignment methods, we set the maximum average utility of
any policy in the KL-ball as the benchmark. For fixed m, D, β, d, µ and correlation between pairwise
comparisons, the distortion of alignment method f is
                                                             maxπ∈Bτ (πref ) AvgUtil(π)
                      dist(f ) = supD,πref ,τ lim supn→∞       AvgUtiln (f (·,πref ,τ )) .
\end{ReviewVerbatim}



\paragraph{Lean-expanded paper semantic target}
\begin{ReviewVerbatim}
type:
{Alternative : Type u_1} →
  [Fintype Alternative] → [DecidableEq Alternative] → (Alternative → ℝ) → PMF Alternative → PMF Alternative → ℝ → Prop

value:
fun {Alternative : Type u_1} [Fintype Alternative] [DecidableEq Alternative] (averageUtility : Alternative → ℝ)
    (reference benchmark : PMF Alternative) (klBudget : ℝ) =>
  AppliedModelingLib.Alignment.Welfare.InFiniteKLBall reference klBudget benchmark ∧
    ∀ (policy : PMF Alternative),
      AppliedModelingLib.Alignment.Welfare.InFiniteKLBall reference klBudget policy →
        AppliedModelingLib.pmfExp policy averageUtility ≤ AppliedModelingLib.pmfExp benchmark averageUtility
\end{ReviewVerbatim}

\paragraph{Direct retained dependencies}
\begin{itemize}\raggedright
\item \hyperlink{library-prerequisite-1d625de57bf58a6a}{Applied\allowbreak{}Modeling\allowbreak{}Lib.\allowbreak{}Alignment.\allowbreak{}Welfare.\allowbreak{}In\allowbreak{}Finite\allowbreak{}KL\allowbreak{}Ball}
\item \hyperlink{governing-declaration-8e4b4389a4f6454b}{Applied\allowbreak{}Modeling\allowbreak{}Lib.\allowbreak{}pmf\allowbreak{}Exp}
\end{itemize}
\paragraph{Recorded source-to-declaration screening}
\textbf{Verdict:} \texttt{matches}
\\
{\raggedright
\textbf{Reason:} This selects a welfare-\allowbreak{}maximizing policy in the stated KL ball.\allowbreak{}
\par}
\reviewmatch{Matches source input}
\noindent\textbf{Reviewer annotation}\par
\reviewerbox
\clearpage
\hypertarget{paper-prerequisite-dbd778917da393ac}{}
\subsection*{\small\ttfamily\raggedright Golz\allowbreak{}Haghtalab\allowbreak{}Yang2025\allowbreak{}Distortion.\allowbreak{}source\allowbreak{}Maximal\allowbreak{}Lottery\allowbreak{}Margin}
\noindent\textbf{Paper-local declaration:}\par
{\footnotesize\ttfamily\raggedright Golz\allowbreak{}Haghtalab\allowbreak{}Yang2025\allowbreak{}Distortion.\allowbreak{}source\allowbreak{}Maximal\allowbreak{}Lottery\allowbreak{}Margin\par}
\textbf{Source locator:} {\footnotesize\raggedright cited publication, lines 230-\allowbreak{}236\par}
\paragraph{Verbatim paper-source connection}
\begin{ReviewVerbatim}
the winner uniformly among all alternatives with maximum Borda score. The Maximal Lotteries vot-
ing rule first computes the margin matrix M ∈ Rm×m , where Mx,y = #(x[U+001F control character]y)−#(y[U+001F control character]x)
                                                                          #(x[U+001F control character]y)+#(y[U+001F control character]x) . It then con-
siders a symmetric two-player zero-sum game in which player 1 selects alternative x1 , player 2 selects
alternative x2 , and the payoffs are Mx1 ,x2 for player 1 and Mx2 ,x1 = −Mx1 ,x2 for player 2. The max-
imal lotteries rule returns a distribution π ∈ argmaxπ1 ∈∆(A) minπ2 ∈∆(A) Ex1 ∼π1 ,x2 ∼π2 [Mx1 ,x2 ],
i.e., a mixed strategy in Nash equilibrium. When several such π exist, our results hold for any choice.
\end{ReviewVerbatim}



\paragraph{Lean-expanded paper semantic target}
\begin{ReviewVerbatim}
type:
{Alternative : Type u_1} →
  (wins : Alternative → Alternative → ℕ) → (first second : Alternative) → wins first second + wins second first ≠ 0 → ℝ

value:
fun {Alternative : Type u_1} (wins : Alternative → Alternative → ℕ) (first second : Alternative)
    (_hdenominator : wins first second + wins second first ≠ 0) =>
  (↑(wins first second) - ↑(wins second first)) / (↑(wins first second) + ↑(wins second first))
\end{ReviewVerbatim}


\paragraph{Recorded source-to-declaration screening}
\textbf{Verdict:} \texttt{matches}
\\
{\raggedright
\textbf{Reason:} This is the literal normalized win-\allowbreak{}loss margin, explicitly restricted to its source nonzero-\allowbreak{}denominator domain.\allowbreak{}
\par}
\reviewmatch{Matches source input}
\noindent\textbf{Reviewer annotation}\par
\reviewerbox
\clearpage
\hypertarget{paper-prerequisite-bbe7079711c4e442}{}
\subsection*{\small\ttfamily\raggedright Golz\allowbreak{}Haghtalab\allowbreak{}Yang2025\allowbreak{}Distortion.\allowbreak{}source\allowbreak{}Maximal\allowbreak{}Lottery\allowbreak{}Value}
\noindent\textbf{Paper-local declaration:}\par
{\footnotesize\ttfamily\raggedright Golz\allowbreak{}Haghtalab\allowbreak{}Yang2025\allowbreak{}Distortion.\allowbreak{}source\allowbreak{}Maximal\allowbreak{}Lottery\allowbreak{}Value\par}
\textbf{Source locator:} {\footnotesize\raggedright cited publication, lines 230-\allowbreak{}236\par}
\paragraph{Verbatim paper-source connection}
\begin{ReviewVerbatim}
the winner uniformly among all alternatives with maximum Borda score. The Maximal Lotteries vot-
ing rule first computes the margin matrix M ∈ Rm×m , where Mx,y = #(x[U+001F control character]y)−#(y[U+001F control character]x)
                                                                          #(x[U+001F control character]y)+#(y[U+001F control character]x) . It then con-
siders a symmetric two-player zero-sum game in which player 1 selects alternative x1 , player 2 selects
alternative x2 , and the payoffs are Mx1 ,x2 for player 1 and Mx2 ,x1 = −Mx1 ,x2 for player 2. The max-
imal lotteries rule returns a distribution π ∈ argmaxπ1 ∈∆(A) minπ2 ∈∆(A) Ex1 ∼π1 ,x2 ∼π2 [Mx1 ,x2 ],
i.e., a mixed strategy in Nash equilibrium. When several such π exist, our results hold for any choice.
\end{ReviewVerbatim}



\paragraph{Lean-expanded paper semantic target}
\begin{ReviewVerbatim}
type:
{Alternative : Type u_1} →
  [Fintype Alternative] →
    [DecidableEq Alternative] →
      (wins : Alternative → Alternative → ℕ) →
        (∀ (first second : Alternative), first ≠ second → wins first second + wins second first ≠ 0) →
          PMF Alternative → ℝ

value:
fun {Alternative : Type u_1} [Fintype Alternative] [DecidableEq Alternative] (wins : Alternative → Alternative → ℕ)
    (hdenominator : ∀ (first second : Alternative), first ≠ second → wins first second + wins second first ≠ 0)
    (candidate : PMF Alternative) =>
  sInf
    (Set.range fun (opponent : PMF Alternative) =>
      AppliedModelingLib.pmfExp candidate fun (first : Alternative) =>
        AppliedModelingLib.pmfExp opponent fun (second : Alternative) =>
          if h : first = second then 0
          else
            GolzHaghtalabYang2025Distortion.sourceMaximalLotteryMargin wins first second (hdenominator first second h))
\end{ReviewVerbatim}

\paragraph{Direct retained dependencies}
\begin{itemize}\raggedright
\item \hyperlink{governing-declaration-8e4b4389a4f6454b}{Applied\allowbreak{}Modeling\allowbreak{}Lib.\allowbreak{}pmf\allowbreak{}Exp}
\item \hyperlink{paper-prerequisite-dbd778917da393ac}{Golz\allowbreak{}Haghtalab\allowbreak{}Yang2025\allowbreak{}Distortion.\allowbreak{}source\allowbreak{}Maximal\allowbreak{}Lottery\allowbreak{}Margin}
\end{itemize}
\paragraph{Recorded source-to-declaration screening}
\textbf{Verdict:} \texttt{matches}
\\
{\raggedright
\textbf{Reason:} This is the minimum expected source margin over opponent lotteries, with diagonal zero and the same off-\allowbreak{}diagonal domain.\allowbreak{}
\par}
\reviewmatch{Matches source input}
\noindent\textbf{Reviewer annotation}\par
\reviewerbox
\clearpage
\hypertarget{paper-prerequisite-91b5f077c58b43c1}{}
\subsection*{\small\ttfamily\raggedright Golz\allowbreak{}Haghtalab\allowbreak{}Yang2025\allowbreak{}Distortion.\allowbreak{}Is\allowbreak{}Source\allowbreak{}Maximal\allowbreak{}Lottery}
\noindent\textbf{Paper-local declaration:}\par
{\footnotesize\ttfamily\raggedright Golz\allowbreak{}Haghtalab\allowbreak{}Yang2025\allowbreak{}Distortion.\allowbreak{}Is\allowbreak{}Source\allowbreak{}Maximal\allowbreak{}Lottery\par}
\textbf{Source locator:} {\footnotesize\raggedright cited publication, lines 230-\allowbreak{}236\par}
\paragraph{Verbatim paper-source connection}
\begin{ReviewVerbatim}
the winner uniformly among all alternatives with maximum Borda score. The Maximal Lotteries vot-
ing rule first computes the margin matrix M ∈ Rm×m , where Mx,y = #(x[U+001F control character]y)−#(y[U+001F control character]x)
                                                                          #(x[U+001F control character]y)+#(y[U+001F control character]x) . It then con-
siders a symmetric two-player zero-sum game in which player 1 selects alternative x1 , player 2 selects
alternative x2 , and the payoffs are Mx1 ,x2 for player 1 and Mx2 ,x1 = −Mx1 ,x2 for player 2. The max-
imal lotteries rule returns a distribution π ∈ argmaxπ1 ∈∆(A) minπ2 ∈∆(A) Ex1 ∼π1 ,x2 ∼π2 [Mx1 ,x2 ],
i.e., a mixed strategy in Nash equilibrium. When several such π exist, our results hold for any choice.
\end{ReviewVerbatim}



\paragraph{Lean-expanded paper semantic target}
\begin{ReviewVerbatim}
type:
{Alternative : Type u_1} →
  [Fintype Alternative] →
    [DecidableEq Alternative] →
      (wins : Alternative → Alternative → ℕ) →
        (∀ (first second : Alternative), first ≠ second → wins first second + wins second first ≠ 0) →
          PMF Alternative → Prop

value:
fun {Alternative : Type u_1} [Fintype Alternative] [DecidableEq Alternative] (wins : Alternative → Alternative → ℕ)
    (hdenominator : ∀ (first second : Alternative), first ≠ second → wins first second + wins second first ≠ 0)
    (lottery : PMF Alternative) =>
  ∀ (candidate : PMF Alternative),
    GolzHaghtalabYang2025Distortion.sourceMaximalLotteryValue wins hdenominator candidate ≤
      GolzHaghtalabYang2025Distortion.sourceMaximalLotteryValue wins hdenominator lottery
\end{ReviewVerbatim}

\paragraph{Direct retained dependencies}
\begin{itemize}\raggedright
\item \hyperlink{paper-prerequisite-bbe7079711c4e442}{Golz\allowbreak{}Haghtalab\allowbreak{}Yang2025\allowbreak{}Distortion.\allowbreak{}source\allowbreak{}Maximal\allowbreak{}Lottery\allowbreak{}Value}
\end{itemize}
\paragraph{Recorded source-to-declaration screening}
\textbf{Verdict:} \texttt{matches}
\\
{\raggedright
\textbf{Reason:} This is the source maximin condition on its raw nonzero-\allowbreak{}denominator margin domain.\allowbreak{}
\par}
\reviewmatch{Matches source input}
\noindent\textbf{Reviewer annotation}\par
\reviewerbox
\clearpage
\hypertarget{paper-prerequisite-ceb3f23b84b56503}{}
\subsection*{\small\ttfamily\raggedright Golz\allowbreak{}Haghtalab\allowbreak{}Yang2025\allowbreak{}Distortion.\allowbreak{}Is\allowbreak{}Source\allowbreak{}Social\allowbreak{}Choice\allowbreak{}Optimal\allowbreak{}Average\allowbreak{}Utility}
\noindent\textbf{Paper-local declaration:}\par
{\footnotesize\ttfamily\raggedright Golz\allowbreak{}Haghtalab\allowbreak{}Yang2025\allowbreak{}Distortion.\allowbreak{}Is\allowbreak{}Source\allowbreak{}Social\allowbreak{}Choice\allowbreak{}Optimal\allowbreak{}Average\allowbreak{}Utility\par}
\textbf{Source locator:} {\footnotesize\raggedright cited publication, lines 202-\allowbreak{}215\par}
\paragraph{Verbatim paper-source connection}
\begin{ReviewVerbatim}
Social Choice Setting. A voting rule f observes the sampled pairwise comparisons {xji [U+001F control character]i
yij }i∈[n],j∈[d] and maps them to a probability distribution over alternatives. For some m, D, β, d, µ,
and
  [U+0002 control character] the correlation between comparisons,     the average utility of f for n samples is AvgUtiln (f ) :=
E AvgUtil f ({xji [U+001F control character]i yij }i,j ) , where the expectation is taken over the pairwise comparisons. The
                               [U+0001 control character][U+0003 control character]

distortion of f on D is the competitive ratio between AvgUtiln (f ) and the optimal average utility
maxx∈A AvgUtil(x) in the limit of n → ∞ samples, and the distortion of f the worst-case distortion
over all D:

            dist(f, D) := lim supn→∞ maxAvgUtil
                                        x∈A AvgUtil(x)
                                                (f )   ,               dist(f ) = supD dist(f, D).
\end{ReviewVerbatim}



\paragraph{Lean-expanded paper semantic target}
\begin{ReviewVerbatim}
type:
{Alternative : Type u_1} → [Fintype Alternative] → (Alternative → ℝ) → ℝ → Prop

value:
fun {Alternative : Type u_1} [Fintype Alternative] (averageUtility : Alternative → ℝ) (optimalAverageUtility : ℝ) =>
  (∃ (alternative : Alternative), averageUtility alternative = optimalAverageUtility) ∧
    ∀ (alternative : Alternative), averageUtility alternative ≤ optimalAverageUtility
\end{ReviewVerbatim}


\paragraph{Recorded source-to-declaration screening}
\textbf{Verdict:} \texttt{matches}
\\
{\raggedright
\textbf{Reason:} This captures attainment and dominance of the finite optimal social-\allowbreak{}choice welfare benchmark.\allowbreak{}
\par}
\reviewmatch{Matches source input}
\noindent\textbf{Reviewer annotation}\par
\reviewerbox
\clearpage
\hypertarget{paper-prerequisite-54cac251e94c8e83}{}
\subsection*{\small\ttfamily\raggedright Golz\allowbreak{}Haghtalab\allowbreak{}Yang2025\allowbreak{}Distortion.\allowbreak{}corollary4\allowbreak{}Iid\allowbreak{}User\allowbreak{}Empirical\allowbreak{}Maximal\allowbreak{}Lottery}
\noindent\textbf{Paper-local declaration:}\par
{\footnotesize\ttfamily\raggedright Golz\allowbreak{}Haghtalab\allowbreak{}Yang2025\allowbreak{}Distortion.\allowbreak{}corollary4\allowbreak{}Iid\allowbreak{}User\allowbreak{}Empirical\allowbreak{}Maximal\allowbreak{}Lottery\par}
\textbf{Source locator:} {\footnotesize\raggedright cited publication, lines 230-\allowbreak{}236\par}
\paragraph{Verbatim paper-source connection}
\begin{ReviewVerbatim}
the winner uniformly among all alternatives with maximum Borda score. The Maximal Lotteries vot-
ing rule first computes the margin matrix M ∈ Rm×m , where Mx,y = #(x[U+001F control character]y)−#(y[U+001F control character]x)
                                                                          #(x[U+001F control character]y)+#(y[U+001F control character]x) . It then con-
siders a symmetric two-player zero-sum game in which player 1 selects alternative x1 , player 2 selects
alternative x2 , and the payoffs are Mx1 ,x2 for player 1 and Mx2 ,x1 = −Mx1 ,x2 for player 2. The max-
imal lotteries rule returns a distribution π ∈ argmaxπ1 ∈∆(A) minπ2 ∈∆(A) Ex1 ∼π1 ,x2 ∼π2 [Mx1 ,x2 ],
i.e., a mixed strategy in Nash equilibrium. When several such π exist, our results hold for any choice.
\end{ReviewVerbatim}



\paragraph{Lean-expanded paper semantic target}
\begin{ReviewVerbatim}
type:
{Alternative : Type u_1} →
  [Fintype Alternative] →
    [DecidableEq Alternative] →
      [Nonempty Alternative] →
        {users comparisonsPerUser : ℕ} →
          (Fin users →
              GolzHaghtalabYang2025Distortion.theorem2UserResponseTable Alternative ×
                (Fin comparisonsPerUser → Alternative × Alternative)) →
            PMF Alternative

value:
fun {Alternative : Type u_1} [Fintype Alternative] [DecidableEq Alternative] [Nonempty Alternative]
    {users comparisonsPerUser : ℕ}
    (sample :
      Fin users →
        GolzHaghtalabYang2025Distortion.theorem2UserResponseTable Alternative ×
          (Fin comparisonsPerUser → Alternative × Alternative)) =>
  GolzHaghtalabYang2025Distortion.theorem7IidUserEmpiricalPolicy (AppliedModelingLib.uniformPMF Alternative)
    (Real.log ↑(Fintype.card Alternative))
    (GolzHaghtalabYang2025Distortion.corollary4_uniformFullSimplexBudget_nonneg Alternative) sample
\end{ReviewVerbatim}

\paragraph{Direct retained dependencies}
\begin{itemize}\raggedright
\item \hyperlink{governing-declaration-48f10f2715c04df3}{Applied\allowbreak{}Modeling\allowbreak{}Lib.\allowbreak{}uniform\allowbreak{}PMF}
\item \hyperlink{governing-declaration-c5d0273cbd7ddcc9}{Golz\allowbreak{}Haghtalab\allowbreak{}Yang2025\allowbreak{}Distortion.\allowbreak{}theorem2\allowbreak{}User\allowbreak{}Response\allowbreak{}Table}
\item \hyperlink{governing-declaration-71b9e9aed968af43}{Golz\allowbreak{}Haghtalab\allowbreak{}Yang2025\allowbreak{}Distortion.\allowbreak{}theorem7\allowbreak{}Iid\allowbreak{}User\allowbreak{}Empirical\allowbreak{}Policy}
\end{itemize}
\paragraph{Recorded source-to-declaration screening}
\textbf{Verdict:} \texttt{matches}
\\
{\raggedright
\textbf{Reason:} The uniform reference with radius log |A| contains the simplex;\allowbreak{} observed pairs recover the paper ratios and unobserved off-\allowbreak{}diagonal pairs have the explicitly disclosed neutral 0/\allowbreak{}0-\allowbreak{}to-\allowbreak{}zero margin.\allowbreak{}
\par}
\reviewmatch{Matches source input}
\noindent\textbf{Reviewer annotation}\par
\reviewerbox
\clearpage
\hypertarget{paper-prerequisite-674d1bc00409609f}{}
\subsection*{\small\ttfamily\raggedright Golz\allowbreak{}Haghtalab\allowbreak{}Yang2025\allowbreak{}Distortion.\allowbreak{}empirical\allowbreak{}Borda\allowbreak{}Definition\allowbreak{}Spec}
\noindent\textbf{Paper-local declaration:}\par
{\footnotesize\ttfamily\raggedright Golz\allowbreak{}Haghtalab\allowbreak{}Yang2025\allowbreak{}Distortion.\allowbreak{}empirical\allowbreak{}Borda\allowbreak{}Definition\allowbreak{}Spec\par}
\textbf{Source locator:} {\footnotesize\raggedright cited publication, lines 220-\allowbreak{}230\par}
\paragraph{Verbatim paper-source connection}
\begin{ReviewVerbatim}
For alternatives x, y, let #(x [U+001F control character] y) := {(i, j) | xji = x, yij = y} denote the number of pairwise
comparisons in which x beat y. The (normalized) Borda score [45] of alternative x is
                                                       P
                                                         y∈A #(x[U+001F control character]y)
                                   BC(x) := P         #(x[U+001F control character]y)+
                                                              P             ,
                                                  y∈A           y6=x #(y[U+001F control character]x)


i.e., the fraction of pairwise comparisons involving x in which it wins. The Borda voting rule chooses
the winner uniformly among all alternatives with maximum Borda score. The Maximal Lotteries vot-
\end{ReviewVerbatim}



\paragraph{Lean-expanded paper semantic target}
\begin{ReviewVerbatim}
type:
{Report : Type u_1} →
  {Alternative : Type u_2} →
    [Fintype Alternative] →
      (Alternative → Report → ℝ) →
        (Alternative → Report → ℝ) → ((horizon : ℕ) → (Fin horizon → Report) → PMF Alternative) → Prop

value:
fun {Report : Type u_1} {Alternative : Type u_2} [Fintype Alternative] (wins incidences : Alternative → Report → ℝ)
    (rule : (horizon : ℕ) → (Fin horizon → Report) → PMF Alternative) =>
  ∀ (horizon : ℕ) (sample : Fin horizon → Report),
    have score : Alternative → ℝ := fun (alternative : Alternative) =>
      if AppliedModelingLib.Probability.finiteIidScoreSum (incidences alternative) sample = 0 then 0
      else
        AppliedModelingLib.Probability.finiteIidScoreSum (wins alternative) sample /
          AppliedModelingLib.Probability.finiteIidScoreSum (incidences alternative) sample;
    (∀ (alternative : Alternative),
        0 < ((rule horizon sample) alternative).toReal ↔ ∀ (other : Alternative), score other ≤ score alternative) ∧
      ∀ (first second : Alternative),
        (∀ (other : Alternative), score other ≤ score first) →
          (∀ (other : Alternative), score other ≤ score second) →
            ((rule horizon sample) first).toReal = ((rule horizon sample) second).toReal
\end{ReviewVerbatim}

\paragraph{Direct retained dependencies}
\begin{itemize}\raggedright
\item \hyperlink{governing-declaration-65283b6bc1ebd8c0}{Applied\allowbreak{}Modeling\allowbreak{}Lib.\allowbreak{}Probability.\allowbreak{}finite\allowbreak{}Iid\allowbreak{}Score\allowbreak{}Sum}
\end{itemize}
\paragraph{Recorded source-to-declaration screening}
\textbf{Verdict:} \texttt{matches}
\\
{\raggedright
\textbf{Reason:} This retains the source positive-\allowbreak{}incidence quotient and uniform maximizer selection.\allowbreak{} Score zero is limited to the otherwise undefined zero-\allowbreak{}incidence evaluator branch and is not a model premise.\allowbreak{}
\par}
\reviewmatch{Matches source input}
\noindent\textbf{Reviewer annotation}\par
\reviewerbox
\clearpage
\hypertarget{paper-prerequisite-acf649baada45fd3}{}
\subsection*{\small\ttfamily\raggedright Golz\allowbreak{}Haghtalab\allowbreak{}Yang2025\allowbreak{}Distortion.\allowbreak{}nlhf\allowbreak{}Definition\allowbreak{}Spec}
\noindent\textbf{Paper-local declaration:}\par
{\footnotesize\ttfamily\raggedright Golz\allowbreak{}Haghtalab\allowbreak{}Yang2025\allowbreak{}Distortion.\allowbreak{}nlhf\allowbreak{}Definition\allowbreak{}Spec\par}
\textbf{Source locator:} {\footnotesize\raggedright cited publication, lines 277-\allowbreak{}280\par}
\paragraph{Verbatim paper-source connection}
\begin{ReviewVerbatim}
The alignment method NLHF was inspired in part by a desire to better align with the
preferences of a heterogeneous group [36]. NLHF naturally adapts the definition of maxi-
mal lotteries by constraining both players’ mixed strategies to the KL-ball, i.e., πNLHF :=
argmaxπ1 ∈Bτ (πref ) minπ2 ∈Bτ (πref ) Ex1 ∼π1 ,x2 ∼π2 [Mx1 ,x2 ].

[Next verbatim source excerpt]

      If τ = 0, the only feasible policy is πref , and the claim clearly holds for λ = ∞. We also assume that
πref (a) > 0 for all a ∈ M . Otherwise if πref (a) = 0 for some a ∈ M , then both the regularized and constrained
versions forbid any policy to put nonzero probability on a, which leads to an effectively smaller candidate set.
\end{ReviewVerbatim}



\paragraph{Lean-expanded paper semantic target}
\begin{ReviewVerbatim}
type:
{User Alternative : Type} →
  [Fintype User] →
    [DecidableEq User] →
      [Fintype Alternative] →
        [DecidableEq Alternative] →
          PMF User →
            AppliedModelingLib.Alignment.Welfare.FiniteUtilityProfile User Alternative →
              ℝ → PMF Alternative → ℝ → PMF Alternative → Prop

value:
fun {User Alternative : Type} [Fintype User] [DecidableEq User] [Fintype Alternative] [DecidableEq Alternative]
    (population : PMF User) (utility : AppliedModelingLib.Alignment.Welfare.FiniteUtilityProfile User Alternative)
    (btScale : ℝ) (reference : PMF Alternative) (klBudget : ℝ) (policy : PMF Alternative) =>
  AppliedModelingLib.PMFFullSupport reference ∧
    AppliedModelingLib.finiteKLDivergence policy reference ≤ klBudget ∧
      ∃ (worst : PMF Alternative),
        AppliedModelingLib.finiteKLDivergence worst reference ≤ klBudget ∧
          (∀ (opponent : PMF Alternative),
              AppliedModelingLib.finiteKLDivergence opponent reference ≤ klBudget →
                AppliedModelingLib.GameTheory.PreferenceGame.contextFreeCenteredPreferencePayoff
                    (AppliedModelingLib.Alignment.Welfare.populationBradleyTerryPreference population utility btScale)
                    policy worst ≤
                  AppliedModelingLib.GameTheory.PreferenceGame.contextFreeCenteredPreferencePayoff
                    (AppliedModelingLib.Alignment.Welfare.populationBradleyTerryPreference population utility btScale)
                    policy opponent) ∧
            ∀ (candidate : PMF Alternative),
              AppliedModelingLib.finiteKLDivergence candidate reference ≤ klBudget →
                AppliedModelingLib.GameTheory.PreferenceGame.contextFreeCenteredPreferencePayoff
                    (AppliedModelingLib.Alignment.Welfare.populationBradleyTerryPreference population utility btScale)
                    candidate worst ≤
                  AppliedModelingLib.GameTheory.PreferenceGame.contextFreeCenteredPreferencePayoff
                    (AppliedModelingLib.Alignment.Welfare.populationBradleyTerryPreference population utility btScale)
                    policy worst
\end{ReviewVerbatim}

\paragraph{Direct retained dependencies}
\begin{itemize}\raggedright
\item \hyperlink{library-prerequisite-56c0734f8877ddbe}{Applied\allowbreak{}Modeling\allowbreak{}Lib.\allowbreak{}Alignment.\allowbreak{}Welfare.\allowbreak{}Finite\allowbreak{}Utility\allowbreak{}Profile}
\item \hyperlink{library-prerequisite-cf06d15d7f8b7c31}{Applied\allowbreak{}Modeling\allowbreak{}Lib.\allowbreak{}Alignment.\allowbreak{}Welfare.\allowbreak{}population\allowbreak{}Bradley\allowbreak{}Terry\allowbreak{}Preference}
\item \hyperlink{governing-declaration-44fb06319bbd7683}{Applied\allowbreak{}Modeling\allowbreak{}Lib.\allowbreak{}Game\allowbreak{}Theory.\allowbreak{}Preference\allowbreak{}Game.\allowbreak{}context\allowbreak{}Free\allowbreak{}Centered\allowbreak{}Preference\allowbreak{}Payoff}
\item \hyperlink{library-prerequisite-268c959ee7d2d148}{Applied\allowbreak{}Modeling\allowbreak{}Lib.\allowbreak{}PMF\allowbreak{}Full\allowbreak{}Support}
\item \hyperlink{library-prerequisite-d3573881ae45d4c6}{Applied\allowbreak{}Modeling\allowbreak{}Lib.\allowbreak{}finite\allowbreak{}KL\allowbreak{}Divergence}
\end{itemize}
\paragraph{Recorded source-to-declaration screening}
\textbf{Verdict:} \texttt{matches}
\\
{\raggedright
\textbf{Reason:} Both players use the same KL ball and centered population Bradley-\allowbreak{}Terry payoff;\allowbreak{} reference-\allowbreak{}support scope is covered by the source-\allowbreak{}supported positive-\allowbreak{}support/\allowbreak{}reduced-\allowbreak{}candidate-\allowbreak{}set clarification.\allowbreak{}
\par}
\reviewmatch{Matches source input}
\noindent\textbf{Reviewer annotation}\par
\reviewerbox
\clearpage
\hypertarget{paper-prerequisite-8df126bd40665caf}{}
\subsection*{\small\ttfamily\raggedright Golz\allowbreak{}Haghtalab\allowbreak{}Yang2025\allowbreak{}Distortion.\allowbreak{}population\allowbreak{}Borda\allowbreak{}Definition\allowbreak{}Spec}
\noindent\textbf{Paper-local declaration:}\par
{\footnotesize\ttfamily\raggedright Golz\allowbreak{}Haghtalab\allowbreak{}Yang2025\allowbreak{}Distortion.\allowbreak{}population\allowbreak{}Borda\allowbreak{}Definition\allowbreak{}Spec\par}
\textbf{Source locator:} {\footnotesize\raggedright cited publication, lines 372-\allowbreak{}375\par}
\paragraph{Verbatim paper-source connection}
\begin{ReviewVerbatim}
for n → ∞, assuming
Proof sketch (full proof in Appendix E.1). For exposition, we sketch this proofP
that each alternative’s Borda count has converged to its expectation BC? (x) := y∈A µ(y)·p(x [U+001F control character] y).
Let x̂ = argmaxx∈A BC? (x) be the Borda winner in the limit, and x? = argmaxx∈A AvgUtil(x) be
\end{ReviewVerbatim}



\paragraph{Lean-expanded paper semantic target}
\begin{ReviewVerbatim}
type:
{Alternative : Type u_1} →
  [Fintype Alternative] →
    [DecidableEq Alternative] →
      PMF Alternative →
        AppliedModelingLib.Learning.HumanFeedback.PairwisePreference PUnit.{1} Alternative → Alternative → Prop

value:
fun {Alternative : Type u_1} [Fintype Alternative] [DecidableEq Alternative] (sampling : PMF Alternative)
    (preference : AppliedModelingLib.Learning.HumanFeedback.PairwisePreference PUnit.{1} Alternative)
    (winner : Alternative) =>
  ∀ (alternative : Alternative),
    (AppliedModelingLib.pmfExp sampling fun (opponent : Alternative) =>
        preference.prob PUnit.unit alternative opponent) ≤
      AppliedModelingLib.pmfExp sampling fun (opponent : Alternative) => preference.prob PUnit.unit winner opponent
\end{ReviewVerbatim}

\paragraph{Direct retained dependencies}
\begin{itemize}\raggedright
\item \hyperlink{library-prerequisite-7e2db20cf087d9ba}{Applied\allowbreak{}Modeling\allowbreak{}Lib.\allowbreak{}Learning.\allowbreak{}Human\allowbreak{}Feedback.\allowbreak{}Pairwise\allowbreak{}Preference}
\item \hyperlink{governing-declaration-8e4b4389a4f6454b}{Applied\allowbreak{}Modeling\allowbreak{}Lib.\allowbreak{}pmf\allowbreak{}Exp}
\end{itemize}
\paragraph{Recorded source-to-declaration screening}
\textbf{Verdict:} \texttt{matches}
\\
{\raggedright
\textbf{Reason:} This uses the source sampling-\allowbreak{}weighted population pairwise score and a maximizing alternative.\allowbreak{}
\par}
\reviewmatch{Matches source input}
\noindent\textbf{Reviewer annotation}\par
\reviewerbox
\clearpage
\hypertarget{paper-prerequisite-58ee6f1d89c4f942}{}
\subsection*{\small\ttfamily\raggedright Golz\allowbreak{}Haghtalab\allowbreak{}Yang2025\allowbreak{}Distortion.\allowbreak{}rlhf\allowbreak{}Definition\allowbreak{}Spec}
\noindent\textbf{Paper-local declaration:}\par
{\footnotesize\ttfamily\raggedright Golz\allowbreak{}Haghtalab\allowbreak{}Yang2025\allowbreak{}Distortion.\allowbreak{}rlhf\allowbreak{}Definition\allowbreak{}Spec\par}
\textbf{Source locator:} {\footnotesize\raggedright cited publication, lines 269-\allowbreak{}276\par}
\paragraph{Verbatim paper-source connection}
\begin{ReviewVerbatim}
The RLHF method first estimates rewards for each alternative, using maximum likelihood estimation
assuming that comparisons were generated by a single Bradley–Terry model:
                  r := argmaxr∈Rm 1≤i≤n,1≤j≤d log σ(r(xji ) − r(yij )) .8
                                     P                                     [U+0001 control character]

Next, RLHF uses PPO [44] to compute the policy πRLHF := argmaxπ∈Bτ (πref ) Ex∼π [r(x)] with
maximum expected reward within the KL-ball. In our setup, DPO is equivalent to RLHF (see
Appendix F.3) and thus has the same distortion.

[Next verbatim source excerpt]

      If τ = 0, the only feasible policy is πref , and the claim clearly holds for λ = ∞. We also assume that
πref (a) > 0 for all a ∈ M . Otherwise if πref (a) = 0 for some a ∈ M , then both the regularized and constrained
versions forbid any policy to put nonzero probability on a, which leads to an effectively smaller candidate set.
\end{ReviewVerbatim}



\paragraph{Lean-expanded paper semantic target}
\begin{ReviewVerbatim}
type:
{Alternative : Type u_1} →
  [Fintype Alternative] →
    [DecidableEq Alternative] →
      {horizon : ℕ} →
        (Fin horizon → AppliedModelingLib.Learning.HumanFeedback.BinaryPairwiseReport Alternative) →
          Alternative →
            PMF Alternative →
              PMF Alternative →
                ℝ → AppliedModelingLib.Learning.HumanFeedback.PairwiseCountDataset.ScoreVector Alternative → Prop

value:
fun {Alternative : Type u_1} [Fintype Alternative] [DecidableEq Alternative] {horizon : ℕ}
    (sample : Fin horizon → AppliedModelingLib.Learning.HumanFeedback.BinaryPairwiseReport Alternative)
    (scoreReference : Alternative) (reference policy : PMF Alternative) (klBudget : ℝ)
    (score : AppliedModelingLib.Learning.HumanFeedback.PairwiseCountDataset.ScoreVector Alternative) =>
  AppliedModelingLib.PMFFullSupport reference ∧
    (score scoreReference = 0 ∧
        IsMaxOn
          ((AppliedModelingLib.Learning.HumanFeedback.PairwiseCountDataset.ofBinaryReports sample).pairwiseLogLikelihood
            Real.sigmoid)
          {candidate : AppliedModelingLib.Learning.HumanFeedback.PairwiseCountDataset.ScoreVector Alternative |
            candidate scoreReference = 0}
          score) ∧
      AppliedModelingLib.Learning.HumanFeedback.contextAveragedPolicyKLDivergence (PMF.pure PUnit.unit)
            (AppliedModelingLib.Alignment.Welfare.contextFreeAlternativePolicy policy)
            (AppliedModelingLib.Alignment.Welfare.contextFreeAlternativePolicy reference) ≤
          klBudget ∧
        ∀ (other : AppliedModelingLib.Learning.HumanFeedback.FinitePolicy PUnit.{1} Alternative),
          AppliedModelingLib.Learning.HumanFeedback.contextAveragedPolicyKLDivergence (PMF.pure PUnit.unit) other
                (AppliedModelingLib.Alignment.Welfare.contextFreeAlternativePolicy reference) ≤
              klBudget →
            (AppliedModelingLib.Learning.HumanFeedback.policyExpectedScore (PMF.pure PUnit.unit) other
                fun (x : PUnit.{1}) => score) ≤
              AppliedModelingLib.Learning.HumanFeedback.policyExpectedScore (PMF.pure PUnit.unit)
                (AppliedModelingLib.Alignment.Welfare.contextFreeAlternativePolicy policy) fun (x : PUnit.{1}) => score
\end{ReviewVerbatim}

\paragraph{Direct retained dependencies}
\begin{itemize}\raggedright
\item \hyperlink{governing-declaration-b7cf9e04e23161e5}{Applied\allowbreak{}Modeling\allowbreak{}Lib.\allowbreak{}Alignment.\allowbreak{}Welfare.\allowbreak{}context\allowbreak{}Free\allowbreak{}Alternative\allowbreak{}Policy}
\item \hyperlink{governing-declaration-bc3d6ca8e7e0ac9d}{Applied\allowbreak{}Modeling\allowbreak{}Lib.\allowbreak{}Learning.\allowbreak{}Human\allowbreak{}Feedback.\allowbreak{}Binary\allowbreak{}Pairwise\allowbreak{}Report}
\item \hyperlink{governing-declaration-577316c60f45ef11}{Applied\allowbreak{}Modeling\allowbreak{}Lib.\allowbreak{}Learning.\allowbreak{}Human\allowbreak{}Feedback.\allowbreak{}Finite\allowbreak{}Policy}
\item \hyperlink{governing-declaration-87eda978eb1b0c9d}{Applied\allowbreak{}Modeling\allowbreak{}Lib.\allowbreak{}Learning.\allowbreak{}Human\allowbreak{}Feedback.\allowbreak{}Pairwise\allowbreak{}Count\allowbreak{}Dataset.\allowbreak{}Score\allowbreak{}Vector}
\item \hyperlink{governing-declaration-449f87e91bf37ce0}{Applied\allowbreak{}Modeling\allowbreak{}Lib.\allowbreak{}Learning.\allowbreak{}Human\allowbreak{}Feedback.\allowbreak{}Pairwise\allowbreak{}Count\allowbreak{}Dataset.\allowbreak{}of\allowbreak{}Binary\allowbreak{}Reports}
\item \hyperlink{governing-declaration-10adb9e7c1ff9f4c}{Applied\allowbreak{}Modeling\allowbreak{}Lib.\allowbreak{}Learning.\allowbreak{}Human\allowbreak{}Feedback.\allowbreak{}Pairwise\allowbreak{}Count\allowbreak{}Dataset.\allowbreak{}pairwise\allowbreak{}Log\allowbreak{}Likelihood}
\item \hyperlink{governing-declaration-c528a6568d75fffd}{Applied\allowbreak{}Modeling\allowbreak{}Lib.\allowbreak{}Learning.\allowbreak{}Human\allowbreak{}Feedback.\allowbreak{}context\allowbreak{}Averaged\allowbreak{}Policy\allowbreak{}KL\allowbreak{}Divergence}
\item \hyperlink{governing-declaration-4daa53fb741a5ea7}{Applied\allowbreak{}Modeling\allowbreak{}Lib.\allowbreak{}Learning.\allowbreak{}Human\allowbreak{}Feedback.\allowbreak{}policy\allowbreak{}Expected\allowbreak{}Score}
\item \hyperlink{library-prerequisite-268c959ee7d2d148}{Applied\allowbreak{}Modeling\allowbreak{}Lib.\allowbreak{}PMF\allowbreak{}Full\allowbreak{}Support}
\end{itemize}
\paragraph{Recorded source-to-declaration screening}
\textbf{Verdict:} \texttt{matches}
\\
{\raggedright
\textbf{Reason:} This faithfully specifies finite logistic MLE up to harmless reference-\allowbreak{}score normalization, followed by reward maximization in the same KL ball.\allowbreak{}
\par}
\reviewmatch{Matches source input}
\noindent\textbf{Reviewer annotation}\par
\reviewerbox
\clearpage
\hypertarget{paper-prerequisite-e2a06898e9a1c2db}{}
\subsection*{\small\ttfamily\raggedright Golz\allowbreak{}Haghtalab\allowbreak{}Yang2025\allowbreak{}Distortion.\allowbreak{}source\allowbreak{}Alignment\allowbreak{}Average\allowbreak{}Utility}
\noindent\textbf{Paper-local declaration:}\par
{\footnotesize\ttfamily\raggedright Golz\allowbreak{}Haghtalab\allowbreak{}Yang2025\allowbreak{}Distortion.\allowbreak{}source\allowbreak{}Alignment\allowbreak{}Average\allowbreak{}Utility\par}
\textbf{Source locator:} {\footnotesize\raggedright cited publication, lines 281-\allowbreak{}285\par}
\paragraph{Verbatim paper-source connection}
\begin{ReviewVerbatim}
To generalize the definition of distortion to alignment methods, we set the maximum average utility of
any policy in the KL-ball as the benchmark. For fixed m, D, β, d, µ and correlation between pairwise
comparisons, the distortion of alignment method f is
                                                             maxπ∈Bτ (πref ) AvgUtil(π)
                      dist(f ) = supD,πref ,τ lim supn→∞       AvgUtiln (f (·,πref ,τ )) .
\end{ReviewVerbatim}



\paragraph{Lean-expanded paper semantic target}
\begin{ReviewVerbatim}
type:
{Alternative : Type u_1} →
  {Observation : Type u_2} →
    [Fintype Alternative] →
      [Fintype Observation] →
        (Alternative → ℝ) →
          ((horizon : ℕ) → PMF (Fin horizon → Observation)) →
            ((horizon : ℕ) → (Fin horizon → Observation) → PMF Alternative → ℝ → PMF Alternative) →
              PMF Alternative → ℝ → ℕ → ℝ

value:
fun {Alternative : Type u_1} {Observation : Type u_2} [Fintype Alternative] [Fintype Observation]
    (averageUtility : Alternative → ℝ) (observationLaw : (horizon : ℕ) → PMF (Fin horizon → Observation))
    (method : (horizon : ℕ) → (Fin horizon → Observation) → PMF Alternative → ℝ → PMF Alternative)
    (reference : PMF Alternative) (klBudget : ℝ) (horizon : ℕ) =>
  AppliedModelingLib.pmfExp (observationLaw horizon) fun (observations : Fin horizon → Observation) =>
    AppliedModelingLib.pmfExp (method horizon observations reference klBudget) averageUtility
\end{ReviewVerbatim}

\paragraph{Direct retained dependencies}
\begin{itemize}\raggedright
\item \hyperlink{governing-declaration-8e4b4389a4f6454b}{Applied\allowbreak{}Modeling\allowbreak{}Lib.\allowbreak{}pmf\allowbreak{}Exp}
\end{itemize}
\paragraph{Recorded source-to-declaration screening}
\textbf{Verdict:} \texttt{matches}
\\
{\raggedright
\textbf{Reason:} This is expected policy welfare over the full report law, parameterized by reference and KL budget.\allowbreak{}
\par}
\reviewmatch{Matches source input}
\noindent\textbf{Reviewer annotation}\par
\reviewerbox
\clearpage
\hypertarget{paper-prerequisite-44ab9adb1cd36c3c}{}
\subsection*{\small\ttfamily\raggedright Golz\allowbreak{}Haghtalab\allowbreak{}Yang2025\allowbreak{}Distortion.\allowbreak{}source\allowbreak{}Alignment\allowbreak{}Distortion}
\noindent\textbf{Paper-local declaration:}\par
{\footnotesize\ttfamily\raggedright Golz\allowbreak{}Haghtalab\allowbreak{}Yang2025\allowbreak{}Distortion.\allowbreak{}source\allowbreak{}Alignment\allowbreak{}Distortion\par}
\textbf{Source locator:} {\footnotesize\raggedright cited publication, lines 281-\allowbreak{}285\par}
\paragraph{Verbatim paper-source connection}
\begin{ReviewVerbatim}
To generalize the definition of distortion to alignment methods, we set the maximum average utility of
any policy in the KL-ball as the benchmark. For fixed m, D, β, d, µ and correlation between pairwise
comparisons, the distortion of alignment method f is
                                                             maxπ∈Bτ (πref ) AvgUtil(π)
                      dist(f ) = supD,πref ,τ lim supn→∞       AvgUtiln (f (·,πref ,τ )) .
\end{ReviewVerbatim}



\paragraph{Lean-expanded paper semantic target}
\begin{ReviewVerbatim}
type:
{Alternative : Type u_1} →
  {Observation : Type u_2} →
    [Fintype Alternative] →
      [DecidableEq Alternative] →
        [Fintype Observation] →
          (Alternative → ℝ) →
            ℝ →
              ((horizon : ℕ) → PMF (Fin horizon → Observation)) →
                ((horizon : ℕ) → (Fin horizon → Observation) → PMF Alternative → ℝ → PMF Alternative) →
                  PMF Alternative → ℝ → ℝ

value:
fun {Alternative : Type u_1} {Observation : Type u_2} [Fintype Alternative] [DecidableEq Alternative]
    [Fintype Observation] (averageUtility : Alternative → ℝ) (benchmarkUtility : ℝ)
    (observationLaw : (horizon : ℕ) → PMF (Fin horizon → Observation))
    (method : (horizon : ℕ) → (Fin horizon → Observation) → PMF Alternative → ℝ → PMF Alternative)
    (reference : PMF Alternative) (klBudget : ℝ) =>
  Filter.limsup
    (fun (horizon : ℕ) =>
      benchmarkUtility /
        GolzHaghtalabYang2025Distortion.sourceAlignmentAverageUtility averageUtility observationLaw method reference
          klBudget horizon)
    Filter.atTop
\end{ReviewVerbatim}

\paragraph{Direct retained dependencies}
\begin{itemize}\raggedright
\item \hyperlink{paper-prerequisite-e2a06898e9a1c2db}{Golz\allowbreak{}Haghtalab\allowbreak{}Yang2025\allowbreak{}Distortion.\allowbreak{}source\allowbreak{}Alignment\allowbreak{}Average\allowbreak{}Utility}
\end{itemize}
\paragraph{Recorded source-to-declaration screening}
\textbf{Verdict:} \texttt{matches}
\\
{\raggedright
\textbf{Reason:} This is the fixed-\allowbreak{}instance limsup of the KL-\allowbreak{}ball welfare benchmark over method welfare.\allowbreak{}
\par}
\reviewmatch{Matches source input}
\noindent\textbf{Reviewer annotation}\par
\reviewerbox
\clearpage
\hypertarget{paper-prerequisite-d7683af95771516b}{}
\subsection*{\small\ttfamily\raggedright Golz\allowbreak{}Haghtalab\allowbreak{}Yang2025\allowbreak{}Distortion.\allowbreak{}source\allowbreak{}Social\allowbreak{}Choice\allowbreak{}Average\allowbreak{}Utility}
\noindent\textbf{Paper-local declaration:}\par
{\footnotesize\ttfamily\raggedright Golz\allowbreak{}Haghtalab\allowbreak{}Yang2025\allowbreak{}Distortion.\allowbreak{}source\allowbreak{}Social\allowbreak{}Choice\allowbreak{}Average\allowbreak{}Utility\par}
\textbf{Source locator:} {\footnotesize\raggedright cited publication, lines 202-\allowbreak{}215\par}
\paragraph{Verbatim paper-source connection}
\begin{ReviewVerbatim}
Social Choice Setting. A voting rule f observes the sampled pairwise comparisons {xji [U+001F control character]i
yij }i∈[n],j∈[d] and maps them to a probability distribution over alternatives. For some m, D, β, d, µ,
and
  [U+0002 control character] the correlation between comparisons,     the average utility of f for n samples is AvgUtiln (f ) :=
E AvgUtil f ({xji [U+001F control character]i yij }i,j ) , where the expectation is taken over the pairwise comparisons. The
                               [U+0001 control character][U+0003 control character]

distortion of f on D is the competitive ratio between AvgUtiln (f ) and the optimal average utility
maxx∈A AvgUtil(x) in the limit of n → ∞ samples, and the distortion of f the worst-case distortion
over all D:

            dist(f, D) := lim supn→∞ maxAvgUtil
                                        x∈A AvgUtil(x)
                                                (f )   ,               dist(f ) = supD dist(f, D).
\end{ReviewVerbatim}



\paragraph{Lean-expanded paper semantic target}
\begin{ReviewVerbatim}
type:
{Alternative : Type u_1} →
  {Observation : Type u_2} →
    [Fintype Alternative] →
      [Fintype Observation] →
        (Alternative → ℝ) →
          ((horizon : ℕ) → PMF (Fin horizon → Observation)) →
            ((horizon : ℕ) → (Fin horizon → Observation) → PMF Alternative) → ℕ → ℝ

value:
fun {Alternative : Type u_1} {Observation : Type u_2} [Fintype Alternative] [Fintype Observation]
    (averageUtility : Alternative → ℝ) (observationLaw : (horizon : ℕ) → PMF (Fin horizon → Observation))
    (rule : (horizon : ℕ) → (Fin horizon → Observation) → PMF Alternative) (horizon : ℕ) =>
  AppliedModelingLib.pmfExp (observationLaw horizon) fun (observations : Fin horizon → Observation) =>
    AppliedModelingLib.pmfExp (rule horizon observations) averageUtility
\end{ReviewVerbatim}

\paragraph{Direct retained dependencies}
\begin{itemize}\raggedright
\item \hyperlink{governing-declaration-8e4b4389a4f6454b}{Applied\allowbreak{}Modeling\allowbreak{}Lib.\allowbreak{}pmf\allowbreak{}Exp}
\end{itemize}
\paragraph{Recorded source-to-declaration screening}
\textbf{Verdict:} \texttt{matches}
\\
{\raggedright
\textbf{Reason:} This is the exact expected social-\allowbreak{}choice welfare definition and preserves the source's lack of a false within-\allowbreak{}user comparison-\allowbreak{}independence assumption.\allowbreak{}
\par}
\reviewmatch{Matches source input}
\noindent\textbf{Reviewer annotation}\par
\reviewerbox
\clearpage
\hypertarget{paper-prerequisite-e6860339dbf51649}{}
\subsection*{\small\ttfamily\raggedright Golz\allowbreak{}Haghtalab\allowbreak{}Yang2025\allowbreak{}Distortion.\allowbreak{}source\allowbreak{}Social\allowbreak{}Choice\allowbreak{}Distortion}
\noindent\textbf{Paper-local declaration:}\par
{\footnotesize\ttfamily\raggedright Golz\allowbreak{}Haghtalab\allowbreak{}Yang2025\allowbreak{}Distortion.\allowbreak{}source\allowbreak{}Social\allowbreak{}Choice\allowbreak{}Distortion\par}
\textbf{Source locator:} {\footnotesize\raggedright cited publication, lines 202-\allowbreak{}215\par}
\paragraph{Verbatim paper-source connection}
\begin{ReviewVerbatim}
Social Choice Setting. A voting rule f observes the sampled pairwise comparisons {xji [U+001F control character]i
yij }i∈[n],j∈[d] and maps them to a probability distribution over alternatives. For some m, D, β, d, µ,
and
  [U+0002 control character] the correlation between comparisons,     the average utility of f for n samples is AvgUtiln (f ) :=
E AvgUtil f ({xji [U+001F control character]i yij }i,j ) , where the expectation is taken over the pairwise comparisons. The
                               [U+0001 control character][U+0003 control character]

distortion of f on D is the competitive ratio between AvgUtiln (f ) and the optimal average utility
maxx∈A AvgUtil(x) in the limit of n → ∞ samples, and the distortion of f the worst-case distortion
over all D:

            dist(f, D) := lim supn→∞ maxAvgUtil
                                        x∈A AvgUtil(x)
                                                (f )   ,               dist(f ) = supD dist(f, D).
\end{ReviewVerbatim}



\paragraph{Lean-expanded paper semantic target}
\begin{ReviewVerbatim}
type:
{Alternative : Type u_1} →
  {Observation : Type u_2} →
    [Fintype Alternative] →
      [Fintype Observation] →
        (Alternative → ℝ) →
          ℝ →
            ((horizon : ℕ) → PMF (Fin horizon → Observation)) →
              ((horizon : ℕ) → (Fin horizon → Observation) → PMF Alternative) → ℝ

value:
fun {Alternative : Type u_1} {Observation : Type u_2} [Fintype Alternative] [Fintype Observation]
    (averageUtility : Alternative → ℝ) (optimalAverageUtility : ℝ)
    (observationLaw : (horizon : ℕ) → PMF (Fin horizon → Observation))
    (rule : (horizon : ℕ) → (Fin horizon → Observation) → PMF Alternative) =>
  Filter.limsup
    (fun (horizon : ℕ) =>
      optimalAverageUtility /
        GolzHaghtalabYang2025Distortion.sourceSocialChoiceAverageUtility averageUtility observationLaw rule horizon)
    Filter.atTop
\end{ReviewVerbatim}

\paragraph{Direct retained dependencies}
\begin{itemize}\raggedright
\item \hyperlink{paper-prerequisite-d7683af95771516b}{Golz\allowbreak{}Haghtalab\allowbreak{}Yang2025\allowbreak{}Distortion.\allowbreak{}source\allowbreak{}Social\allowbreak{}Choice\allowbreak{}Average\allowbreak{}Utility}
\end{itemize}
\paragraph{Recorded source-to-declaration screening}
\textbf{Verdict:} \texttt{matches}
\\
{\raggedright
\textbf{Reason:} This is the fixed-\allowbreak{}instance limsup of optimal alternative welfare over expected rule welfare.\allowbreak{}
\par}
\reviewmatch{Matches source input}
\noindent\textbf{Reviewer annotation}\par
\reviewerbox
\clearpage
\hypertarget{paper-prerequisite-a9bfbfb57b04c106}{}
\subsection*{\small\ttfamily\raggedright Golz\allowbreak{}Haghtalab\allowbreak{}Yang2025\allowbreak{}Distortion.\allowbreak{}source\allowbreak{}Worst\allowbreak{}Case\allowbreak{}Alignment\allowbreak{}Distortion}
\noindent\textbf{Paper-local declaration:}\par
{\footnotesize\ttfamily\raggedright Golz\allowbreak{}Haghtalab\allowbreak{}Yang2025\allowbreak{}Distortion.\allowbreak{}source\allowbreak{}Worst\allowbreak{}Case\allowbreak{}Alignment\allowbreak{}Distortion\par}
\textbf{Source locator:} {\footnotesize\raggedright cited publication, lines 281-\allowbreak{}285\par}
\paragraph{Verbatim paper-source connection}
\begin{ReviewVerbatim}
To generalize the definition of distortion to alignment methods, we set the maximum average utility of
any policy in the KL-ball as the benchmark. For fixed m, D, β, d, µ and correlation between pairwise
comparisons, the distortion of alignment method f is
                                                             maxπ∈Bτ (πref ) AvgUtil(π)
                      dist(f ) = supD,πref ,τ lim supn→∞       AvgUtiln (f (·,πref ,τ )) .
\end{ReviewVerbatim}



\paragraph{Lean-expanded paper semantic target}
\begin{ReviewVerbatim}
type:
{Instance : Type u_1} → (Instance → ℝ) → ℝ

value:
fun {Instance : Type u_1} (distortion : Instance → ℝ) => sSup (Set.range distortion)
\end{ReviewVerbatim}


\paragraph{Recorded source-to-declaration screening}
\textbf{Verdict:} \texttt{matches}
\\
{\raggedright
\textbf{Reason:} This takes the supremum over the source alignment-\allowbreak{}instance scope, including utility distribution, reference, and KL budget.\allowbreak{}
\par}
\reviewmatch{Matches source input}
\noindent\textbf{Reviewer annotation}\par
\reviewerbox
\clearpage
\hypertarget{paper-prerequisite-43538553aea8c3ab}{}
\subsection*{\small\ttfamily\raggedright Golz\allowbreak{}Haghtalab\allowbreak{}Yang2025\allowbreak{}Distortion.\allowbreak{}source\allowbreak{}Worst\allowbreak{}Case\allowbreak{}Social\allowbreak{}Choice\allowbreak{}Distortion}
\noindent\textbf{Paper-local declaration:}\par
{\footnotesize\ttfamily\raggedright Golz\allowbreak{}Haghtalab\allowbreak{}Yang2025\allowbreak{}Distortion.\allowbreak{}source\allowbreak{}Worst\allowbreak{}Case\allowbreak{}Social\allowbreak{}Choice\allowbreak{}Distortion\par}
\textbf{Source locator:} {\footnotesize\raggedright cited publication, lines 202-\allowbreak{}215\par}
\paragraph{Verbatim paper-source connection}
\begin{ReviewVerbatim}
Social Choice Setting. A voting rule f observes the sampled pairwise comparisons {xji [U+001F control character]i
yij }i∈[n],j∈[d] and maps them to a probability distribution over alternatives. For some m, D, β, d, µ,
and
  [U+0002 control character] the correlation between comparisons,     the average utility of f for n samples is AvgUtiln (f ) :=
E AvgUtil f ({xji [U+001F control character]i yij }i,j ) , where the expectation is taken over the pairwise comparisons. The
                               [U+0001 control character][U+0003 control character]

distortion of f on D is the competitive ratio between AvgUtiln (f ) and the optimal average utility
maxx∈A AvgUtil(x) in the limit of n → ∞ samples, and the distortion of f the worst-case distortion
over all D:

            dist(f, D) := lim supn→∞ maxAvgUtil
                                        x∈A AvgUtil(x)
                                                (f )   ,               dist(f ) = supD dist(f, D).
\end{ReviewVerbatim}



\paragraph{Lean-expanded paper semantic target}
\begin{ReviewVerbatim}
type:
{Instance : Type u_1} → (Instance → ℝ) → ℝ

value:
fun {Instance : Type u_1} (distortion : Instance → ℝ) => sSup (Set.range distortion)
\end{ReviewVerbatim}


\paragraph{Recorded source-to-declaration screening}
\textbf{Verdict:} \texttt{matches}
\\
{\raggedright
\textbf{Reason:} This takes the supremum over the source social-\allowbreak{}choice instance scope.\allowbreak{}
\par}
\reviewmatch{Matches source input}
\noindent\textbf{Reviewer annotation}\par
\reviewerbox

\clearpage
\hypertarget{source-claim-979ce61abfc9e53f}{}
\hypertarget{source-section-f10b5f68e4acf3dc}{}
\section*{Source claims}
\subsection*{Review row 1}
\textbf{Source locator:} {\footnotesize\raggedright cited publication, lines 338-\allowbreak{}344\par}
\paragraph{Verbatim source input}
\begin{ReviewVerbatim}
Lemma 1 (Linearization of Expected Win-Rates). Let L := σ 0 (0) = 1/4 and `β :=                β       =
 1   1−e−β
2β · 1+e−β . For any pair of alternatives x, y ∈ A, we have

β · (`β · AvgUtil(x) − L · AvgUtil(y)) ≤ p(x [U+001F control character] y) − 12 ≤ β · (L · AvgUtil(x) − `β · AvgUtil(y)) .

3.1   Upper Bound on Borda Distortion

[Next verbatim source excerpt]

Let A = {1, . . . , m} be a finite set of alternatives. The population of users is described by a
probability distribution D over utility vectors u = (u(1), . . . , u(m)), whose entries 0 ≤ u(x) ≤ 1
indicate a user’s utility for alternative x.3 The objective is to find an alternative x such that its average
utility AvgUtil(x) := Eu∼D [u(x)] across the user population (also known as the utilitarian social
welfare) is as high as possible. We extend this notation to probability distributions π over alternatives
by setting AvgUtil(π) := Ex∼π [AvgUtil(x)].
Both voting rules and alignment methods observe comparisons from n users. We model each
i = 1, . . . , n as a fresh user with independently drawn utility vector ui ∼ D. For exposition, we
assume that each user i provides an equal number d ≥ 1 of pairwise comparisons. For each i,
and for j = 1, . . . , d, we independently draw alternatives xji , yij from a fixed distribution µ over

    3
     Our assumption that utilities be in [0, 1] is weaker than any of the three assumptions — unit-sum, unit-range,
or approval [23] — made in classic distortion to avoid trivial infinite lower bounds.


                                                           3

[form-feed in source extraction]
alternatives, in which the minimum probability mass µmin := minx∈A µ(x) is positive.4 User i
then compares each pair {xji , yij } (for j = 1, . . . , d) through a Bradley-Terry model based on i’s
utilities: i prefers xji over yij (written “xji [U+001F control character]i yij ”) with probability σ β · (ui (xji ) − ui (yij )) , where
                                                                                                        [U+0001 control character]

σ(t) := 1/(1+e−t ) is the logistic sigmoid function and β > 0 is a temperature parameter, and prefers
yij over xji (“yij [U+001F control character]i xji ”) otherwise.5 Whereas this specifies the marginal probability of each pairwise
comparison, we make no assumption about the correlation between i’s choices. For example, i might
derive the pairwise comparisons from a Plackett-Luce ranking, ensuring that      [U+0002 control character] the user’s comparisons  [U+0001 control character][U+0003 control character]
are always consistent.6 We set p(x [U+001F control character] y) for the expected win rate Eu∼D σ β · (u(x) − u(y)) .
\end{ReviewVerbatim}



\paragraph{Expanded PaperInterface specification}
\begin{ReviewVerbatim}
∀ {User Alternative : Type} [inst : Fintype User] [inst_1 : DecidableEq User] (population : PMF User)
  (utility : AppliedModelingLib.Alignment.Welfare.FiniteUtilityProfile User Alternative),
  AppliedModelingLib.Alignment.Welfare.UnitIntervalUtilityProfile utility →
    ∀ {btScale : ℝ},
      0 < btScale →
        ∀ (first second : Alternative),
          btScale *
                (AppliedModelingLib.Alignment.Welfare.sigmoidChordSlope btScale *
                    AppliedModelingLib.Alignment.Welfare.populationAverageUtility population utility first -
                  1 / 4 * AppliedModelingLib.Alignment.Welfare.populationAverageUtility population utility second) ≤
              (AppliedModelingLib.Alignment.Welfare.populationBradleyTerryPreference population utility btScale).prob
                  PUnit.unit first second -
                1 / 2 ∧
            (AppliedModelingLib.Alignment.Welfare.populationBradleyTerryPreference population utility btScale).prob
                  PUnit.unit first second -
                1 / 2 ≤
              btScale *
                (1 / 4 * AppliedModelingLib.Alignment.Welfare.populationAverageUtility population utility first -
                  AppliedModelingLib.Alignment.Welfare.sigmoidChordSlope btScale *
                    AppliedModelingLib.Alignment.Welfare.populationAverageUtility population utility second)
\end{ReviewVerbatim}

\paragraph{Retained governing declarations}
\begin{itemize}\raggedright
\item \hyperlink{library-prerequisite-56c0734f8877ddbe}{Applied\allowbreak{}Modeling\allowbreak{}Lib.\allowbreak{}Alignment.\allowbreak{}Welfare.\allowbreak{}Finite\allowbreak{}Utility\allowbreak{}Profile}
\item \hyperlink{library-prerequisite-cbe8afa0d6d1ba78}{Applied\allowbreak{}Modeling\allowbreak{}Lib.\allowbreak{}Alignment.\allowbreak{}Welfare.\allowbreak{}Unit\allowbreak{}Interval\allowbreak{}Utility\allowbreak{}Profile}
\item \hyperlink{library-prerequisite-15f4b442749deab2}{Applied\allowbreak{}Modeling\allowbreak{}Lib.\allowbreak{}Alignment.\allowbreak{}Welfare.\allowbreak{}population\allowbreak{}Average\allowbreak{}Utility}
\item \hyperlink{library-prerequisite-cf06d15d7f8b7c31}{Applied\allowbreak{}Modeling\allowbreak{}Lib.\allowbreak{}Alignment.\allowbreak{}Welfare.\allowbreak{}population\allowbreak{}Bradley\allowbreak{}Terry\allowbreak{}Preference}
\item \hyperlink{governing-declaration-a4d623b8c608b566}{Applied\allowbreak{}Modeling\allowbreak{}Lib.\allowbreak{}Alignment.\allowbreak{}Welfare.\allowbreak{}sigmoid\allowbreak{}Chord\allowbreak{}Slope}
\item \hyperlink{library-prerequisite-7e2db20cf087d9ba}{Applied\allowbreak{}Modeling\allowbreak{}Lib.\allowbreak{}Learning.\allowbreak{}Human\allowbreak{}Feedback.\allowbreak{}Pairwise\allowbreak{}Preference}
\end{itemize}
\paragraph{Lean-elaborated claim roles}\begin{ReviewVerbatim}
1. parameter: Type
2. parameter: Type
3. parameter: Fintype User
4. parameter: DecidableEq User
5. parameter: PMF User
6. parameter: AppliedModelingLib.Alignment.Welfare.FiniteUtilityProfile User Alternative
7. assumption: AppliedModelingLib.Alignment.Welfare.UnitIntervalUtilityProfile utility
8. parameter: ℝ
9. assumption: 0 < btScale
10. parameter: Alternative
11. parameter: Alternative
12. conclusion: btScale *
      (AppliedModelingLib.Alignment.Welfare.sigmoidChordSlope btScale *
          AppliedModelingLib.Alignment.Welfare.populationAverageUtility population utility first -
        1 / 4 * AppliedModelingLib.Alignment.Welfare.populationAverageUtility population utility second) ≤
    (AppliedModelingLib.Alignment.Welfare.populationBradleyTerryPreference population utility btScale).prob PUnit.unit
        first second -
      1 / 2 ∧
  (AppliedModelingLib.Alignment.Welfare.populationBradleyTerryPreference population utility btScale).prob PUnit.unit
        first second -
      1 / 2 ≤
    btScale *
      (1 / 4 * AppliedModelingLib.Alignment.Welfare.populationAverageUtility population utility first -
        AppliedModelingLib.Alignment.Welfare.sigmoidChordSlope btScale *
          AppliedModelingLib.Alignment.Welfare.populationAverageUtility population utility second)
\end{ReviewVerbatim}

\paragraph{Lean proof endpoint (built separately)}\begin{ReviewVerbatim}
GolzHaghtalabYang2025Distortion.lemma1_linearization
\end{ReviewVerbatim}

\paragraph{Recorded source-to-Spec screening}
\textbf{Verdict:} \texttt{matches}
\\
{\raggedright
\textbf{Reason:} This is the source two-\allowbreak{}sided Bradley-\allowbreak{}Terry linearization using L = 1/\allowbreak{}4 and the source chord coefficient.\allowbreak{}
\par}
\reviewmatch{Matches source input}
\noindent\textbf{Reviewer annotation}\par
\reviewerbox
\clearpage
\hypertarget{source-claim-43099ab484dcdf43}{}
\section*{Review row 2}
\textbf{Source locator:} {\footnotesize\raggedright cited publication, lines 1358-\allowbreak{}1370\par}
\paragraph{Verbatim source input}
\begin{ReviewVerbatim}
Lemma 10. For any instance D with any number of alternatives m, any distribution µ over alter-
natives with µmin = minx∈A µ(x), and n i.i.d. users sampled from D where each user labels d
comparison pairs following the Bradley-Terry[U+0011 control character] model with temperature β, we have that with probability
                                    ndµ2min
                                [U+0010 control character]
                          2                                                                #(x[U+001F control character]y)
at least 1 − δ where δ ≥ m exp − 8           , the empirical win rates pn (x [U+001F control character] y) := #(x[U+001F control character]y)+#(y[U+001F control character]x)
satisfies that:
                                                     s                                      !
                                                              log(m/δ)            log(m/δ)
        ∀x, y ∈ A, |pn (x [U+001F control character] y) − p(x [U+001F control character] y)| ≤ O                                  +             .
                                                         n · min{1, d · µ2min }     nµ2min


[Next verbatim source excerpt]

Let A = {1, . . . , m} be a finite set of alternatives. The population of users is described by a
probability distribution D over utility vectors u = (u(1), . . . , u(m)), whose entries 0 ≤ u(x) ≤ 1
indicate a user’s utility for alternative x.3 The objective is to find an alternative x such that its average
utility AvgUtil(x) := Eu∼D [u(x)] across the user population (also known as the utilitarian social
welfare) is as high as possible. We extend this notation to probability distributions π over alternatives
by setting AvgUtil(π) := Ex∼π [AvgUtil(x)].
Both voting rules and alignment methods observe comparisons from n users. We model each
i = 1, . . . , n as a fresh user with independently drawn utility vector ui ∼ D. For exposition, we
assume that each user i provides an equal number d ≥ 1 of pairwise comparisons. For each i,
and for j = 1, . . . , d, we independently draw alternatives xji , yij from a fixed distribution µ over

    3
     Our assumption that utilities be in [0, 1] is weaker than any of the three assumptions — unit-sum, unit-range,
or approval [23] — made in classic distortion to avoid trivial infinite lower bounds.


                                                           3

[form-feed in source extraction]
alternatives, in which the minimum probability mass µmin := minx∈A µ(x) is positive.4 User i
then compares each pair {xji , yij } (for j = 1, . . . , d) through a Bradley-Terry model based on i’s
utilities: i prefers xji over yij (written “xji [U+001F control character]i yij ”) with probability σ β · (ui (xji ) − ui (yij )) , where
                                                                                                        [U+0001 control character]

σ(t) := 1/(1+e−t ) is the logistic sigmoid function and β > 0 is a temperature parameter, and prefers
yij over xji (“yij [U+001F control character]i xji ”) otherwise.5 Whereas this specifies the marginal probability of each pairwise
comparison, we make no assumption about the correlation between i’s choices. For example, i might
derive the pairwise comparisons from a Plackett-Luce ranking, ensuring that      [U+0002 control character] the user’s comparisons  [U+0001 control character][U+0003 control character]
are always consistent.6 We set p(x [U+001F control character] y) for the expected win rate Eu∼D σ β · (u(x) − u(y)) .
\end{ReviewVerbatim}



\paragraph{Expanded PaperInterface specification}
\begin{ReviewVerbatim}
∀ {Alternative : Type} [inst : Fintype Alternative] [inst_1 : DecidableEq Alternative] [Nonempty Alternative]
  {responseLaw : PMF (GolzHaghtalabYang2025Distortion.theorem2UserResponseTable Alternative)}
  {preference : AppliedModelingLib.Learning.HumanFeedback.PairwisePreference PUnit.{1} Alternative},
  GolzHaghtalabYang2025Distortion.theorem2UserResponseCalibrated responseLaw preference →
    ∀ (sampling : PMF Alternative) (users comparisonsPerUser : ℕ) (minimumMass delta : ℝ),
      0 < users →
        0 < comparisonsPerUser →
          0 < minimumMass →
            (∀ (alternative : Alternative), minimumMass ≤ (sampling alternative).toReal) →
              0 < delta →
                delta ≤ 1 →
                  ↑(Fintype.card Alternative) ^ 2 * Real.exp (-↑(users * comparisonsPerUser) * minimumMass ^ 2 / 8) ≤
                      delta →
                    have failureBudget : ℝ := delta / ↑(Fintype.card Alternative) ^ 2;
                    (AppliedModelingLib.pmfProb
                        (AppliedModelingLib.pmfProduct (Fin users)
                          (GolzHaghtalabYang2025Distortion.theorem2UserResponseTable Alternative ×
                            (Fin comparisonsPerUser → Alternative × Alternative))
                          (AppliedModelingLib.pmfProd responseLaw
                            (AppliedModelingLib.pmfProduct (Fin comparisonsPerUser) (Alternative × Alternative)
                              (GolzHaghtalabYang2025Distortion.theorem2UserPairLabelLaw sampling))))
                        fun
                          (sample :
                            Fin users →
                              GolzHaghtalabYang2025Distortion.theorem2UserResponseTable Alternative ×
                                (Fin comparisonsPerUser → Alternative × Alternative)) =>
                        ∃ (pair : Alternative × Alternative),
                          GolzHaghtalabYang2025Distortion.theorem10SourceRate users comparisonsPerUser minimumMass
                              failureBudget <
                            |GolzHaghtalabYang2025Distortion.theorem10IidUserPaperWinRate pair.1 pair.2 sample -
                                preference.prob PUnit.unit pair.1 pair.2|) ≤
                      delta
\end{ReviewVerbatim}

\paragraph{Retained governing declarations}
\begin{itemize}\raggedright
\item \hyperlink{library-prerequisite-7e2db20cf087d9ba}{Applied\allowbreak{}Modeling\allowbreak{}Lib.\allowbreak{}Learning.\allowbreak{}Human\allowbreak{}Feedback.\allowbreak{}Pairwise\allowbreak{}Preference}
\item \hyperlink{governing-declaration-9c38d53b9d664221}{Applied\allowbreak{}Modeling\allowbreak{}Lib.\allowbreak{}pmf\allowbreak{}Prob}
\item \hyperlink{library-prerequisite-de8e2a8bfe221a38}{Applied\allowbreak{}Modeling\allowbreak{}Lib.\allowbreak{}pmf\allowbreak{}Prod}
\item \hyperlink{governing-declaration-5ac12d0a52b69068}{Applied\allowbreak{}Modeling\allowbreak{}Lib.\allowbreak{}pmf\allowbreak{}Product}
\item \hyperlink{governing-declaration-a1441da0ef6b2fa2}{Golz\allowbreak{}Haghtalab\allowbreak{}Yang2025\allowbreak{}Distortion.\allowbreak{}theorem10\allowbreak{}Iid\allowbreak{}User\allowbreak{}Paper\allowbreak{}Win\allowbreak{}Rate}
\item \hyperlink{governing-declaration-aaf6dc2063f9888f}{Golz\allowbreak{}Haghtalab\allowbreak{}Yang2025\allowbreak{}Distortion.\allowbreak{}theorem10\allowbreak{}Source\allowbreak{}Rate}
\item \hyperlink{governing-declaration-98c3ee9a11a59ef9}{Golz\allowbreak{}Haghtalab\allowbreak{}Yang2025\allowbreak{}Distortion.\allowbreak{}theorem2\allowbreak{}User\allowbreak{}Pair\allowbreak{}Label\allowbreak{}Law}
\item \hyperlink{governing-declaration-080b9388e51a6018}{Golz\allowbreak{}Haghtalab\allowbreak{}Yang2025\allowbreak{}Distortion.\allowbreak{}theorem2\allowbreak{}User\allowbreak{}Response\allowbreak{}Calibrated}
\item \hyperlink{governing-declaration-c5d0273cbd7ddcc9}{Golz\allowbreak{}Haghtalab\allowbreak{}Yang2025\allowbreak{}Distortion.\allowbreak{}theorem2\allowbreak{}User\allowbreak{}Response\allowbreak{}Table}
\end{itemize}
\paragraph{Lean-elaborated claim roles}\begin{ReviewVerbatim}
1. parameter: Type
2. parameter: Fintype Alternative
3. parameter: DecidableEq Alternative
4. assumption: Nonempty Alternative
5. parameter: PMF (GolzHaghtalabYang2025Distortion.theorem2UserResponseTable Alternative)
6. parameter: AppliedModelingLib.Learning.HumanFeedback.PairwisePreference PUnit.{1} Alternative
7. assumption: GolzHaghtalabYang2025Distortion.theorem2UserResponseCalibrated responseLaw preference
8. parameter: PMF Alternative
9. parameter: ℕ
10. parameter: ℕ
11. parameter: ℝ
12. parameter: ℝ
13. assumption: 0 < users
14. assumption: 0 < comparisonsPerUser
15. assumption: 0 < minimumMass
16. assumption: ∀ (alternative : Alternative), minimumMass ≤ (sampling alternative).toReal
17. assumption: 0 < delta
18. assumption: delta ≤ 1
19. assumption: ↑(Fintype.card Alternative) ^ 2 * Real.exp (-↑(users * comparisonsPerUser) * minimumMass ^ 2 / 8) ≤ delta
20. conclusion: (AppliedModelingLib.pmfProb
    (AppliedModelingLib.pmfProduct (Fin users)
      (GolzHaghtalabYang2025Distortion.theorem2UserResponseTable Alternative ×
        (Fin comparisonsPerUser → Alternative × Alternative))
      (AppliedModelingLib.pmfProd responseLaw
        (AppliedModelingLib.pmfProduct (Fin comparisonsPerUser) (Alternative × Alternative)
          (GolzHaghtalabYang2025Distortion.theorem2UserPairLabelLaw sampling))))
    fun
      (sample :
        Fin users →
          GolzHaghtalabYang2025Distortion.theorem2UserResponseTable Alternative ×
            (Fin comparisonsPerUser → Alternative × Alternative)) =>
    ∃ (pair : Alternative × Alternative),
      GolzHaghtalabYang2025Distortion.theorem10SourceRate users comparisonsPerUser minimumMass
          (delta / ↑(Fintype.card Alternative) ^ 2) <
        |GolzHaghtalabYang2025Distortion.theorem10IidUserPaperWinRate pair.1 pair.2 sample -
            preference.prob PUnit.unit pair.1 pair.2|) ≤
  delta
\end{ReviewVerbatim}

\paragraph{Lean proof endpoint (built separately)}\begin{ReviewVerbatim}
GolzHaghtalabYang2025Distortion.lemma10_empirical_win_rate_concentration
\end{ReviewVerbatim}

\paragraph{Recorded source-to-Spec screening}
\textbf{Verdict:} \texttt{matches}
\\
{\raggedright
\textbf{Reason:} This is the source uniform empirical pairwise-\allowbreak{}rate concentration statement under iid users, positive sampling mass, and its stated exponential confidence regime.\allowbreak{}
\par}
\reviewmatch{Matches source input}
\noindent\textbf{Reviewer annotation}\par
\reviewerbox
\clearpage
\hypertarget{source-claim-069675611a64cc9c}{}
\section*{Review row 3}
\textbf{Source locator:} {\footnotesize\raggedright cited publication, lines 1464-\allowbreak{}1478\par}
\paragraph{Verbatim source input}
\begin{ReviewVerbatim}
Lemma 11. For any instance D with any number of alternatives m, any distribution µ over alter-
natives with µmin = minx∈A µ(x), and n i.i.d. users sampled from D where each user labels d
comparison pairs, the normalized Borda score BCn (x) of any alternative x ∈ A satisfies that with
probability at least 1 − δ where δ ≥ 2m exp(− ndµ8min ),
                                              s                                      !
                                                        log(m/δ)          m log(m/δ)
        ∀x ∈ A, |BCn (x) − BC? (x)| ≤ O                                 +              ,      (5)
                                                 n · min{1, d · µ2min }      nµmin

where BC? (x) is the limiting normalized Borda score of candidate x, defined as
                                X                              1        X
                 BC? (x) :=          µ(y) · p(x [U+001F control character] y) =           µ(x) +   µ(y) · p(x [U+001F control character] y).                     (6)
                                                               2
                               y∈A                                     y6=x


[Next verbatim source excerpt]

Let A = {1, . . . , m} be a finite set of alternatives. The population of users is described by a
probability distribution D over utility vectors u = (u(1), . . . , u(m)), whose entries 0 ≤ u(x) ≤ 1
indicate a user’s utility for alternative x.3 The objective is to find an alternative x such that its average
utility AvgUtil(x) := Eu∼D [u(x)] across the user population (also known as the utilitarian social
welfare) is as high as possible. We extend this notation to probability distributions π over alternatives
by setting AvgUtil(π) := Ex∼π [AvgUtil(x)].
Both voting rules and alignment methods observe comparisons from n users. We model each
i = 1, . . . , n as a fresh user with independently drawn utility vector ui ∼ D. For exposition, we
assume that each user i provides an equal number d ≥ 1 of pairwise comparisons. For each i,
and for j = 1, . . . , d, we independently draw alternatives xji , yij from a fixed distribution µ over

    3
     Our assumption that utilities be in [0, 1] is weaker than any of the three assumptions — unit-sum, unit-range,
or approval [23] — made in classic distortion to avoid trivial infinite lower bounds.


                                                           3

[form-feed in source extraction]
alternatives, in which the minimum probability mass µmin := minx∈A µ(x) is positive.4 User i
then compares each pair {xji , yij } (for j = 1, . . . , d) through a Bradley-Terry model based on i’s
utilities: i prefers xji over yij (written “xji [U+001F control character]i yij ”) with probability σ β · (ui (xji ) − ui (yij )) , where
                                                                                                        [U+0001 control character]

σ(t) := 1/(1+e−t ) is the logistic sigmoid function and β > 0 is a temperature parameter, and prefers
yij over xji (“yij [U+001F control character]i xji ”) otherwise.5 Whereas this specifies the marginal probability of each pairwise
comparison, we make no assumption about the correlation between i’s choices. For example, i might
derive the pairwise comparisons from a Plackett-Luce ranking, ensuring that      [U+0002 control character] the user’s comparisons  [U+0001 control character][U+0003 control character]
are always consistent.6 We set p(x [U+001F control character] y) for the expected win rate Eu∼D σ β · (u(x) − u(y)) .

[Next verbatim source excerpt]

Social Choice Setting. A voting rule f observes the sampled pairwise comparisons {xji [U+001F control character]i
yij }i∈[n],j∈[d] and maps them to a probability distribution over alternatives. For some m, D, β, d, µ,
and
  [U+0002 control character] the correlation between comparisons,     the average utility of f for n samples is AvgUtiln (f ) :=
E AvgUtil f ({xji [U+001F control character]i yij }i,j ) , where the expectation is taken over the pairwise comparisons. The
                               [U+0001 control character][U+0003 control character]

distortion of f on D is the competitive ratio between AvgUtiln (f ) and the optimal average utility
maxx∈A AvgUtil(x) in the limit of n → ∞ samples, and the distortion of f the worst-case distortion
over all D:

            dist(f, D) := lim supn→∞ maxAvgUtil
                                        x∈A AvgUtil(x)
                                                (f )   ,               dist(f ) = supD dist(f, D).
                                                       n



For alternatives x, y, let #(x [U+001F control character] y) := {(i, j) | xji = x, yij = y} denote the number of pairwise
comparisons in which x beat y. The (normalized) Borda score [45] of alternative x is
                                                       P
                                                         y∈A #(x[U+001F control character]y)
                                   BC(x) := P         #(x[U+001F control character]y)+
                                                              P             ,
                                                  y∈A           y6=x #(y[U+001F control character]x)


i.e., the fraction of pairwise comparisons involving x in which it wins. The Borda voting rule chooses
the winner uniformly among all alternatives with maximum Borda score. The Maximal Lotteries vot-
ing rule first computes the margin matrix M ∈ Rm×m , where Mx,y = #(x[U+001F control character]y)−#(y[U+001F control character]x)
                                                                          #(x[U+001F control character]y)+#(y[U+001F control character]x) . It then con-
siders a symmetric two-player zero-sum game in which player 1 selects alternative x1 , player 2 selects
alternative x2 , and the payoffs are Mx1 ,x2 for player 1 and Mx2 ,x1 = −Mx1 ,x2 for player 2. The max-
imal lotteries rule returns a distribution π ∈ argmaxπ1 ∈∆(A) minπ2 ∈∆(A) Ex1 ∼π1 ,x2 ∼π2 [Mx1 ,x2 ],
i.e., a mixed strategy in Nash equilibrium. When several such π exist, our results hold for any choice.
\end{ReviewVerbatim}



\paragraph{Expanded PaperInterface specification}
\begin{ReviewVerbatim}
∀ {Alternative : Type} [inst : Fintype Alternative] [inst_1 : DecidableEq Alternative] [Nonempty Alternative]
  (responseLaw : PMF (GolzHaghtalabYang2025Distortion.theorem2UserResponseTable Alternative))
  (sampling : PMF Alternative)
  (preference : AppliedModelingLib.Learning.HumanFeedback.PairwisePreference PUnit.{1} Alternative)
  (users comparisonsPerUser : ℕ),
  0 < users →
    0 < comparisonsPerUser →
      GolzHaghtalabYang2025Distortion.theorem2UserResponseCalibrated responseLaw preference →
        ∀ (minimumMass delta : ℝ),
          0 < minimumMass →
            (∀ (alternative : Alternative), minimumMass ≤ (sampling alternative).toReal) →
              0 < delta →
                delta ≤ 1 →
                  2 * ↑(Fintype.card Alternative) * Real.exp (-↑(users * comparisonsPerUser) * minimumMass / 8) ≤
                      delta →
                    (AppliedModelingLib.pmfProb
                        (AppliedModelingLib.pmfProduct (Fin users)
                          (GolzHaghtalabYang2025Distortion.theorem2UserResponseTable Alternative ×
                            (Fin comparisonsPerUser → Alternative × Alternative))
                          (AppliedModelingLib.pmfProd responseLaw
                            (AppliedModelingLib.pmfProduct (Fin comparisonsPerUser) (Alternative × Alternative)
                              (GolzHaghtalabYang2025Distortion.theorem2UserPairLabelLaw sampling))))
                        fun
                          (sample :
                            Fin users →
                              GolzHaghtalabYang2025Distortion.theorem2UserResponseTable Alternative ×
                                (Fin comparisonsPerUser → Alternative × Alternative)) =>
                        have logLevel : ℝ := Real.log (8 * ↑(Fintype.card Alternative) ^ 2 / delta);
                        ∃ (alternative : Alternative),
                          20 * √(↑users * logLevel / min 1 (↑comparisonsPerUser * minimumMass ^ 2)) / ↑users +
                              8 * (↑(Fintype.card Alternative) + 1) * logLevel / (↑users * minimumMass) <
                            |GolzHaghtalabYang2025Distortion.theorem11IidUserLatentBordaScore alternative sample -
                                AppliedModelingLib.Alignment.Welfare.pairwiseBordaScore sampling preference
                                  alternative|) ≤
                      delta
\end{ReviewVerbatim}

\paragraph{Retained governing declarations}
\begin{itemize}\raggedright
\item \hyperlink{governing-declaration-3808cf59ea39b435}{Applied\allowbreak{}Modeling\allowbreak{}Lib.\allowbreak{}Alignment.\allowbreak{}Welfare.\allowbreak{}pairwise\allowbreak{}Borda\allowbreak{}Score}
\item \hyperlink{library-prerequisite-7e2db20cf087d9ba}{Applied\allowbreak{}Modeling\allowbreak{}Lib.\allowbreak{}Learning.\allowbreak{}Human\allowbreak{}Feedback.\allowbreak{}Pairwise\allowbreak{}Preference}
\item \hyperlink{governing-declaration-9c38d53b9d664221}{Applied\allowbreak{}Modeling\allowbreak{}Lib.\allowbreak{}pmf\allowbreak{}Prob}
\item \hyperlink{library-prerequisite-de8e2a8bfe221a38}{Applied\allowbreak{}Modeling\allowbreak{}Lib.\allowbreak{}pmf\allowbreak{}Prod}
\item \hyperlink{governing-declaration-5ac12d0a52b69068}{Applied\allowbreak{}Modeling\allowbreak{}Lib.\allowbreak{}pmf\allowbreak{}Product}
\item \hyperlink{governing-declaration-7506044331f97cc9}{Golz\allowbreak{}Haghtalab\allowbreak{}Yang2025\allowbreak{}Distortion.\allowbreak{}theorem11\allowbreak{}Iid\allowbreak{}User\allowbreak{}Latent\allowbreak{}Borda\allowbreak{}Score}
\item \hyperlink{governing-declaration-98c3ee9a11a59ef9}{Golz\allowbreak{}Haghtalab\allowbreak{}Yang2025\allowbreak{}Distortion.\allowbreak{}theorem2\allowbreak{}User\allowbreak{}Pair\allowbreak{}Label\allowbreak{}Law}
\item \hyperlink{governing-declaration-080b9388e51a6018}{Golz\allowbreak{}Haghtalab\allowbreak{}Yang2025\allowbreak{}Distortion.\allowbreak{}theorem2\allowbreak{}User\allowbreak{}Response\allowbreak{}Calibrated}
\item \hyperlink{governing-declaration-c5d0273cbd7ddcc9}{Golz\allowbreak{}Haghtalab\allowbreak{}Yang2025\allowbreak{}Distortion.\allowbreak{}theorem2\allowbreak{}User\allowbreak{}Response\allowbreak{}Table}
\end{itemize}
\paragraph{Lean-elaborated claim roles}\begin{ReviewVerbatim}
1. parameter: Type
2. parameter: Fintype Alternative
3. parameter: DecidableEq Alternative
4. assumption: Nonempty Alternative
5. parameter: PMF (GolzHaghtalabYang2025Distortion.theorem2UserResponseTable Alternative)
6. parameter: PMF Alternative
7. parameter: AppliedModelingLib.Learning.HumanFeedback.PairwisePreference PUnit.{1} Alternative
8. parameter: ℕ
9. parameter: ℕ
10. assumption: 0 < users
11. assumption: 0 < comparisonsPerUser
12. assumption: GolzHaghtalabYang2025Distortion.theorem2UserResponseCalibrated responseLaw preference
13. parameter: ℝ
14. parameter: ℝ
15. assumption: 0 < minimumMass
16. assumption: ∀ (alternative : Alternative), minimumMass ≤ (sampling alternative).toReal
17. assumption: 0 < delta
18. assumption: delta ≤ 1
19. assumption: 2 * ↑(Fintype.card Alternative) * Real.exp (-↑(users * comparisonsPerUser) * minimumMass / 8) ≤ delta
20. conclusion: (AppliedModelingLib.pmfProb
    (AppliedModelingLib.pmfProduct (Fin users)
      (GolzHaghtalabYang2025Distortion.theorem2UserResponseTable Alternative ×
        (Fin comparisonsPerUser → Alternative × Alternative))
      (AppliedModelingLib.pmfProd responseLaw
        (AppliedModelingLib.pmfProduct (Fin comparisonsPerUser) (Alternative × Alternative)
          (GolzHaghtalabYang2025Distortion.theorem2UserPairLabelLaw sampling))))
    fun
      (sample :
        Fin users →
          GolzHaghtalabYang2025Distortion.theorem2UserResponseTable Alternative ×
            (Fin comparisonsPerUser → Alternative × Alternative)) =>
    have logLevel : ℝ := Real.log (8 * ↑(Fintype.card Alternative) ^ 2 / delta);
    ∃ (alternative : Alternative),
      20 * √(↑users * logLevel / min 1 (↑comparisonsPerUser * minimumMass ^ 2)) / ↑users +
          8 * (↑(Fintype.card Alternative) + 1) * logLevel / (↑users * minimumMass) <
        |GolzHaghtalabYang2025Distortion.theorem11IidUserLatentBordaScore alternative sample -
            AppliedModelingLib.Alignment.Welfare.pairwiseBordaScore sampling preference alternative|) ≤
  delta
\end{ReviewVerbatim}

\paragraph{Lean proof endpoint (built separately)}\begin{ReviewVerbatim}
GolzHaghtalabYang2025Distortion.lemma11_borda_concentration
\end{ReviewVerbatim}

\paragraph{Recorded source-to-Spec screening}
\textbf{Verdict:} \texttt{matches}
\\
{\raggedright
\textbf{Reason:} This retains the simultaneous normalized-\allowbreak{}Borda concentration claim with the stated rate dependence and confidence regime.\allowbreak{}
\par}
\reviewmatch{Matches source input}
\noindent\textbf{Reviewer annotation}\par
\reviewerbox
\clearpage
\hypertarget{source-claim-99f78f35feb19dfc}{}
\section*{Review row 4}
\textbf{Source locator:} {\footnotesize\raggedright cited publication, lines 2312-\allowbreak{}2324\par}
\paragraph{Verbatim source input}
\begin{ReviewVerbatim}
Lemma 15. For any β > 0, there is an infinite sequence a1 , a2 , . . . of alternatives, and a distribution
  D of utility functions over these alternatives such that

                • AvgUtil(a1 ) = 1/3,
                                                                            [U+0010 control character]
                • for all t ≥ 2, 0 < AvgUtil(at ) ≤ AvgUtil(at−1 ) − 32β log 1 + tanh β/4 ·
                                [U+0001 control character]3 [U+0011 control character]
                  AvgUtil(at−1 )     < AvgUtil(at−1 ), and

                • for all t ≥ 2, p(at [U+001F control character] at−1 ) > 1/2.

  Proof. Our population D will be a uniform distribution over three utility vectors, u1 , u2 , and u3 . We
  define the sequence of alternatives and prove the claim by induction over t ≥ 1.

[Next verbatim source excerpt]

Let A = {1, . . . , m} be a finite set of alternatives. The population of users is described by a
probability distribution D over utility vectors u = (u(1), . . . , u(m)), whose entries 0 ≤ u(x) ≤ 1
indicate a user’s utility for alternative x.3 The objective is to find an alternative x such that its average
utility AvgUtil(x) := Eu∼D [u(x)] across the user population (also known as the utilitarian social
welfare) is as high as possible. We extend this notation to probability distributions π over alternatives
by setting AvgUtil(π) := Ex∼π [AvgUtil(x)].
Both voting rules and alignment methods observe comparisons from n users. We model each
i = 1, . . . , n as a fresh user with independently drawn utility vector ui ∼ D. For exposition, we
assume that each user i provides an equal number d ≥ 1 of pairwise comparisons. For each i,
and for j = 1, . . . , d, we independently draw alternatives xji , yij from a fixed distribution µ over

    3
     Our assumption that utilities be in [0, 1] is weaker than any of the three assumptions — unit-sum, unit-range,
or approval [23] — made in classic distortion to avoid trivial infinite lower bounds.


                                                           3

[form-feed in source extraction]
alternatives, in which the minimum probability mass µmin := minx∈A µ(x) is positive.4 User i
then compares each pair {xji , yij } (for j = 1, . . . , d) through a Bradley-Terry model based on i’s
utilities: i prefers xji over yij (written “xji [U+001F control character]i yij ”) with probability σ β · (ui (xji ) − ui (yij )) , where
                                                                                                        [U+0001 control character]

σ(t) := 1/(1+e−t ) is the logistic sigmoid function and β > 0 is a temperature parameter, and prefers
yij over xji (“yij [U+001F control character]i xji ”) otherwise.5 Whereas this specifies the marginal probability of each pairwise
comparison, we make no assumption about the correlation between i’s choices. For example, i might
derive the pairwise comparisons from a Plackett-Luce ranking, ensuring that      [U+0002 control character] the user’s comparisons  [U+0001 control character][U+0003 control character]
are always consistent.6 We set p(x [U+001F control character] y) for the expected win rate Eu∼D σ β · (u(x) − u(y)) .
\end{ReviewVerbatim}



\paragraph{Expanded PaperInterface specification}
\begin{ReviewVerbatim}
∀ {beta : ℝ},
  0 < beta →
    ∃ (population : PMF (Fin 3)) (utility : AppliedModelingLib.Alignment.Welfare.FiniteUtilityProfile (Fin 3) ℕ),
      AppliedModelingLib.Alignment.Welfare.UnitIntervalUtilityProfile utility ∧
        AppliedModelingLib.Alignment.Welfare.populationAverageUtility population utility 0 = 1 / 3 ∧
          ∀ (n : ℕ),
            0 < AppliedModelingLib.Alignment.Welfare.populationAverageUtility population utility (n + 1) ∧
              AppliedModelingLib.Alignment.Welfare.populationAverageUtility population utility (n + 1) ≤
                  AppliedModelingLib.Alignment.Welfare.populationAverageUtility population utility n -
                    2 / (3 * beta) *
                      Real.log
                        (1 +
                          Real.tanh
                              (beta *
                                  AppliedModelingLib.Alignment.Welfare.populationAverageUtility population utility n /
                                4) ^
                            3) ∧
                AppliedModelingLib.Alignment.Welfare.populationAverageUtility population utility (n + 1) <
                    AppliedModelingLib.Alignment.Welfare.populationAverageUtility population utility n ∧
                  1 / 2 <
                    (AppliedModelingLib.Alignment.Welfare.populationBradleyTerryPreference population utility beta).prob
                      PUnit.unit (n + 1) n
\end{ReviewVerbatim}

\paragraph{Retained governing declarations}
\begin{itemize}\raggedright
\item \hyperlink{library-prerequisite-56c0734f8877ddbe}{Applied\allowbreak{}Modeling\allowbreak{}Lib.\allowbreak{}Alignment.\allowbreak{}Welfare.\allowbreak{}Finite\allowbreak{}Utility\allowbreak{}Profile}
\item \hyperlink{library-prerequisite-cbe8afa0d6d1ba78}{Applied\allowbreak{}Modeling\allowbreak{}Lib.\allowbreak{}Alignment.\allowbreak{}Welfare.\allowbreak{}Unit\allowbreak{}Interval\allowbreak{}Utility\allowbreak{}Profile}
\item \hyperlink{library-prerequisite-15f4b442749deab2}{Applied\allowbreak{}Modeling\allowbreak{}Lib.\allowbreak{}Alignment.\allowbreak{}Welfare.\allowbreak{}population\allowbreak{}Average\allowbreak{}Utility}
\item \hyperlink{library-prerequisite-cf06d15d7f8b7c31}{Applied\allowbreak{}Modeling\allowbreak{}Lib.\allowbreak{}Alignment.\allowbreak{}Welfare.\allowbreak{}population\allowbreak{}Bradley\allowbreak{}Terry\allowbreak{}Preference}
\item \hyperlink{library-prerequisite-7e2db20cf087d9ba}{Applied\allowbreak{}Modeling\allowbreak{}Lib.\allowbreak{}Learning.\allowbreak{}Human\allowbreak{}Feedback.\allowbreak{}Pairwise\allowbreak{}Preference}
\end{itemize}
\paragraph{Lean-elaborated claim roles}\begin{ReviewVerbatim}
1. parameter: ℝ
2. assumption: 0 < beta
3. conclusion: ∃ (population : PMF (Fin 3)) (utility : AppliedModelingLib.Alignment.Welfare.FiniteUtilityProfile (Fin 3) ℕ),
  AppliedModelingLib.Alignment.Welfare.UnitIntervalUtilityProfile utility ∧
    AppliedModelingLib.Alignment.Welfare.populationAverageUtility population utility 0 = 1 / 3 ∧
      ∀ (n : ℕ),
        0 < AppliedModelingLib.Alignment.Welfare.populationAverageUtility population utility (n + 1) ∧
          AppliedModelingLib.Alignment.Welfare.populationAverageUtility population utility (n + 1) ≤
              AppliedModelingLib.Alignment.Welfare.populationAverageUtility population utility n -
                2 / (3 * beta) *
                  Real.log
                    (1 +
                      Real.tanh
                          (beta * AppliedModelingLib.Alignment.Welfare.populationAverageUtility population utility n /
                            4) ^
                        3) ∧
            AppliedModelingLib.Alignment.Welfare.populationAverageUtility population utility (n + 1) <
                AppliedModelingLib.Alignment.Welfare.populationAverageUtility population utility n ∧
              1 / 2 <
                (AppliedModelingLib.Alignment.Welfare.populationBradleyTerryPreference population utility beta).prob
                  PUnit.unit (n + 1) n
\end{ReviewVerbatim}

\paragraph{Lean proof endpoint (built separately)}\begin{ReviewVerbatim}
GolzHaghtalabYang2025Distortion.lemma15_infinite_sequence
\end{ReviewVerbatim}

\paragraph{Recorded source-to-Spec screening}
\textbf{Verdict:} \texttt{matches}
\\
{\raggedright
\textbf{Reason:} The infinite sequence, initial welfare one-\allowbreak{}third, stated welfare decay, and adjacent win-\allowbreak{}rate strictly above one-\allowbreak{}half are all retained.\allowbreak{}
\par}
\reviewmatch{Matches source input}
\noindent\textbf{Reviewer annotation}\par
\reviewerbox
\clearpage
\hypertarget{source-claim-c85b480e847e3345}{}
\section*{Review row 5}
\textbf{Source locator:} {\footnotesize\raggedright cited publication, lines 352-\allowbreak{}372\par}
\paragraph{Verbatim source input}
\begin{ReviewVerbatim}
Theorem 2 (Borda Distortion Upper Bound). For any instance D with any number of alternatives m,
                                                                                       1+e−β 2
distribution µ over alternatives, and temperature β, Borda has at most distortion β2 · 1−e
                                                                                               [U+0001 control character]
                                                                                           −β     =
O(β 2 ). In the finite-sample regime, we have that
                                −β [U+0001 control character]2
                                                                   [U+0010 control character] q                              [U+0011 control character]
                                                                           log(mnβ)     m log(mnβ)
AvgUtiln (Borda) ≥ β2 · 1−e 1+e −β    · max x
                                                            ?
                                              ? ∈A AvgUtil(x ) − O
                                                                    β
                                                                     1
                                                                      2 n·min{1,dµ2 }
                                                                                      +      2
                                                                                         n·β µ  2     .
                                                                                   min             min


                                                                               for n → ∞, assuming

[Next verbatim source excerpt]

Let A = {1, . . . , m} be a finite set of alternatives. The population of users is described by a
probability distribution D over utility vectors u = (u(1), . . . , u(m)), whose entries 0 ≤ u(x) ≤ 1
indicate a user’s utility for alternative x.3 The objective is to find an alternative x such that its average
utility AvgUtil(x) := Eu∼D [u(x)] across the user population (also known as the utilitarian social
welfare) is as high as possible. We extend this notation to probability distributions π over alternatives
by setting AvgUtil(π) := Ex∼π [AvgUtil(x)].
Both voting rules and alignment methods observe comparisons from n users. We model each
i = 1, . . . , n as a fresh user with independently drawn utility vector ui ∼ D. For exposition, we
assume that each user i provides an equal number d ≥ 1 of pairwise comparisons. For each i,
and for j = 1, . . . , d, we independently draw alternatives xji , yij from a fixed distribution µ over

    3
     Our assumption that utilities be in [0, 1] is weaker than any of the three assumptions — unit-sum, unit-range,
or approval [23] — made in classic distortion to avoid trivial infinite lower bounds.


                                                           3

[form-feed in source extraction]
alternatives, in which the minimum probability mass µmin := minx∈A µ(x) is positive.4 User i
then compares each pair {xji , yij } (for j = 1, . . . , d) through a Bradley-Terry model based on i’s
utilities: i prefers xji over yij (written “xji [U+001F control character]i yij ”) with probability σ β · (ui (xji ) − ui (yij )) , where
                                                                                                        [U+0001 control character]

σ(t) := 1/(1+e−t ) is the logistic sigmoid function and β > 0 is a temperature parameter, and prefers
yij over xji (“yij [U+001F control character]i xji ”) otherwise.5 Whereas this specifies the marginal probability of each pairwise
comparison, we make no assumption about the correlation between i’s choices. For example, i might
derive the pairwise comparisons from a Plackett-Luce ranking, ensuring that      [U+0002 control character] the user’s comparisons  [U+0001 control character][U+0003 control character]
are always consistent.6 We set p(x [U+001F control character] y) for the expected win rate Eu∼D σ β · (u(x) − u(y)) .

[Next verbatim source excerpt]

Social Choice Setting. A voting rule f observes the sampled pairwise comparisons {xji [U+001F control character]i
yij }i∈[n],j∈[d] and maps them to a probability distribution over alternatives. For some m, D, β, d, µ,
and
  [U+0002 control character] the correlation between comparisons,     the average utility of f for n samples is AvgUtiln (f ) :=
E AvgUtil f ({xji [U+001F control character]i yij }i,j ) , where the expectation is taken over the pairwise comparisons. The
                               [U+0001 control character][U+0003 control character]

distortion of f on D is the competitive ratio between AvgUtiln (f ) and the optimal average utility
maxx∈A AvgUtil(x) in the limit of n → ∞ samples, and the distortion of f the worst-case distortion
over all D:

            dist(f, D) := lim supn→∞ maxAvgUtil
                                        x∈A AvgUtil(x)
                                                (f )   ,               dist(f ) = supD dist(f, D).
                                                       n



For alternatives x, y, let #(x [U+001F control character] y) := {(i, j) | xji = x, yij = y} denote the number of pairwise
comparisons in which x beat y. The (normalized) Borda score [45] of alternative x is
                                                       P
                                                         y∈A #(x[U+001F control character]y)
                                   BC(x) := P         #(x[U+001F control character]y)+
                                                              P             ,
                                                  y∈A           y6=x #(y[U+001F control character]x)


i.e., the fraction of pairwise comparisons involving x in which it wins. The Borda voting rule chooses
the winner uniformly among all alternatives with maximum Borda score. The Maximal Lotteries vot-
ing rule first computes the margin matrix M ∈ Rm×m , where Mx,y = #(x[U+001F control character]y)−#(y[U+001F control character]x)
                                                                          #(x[U+001F control character]y)+#(y[U+001F control character]x) . It then con-
siders a symmetric two-player zero-sum game in which player 1 selects alternative x1 , player 2 selects
alternative x2 , and the payoffs are Mx1 ,x2 for player 1 and Mx2 ,x1 = −Mx1 ,x2 for player 2. The max-
imal lotteries rule returns a distribution π ∈ argmaxπ1 ∈∆(A) minπ2 ∈∆(A) Ex1 ∼π1 ,x2 ∼π2 [Mx1 ,x2 ],
i.e., a mixed strategy in Nash equilibrium. When several such π exist, our results hold for any choice.
\end{ReviewVerbatim}



\paragraph{Expanded PaperInterface specification}
\begin{ReviewVerbatim}
(∀ {User Alternative : Type} [inst : Fintype User] [inst_1 : DecidableEq User] [inst_2 : Fintype Alternative]
    [inst_3 : DecidableEq Alternative] (population : PMF User)
    (utility : AppliedModelingLib.Alignment.Welfare.FiniteUtilityProfile User Alternative),
    AppliedModelingLib.Alignment.Welfare.UnitIntervalUtilityProfile utility →
      ∀ (sampling : PMF Alternative) {btScale : ℝ},
        0 < btScale →
          ∀ (winner : Alternative),
            (∀ (alternative : Alternative),
                AppliedModelingLib.Alignment.Welfare.pairwiseBordaScore sampling
                    (AppliedModelingLib.Alignment.Welfare.populationBradleyTerryPreference population utility btScale)
                    alternative ≤
                  AppliedModelingLib.Alignment.Welfare.pairwiseBordaScore sampling
                    (AppliedModelingLib.Alignment.Welfare.populationBradleyTerryPreference population utility btScale)
                    winner) →
              ∀ (alternative : Alternative),
                (AppliedModelingLib.Alignment.Welfare.sigmoidChordSlope btScale / (1 / 4)) ^ 2 *
                    AppliedModelingLib.Alignment.Welfare.populationAverageUtility population utility alternative ≤
                  AppliedModelingLib.Alignment.Welfare.populationAverageUtility population utility winner) ∧
  ∀ {User Alternative : Type} [inst : Fintype User] [inst_1 : DecidableEq User] [inst_2 : Fintype Alternative]
    [inst_3 : DecidableEq Alternative] [Nonempty Alternative] (population : PMF User)
    (utility : AppliedModelingLib.Alignment.Welfare.FiniteUtilityProfile User Alternative),
    AppliedModelingLib.Alignment.Welfare.UnitIntervalUtilityProfile utility →
      ∀ (sampling : PMF Alternative) {btScale : ℝ},
        0 < btScale →
          ∀ (responseLaw : PMF (GolzHaghtalabYang2025Distortion.theorem2UserResponseTable Alternative)),
            GolzHaghtalabYang2025Distortion.theorem2UserResponseCalibrated responseLaw
                (AppliedModelingLib.Alignment.Welfare.populationBradleyTerryPreference population utility btScale) →
              ∀ (users comparisonsPerUser : ℕ),
                0 < users →
                  0 < comparisonsPerUser →
                    ∀ (minimumMass delta : ℝ),
                      0 < minimumMass →
                        (∀ (alternative : Alternative), minimumMass ≤ (sampling alternative).toReal) →
                          0 < delta →
                            delta ≤ 1 →
                              2 * ↑(Fintype.card Alternative) *
                                    Real.exp (-↑(users * comparisonsPerUser) * minimumMass / 8) ≤
                                  delta →
                                ∀
                                  (winner :
                                    (Fin users →
                                        GolzHaghtalabYang2025Distortion.theorem2UserResponseTable Alternative ×
                                          (Fin comparisonsPerUser → Alternative × Alternative)) →
                                      Alternative),
                                  (∀
                                      (sample :
                                        Fin users →
                                          GolzHaghtalabYang2025Distortion.theorem2UserResponseTable Alternative ×
                                            (Fin comparisonsPerUser → Alternative × Alternative))
                                      (alternative : Alternative),
                                      GolzHaghtalabYang2025Distortion.theorem11IidUserLatentBordaScore alternative
                                          sample ≤
                                        GolzHaghtalabYang2025Distortion.theorem11IidUserLatentBordaScore (winner sample)
                                          sample) →
                                    ∀ (benchmark : Alternative),
                                      (1 - delta) *
                                          ((AppliedModelingLib.Alignment.Welfare.sigmoidChordSlope btScale / (1 / 4)) ^
                                                2 *
                                              AppliedModelingLib.Alignment.Welfare.populationAverageUtility population
                                                utility benchmark -
                                            2 *
                                                GolzHaghtalabYang2025Distortion.theorem2SourceRate
                                                  (Fintype.card Alternative) users comparisonsPerUser minimumMass
                                                  delta /
                                              (btScale * (1 / 4))) ≤
                                        AppliedModelingLib.pmfExp
                                          (AppliedModelingLib.pmfProduct (Fin users)
                                            (GolzHaghtalabYang2025Distortion.theorem2UserResponseTable Alternative ×
                                              (Fin comparisonsPerUser → Alternative × Alternative))
                                            (AppliedModelingLib.pmfProd responseLaw
                                              (AppliedModelingLib.pmfProduct (Fin comparisonsPerUser)
                                                (Alternative × Alternative)
                                                (GolzHaghtalabYang2025Distortion.theorem2UserPairLabelLaw sampling))))
                                          fun
                                            (sample :
                                              Fin users →
                                                GolzHaghtalabYang2025Distortion.theorem2UserResponseTable Alternative ×
                                                  (Fin comparisonsPerUser → Alternative × Alternative)) =>
                                          AppliedModelingLib.Alignment.Welfare.populationAverageUtility population
                                            utility (winner sample)
\end{ReviewVerbatim}

\paragraph{Retained governing declarations}
\begin{itemize}\raggedright
\item \hyperlink{library-prerequisite-56c0734f8877ddbe}{Applied\allowbreak{}Modeling\allowbreak{}Lib.\allowbreak{}Alignment.\allowbreak{}Welfare.\allowbreak{}Finite\allowbreak{}Utility\allowbreak{}Profile}
\item \hyperlink{library-prerequisite-cbe8afa0d6d1ba78}{Applied\allowbreak{}Modeling\allowbreak{}Lib.\allowbreak{}Alignment.\allowbreak{}Welfare.\allowbreak{}Unit\allowbreak{}Interval\allowbreak{}Utility\allowbreak{}Profile}
\item \hyperlink{governing-declaration-3808cf59ea39b435}{Applied\allowbreak{}Modeling\allowbreak{}Lib.\allowbreak{}Alignment.\allowbreak{}Welfare.\allowbreak{}pairwise\allowbreak{}Borda\allowbreak{}Score}
\item \hyperlink{library-prerequisite-15f4b442749deab2}{Applied\allowbreak{}Modeling\allowbreak{}Lib.\allowbreak{}Alignment.\allowbreak{}Welfare.\allowbreak{}population\allowbreak{}Average\allowbreak{}Utility}
\item \hyperlink{library-prerequisite-cf06d15d7f8b7c31}{Applied\allowbreak{}Modeling\allowbreak{}Lib.\allowbreak{}Alignment.\allowbreak{}Welfare.\allowbreak{}population\allowbreak{}Bradley\allowbreak{}Terry\allowbreak{}Preference}
\item \hyperlink{governing-declaration-a4d623b8c608b566}{Applied\allowbreak{}Modeling\allowbreak{}Lib.\allowbreak{}Alignment.\allowbreak{}Welfare.\allowbreak{}sigmoid\allowbreak{}Chord\allowbreak{}Slope}
\item \hyperlink{governing-declaration-8e4b4389a4f6454b}{Applied\allowbreak{}Modeling\allowbreak{}Lib.\allowbreak{}pmf\allowbreak{}Exp}
\item \hyperlink{library-prerequisite-de8e2a8bfe221a38}{Applied\allowbreak{}Modeling\allowbreak{}Lib.\allowbreak{}pmf\allowbreak{}Prod}
\item \hyperlink{governing-declaration-5ac12d0a52b69068}{Applied\allowbreak{}Modeling\allowbreak{}Lib.\allowbreak{}pmf\allowbreak{}Product}
\item \hyperlink{governing-declaration-7506044331f97cc9}{Golz\allowbreak{}Haghtalab\allowbreak{}Yang2025\allowbreak{}Distortion.\allowbreak{}theorem11\allowbreak{}Iid\allowbreak{}User\allowbreak{}Latent\allowbreak{}Borda\allowbreak{}Score}
\item \hyperlink{governing-declaration-5d4c92feb59536c0}{Golz\allowbreak{}Haghtalab\allowbreak{}Yang2025\allowbreak{}Distortion.\allowbreak{}theorem2\allowbreak{}Source\allowbreak{}Rate}
\item \hyperlink{governing-declaration-98c3ee9a11a59ef9}{Golz\allowbreak{}Haghtalab\allowbreak{}Yang2025\allowbreak{}Distortion.\allowbreak{}theorem2\allowbreak{}User\allowbreak{}Pair\allowbreak{}Label\allowbreak{}Law}
\item \hyperlink{governing-declaration-080b9388e51a6018}{Golz\allowbreak{}Haghtalab\allowbreak{}Yang2025\allowbreak{}Distortion.\allowbreak{}theorem2\allowbreak{}User\allowbreak{}Response\allowbreak{}Calibrated}
\item \hyperlink{governing-declaration-c5d0273cbd7ddcc9}{Golz\allowbreak{}Haghtalab\allowbreak{}Yang2025\allowbreak{}Distortion.\allowbreak{}theorem2\allowbreak{}User\allowbreak{}Response\allowbreak{}Table}
\end{itemize}
\paragraph{Lean-elaborated claim roles}\begin{ReviewVerbatim}
1. conclusion: (∀ {User Alternative : Type} [inst : Fintype User] [inst_1 : DecidableEq User] [inst_2 : Fintype Alternative]
    [inst_3 : DecidableEq Alternative] (population : PMF User)
    (utility : AppliedModelingLib.Alignment.Welfare.FiniteUtilityProfile User Alternative),
    AppliedModelingLib.Alignment.Welfare.UnitIntervalUtilityProfile utility →
      ∀ (sampling : PMF Alternative) {btScale : ℝ},
        0 < btScale →
          ∀ (winner : Alternative),
            (∀ (alternative : Alternative),
                AppliedModelingLib.Alignment.Welfare.pairwiseBordaScore sampling
                    (AppliedModelingLib.Alignment.Welfare.populationBradleyTerryPreference population utility btScale)
                    alternative ≤
                  AppliedModelingLib.Alignment.Welfare.pairwiseBordaScore sampling
                    (AppliedModelingLib.Alignment.Welfare.populationBradleyTerryPreference population utility btScale)
                    winner) →
              ∀ (alternative : Alternative),
                (AppliedModelingLib.Alignment.Welfare.sigmoidChordSlope btScale / (1 / 4)) ^ 2 *
                    AppliedModelingLib.Alignment.Welfare.populationAverageUtility population utility alternative ≤
                  AppliedModelingLib.Alignment.Welfare.populationAverageUtility population utility winner) ∧
  ∀ {User Alternative : Type} [inst : Fintype User] [inst_1 : DecidableEq User] [inst_2 : Fintype Alternative]
    [inst_3 : DecidableEq Alternative] [Nonempty Alternative] (population : PMF User)
    (utility : AppliedModelingLib.Alignment.Welfare.FiniteUtilityProfile User Alternative),
    AppliedModelingLib.Alignment.Welfare.UnitIntervalUtilityProfile utility →
      ∀ (sampling : PMF Alternative) {btScale : ℝ},
        0 < btScale →
          ∀ (responseLaw : PMF (GolzHaghtalabYang2025Distortion.theorem2UserResponseTable Alternative)),
            GolzHaghtalabYang2025Distortion.theorem2UserResponseCalibrated responseLaw
                (AppliedModelingLib.Alignment.Welfare.populationBradleyTerryPreference population utility btScale) →
              ∀ (users comparisonsPerUser : ℕ),
                0 < users →
                  0 < comparisonsPerUser →
                    ∀ (minimumMass delta : ℝ),
                      0 < minimumMass →
                        (∀ (alternative : Alternative), minimumMass ≤ (sampling alternative).toReal) →
                          0 < delta →
                            delta ≤ 1 →
                              2 * ↑(Fintype.card Alternative) *
                                    Real.exp (-↑(users * comparisonsPerUser) * minimumMass / 8) ≤
                                  delta →
                                ∀
                                  (winner :
                                    (Fin users →
                                        GolzHaghtalabYang2025Distortion.theorem2UserResponseTable Alternative ×
                                          (Fin comparisonsPerUser → Alternative × Alternative)) →
                                      Alternative),
                                  (∀
                                      (sample :
                                        Fin users →
                                          GolzHaghtalabYang2025Distortion.theorem2UserResponseTable Alternative ×
                                            (Fin comparisonsPerUser → Alternative × Alternative))
                                      (alternative : Alternative),
                                      GolzHaghtalabYang2025Distortion.theorem11IidUserLatentBordaScore alternative
                                          sample ≤
                                        GolzHaghtalabYang2025Distortion.theorem11IidUserLatentBordaScore (winner sample)
                                          sample) →
                                    ∀ (benchmark : Alternative),
                                      (1 - delta) *
                                          ((AppliedModelingLib.Alignment.Welfare.sigmoidChordSlope btScale / (1 / 4)) ^
                                                2 *
                                              AppliedModelingLib.Alignment.Welfare.populationAverageUtility population
                                                utility benchmark -
                                            2 *
                                                GolzHaghtalabYang2025Distortion.theorem2SourceRate
                                                  (Fintype.card Alternative) users comparisonsPerUser minimumMass
                                                  delta /
                                              (btScale * (1 / 4))) ≤
                                        AppliedModelingLib.pmfExp
                                          (AppliedModelingLib.pmfProduct (Fin users)
                                            (GolzHaghtalabYang2025Distortion.theorem2UserResponseTable Alternative ×
                                              (Fin comparisonsPerUser → Alternative × Alternative))
                                            (AppliedModelingLib.pmfProd responseLaw
                                              (AppliedModelingLib.pmfProduct (Fin comparisonsPerUser)
                                                (Alternative × Alternative)
                                                (GolzHaghtalabYang2025Distortion.theorem2UserPairLabelLaw sampling))))
                                          fun
                                            (sample :
                                              Fin users →
                                                GolzHaghtalabYang2025Distortion.theorem2UserResponseTable Alternative ×
                                                  (Fin comparisonsPerUser → Alternative × Alternative)) =>
                                          AppliedModelingLib.Alignment.Welfare.populationAverageUtility population
                                            utility (winner sample)
\end{ReviewVerbatim}

\paragraph{Lean proof endpoint (built separately)}\begin{ReviewVerbatim}
GolzHaghtalabYang2025Distortion.theorem2_borda_distortion
\end{ReviewVerbatim}

\paragraph{Recorded source-to-Spec screening}
\textbf{Verdict:} \texttt{matches}
\\
{\raggedright
\textbf{Reason:} The population squared chord/\allowbreak{}tangent Borda factor and literal finite-\allowbreak{}user expected-\allowbreak{}welfare source rate are retained.\allowbreak{} The empirical-\allowbreak{}Borda zero-\allowbreak{}incidence branch is now an explicit neutral-\allowbreak{}score finite-\allowbreak{}rule convention, not a model premise.\allowbreak{}
\par}
\reviewmatch{Matches source input}
\noindent\textbf{Reviewer annotation}\par
\reviewerbox
\clearpage
\hypertarget{source-claim-c933c69486621c74}{}
\section*{Review row 6}
\textbf{Source locator:} {\footnotesize\raggedright cited publication, lines 431-\allowbreak{}436\par}
\paragraph{Verbatim source input}
\begin{ReviewVerbatim}
Theorem 3 (Voting Rule-Independent Distortion Lower Bound). Fix any β > 0. If each user provides
                                                                           1+e−β
d=1 comparison, no voting rule can guarantee distortion better than β2 · 1−e   −β for large m. If each

user reports d ≥ 2 pairwise comparisons, any voting rule that puts at most 1/m probability mass on
a Condorcet loser must have at least the above distortion.

[Next verbatim source excerpt]

Let A = {1, . . . , m} be a finite set of alternatives. The population of users is described by a
probability distribution D over utility vectors u = (u(1), . . . , u(m)), whose entries 0 ≤ u(x) ≤ 1
indicate a user’s utility for alternative x.3 The objective is to find an alternative x such that its average
utility AvgUtil(x) := Eu∼D [u(x)] across the user population (also known as the utilitarian social
welfare) is as high as possible. We extend this notation to probability distributions π over alternatives
by setting AvgUtil(π) := Ex∼π [AvgUtil(x)].
Both voting rules and alignment methods observe comparisons from n users. We model each
i = 1, . . . , n as a fresh user with independently drawn utility vector ui ∼ D. For exposition, we
assume that each user i provides an equal number d ≥ 1 of pairwise comparisons. For each i,
and for j = 1, . . . , d, we independently draw alternatives xji , yij from a fixed distribution µ over

    3
     Our assumption that utilities be in [0, 1] is weaker than any of the three assumptions — unit-sum, unit-range,
or approval [23] — made in classic distortion to avoid trivial infinite lower bounds.


                                                           3

[form-feed in source extraction]
alternatives, in which the minimum probability mass µmin := minx∈A µ(x) is positive.4 User i
then compares each pair {xji , yij } (for j = 1, . . . , d) through a Bradley-Terry model based on i’s
utilities: i prefers xji over yij (written “xji [U+001F control character]i yij ”) with probability σ β · (ui (xji ) − ui (yij )) , where
                                                                                                        [U+0001 control character]

σ(t) := 1/(1+e−t ) is the logistic sigmoid function and β > 0 is a temperature parameter, and prefers
yij over xji (“yij [U+001F control character]i xji ”) otherwise.5 Whereas this specifies the marginal probability of each pairwise
comparison, we make no assumption about the correlation between i’s choices. For example, i might
derive the pairwise comparisons from a Plackett-Luce ranking, ensuring that      [U+0002 control character] the user’s comparisons  [U+0001 control character][U+0003 control character]
are always consistent.6 We set p(x [U+001F control character] y) for the expected win rate Eu∼D σ β · (u(x) − u(y)) .

[Next verbatim source excerpt]

Social Choice Setting. A voting rule f observes the sampled pairwise comparisons {xji [U+001F control character]i
yij }i∈[n],j∈[d] and maps them to a probability distribution over alternatives. For some m, D, β, d, µ,
and
  [U+0002 control character] the correlation between comparisons,     the average utility of f for n samples is AvgUtiln (f ) :=
E AvgUtil f ({xji [U+001F control character]i yij }i,j ) , where the expectation is taken over the pairwise comparisons. The
                               [U+0001 control character][U+0003 control character]

distortion of f on D is the competitive ratio between AvgUtiln (f ) and the optimal average utility
maxx∈A AvgUtil(x) in the limit of n → ∞ samples, and the distortion of f the worst-case distortion
over all D:

            dist(f, D) := lim supn→∞ maxAvgUtil
                                        x∈A AvgUtil(x)
                                                (f )   ,               dist(f ) = supD dist(f, D).
                                                       n



For alternatives x, y, let #(x [U+001F control character] y) := {(i, j) | xji = x, yij = y} denote the number of pairwise
comparisons in which x beat y. The (normalized) Borda score [45] of alternative x is
                                                       P
                                                         y∈A #(x[U+001F control character]y)
                                   BC(x) := P         #(x[U+001F control character]y)+
                                                              P             ,
                                                  y∈A           y6=x #(y[U+001F control character]x)


i.e., the fraction of pairwise comparisons involving x in which it wins. The Borda voting rule chooses
the winner uniformly among all alternatives with maximum Borda score. The Maximal Lotteries vot-
ing rule first computes the margin matrix M ∈ Rm×m , where Mx,y = #(x[U+001F control character]y)−#(y[U+001F control character]x)
                                                                          #(x[U+001F control character]y)+#(y[U+001F control character]x) . It then con-
siders a symmetric two-player zero-sum game in which player 1 selects alternative x1 , player 2 selects
alternative x2 , and the payoffs are Mx1 ,x2 for player 1 and Mx2 ,x1 = −Mx1 ,x2 for player 2. The max-
imal lotteries rule returns a distribution π ∈ argmaxπ1 ∈∆(A) minπ2 ∈∆(A) Ex1 ∼π1 ,x2 ∼π2 [Mx1 ,x2 ],
i.e., a mixed strategy in Nash equilibrium. When several such π exist, our results hold for any choice.
\end{ReviewVerbatim}



\paragraph{Expanded PaperInterface specification}
\begin{ReviewVerbatim}
(∀ {beta : ℝ} (hbeta : 0 < beta) (pairSampling : (n : ℕ) → PMF (Fin (n + 2) × Fin (n + 2)))
    (rule : (n horizon : ℕ) → (Fin horizon → (Fin (n + 2) × Fin (n + 2)) × Bool) → PMF (Fin (n + 2))) (delta : ℝ),
    0 < delta →
      ∃ (N : ℕ),
        ∀ n ≥ N,
          ∃ (special : Fin (n + 2)),
            {horizon : ℕ |
                have policy : PMF (Fin (n + 2)) :=
                  GolzHaghtalabYang2025Distortion.theorem3ObservationSelectionLaw
                    (GolzHaghtalabYang2025Distortion.theorem3ManyAlternativeObservationLaw (pairSampling n) beta
                      (GolzHaghtalabYang2025Distortion.theorem3DiagonalEpsilon n) hbeta
                      (GolzHaghtalabYang2025Distortion.theorem3DiagonalEpsilon_pos n) special horizon)
                    (rule n horizon);
                (∀ (alternative : Fin (n + 2)),
                    AppliedModelingLib.Alignment.Welfare.populationAverageUtility
                        (GolzHaghtalabYang2025Distortion.theorem3TwoTypePopulation beta
                          (GolzHaghtalabYang2025Distortion.theorem3DiagonalEpsilon n) hbeta
                          (GolzHaghtalabYang2025Distortion.theorem3DiagonalEpsilon_pos n))
                        (GolzHaghtalabYang2025Distortion.theorem3ManyAlternativeUtility
                          (GolzHaghtalabYang2025Distortion.theorem3DiagonalEpsilon n) 1 special)
                        alternative ≤
                      AppliedModelingLib.Alignment.Welfare.populationAverageUtility
                        (GolzHaghtalabYang2025Distortion.theorem3TwoTypePopulation beta
                          (GolzHaghtalabYang2025Distortion.theorem3DiagonalEpsilon n) hbeta
                          (GolzHaghtalabYang2025Distortion.theorem3DiagonalEpsilon_pos n))
                        (GolzHaghtalabYang2025Distortion.theorem3ManyAlternativeUtility
                          (GolzHaghtalabYang2025Distortion.theorem3DiagonalEpsilon n) 1 special)
                        special) ∧
                  beta / 2 * (1 + Real.exp (-beta)) / (1 - Real.exp (-beta)) - delta <
                    GolzHaghtalabYang2025Distortion.theorem3SpecialTypeMass beta
                        (GolzHaghtalabYang2025Distortion.theorem3DiagonalEpsilon n) /
                      AppliedModelingLib.Alignment.Welfare.policyAverageUtility
                        (GolzHaghtalabYang2025Distortion.theorem3TwoTypePopulation beta
                          (GolzHaghtalabYang2025Distortion.theorem3DiagonalEpsilon n) hbeta
                          (GolzHaghtalabYang2025Distortion.theorem3DiagonalEpsilon_pos n))
                        (GolzHaghtalabYang2025Distortion.theorem3ManyAlternativeUtility
                          (GolzHaghtalabYang2025Distortion.theorem3DiagonalEpsilon n) 1 special)
                        policy}.Infinite) ∧
  ∀ {beta : ℝ} (hbeta : 0 < beta) (d : ℕ),
    2 ≤ d →
      ∀ (pairSampling : (n : ℕ) → PMF (Fin (n + 3) × Fin (n + 3))),
        (∀ (n : ℕ) (ordinary : Fin (n + 3)), ordinary ≠ 0 → 0 < ((pairSampling n) (ordinary, 0)).toReal) →
          ∀ (reportLaw : (n : ℕ) → PMF (Fin d → (Fin (n + 3) × Fin (n + 3)) × Bool)),
            (∀ (n : ℕ) (comparisonIndex : Fin d),
                PMF.map (fun (report : Fin d → (Fin (n + 3) × Fin (n + 3)) × Bool) => report comparisonIndex)
                    (reportLaw n) =
                  GolzHaghtalabYang2025Distortion.theorem3ManyAlternativeOneComparisonLawXi (pairSampling n) beta
                    (GolzHaghtalabYang2025Distortion.theorem3D2DiagonalEpsilon n)
                    (GolzHaghtalabYang2025Distortion.theorem3D2DiagonalXi n) hbeta
                    (GolzHaghtalabYang2025Distortion.theorem3D2DiagonalEpsilon_pos n) 0) →
              ∀
                (rule :
                  (n horizon : ℕ) → (Fin horizon → Fin d → (Fin (n + 3) × Fin (n + 3)) × Bool) → PMF (Fin (n + 3))),
                (∀ (n horizon : ℕ),
                    GolzHaghtalabYang2025Distortion.theorem3ProbabilisticCondorcetLoserCriterion (rule n horizon)
                      fun (observation : Fin horizon → Fin d → (Fin (n + 3) × Fin (n + 3)) × Bool)
                        (alternative : Fin (n + 3)) =>
                      GolzHaghtalabYang2025Distortion.theorem3IidEmpiricalStrictCondorcetLoser
                        (fun (ordinary : Fin (n + 3)) =>
                          GolzHaghtalabYang2025Distortion.theorem3ReportedOrdinarySpecialMargin alternative ordinary)
                        alternative observation) →
                  ∀ (delta : ℝ),
                    0 < delta →
                      ∃ (N : ℕ),
                        ∀ n ≥ N,
                          ∀ᶠ (horizon : ℕ) in Filter.atTop,
                            beta / 2 * (1 + Real.exp (-beta)) / (1 - Real.exp (-beta)) - delta <
                              GolzHaghtalabYang2025Distortion.theorem3SpecialTypeMass beta
                                  (GolzHaghtalabYang2025Distortion.theorem3D2DiagonalEpsilon n) /
                                AppliedModelingLib.Alignment.Welfare.policyAverageUtility
                                  (GolzHaghtalabYang2025Distortion.theorem3TwoTypePopulation beta
                                    (GolzHaghtalabYang2025Distortion.theorem3D2DiagonalEpsilon n) hbeta
                                    (GolzHaghtalabYang2025Distortion.theorem3D2DiagonalEpsilon_pos n))
                                  (GolzHaghtalabYang2025Distortion.theorem3ManyAlternativeUtility
                                    (GolzHaghtalabYang2025Distortion.theorem3D2DiagonalEpsilon n)
                                    (GolzHaghtalabYang2025Distortion.theorem3D2DiagonalXi n) 0)
                                  (GolzHaghtalabYang2025Distortion.theorem3ObservationSelectionLaw
                                    (GolzHaghtalabYang2025Distortion.theorem3IidUserReportBatchLaw (reportLaw n)
                                      horizon)
                                    (rule n horizon))
\end{ReviewVerbatim}

\paragraph{Retained governing declarations}
\begin{itemize}\raggedright
\item \hyperlink{library-prerequisite-184c74a4db4311c2}{Applied\allowbreak{}Modeling\allowbreak{}Lib.\allowbreak{}Alignment.\allowbreak{}Welfare.\allowbreak{}policy\allowbreak{}Average\allowbreak{}Utility}
\item \hyperlink{library-prerequisite-15f4b442749deab2}{Applied\allowbreak{}Modeling\allowbreak{}Lib.\allowbreak{}Alignment.\allowbreak{}Welfare.\allowbreak{}population\allowbreak{}Average\allowbreak{}Utility}
\item \hyperlink{governing-declaration-0f03655f38b9e33e}{Golz\allowbreak{}Haghtalab\allowbreak{}Yang2025\allowbreak{}Distortion.\allowbreak{}theorem3\allowbreak{}D2\allowbreak{}Diagonal\allowbreak{}Epsilon}
\item \hyperlink{governing-declaration-783b316e1c8a0ee2}{Golz\allowbreak{}Haghtalab\allowbreak{}Yang2025\allowbreak{}Distortion.\allowbreak{}theorem3\allowbreak{}D2\allowbreak{}Diagonal\allowbreak{}Xi}
\item \hyperlink{governing-declaration-a7437eb8cad7f1a3}{Golz\allowbreak{}Haghtalab\allowbreak{}Yang2025\allowbreak{}Distortion.\allowbreak{}theorem3\allowbreak{}Diagonal\allowbreak{}Epsilon}
\item \hyperlink{governing-declaration-8d2e05cae4023172}{Golz\allowbreak{}Haghtalab\allowbreak{}Yang2025\allowbreak{}Distortion.\allowbreak{}theorem3\allowbreak{}Iid\allowbreak{}Empirical\allowbreak{}Strict\allowbreak{}Condorcet\allowbreak{}Loser}
\item \hyperlink{governing-declaration-7fc5db381e1a8f71}{Golz\allowbreak{}Haghtalab\allowbreak{}Yang2025\allowbreak{}Distortion.\allowbreak{}theorem3\allowbreak{}Iid\allowbreak{}User\allowbreak{}Report\allowbreak{}Batch\allowbreak{}Law}
\item \hyperlink{governing-declaration-8cd7a7c5a45b332d}{Golz\allowbreak{}Haghtalab\allowbreak{}Yang2025\allowbreak{}Distortion.\allowbreak{}theorem3\allowbreak{}Many\allowbreak{}Alternative\allowbreak{}Observation\allowbreak{}Law}
\item \hyperlink{governing-declaration-94724c60280865f4}{Golz\allowbreak{}Haghtalab\allowbreak{}Yang2025\allowbreak{}Distortion.\allowbreak{}theorem3\allowbreak{}Many\allowbreak{}Alternative\allowbreak{}One\allowbreak{}Comparison\allowbreak{}Law\allowbreak{}Xi}
\item \hyperlink{governing-declaration-bfef2d422eb25497}{Golz\allowbreak{}Haghtalab\allowbreak{}Yang2025\allowbreak{}Distortion.\allowbreak{}theorem3\allowbreak{}Many\allowbreak{}Alternative\allowbreak{}Utility}
\item \hyperlink{governing-declaration-f48897680b4251cd}{Golz\allowbreak{}Haghtalab\allowbreak{}Yang2025\allowbreak{}Distortion.\allowbreak{}theorem3\allowbreak{}Observation\allowbreak{}Selection\allowbreak{}Law}
\item \hyperlink{governing-declaration-fa9a73ecc3173034}{Golz\allowbreak{}Haghtalab\allowbreak{}Yang2025\allowbreak{}Distortion.\allowbreak{}theorem3\allowbreak{}Probabilistic\allowbreak{}Condorcet\allowbreak{}Loser\allowbreak{}Criterion}
\item \hyperlink{governing-declaration-5e11e0b8b649abae}{Golz\allowbreak{}Haghtalab\allowbreak{}Yang2025\allowbreak{}Distortion.\allowbreak{}theorem3\allowbreak{}Reported\allowbreak{}Ordinary\allowbreak{}Special\allowbreak{}Margin}
\item \hyperlink{governing-declaration-d300ff02fbceb879}{Golz\allowbreak{}Haghtalab\allowbreak{}Yang2025\allowbreak{}Distortion.\allowbreak{}theorem3\allowbreak{}Special\allowbreak{}Type\allowbreak{}Mass}
\item \hyperlink{governing-declaration-5256b014ca2f6f1b}{Golz\allowbreak{}Haghtalab\allowbreak{}Yang2025\allowbreak{}Distortion.\allowbreak{}theorem3\allowbreak{}Two\allowbreak{}Type\allowbreak{}Population}
\end{itemize}
\paragraph{Lean-elaborated claim roles}\begin{ReviewVerbatim}
1. conclusion: (∀ {beta : ℝ} (hbeta : 0 < beta) (pairSampling : (n : ℕ) → PMF (Fin (n + 2) × Fin (n + 2)))
    (rule : (n horizon : ℕ) → (Fin horizon → (Fin (n + 2) × Fin (n + 2)) × Bool) → PMF (Fin (n + 2))) (delta : ℝ),
    0 < delta →
      ∃ (N : ℕ),
        ∀ n ≥ N,
          ∃ (special : Fin (n + 2)),
            {horizon : ℕ |
                have policy : PMF (Fin (n + 2)) :=
                  GolzHaghtalabYang2025Distortion.theorem3ObservationSelectionLaw
                    (GolzHaghtalabYang2025Distortion.theorem3ManyAlternativeObservationLaw (pairSampling n) beta
                      (GolzHaghtalabYang2025Distortion.theorem3DiagonalEpsilon n) hbeta
                      (GolzHaghtalabYang2025Distortion.theorem3DiagonalEpsilon_pos n) special horizon)
                    (rule n horizon);
                (∀ (alternative : Fin (n + 2)),
                    AppliedModelingLib.Alignment.Welfare.populationAverageUtility
                        (GolzHaghtalabYang2025Distortion.theorem3TwoTypePopulation beta
                          (GolzHaghtalabYang2025Distortion.theorem3DiagonalEpsilon n) hbeta
                          (GolzHaghtalabYang2025Distortion.theorem3DiagonalEpsilon_pos n))
                        (GolzHaghtalabYang2025Distortion.theorem3ManyAlternativeUtility
                          (GolzHaghtalabYang2025Distortion.theorem3DiagonalEpsilon n) 1 special)
                        alternative ≤
                      AppliedModelingLib.Alignment.Welfare.populationAverageUtility
                        (GolzHaghtalabYang2025Distortion.theorem3TwoTypePopulation beta
                          (GolzHaghtalabYang2025Distortion.theorem3DiagonalEpsilon n) hbeta
                          (GolzHaghtalabYang2025Distortion.theorem3DiagonalEpsilon_pos n))
                        (GolzHaghtalabYang2025Distortion.theorem3ManyAlternativeUtility
                          (GolzHaghtalabYang2025Distortion.theorem3DiagonalEpsilon n) 1 special)
                        special) ∧
                  beta / 2 * (1 + Real.exp (-beta)) / (1 - Real.exp (-beta)) - delta <
                    GolzHaghtalabYang2025Distortion.theorem3SpecialTypeMass beta
                        (GolzHaghtalabYang2025Distortion.theorem3DiagonalEpsilon n) /
                      AppliedModelingLib.Alignment.Welfare.policyAverageUtility
                        (GolzHaghtalabYang2025Distortion.theorem3TwoTypePopulation beta
                          (GolzHaghtalabYang2025Distortion.theorem3DiagonalEpsilon n) hbeta
                          (GolzHaghtalabYang2025Distortion.theorem3DiagonalEpsilon_pos n))
                        (GolzHaghtalabYang2025Distortion.theorem3ManyAlternativeUtility
                          (GolzHaghtalabYang2025Distortion.theorem3DiagonalEpsilon n) 1 special)
                        policy}.Infinite) ∧
  ∀ {beta : ℝ} (hbeta : 0 < beta) (d : ℕ),
    2 ≤ d →
      ∀ (pairSampling : (n : ℕ) → PMF (Fin (n + 3) × Fin (n + 3))),
        (∀ (n : ℕ) (ordinary : Fin (n + 3)), ordinary ≠ 0 → 0 < ((pairSampling n) (ordinary, 0)).toReal) →
          ∀ (reportLaw : (n : ℕ) → PMF (Fin d → (Fin (n + 3) × Fin (n + 3)) × Bool)),
            (∀ (n : ℕ) (comparisonIndex : Fin d),
                PMF.map (fun (report : Fin d → (Fin (n + 3) × Fin (n + 3)) × Bool) => report comparisonIndex)
                    (reportLaw n) =
                  GolzHaghtalabYang2025Distortion.theorem3ManyAlternativeOneComparisonLawXi (pairSampling n) beta
                    (GolzHaghtalabYang2025Distortion.theorem3D2DiagonalEpsilon n)
                    (GolzHaghtalabYang2025Distortion.theorem3D2DiagonalXi n) hbeta
                    (GolzHaghtalabYang2025Distortion.theorem3D2DiagonalEpsilon_pos n) 0) →
              ∀
                (rule :
                  (n horizon : ℕ) → (Fin horizon → Fin d → (Fin (n + 3) × Fin (n + 3)) × Bool) → PMF (Fin (n + 3))),
                (∀ (n horizon : ℕ),
                    GolzHaghtalabYang2025Distortion.theorem3ProbabilisticCondorcetLoserCriterion (rule n horizon)
                      fun (observation : Fin horizon → Fin d → (Fin (n + 3) × Fin (n + 3)) × Bool)
                        (alternative : Fin (n + 3)) =>
                      GolzHaghtalabYang2025Distortion.theorem3IidEmpiricalStrictCondorcetLoser
                        (fun (ordinary : Fin (n + 3)) =>
                          GolzHaghtalabYang2025Distortion.theorem3ReportedOrdinarySpecialMargin alternative ordinary)
                        alternative observation) →
                  ∀ (delta : ℝ),
                    0 < delta →
                      ∃ (N : ℕ),
                        ∀ n ≥ N,
                          ∀ᶠ (horizon : ℕ) in Filter.atTop,
                            beta / 2 * (1 + Real.exp (-beta)) / (1 - Real.exp (-beta)) - delta <
                              GolzHaghtalabYang2025Distortion.theorem3SpecialTypeMass beta
                                  (GolzHaghtalabYang2025Distortion.theorem3D2DiagonalEpsilon n) /
                                AppliedModelingLib.Alignment.Welfare.policyAverageUtility
                                  (GolzHaghtalabYang2025Distortion.theorem3TwoTypePopulation beta
                                    (GolzHaghtalabYang2025Distortion.theorem3D2DiagonalEpsilon n) hbeta
                                    (GolzHaghtalabYang2025Distortion.theorem3D2DiagonalEpsilon_pos n))
                                  (GolzHaghtalabYang2025Distortion.theorem3ManyAlternativeUtility
                                    (GolzHaghtalabYang2025Distortion.theorem3D2DiagonalEpsilon n)
                                    (GolzHaghtalabYang2025Distortion.theorem3D2DiagonalXi n) 0)
                                  (GolzHaghtalabYang2025Distortion.theorem3ObservationSelectionLaw
                                    (GolzHaghtalabYang2025Distortion.theorem3IidUserReportBatchLaw (reportLaw n)
                                      horizon)
                                    (rule n horizon))
\end{ReviewVerbatim}

\paragraph{Lean proof endpoint (built separately)}\begin{ReviewVerbatim}
GolzHaghtalabYang2025Distortion.theorem3_voting_rule_lower_bound
\end{ReviewVerbatim}

\paragraph{Recorded source-to-Spec screening}
\textbf{Verdict:} \texttt{matches}
\\
{\raggedright
\textbf{Reason:} Both the d = 1 arbitrary-\allowbreak{}rule and d >= 2 Condorcet-\allowbreak{}loser lower-\allowbreak{}bound branches are covered.\allowbreak{}
\par}
\reviewmatch{Matches source input}
\noindent\textbf{Reviewer annotation}\par
\reviewerbox
\clearpage
\hypertarget{source-claim-27fd47f63d6ea396}{}
\section*{Review row 7}
\textbf{Source locator:} {\footnotesize\raggedright cited publication, lines 450-\allowbreak{}454\par}
\paragraph{Verbatim source input}
\begin{ReviewVerbatim}
Theorem 5 (Borda Distortion Lower Bound, Informally). For any β > 0 and m ≥ 3, the distortion
                                                                                          1+e−β
guaranteed by Borda (and, hence, RLHF) is greater than and bounded away from β2 1−e           −β . In

particular, as β → ∞, this distortion guarantee is at least (1 − o(1)) · β.

[Next verbatim source excerpt]

Let A = {1, . . . , m} be a finite set of alternatives. The population of users is described by a
probability distribution D over utility vectors u = (u(1), . . . , u(m)), whose entries 0 ≤ u(x) ≤ 1
indicate a user’s utility for alternative x.3 The objective is to find an alternative x such that its average
utility AvgUtil(x) := Eu∼D [u(x)] across the user population (also known as the utilitarian social
welfare) is as high as possible. We extend this notation to probability distributions π over alternatives
by setting AvgUtil(π) := Ex∼π [AvgUtil(x)].
Both voting rules and alignment methods observe comparisons from n users. We model each
i = 1, . . . , n as a fresh user with independently drawn utility vector ui ∼ D. For exposition, we
assume that each user i provides an equal number d ≥ 1 of pairwise comparisons. For each i,
and for j = 1, . . . , d, we independently draw alternatives xji , yij from a fixed distribution µ over

    3
     Our assumption that utilities be in [0, 1] is weaker than any of the three assumptions — unit-sum, unit-range,
or approval [23] — made in classic distortion to avoid trivial infinite lower bounds.


                                                           3

[form-feed in source extraction]
alternatives, in which the minimum probability mass µmin := minx∈A µ(x) is positive.4 User i
then compares each pair {xji , yij } (for j = 1, . . . , d) through a Bradley-Terry model based on i’s
utilities: i prefers xji over yij (written “xji [U+001F control character]i yij ”) with probability σ β · (ui (xji ) − ui (yij )) , where
                                                                                                        [U+0001 control character]

σ(t) := 1/(1+e−t ) is the logistic sigmoid function and β > 0 is a temperature parameter, and prefers
yij over xji (“yij [U+001F control character]i xji ”) otherwise.5 Whereas this specifies the marginal probability of each pairwise
comparison, we make no assumption about the correlation between i’s choices. For example, i might
derive the pairwise comparisons from a Plackett-Luce ranking, ensuring that      [U+0002 control character] the user’s comparisons  [U+0001 control character][U+0003 control character]
are always consistent.6 We set p(x [U+001F control character] y) for the expected win rate Eu∼D σ β · (u(x) − u(y)) .

[Next verbatim source excerpt]

Social Choice Setting. A voting rule f observes the sampled pairwise comparisons {xji [U+001F control character]i
yij }i∈[n],j∈[d] and maps them to a probability distribution over alternatives. For some m, D, β, d, µ,
and
  [U+0002 control character] the correlation between comparisons,     the average utility of f for n samples is AvgUtiln (f ) :=
E AvgUtil f ({xji [U+001F control character]i yij }i,j ) , where the expectation is taken over the pairwise comparisons. The
                               [U+0001 control character][U+0003 control character]

distortion of f on D is the competitive ratio between AvgUtiln (f ) and the optimal average utility
maxx∈A AvgUtil(x) in the limit of n → ∞ samples, and the distortion of f the worst-case distortion
over all D:

            dist(f, D) := lim supn→∞ maxAvgUtil
                                        x∈A AvgUtil(x)
                                                (f )   ,               dist(f ) = supD dist(f, D).
                                                       n



For alternatives x, y, let #(x [U+001F control character] y) := {(i, j) | xji = x, yij = y} denote the number of pairwise
comparisons in which x beat y. The (normalized) Borda score [45] of alternative x is
                                                       P
                                                         y∈A #(x[U+001F control character]y)
                                   BC(x) := P         #(x[U+001F control character]y)+
                                                              P             ,
                                                  y∈A           y6=x #(y[U+001F control character]x)


i.e., the fraction of pairwise comparisons involving x in which it wins. The Borda voting rule chooses
the winner uniformly among all alternatives with maximum Borda score. The Maximal Lotteries vot-
ing rule first computes the margin matrix M ∈ Rm×m , where Mx,y = #(x[U+001F control character]y)−#(y[U+001F control character]x)
                                                                          #(x[U+001F control character]y)+#(y[U+001F control character]x) . It then con-
siders a symmetric two-player zero-sum game in which player 1 selects alternative x1 , player 2 selects
alternative x2 , and the payoffs are Mx1 ,x2 for player 1 and Mx2 ,x1 = −Mx1 ,x2 for player 2. The max-
imal lotteries rule returns a distribution π ∈ argmaxπ1 ∈∆(A) minπ2 ∈∆(A) Ex1 ∼π1 ,x2 ∼π2 [Mx1 ,x2 ],
i.e., a mixed strategy in Nash equilibrium. When several such π exist, our results hold for any choice.
\end{ReviewVerbatim}



\paragraph{Expanded PaperInterface specification}
\begin{ReviewVerbatim}
(∀ {beta : ℝ} (hbeta : 0 < beta),
    have gamma : ℝ := GolzHaghtalabYang2025Distortion.theorem12LogChoice beta;
    ∃ (hgamma : 0 < gamma) (hgamma_lt_one : gamma < 1),
      beta / 2 * (1 + Real.exp (-beta)) / (1 - Real.exp (-beta)) <
          GolzHaghtalabYang2025Distortion.theorem12Eq10Coefficient beta gamma ∧
        ∀ q < GolzHaghtalabYang2025Distortion.theorem12Eq10Coefficient beta gamma,
          ∀ (m : ℕ),
            3 ≤ m →
              ∃ (epsilon : ℝ) (hepsilon : 0 < epsilon) (hepsilon_lt_half : epsilon < 1 / 2) (sampling :
                PMF (Fin 3 ⊕ Fin (m - 3))),
                epsilon ^ 2 < 1 - epsilon ∧
                  Fintype.card (Fin 3 ⊕ Fin (m - 3)) = m ∧
                    have collapse : Fin 3 ⊕ Fin (m - 3) → Fin 3 :=
                      GolzHaghtalabYang2025Distortion.theorem12CloneCollapse 1 (m - 3);
                    have splitUtility :
                      AppliedModelingLib.Alignment.Welfare.FiniteUtilityProfile (Fin 3) (Fin 3 ⊕ Fin (m - 3)) :=
                      GolzHaghtalabYang2025Distortion.theorem12SplitUtility
                        (GolzHaghtalabYang2025Distortion.theorem12Utility epsilon (epsilon ^ 2) gamma) collapse;
                    have population : PMF (Fin 3) :=
                      GolzHaghtalabYang2025Distortion.theorem12Population beta epsilon gamma hbeta hepsilon
                        (GolzHaghtalabYang2025Distortion.theorem5Spec._proof_2 epsilon hepsilon hepsilon_lt_half) hgamma
                        hgamma_lt_one;
                    have preference :
                      AppliedModelingLib.Learning.HumanFeedback.PairwisePreference PUnit.{1} (Fin 3 ⊕ Fin (m - 3)) :=
                      AppliedModelingLib.Alignment.Welfare.populationBradleyTerryPreference population splitUtility
                        beta;
                    AppliedModelingLib.Alignment.Welfare.UnitIntervalUtilityProfile splitUtility ∧
                      AppliedModelingLib.PMFFullSupport sampling ∧
                        q <
                            AppliedModelingLib.Alignment.Welfare.populationAverageUtility population splitUtility
                                (Sum.inl 0) /
                              AppliedModelingLib.Alignment.Welfare.populationAverageUtility population splitUtility
                                (Sum.inl 2) ∧
                          Filter.Tendsto
                            (fun (horizon : ℕ) =>
                              AppliedModelingLib.pmfProb
                                (GolzHaghtalabYang2025Distortion.theorem12IidBordaReportBatchLaw
                                  (GolzHaghtalabYang2025Distortion.theorem12OneComparisonReportLaw sampling preference)
                                  horizon)
                                fun (sample : Fin horizon → ((Fin 3 ⊕ Fin (m - 3)) × (Fin 3 ⊕ Fin (m - 3))) × Bool) =>
                                GolzHaghtalabYang2025Distortion.theorem12IidEmpiricalBordaFailure
                                  GolzHaghtalabYang2025Distortion.theorem12OneComparisonWins
                                  GolzHaghtalabYang2025Distortion.theorem12OneComparisonIncidences (Sum.inl 2) sample)
                            Filter.atTop (nhds 0)) ∧
  Filter.Tendsto
    (fun (beta : ℝ) =>
      GolzHaghtalabYang2025Distortion.theorem12Eq10Coefficient beta
          (GolzHaghtalabYang2025Distortion.theorem12LogChoice beta) /
        beta)
    Filter.atTop (nhds 1)
\end{ReviewVerbatim}

\paragraph{Retained governing declarations}
\begin{itemize}\raggedright
\item \hyperlink{library-prerequisite-56c0734f8877ddbe}{Applied\allowbreak{}Modeling\allowbreak{}Lib.\allowbreak{}Alignment.\allowbreak{}Welfare.\allowbreak{}Finite\allowbreak{}Utility\allowbreak{}Profile}
\item \hyperlink{library-prerequisite-cbe8afa0d6d1ba78}{Applied\allowbreak{}Modeling\allowbreak{}Lib.\allowbreak{}Alignment.\allowbreak{}Welfare.\allowbreak{}Unit\allowbreak{}Interval\allowbreak{}Utility\allowbreak{}Profile}
\item \hyperlink{library-prerequisite-15f4b442749deab2}{Applied\allowbreak{}Modeling\allowbreak{}Lib.\allowbreak{}Alignment.\allowbreak{}Welfare.\allowbreak{}population\allowbreak{}Average\allowbreak{}Utility}
\item \hyperlink{library-prerequisite-cf06d15d7f8b7c31}{Applied\allowbreak{}Modeling\allowbreak{}Lib.\allowbreak{}Alignment.\allowbreak{}Welfare.\allowbreak{}population\allowbreak{}Bradley\allowbreak{}Terry\allowbreak{}Preference}
\item \hyperlink{library-prerequisite-7e2db20cf087d9ba}{Applied\allowbreak{}Modeling\allowbreak{}Lib.\allowbreak{}Learning.\allowbreak{}Human\allowbreak{}Feedback.\allowbreak{}Pairwise\allowbreak{}Preference}
\item \hyperlink{library-prerequisite-268c959ee7d2d148}{Applied\allowbreak{}Modeling\allowbreak{}Lib.\allowbreak{}PMF\allowbreak{}Full\allowbreak{}Support}
\item \hyperlink{governing-declaration-9c38d53b9d664221}{Applied\allowbreak{}Modeling\allowbreak{}Lib.\allowbreak{}pmf\allowbreak{}Prob}
\item \hyperlink{governing-declaration-ec899679d7cb6165}{Golz\allowbreak{}Haghtalab\allowbreak{}Yang2025\allowbreak{}Distortion.\allowbreak{}theorem12\allowbreak{}Clone\allowbreak{}Collapse}
\item \hyperlink{governing-declaration-b2a8d113fc3130c2}{Golz\allowbreak{}Haghtalab\allowbreak{}Yang2025\allowbreak{}Distortion.\allowbreak{}theorem12\allowbreak{}Empirical\allowbreak{}Borda\allowbreak{}Score}
\item \hyperlink{governing-declaration-aec5a9f50c9e967a}{Golz\allowbreak{}Haghtalab\allowbreak{}Yang2025\allowbreak{}Distortion.\allowbreak{}theorem12\allowbreak{}Eq10\allowbreak{}Coefficient}
\item \hyperlink{governing-declaration-e5cae00c98b32d81}{Golz\allowbreak{}Haghtalab\allowbreak{}Yang2025\allowbreak{}Distortion.\allowbreak{}theorem12\allowbreak{}Iid\allowbreak{}Borda\allowbreak{}Report\allowbreak{}Batch\allowbreak{}Law}
\item \hyperlink{governing-declaration-3940e2eb509393ea}{Golz\allowbreak{}Haghtalab\allowbreak{}Yang2025\allowbreak{}Distortion.\allowbreak{}theorem12\allowbreak{}Iid\allowbreak{}Empirical\allowbreak{}Borda\allowbreak{}Failure}
\item \hyperlink{governing-declaration-df5d365b55936500}{Golz\allowbreak{}Haghtalab\allowbreak{}Yang2025\allowbreak{}Distortion.\allowbreak{}theorem12\allowbreak{}Log\allowbreak{}Choice}
\item \hyperlink{governing-declaration-baa615edc2baf254}{Golz\allowbreak{}Haghtalab\allowbreak{}Yang2025\allowbreak{}Distortion.\allowbreak{}theorem12\allowbreak{}One\allowbreak{}Comparison\allowbreak{}Incidences}
\item \hyperlink{governing-declaration-9f5f3a71633fcfda}{Golz\allowbreak{}Haghtalab\allowbreak{}Yang2025\allowbreak{}Distortion.\allowbreak{}theorem12\allowbreak{}One\allowbreak{}Comparison\allowbreak{}Report\allowbreak{}Law}
\item \hyperlink{governing-declaration-0d61e98c23f20a22}{Golz\allowbreak{}Haghtalab\allowbreak{}Yang2025\allowbreak{}Distortion.\allowbreak{}theorem12\allowbreak{}One\allowbreak{}Comparison\allowbreak{}Wins}
\item \hyperlink{governing-declaration-787bb3a4dec78d2f}{Golz\allowbreak{}Haghtalab\allowbreak{}Yang2025\allowbreak{}Distortion.\allowbreak{}theorem12\allowbreak{}Population}
\item \hyperlink{governing-declaration-35d01f4df2faa71d}{Golz\allowbreak{}Haghtalab\allowbreak{}Yang2025\allowbreak{}Distortion.\allowbreak{}theorem12\allowbreak{}Split\allowbreak{}Utility}
\item \hyperlink{governing-declaration-a73ce8ecb43a8abd}{Golz\allowbreak{}Haghtalab\allowbreak{}Yang2025\allowbreak{}Distortion.\allowbreak{}theorem12\allowbreak{}Utility}
\end{itemize}
\paragraph{Lean-elaborated claim roles}\begin{ReviewVerbatim}
1. conclusion: (∀ {beta : ℝ} (hbeta : 0 < beta),
    have gamma : ℝ := GolzHaghtalabYang2025Distortion.theorem12LogChoice beta;
    ∃ (hgamma : 0 < gamma) (hgamma_lt_one : gamma < 1),
      beta / 2 * (1 + Real.exp (-beta)) / (1 - Real.exp (-beta)) <
          GolzHaghtalabYang2025Distortion.theorem12Eq10Coefficient beta gamma ∧
        ∀ q < GolzHaghtalabYang2025Distortion.theorem12Eq10Coefficient beta gamma,
          ∀ (m : ℕ),
            3 ≤ m →
              ∃ (epsilon : ℝ) (hepsilon : 0 < epsilon) (hepsilon_lt_half : epsilon < 1 / 2) (sampling :
                PMF (Fin 3 ⊕ Fin (m - 3))),
                epsilon ^ 2 < 1 - epsilon ∧
                  Fintype.card (Fin 3 ⊕ Fin (m - 3)) = m ∧
                    have collapse : Fin 3 ⊕ Fin (m - 3) → Fin 3 :=
                      GolzHaghtalabYang2025Distortion.theorem12CloneCollapse 1 (m - 3);
                    have splitUtility :
                      AppliedModelingLib.Alignment.Welfare.FiniteUtilityProfile (Fin 3) (Fin 3 ⊕ Fin (m - 3)) :=
                      GolzHaghtalabYang2025Distortion.theorem12SplitUtility
                        (GolzHaghtalabYang2025Distortion.theorem12Utility epsilon (epsilon ^ 2) gamma) collapse;
                    have population : PMF (Fin 3) :=
                      GolzHaghtalabYang2025Distortion.theorem12Population beta epsilon gamma hbeta hepsilon
                        (GolzHaghtalabYang2025Distortion.theorem5Spec._proof_2 epsilon hepsilon hepsilon_lt_half) hgamma
                        hgamma_lt_one;
                    have preference :
                      AppliedModelingLib.Learning.HumanFeedback.PairwisePreference PUnit.{1} (Fin 3 ⊕ Fin (m - 3)) :=
                      AppliedModelingLib.Alignment.Welfare.populationBradleyTerryPreference population splitUtility
                        beta;
                    AppliedModelingLib.Alignment.Welfare.UnitIntervalUtilityProfile splitUtility ∧
                      AppliedModelingLib.PMFFullSupport sampling ∧
                        q <
                            AppliedModelingLib.Alignment.Welfare.populationAverageUtility population splitUtility
                                (Sum.inl 0) /
                              AppliedModelingLib.Alignment.Welfare.populationAverageUtility population splitUtility
                                (Sum.inl 2) ∧
                          Filter.Tendsto
                            (fun (horizon : ℕ) =>
                              AppliedModelingLib.pmfProb
                                (GolzHaghtalabYang2025Distortion.theorem12IidBordaReportBatchLaw
                                  (GolzHaghtalabYang2025Distortion.theorem12OneComparisonReportLaw sampling preference)
                                  horizon)
                                fun (sample : Fin horizon → ((Fin 3 ⊕ Fin (m - 3)) × (Fin 3 ⊕ Fin (m - 3))) × Bool) =>
                                GolzHaghtalabYang2025Distortion.theorem12IidEmpiricalBordaFailure
                                  GolzHaghtalabYang2025Distortion.theorem12OneComparisonWins
                                  GolzHaghtalabYang2025Distortion.theorem12OneComparisonIncidences (Sum.inl 2) sample)
                            Filter.atTop (nhds 0)) ∧
  Filter.Tendsto
    (fun (beta : ℝ) =>
      GolzHaghtalabYang2025Distortion.theorem12Eq10Coefficient beta
          (GolzHaghtalabYang2025Distortion.theorem12LogChoice beta) /
        beta)
    Filter.atTop (nhds 1)
\end{ReviewVerbatim}

\paragraph{Lean proof endpoint (built separately)}\begin{ReviewVerbatim}
GolzHaghtalabYang2025Distortion.theorem5_borda_lower_bound
\end{ReviewVerbatim}

\paragraph{Recorded source-to-Spec screening}
\textbf{Verdict:} \texttt{matches}
\\
{\raggedright
\textbf{Reason:} This realizes the source Borda/\allowbreak{}RLHF lower bound strictly above the universal coefficient and asymptotic to beta.\allowbreak{}
\par}
\reviewmatch{Matches source input}
\noindent\textbf{Reviewer annotation}\par
\reviewerbox
\clearpage
\hypertarget{source-claim-8c95df698cbd00a0}{}
\section*{Review row 8}
\textbf{Source locator:} {\footnotesize\raggedright cited publication, lines 1728-\allowbreak{}1734\par}
\paragraph{Verbatim source input}
\begin{ReviewVerbatim}
Theorem 12 (Lower Bound for Borda; Formal Version of Theorem 5). For any β > 0 and m ≥ 3, the
                                                                                           1+e−β
Borda voting rule (and, hence, RHLF) cannot guarantee a distortion better than max0<γ<1 β2 1−e −β ·
          σ(βγ)−1/2  [U+0001 control character]
 1 − γ + σ(β)−1/2 . This bound is strictly higher than the voting-rule independent lower bound
β 1+e−β
2 1−e−β for all β and is at least (1 − o(1)) β as β → ∞.

[Next verbatim source excerpt]

Let A = {1, . . . , m} be a finite set of alternatives. The population of users is described by a
probability distribution D over utility vectors u = (u(1), . . . , u(m)), whose entries 0 ≤ u(x) ≤ 1
indicate a user’s utility for alternative x.3 The objective is to find an alternative x such that its average
utility AvgUtil(x) := Eu∼D [u(x)] across the user population (also known as the utilitarian social
welfare) is as high as possible. We extend this notation to probability distributions π over alternatives
by setting AvgUtil(π) := Ex∼π [AvgUtil(x)].
Both voting rules and alignment methods observe comparisons from n users. We model each
i = 1, . . . , n as a fresh user with independently drawn utility vector ui ∼ D. For exposition, we
assume that each user i provides an equal number d ≥ 1 of pairwise comparisons. For each i,
and for j = 1, . . . , d, we independently draw alternatives xji , yij from a fixed distribution µ over

    3
     Our assumption that utilities be in [0, 1] is weaker than any of the three assumptions — unit-sum, unit-range,
or approval [23] — made in classic distortion to avoid trivial infinite lower bounds.


                                                           3

[form-feed in source extraction]
alternatives, in which the minimum probability mass µmin := minx∈A µ(x) is positive.4 User i
then compares each pair {xji , yij } (for j = 1, . . . , d) through a Bradley-Terry model based on i’s
utilities: i prefers xji over yij (written “xji [U+001F control character]i yij ”) with probability σ β · (ui (xji ) − ui (yij )) , where
                                                                                                        [U+0001 control character]

σ(t) := 1/(1+e−t ) is the logistic sigmoid function and β > 0 is a temperature parameter, and prefers
yij over xji (“yij [U+001F control character]i xji ”) otherwise.5 Whereas this specifies the marginal probability of each pairwise
comparison, we make no assumption about the correlation between i’s choices. For example, i might
derive the pairwise comparisons from a Plackett-Luce ranking, ensuring that      [U+0002 control character] the user’s comparisons  [U+0001 control character][U+0003 control character]
are always consistent.6 We set p(x [U+001F control character] y) for the expected win rate Eu∼D σ β · (u(x) − u(y)) .

[Next verbatim source excerpt]

Social Choice Setting. A voting rule f observes the sampled pairwise comparisons {xji [U+001F control character]i
yij }i∈[n],j∈[d] and maps them to a probability distribution over alternatives. For some m, D, β, d, µ,
and
  [U+0002 control character] the correlation between comparisons,     the average utility of f for n samples is AvgUtiln (f ) :=
E AvgUtil f ({xji [U+001F control character]i yij }i,j ) , where the expectation is taken over the pairwise comparisons. The
                               [U+0001 control character][U+0003 control character]

distortion of f on D is the competitive ratio between AvgUtiln (f ) and the optimal average utility
maxx∈A AvgUtil(x) in the limit of n → ∞ samples, and the distortion of f the worst-case distortion
over all D:

            dist(f, D) := lim supn→∞ maxAvgUtil
                                        x∈A AvgUtil(x)
                                                (f )   ,               dist(f ) = supD dist(f, D).
                                                       n



For alternatives x, y, let #(x [U+001F control character] y) := {(i, j) | xji = x, yij = y} denote the number of pairwise
comparisons in which x beat y. The (normalized) Borda score [45] of alternative x is
                                                       P
                                                         y∈A #(x[U+001F control character]y)
                                   BC(x) := P         #(x[U+001F control character]y)+
                                                              P             ,
                                                  y∈A           y6=x #(y[U+001F control character]x)


i.e., the fraction of pairwise comparisons involving x in which it wins. The Borda voting rule chooses
the winner uniformly among all alternatives with maximum Borda score. The Maximal Lotteries vot-
ing rule first computes the margin matrix M ∈ Rm×m , where Mx,y = #(x[U+001F control character]y)−#(y[U+001F control character]x)
                                                                          #(x[U+001F control character]y)+#(y[U+001F control character]x) . It then con-
siders a symmetric two-player zero-sum game in which player 1 selects alternative x1 , player 2 selects
alternative x2 , and the payoffs are Mx1 ,x2 for player 1 and Mx2 ,x1 = −Mx1 ,x2 for player 2. The max-
imal lotteries rule returns a distribution π ∈ argmaxπ1 ∈∆(A) minπ2 ∈∆(A) Ex1 ∼π1 ,x2 ∼π2 [Mx1 ,x2 ],
i.e., a mixed strategy in Nash equilibrium. When several such π exist, our results hold for any choice.
\end{ReviewVerbatim}



\paragraph{Expanded PaperInterface specification}
\begin{ReviewVerbatim}
(∀ {beta gamma q : ℝ} (hbeta : 0 < beta) (hgamma : 0 < gamma) (hgamma_lt_one : gamma < 1),
    q < GolzHaghtalabYang2025Distortion.theorem12Eq10Coefficient beta gamma →
      ∀ (m : ℕ),
        3 ≤ m →
          ∃ (epsilon : ℝ) (hepsilon : 0 < epsilon) (hepsilon_lt_half : epsilon < 1 / 2) (sampling :
            PMF (Fin 3 ⊕ Fin (m - 3))),
            epsilon ^ 2 < 1 - epsilon ∧
              Fintype.card (Fin 3 ⊕ Fin (m - 3)) = m ∧
                have collapse : Fin 3 ⊕ Fin (m - 3) → Fin 3 :=
                  GolzHaghtalabYang2025Distortion.theorem12CloneCollapse 1 (m - 3);
                have splitUtility :
                  AppliedModelingLib.Alignment.Welfare.FiniteUtilityProfile (Fin 3) (Fin 3 ⊕ Fin (m - 3)) :=
                  GolzHaghtalabYang2025Distortion.theorem12SplitUtility
                    (GolzHaghtalabYang2025Distortion.theorem12Utility epsilon (epsilon ^ 2) gamma) collapse;
                have population : PMF (Fin 3) :=
                  GolzHaghtalabYang2025Distortion.theorem12Population beta epsilon gamma hbeta hepsilon
                    (GolzHaghtalabYang2025Distortion.theorem5Spec._proof_2 epsilon hepsilon hepsilon_lt_half) hgamma
                    hgamma_lt_one;
                have preference :
                  AppliedModelingLib.Learning.HumanFeedback.PairwisePreference PUnit.{1} (Fin 3 ⊕ Fin (m - 3)) :=
                  AppliedModelingLib.Alignment.Welfare.populationBradleyTerryPreference population splitUtility beta;
                AppliedModelingLib.Alignment.Welfare.UnitIntervalUtilityProfile splitUtility ∧
                  AppliedModelingLib.PMFFullSupport sampling ∧
                    q <
                        AppliedModelingLib.Alignment.Welfare.populationAverageUtility population splitUtility
                            (Sum.inl 0) /
                          AppliedModelingLib.Alignment.Welfare.populationAverageUtility population splitUtility
                            (Sum.inl 2) ∧
                      Filter.Tendsto
                        (fun (horizon : ℕ) =>
                          AppliedModelingLib.pmfProb
                            (GolzHaghtalabYang2025Distortion.theorem12IidBordaReportBatchLaw
                              (GolzHaghtalabYang2025Distortion.theorem12OneComparisonReportLaw sampling preference)
                              horizon)
                            fun (sample : Fin horizon → ((Fin 3 ⊕ Fin (m - 3)) × (Fin 3 ⊕ Fin (m - 3))) × Bool) =>
                            GolzHaghtalabYang2025Distortion.theorem12IidEmpiricalBordaFailure
                              GolzHaghtalabYang2025Distortion.theorem12OneComparisonWins
                              GolzHaghtalabYang2025Distortion.theorem12OneComparisonIncidences (Sum.inl 2) sample)
                        Filter.atTop (nhds 0)) ∧
  (∀ {beta gamma : ℝ},
      0 < beta →
        0 < gamma →
          gamma < 1 →
            beta / 2 * (1 + Real.exp (-beta)) / (1 - Real.exp (-beta)) <
              GolzHaghtalabYang2025Distortion.theorem12Eq10Coefficient beta gamma) ∧
    Filter.Tendsto
      (fun (beta : ℝ) =>
        GolzHaghtalabYang2025Distortion.theorem12Eq10Coefficient beta
            (GolzHaghtalabYang2025Distortion.theorem12LogChoice beta) /
          beta)
      Filter.atTop (nhds 1)
\end{ReviewVerbatim}

\paragraph{Retained governing declarations}
\begin{itemize}\raggedright
\item \hyperlink{library-prerequisite-56c0734f8877ddbe}{Applied\allowbreak{}Modeling\allowbreak{}Lib.\allowbreak{}Alignment.\allowbreak{}Welfare.\allowbreak{}Finite\allowbreak{}Utility\allowbreak{}Profile}
\item \hyperlink{library-prerequisite-cbe8afa0d6d1ba78}{Applied\allowbreak{}Modeling\allowbreak{}Lib.\allowbreak{}Alignment.\allowbreak{}Welfare.\allowbreak{}Unit\allowbreak{}Interval\allowbreak{}Utility\allowbreak{}Profile}
\item \hyperlink{library-prerequisite-15f4b442749deab2}{Applied\allowbreak{}Modeling\allowbreak{}Lib.\allowbreak{}Alignment.\allowbreak{}Welfare.\allowbreak{}population\allowbreak{}Average\allowbreak{}Utility}
\item \hyperlink{library-prerequisite-cf06d15d7f8b7c31}{Applied\allowbreak{}Modeling\allowbreak{}Lib.\allowbreak{}Alignment.\allowbreak{}Welfare.\allowbreak{}population\allowbreak{}Bradley\allowbreak{}Terry\allowbreak{}Preference}
\item \hyperlink{library-prerequisite-7e2db20cf087d9ba}{Applied\allowbreak{}Modeling\allowbreak{}Lib.\allowbreak{}Learning.\allowbreak{}Human\allowbreak{}Feedback.\allowbreak{}Pairwise\allowbreak{}Preference}
\item \hyperlink{library-prerequisite-268c959ee7d2d148}{Applied\allowbreak{}Modeling\allowbreak{}Lib.\allowbreak{}PMF\allowbreak{}Full\allowbreak{}Support}
\item \hyperlink{governing-declaration-9c38d53b9d664221}{Applied\allowbreak{}Modeling\allowbreak{}Lib.\allowbreak{}pmf\allowbreak{}Prob}
\item \hyperlink{governing-declaration-ec899679d7cb6165}{Golz\allowbreak{}Haghtalab\allowbreak{}Yang2025\allowbreak{}Distortion.\allowbreak{}theorem12\allowbreak{}Clone\allowbreak{}Collapse}
\item \hyperlink{governing-declaration-b2a8d113fc3130c2}{Golz\allowbreak{}Haghtalab\allowbreak{}Yang2025\allowbreak{}Distortion.\allowbreak{}theorem12\allowbreak{}Empirical\allowbreak{}Borda\allowbreak{}Score}
\item \hyperlink{governing-declaration-aec5a9f50c9e967a}{Golz\allowbreak{}Haghtalab\allowbreak{}Yang2025\allowbreak{}Distortion.\allowbreak{}theorem12\allowbreak{}Eq10\allowbreak{}Coefficient}
\item \hyperlink{governing-declaration-e5cae00c98b32d81}{Golz\allowbreak{}Haghtalab\allowbreak{}Yang2025\allowbreak{}Distortion.\allowbreak{}theorem12\allowbreak{}Iid\allowbreak{}Borda\allowbreak{}Report\allowbreak{}Batch\allowbreak{}Law}
\item \hyperlink{governing-declaration-3940e2eb509393ea}{Golz\allowbreak{}Haghtalab\allowbreak{}Yang2025\allowbreak{}Distortion.\allowbreak{}theorem12\allowbreak{}Iid\allowbreak{}Empirical\allowbreak{}Borda\allowbreak{}Failure}
\item \hyperlink{governing-declaration-df5d365b55936500}{Golz\allowbreak{}Haghtalab\allowbreak{}Yang2025\allowbreak{}Distortion.\allowbreak{}theorem12\allowbreak{}Log\allowbreak{}Choice}
\item \hyperlink{governing-declaration-baa615edc2baf254}{Golz\allowbreak{}Haghtalab\allowbreak{}Yang2025\allowbreak{}Distortion.\allowbreak{}theorem12\allowbreak{}One\allowbreak{}Comparison\allowbreak{}Incidences}
\item \hyperlink{governing-declaration-9f5f3a71633fcfda}{Golz\allowbreak{}Haghtalab\allowbreak{}Yang2025\allowbreak{}Distortion.\allowbreak{}theorem12\allowbreak{}One\allowbreak{}Comparison\allowbreak{}Report\allowbreak{}Law}
\item \hyperlink{governing-declaration-0d61e98c23f20a22}{Golz\allowbreak{}Haghtalab\allowbreak{}Yang2025\allowbreak{}Distortion.\allowbreak{}theorem12\allowbreak{}One\allowbreak{}Comparison\allowbreak{}Wins}
\item \hyperlink{governing-declaration-787bb3a4dec78d2f}{Golz\allowbreak{}Haghtalab\allowbreak{}Yang2025\allowbreak{}Distortion.\allowbreak{}theorem12\allowbreak{}Population}
\item \hyperlink{governing-declaration-35d01f4df2faa71d}{Golz\allowbreak{}Haghtalab\allowbreak{}Yang2025\allowbreak{}Distortion.\allowbreak{}theorem12\allowbreak{}Split\allowbreak{}Utility}
\item \hyperlink{governing-declaration-a73ce8ecb43a8abd}{Golz\allowbreak{}Haghtalab\allowbreak{}Yang2025\allowbreak{}Distortion.\allowbreak{}theorem12\allowbreak{}Utility}
\end{itemize}
\paragraph{Lean-elaborated claim roles}\begin{ReviewVerbatim}
1. conclusion: (∀ {beta gamma q : ℝ} (hbeta : 0 < beta) (hgamma : 0 < gamma) (hgamma_lt_one : gamma < 1),
    q < GolzHaghtalabYang2025Distortion.theorem12Eq10Coefficient beta gamma →
      ∀ (m : ℕ),
        3 ≤ m →
          ∃ (epsilon : ℝ) (hepsilon : 0 < epsilon) (hepsilon_lt_half : epsilon < 1 / 2) (sampling :
            PMF (Fin 3 ⊕ Fin (m - 3))),
            epsilon ^ 2 < 1 - epsilon ∧
              Fintype.card (Fin 3 ⊕ Fin (m - 3)) = m ∧
                have collapse : Fin 3 ⊕ Fin (m - 3) → Fin 3 :=
                  GolzHaghtalabYang2025Distortion.theorem12CloneCollapse 1 (m - 3);
                have splitUtility :
                  AppliedModelingLib.Alignment.Welfare.FiniteUtilityProfile (Fin 3) (Fin 3 ⊕ Fin (m - 3)) :=
                  GolzHaghtalabYang2025Distortion.theorem12SplitUtility
                    (GolzHaghtalabYang2025Distortion.theorem12Utility epsilon (epsilon ^ 2) gamma) collapse;
                have population : PMF (Fin 3) :=
                  GolzHaghtalabYang2025Distortion.theorem12Population beta epsilon gamma hbeta hepsilon
                    (GolzHaghtalabYang2025Distortion.theorem5Spec._proof_2 epsilon hepsilon hepsilon_lt_half) hgamma
                    hgamma_lt_one;
                have preference :
                  AppliedModelingLib.Learning.HumanFeedback.PairwisePreference PUnit.{1} (Fin 3 ⊕ Fin (m - 3)) :=
                  AppliedModelingLib.Alignment.Welfare.populationBradleyTerryPreference population splitUtility beta;
                AppliedModelingLib.Alignment.Welfare.UnitIntervalUtilityProfile splitUtility ∧
                  AppliedModelingLib.PMFFullSupport sampling ∧
                    q <
                        AppliedModelingLib.Alignment.Welfare.populationAverageUtility population splitUtility
                            (Sum.inl 0) /
                          AppliedModelingLib.Alignment.Welfare.populationAverageUtility population splitUtility
                            (Sum.inl 2) ∧
                      Filter.Tendsto
                        (fun (horizon : ℕ) =>
                          AppliedModelingLib.pmfProb
                            (GolzHaghtalabYang2025Distortion.theorem12IidBordaReportBatchLaw
                              (GolzHaghtalabYang2025Distortion.theorem12OneComparisonReportLaw sampling preference)
                              horizon)
                            fun (sample : Fin horizon → ((Fin 3 ⊕ Fin (m - 3)) × (Fin 3 ⊕ Fin (m - 3))) × Bool) =>
                            GolzHaghtalabYang2025Distortion.theorem12IidEmpiricalBordaFailure
                              GolzHaghtalabYang2025Distortion.theorem12OneComparisonWins
                              GolzHaghtalabYang2025Distortion.theorem12OneComparisonIncidences (Sum.inl 2) sample)
                        Filter.atTop (nhds 0)) ∧
  (∀ {beta gamma : ℝ},
      0 < beta →
        0 < gamma →
          gamma < 1 →
            beta / 2 * (1 + Real.exp (-beta)) / (1 - Real.exp (-beta)) <
              GolzHaghtalabYang2025Distortion.theorem12Eq10Coefficient beta gamma) ∧
    Filter.Tendsto
      (fun (beta : ℝ) =>
        GolzHaghtalabYang2025Distortion.theorem12Eq10Coefficient beta
            (GolzHaghtalabYang2025Distortion.theorem12LogChoice beta) /
          beta)
      Filter.atTop (nhds 1)
\end{ReviewVerbatim}

\paragraph{Lean proof endpoint (built separately)}\begin{ReviewVerbatim}
GolzHaghtalabYang2025Distortion.theorem12_borda_lower_bound
\end{ReviewVerbatim}

\paragraph{Recorded source-to-Spec screening}
\textbf{Verdict:} \texttt{matches}
\\
{\raggedright
\textbf{Reason:} The all-\allowbreak{}m Borda lower-\allowbreak{}bound construction, its strict separation from the universal factor, and the asymptotic coefficient are retained.\allowbreak{}
\par}
\reviewmatch{Matches source input}
\noindent\textbf{Reviewer annotation}\par
\reviewerbox
\clearpage
\hypertarget{source-claim-b1ccdc13d898700c}{}
\section*{Review row 9}
\textbf{Source locator:} {\footnotesize\raggedright cited publication, lines 494-\allowbreak{}496\par}
\paragraph{Verbatim source input}
\begin{ReviewVerbatim}
Theorem 6 (RLHF Distortion Lower Bound). For m ≥ 3, there is a sequence of alignment problems
on which the distortion of RLHF scales as eΩ(β) in β.


[Next verbatim source excerpt]

Let A = {1, . . . , m} be a finite set of alternatives. The population of users is described by a
probability distribution D over utility vectors u = (u(1), . . . , u(m)), whose entries 0 ≤ u(x) ≤ 1
indicate a user’s utility for alternative x.3 The objective is to find an alternative x such that its average
utility AvgUtil(x) := Eu∼D [u(x)] across the user population (also known as the utilitarian social
welfare) is as high as possible. We extend this notation to probability distributions π over alternatives
by setting AvgUtil(π) := Ex∼π [AvgUtil(x)].
Both voting rules and alignment methods observe comparisons from n users. We model each
i = 1, . . . , n as a fresh user with independently drawn utility vector ui ∼ D. For exposition, we
assume that each user i provides an equal number d ≥ 1 of pairwise comparisons. For each i,
and for j = 1, . . . , d, we independently draw alternatives xji , yij from a fixed distribution µ over

    3
     Our assumption that utilities be in [0, 1] is weaker than any of the three assumptions — unit-sum, unit-range,
or approval [23] — made in classic distortion to avoid trivial infinite lower bounds.


                                                           3

[form-feed in source extraction]
alternatives, in which the minimum probability mass µmin := minx∈A µ(x) is positive.4 User i
then compares each pair {xji , yij } (for j = 1, . . . , d) through a Bradley-Terry model based on i’s
utilities: i prefers xji over yij (written “xji [U+001F control character]i yij ”) with probability σ β · (ui (xji ) − ui (yij )) , where
                                                                                                        [U+0001 control character]

σ(t) := 1/(1+e−t ) is the logistic sigmoid function and β > 0 is a temperature parameter, and prefers
yij over xji (“yij [U+001F control character]i xji ”) otherwise.5 Whereas this specifies the marginal probability of each pairwise
comparison, we make no assumption about the correlation between i’s choices. For example, i might
derive the pairwise comparisons from a Plackett-Luce ranking, ensuring that      [U+0002 control character] the user’s comparisons  [U+0001 control character][U+0003 control character]
are always consistent.6 We set p(x [U+001F control character] y) for the expected win rate Eu∼D σ β · (u(x) − u(y)) .

[Next verbatim source excerpt]

Alignment Setting. The alignment setting generalizes the social choice setting in two ways: first,
user utilities ui (y | x) may depend on a state x; second, the goal in determining a policy π is not
purely to maximize the reward AvgUtil(π | x), but a trade-off between this reward and the goal of
remaining close to a reference policy πref (· | x) ∈ ∆(A) in terms of the KL divergence DKL (· k πref ).
For theoretical tractability, we focus our analysis on a single state x, which we from here on omit
from the notation. Conceptually, this treatment of alignment on a state-by-state basis corresponds to
an assumption that our policy class is expressive enough so that it can take the optimal distribution of
actions at each state7 and abstracts from the generalization problem of estimating the population’s
preference between a pair of alternatives at the given state x based on preferences in similar states x0 .
Having set aside the dependency between states, we focus on how the regularization with respect
to a reference policy impacts the ability to optimize the average utility of three alignment methods:
RLHF, DPO, and NLHF. In addition to the pairwise comparisons, these methods take in a reference
policy πref ∈ ∆(A) and a KL bound τ ≥ 0, and map these inputs to a policy π in the KL-ball
Bτ (πref ) := {π ∈ ∆(A) | DKL (π k πref ) ≤ τ } around the reference policy. RLHF, DPO, and NLHF
are typically implemented with a KL regularization rather than our KL-constrained formulation,
which we adopt to enable a comparison on equal terms. In Section 5, we show that these perspectives
are equivalent, and that distortion upper bounds carry over to regularized alignment methods.

   4
     In Section 5, we discuss how most of our results extend more general distributions over comparison pairs,
which can, for example. capture k-wise comparisons.
   5
     Should we sample the same alternative x = xji = yij twice for a pair, the user is not asked for a pairwise
comparison. We record this as “x [U+001F control character]i x”, in a slight abuse of notation.
   6
     In particular, if we sample the same unordered pair twice for a user, the answers can be perfectly correlated.
   7
     This assumption is common in the literature; see, for example, Rafailov et al. [43]’s application of first-order
optimality conditions of PPO loss minimization.


                                                           4

[form-feed in source extraction]
The RLHF method first estimates rewards for each alternative, using maximum likelihood estimation
assuming that comparisons were generated by a single Bradley–Terry model:
                  r := argmaxr∈Rm 1≤i≤n,1≤j≤d log σ(r(xji ) − r(yij )) .8
                                     P                                     [U+0001 control character]

Next, RLHF uses PPO [44] to compute the policy πRLHF := argmaxπ∈Bτ (πref ) Ex∼π [r(x)] with
maximum expected reward within the KL-ball. In our setup, DPO is equivalent to RLHF (see
Appendix F.3) and thus has the same distortion.
The alignment method NLHF was inspired in part by a desire to better align with the
preferences of a heterogeneous group [36]. NLHF naturally adapts the definition of maxi-
mal lotteries by constraining both players’ mixed strategies to the KL-ball, i.e., πNLHF :=
argmaxπ1 ∈Bτ (πref ) minπ2 ∈Bτ (πref ) Ex1 ∼π1 ,x2 ∼π2 [Mx1 ,x2 ].
To generalize the definition of distortion to alignment methods, we set the maximum average utility of
any policy in the KL-ball as the benchmark. For fixed m, D, β, d, µ and correlation between pairwise
comparisons, the distortion of alignment method f is
                                                             maxπ∈Bτ (πref ) AvgUtil(π)
                      dist(f ) = supD,πref ,τ lim supn→∞       AvgUtiln (f (·,πref ,τ )) .

[Next verbatim source excerpt]

Theorem 6 (RLHF Distortion Lower Bound). For m ≥ 3, there is a sequence of alignment problems
on which the distortion of RLHF scales as eΩ(β) in β.

Proof of Theorem 6. Suppose that the instance has m alternatives A = {a, b, c1 , . . . , cm−2 }, where
m−2 ≥ 4eβ . Let the data collection distribution be uniform over all candidates, i.e., µ = Uniform(A).
We consider the following distribution D over utility vectors, such that the utility vector of a random
agent i ∼ D satisfies
                                                   (
                                                     (0, 1, 0, . . . , 0) [U+0011 control character] (type I) with probability δ,
   (ui (a), ui (b), ui (c1 ), . . . , ui (cm−2 )) = [U+0010 control character] 1
                                                       β , 0, 1, . . . , 1  (type II) with probability 1 − δ,
             10           −β
where δ = 10+e  β = Θ(e      ). In other words, type I users have a strong preference for candidate b
but only constitute a δ fraction of the population, while type II users have a strong preference for c
and weak preference for a [U+001F control character] b and make up for a 1 − δ fraction of the population.
For the reference policy and the KL budget, we set πref (a) = πref (b) = 1−ε                   ε
                                                                          2 and πref (ci ) = m−2 for
all i ∈ [m − 2]. We leave the choice of ε to be determined later. The KL budget τ is set to be τ = 1.


                                                              32

[form-feed in source extraction]
Analysis of the MLE reward. Now we show that when n → ∞, the MLE reward satisfies
r(b) − r(a) > 0. According to [46, 42], it suffices to show that limn→∞ BCn (b) − BCn (a) > 0,

[Next verbatim source excerpt]


[form-feed in source extraction]
       • For πRLHF , we assume that the estimated reward is shifted such that r(c) = 0 (as a result,
         r(a) < r(b) < 0). Since π 0 = (0, 1, 0) also satisfies the KL constraint, we have r(πRLHF ) ≥
         r(π 0 ). Together with the fact that πRLHF (c) ≤ η, we have
\end{ReviewVerbatim}



\paragraph{Expanded PaperInterface specification}
\begin{ReviewVerbatim}
∀ {beta : ℝ},
  10 ≤ beta →
    ∃ (eta : ℝ) (heta_pos : 0 < eta) (_ : eta ≤ 1),
      3 ≤
          Fintype.card
            (GolzHaghtalabYang2025Distortion.theorem6Alternative
              (GolzHaghtalabYang2025Distortion.theorem6CopyCount beta)) ∧
        ∀ᵐ (path :
          ℕ →
            AppliedModelingLib.Learning.HumanFeedback.BinaryPairwiseReport
              (GolzHaghtalabYang2025Distortion.theorem6Alternative
                (GolzHaghtalabYang2025Distortion.theorem6CopyCount
                  beta))) ∂AppliedModelingLib.finitePMFIidPathMeasure
            (GolzHaghtalabYang2025Distortion.theorem12OneComparisonReportLaw
              (AppliedModelingLib.uniformPMF
                (GolzHaghtalabYang2025Distortion.theorem6Alternative
                  (GolzHaghtalabYang2025Distortion.theorem6CopyCount beta)))
              (GolzHaghtalabYang2025Distortion.theorem6Preference beta)),
          ∀ᶠ (horizon : ℕ) in Filter.atTop,
            ∀
              (reward :
                AppliedModelingLib.Learning.HumanFeedback.PairwiseCountDataset.ScoreVector
                  (GolzHaghtalabYang2025Distortion.theorem6Alternative
                    (GolzHaghtalabYang2025Distortion.theorem6CopyCount beta))),
              (AppliedModelingLib.Learning.HumanFeedback.PairwiseCountDataset.ofBinaryReports
                      (AppliedModelingLib.finitePMFIidPathPrefix path horizon)).isPairwiseMLE
                  Real.sigmoid GolzHaghtalabYang2025Distortion.theorem6A reward →
                ∃ (rewardPolicy :
                  PMF
                    (GolzHaghtalabYang2025Distortion.theorem6Alternative
                      (GolzHaghtalabYang2025Distortion.theorem6CopyCount beta)))
                  (welfarePolicy :
                  PMF
                    (GolzHaghtalabYang2025Distortion.theorem6Alternative
                      (GolzHaghtalabYang2025Distortion.theorem6CopyCount beta))),
                  AppliedModelingLib.finiteKLDivergence rewardPolicy
                        (GolzHaghtalabYang2025Distortion.theorem6Reference
                          (GolzHaghtalabYang2025Distortion.theorem6Epsilon eta)
                          (GolzHaghtalabYang2025Distortion.theorem6Epsilon_pos eta)
                          (GolzHaghtalabYang2025Distortion.theorem6Epsilon_lt_one heta_pos)
                          (GolzHaghtalabYang2025Distortion.theorem6CopyCount_pos beta)) ≤
                      1 ∧
                    (∀
                        (other :
                          PMF
                            (GolzHaghtalabYang2025Distortion.theorem6Alternative
                              (GolzHaghtalabYang2025Distortion.theorem6CopyCount beta))),
                        AppliedModelingLib.finiteKLDivergence other
                              (GolzHaghtalabYang2025Distortion.theorem6Reference
                                (GolzHaghtalabYang2025Distortion.theorem6Epsilon eta)
                                (GolzHaghtalabYang2025Distortion.theorem6Epsilon_pos eta)
                                (GolzHaghtalabYang2025Distortion.theorem6Epsilon_lt_one heta_pos)
                                (GolzHaghtalabYang2025Distortion.theorem6CopyCount_pos beta)) ≤
                            1 →
                          AppliedModelingLib.pmfExp other reward ≤ AppliedModelingLib.pmfExp rewardPolicy reward) ∧
                      AppliedModelingLib.finiteKLDivergence welfarePolicy
                            (GolzHaghtalabYang2025Distortion.theorem6Reference
                              (GolzHaghtalabYang2025Distortion.theorem6Epsilon eta)
                              (GolzHaghtalabYang2025Distortion.theorem6Epsilon_pos eta)
                              (GolzHaghtalabYang2025Distortion.theorem6Epsilon_lt_one heta_pos)
                              (GolzHaghtalabYang2025Distortion.theorem6CopyCount_pos beta)) ≤
                          1 ∧
                        (∀
                            (other :
                              PMF
                                (GolzHaghtalabYang2025Distortion.theorem6Alternative
                                  (GolzHaghtalabYang2025Distortion.theorem6CopyCount beta))),
                            AppliedModelingLib.finiteKLDivergence other
                                  (GolzHaghtalabYang2025Distortion.theorem6Reference
                                    (GolzHaghtalabYang2025Distortion.theorem6Epsilon eta)
                                    (GolzHaghtalabYang2025Distortion.theorem6Epsilon_pos eta)
                                    (GolzHaghtalabYang2025Distortion.theorem6Epsilon_lt_one heta_pos)
                                    (GolzHaghtalabYang2025Distortion.theorem6CopyCount_pos beta)) ≤
                                1 →
                              AppliedModelingLib.pmfExp other (GolzHaghtalabYang2025Distortion.theorem6Welfare beta) ≤
                                AppliedModelingLib.pmfExp welfarePolicy
                                  (GolzHaghtalabYang2025Distortion.theorem6Welfare beta)) ∧
                          Real.exp beta / (44 * beta) ≤
                            AppliedModelingLib.pmfExp welfarePolicy
                                (GolzHaghtalabYang2025Distortion.theorem6Welfare beta) /
                              AppliedModelingLib.pmfExp rewardPolicy
                                (GolzHaghtalabYang2025Distortion.theorem6Welfare beta)
\end{ReviewVerbatim}

\paragraph{Retained governing declarations}
\begin{itemize}\raggedright
\item \hyperlink{governing-declaration-bc3d6ca8e7e0ac9d}{Applied\allowbreak{}Modeling\allowbreak{}Lib.\allowbreak{}Learning.\allowbreak{}Human\allowbreak{}Feedback.\allowbreak{}Binary\allowbreak{}Pairwise\allowbreak{}Report}
\item \hyperlink{governing-declaration-87eda978eb1b0c9d}{Applied\allowbreak{}Modeling\allowbreak{}Lib.\allowbreak{}Learning.\allowbreak{}Human\allowbreak{}Feedback.\allowbreak{}Pairwise\allowbreak{}Count\allowbreak{}Dataset.\allowbreak{}Score\allowbreak{}Vector}
\item \hyperlink{governing-declaration-cba21d23e0e5aa10}{Applied\allowbreak{}Modeling\allowbreak{}Lib.\allowbreak{}Learning.\allowbreak{}Human\allowbreak{}Feedback.\allowbreak{}Pairwise\allowbreak{}Count\allowbreak{}Dataset.\allowbreak{}is\allowbreak{}Pairwise\allowbreak{}MLE}
\item \hyperlink{governing-declaration-449f87e91bf37ce0}{Applied\allowbreak{}Modeling\allowbreak{}Lib.\allowbreak{}Learning.\allowbreak{}Human\allowbreak{}Feedback.\allowbreak{}Pairwise\allowbreak{}Count\allowbreak{}Dataset.\allowbreak{}of\allowbreak{}Binary\allowbreak{}Reports}
\item \hyperlink{library-prerequisite-d3573881ae45d4c6}{Applied\allowbreak{}Modeling\allowbreak{}Lib.\allowbreak{}finite\allowbreak{}KL\allowbreak{}Divergence}
\item \hyperlink{governing-declaration-779fe1555fbbcea5}{Applied\allowbreak{}Modeling\allowbreak{}Lib.\allowbreak{}finite\allowbreak{}PMF\allowbreak{}Iid\allowbreak{}Path\allowbreak{}Measure}
\item \hyperlink{governing-declaration-89219e02ea5145c7}{Applied\allowbreak{}Modeling\allowbreak{}Lib.\allowbreak{}finite\allowbreak{}PMF\allowbreak{}Iid\allowbreak{}Path\allowbreak{}Prefix}
\item \hyperlink{governing-declaration-8e4b4389a4f6454b}{Applied\allowbreak{}Modeling\allowbreak{}Lib.\allowbreak{}pmf\allowbreak{}Exp}
\item \hyperlink{governing-declaration-48f10f2715c04df3}{Applied\allowbreak{}Modeling\allowbreak{}Lib.\allowbreak{}uniform\allowbreak{}PMF}
\item \hyperlink{governing-declaration-9f5f3a71633fcfda}{Golz\allowbreak{}Haghtalab\allowbreak{}Yang2025\allowbreak{}Distortion.\allowbreak{}theorem12\allowbreak{}One\allowbreak{}Comparison\allowbreak{}Report\allowbreak{}Law}
\item \hyperlink{governing-declaration-7b9a203f54b7d337}{Golz\allowbreak{}Haghtalab\allowbreak{}Yang2025\allowbreak{}Distortion.\allowbreak{}theorem6\allowbreak{}A}
\item \hyperlink{governing-declaration-3451d41a5ff413f7}{Golz\allowbreak{}Haghtalab\allowbreak{}Yang2025\allowbreak{}Distortion.\allowbreak{}theorem6\allowbreak{}Alternative}
\item \hyperlink{governing-declaration-a14b7af0e1436262}{Golz\allowbreak{}Haghtalab\allowbreak{}Yang2025\allowbreak{}Distortion.\allowbreak{}theorem6\allowbreak{}Copy\allowbreak{}Count}
\item \hyperlink{governing-declaration-2108be1d9586e812}{Golz\allowbreak{}Haghtalab\allowbreak{}Yang2025\allowbreak{}Distortion.\allowbreak{}theorem6\allowbreak{}Epsilon}
\item \hyperlink{governing-declaration-b14cd971fad8513c}{Golz\allowbreak{}Haghtalab\allowbreak{}Yang2025\allowbreak{}Distortion.\allowbreak{}theorem6\allowbreak{}Preference}
\item \hyperlink{governing-declaration-0527d6356fcaa0f2}{Golz\allowbreak{}Haghtalab\allowbreak{}Yang2025\allowbreak{}Distortion.\allowbreak{}theorem6\allowbreak{}Reference}
\item \hyperlink{governing-declaration-ea43bde8cd383b24}{Golz\allowbreak{}Haghtalab\allowbreak{}Yang2025\allowbreak{}Distortion.\allowbreak{}theorem6\allowbreak{}Welfare}
\end{itemize}
\paragraph{Lean-elaborated claim roles}\begin{ReviewVerbatim}
1. parameter: ℝ
2. assumption: 10 ≤ beta
3. conclusion: ∃ (eta : ℝ) (heta_pos : 0 < eta) (_ : eta ≤ 1),
  3 ≤
      Fintype.card
        (GolzHaghtalabYang2025Distortion.theorem6Alternative (GolzHaghtalabYang2025Distortion.theorem6CopyCount beta)) ∧
    ∀ᵐ (path :
      ℕ →
        AppliedModelingLib.Learning.HumanFeedback.BinaryPairwiseReport
          (GolzHaghtalabYang2025Distortion.theorem6Alternative
            (GolzHaghtalabYang2025Distortion.theorem6CopyCount
              beta))) ∂AppliedModelingLib.finitePMFIidPathMeasure
        (GolzHaghtalabYang2025Distortion.theorem12OneComparisonReportLaw
          (AppliedModelingLib.uniformPMF
            (GolzHaghtalabYang2025Distortion.theorem6Alternative
              (GolzHaghtalabYang2025Distortion.theorem6CopyCount beta)))
          (GolzHaghtalabYang2025Distortion.theorem6Preference beta)),
      ∀ᶠ (horizon : ℕ) in Filter.atTop,
        ∀
          (reward :
            AppliedModelingLib.Learning.HumanFeedback.PairwiseCountDataset.ScoreVector
              (GolzHaghtalabYang2025Distortion.theorem6Alternative
                (GolzHaghtalabYang2025Distortion.theorem6CopyCount beta))),
          (AppliedModelingLib.Learning.HumanFeedback.PairwiseCountDataset.ofBinaryReports
                  (AppliedModelingLib.finitePMFIidPathPrefix path horizon)).isPairwiseMLE
              Real.sigmoid GolzHaghtalabYang2025Distortion.theorem6A reward →
            ∃ (rewardPolicy :
              PMF
                (GolzHaghtalabYang2025Distortion.theorem6Alternative
                  (GolzHaghtalabYang2025Distortion.theorem6CopyCount beta)))
              (welfarePolicy :
              PMF
                (GolzHaghtalabYang2025Distortion.theorem6Alternative
                  (GolzHaghtalabYang2025Distortion.theorem6CopyCount beta))),
              AppliedModelingLib.finiteKLDivergence rewardPolicy
                    (GolzHaghtalabYang2025Distortion.theorem6Reference
                      (GolzHaghtalabYang2025Distortion.theorem6Epsilon eta)
                      (GolzHaghtalabYang2025Distortion.theorem6Epsilon_pos eta)
                      (GolzHaghtalabYang2025Distortion.theorem6Epsilon_lt_one heta_pos)
                      (GolzHaghtalabYang2025Distortion.theorem6CopyCount_pos beta)) ≤
                  1 ∧
                (∀
                    (other :
                      PMF
                        (GolzHaghtalabYang2025Distortion.theorem6Alternative
                          (GolzHaghtalabYang2025Distortion.theorem6CopyCount beta))),
                    AppliedModelingLib.finiteKLDivergence other
                          (GolzHaghtalabYang2025Distortion.theorem6Reference
                            (GolzHaghtalabYang2025Distortion.theorem6Epsilon eta)
                            (GolzHaghtalabYang2025Distortion.theorem6Epsilon_pos eta)
                            (GolzHaghtalabYang2025Distortion.theorem6Epsilon_lt_one heta_pos)
                            (GolzHaghtalabYang2025Distortion.theorem6CopyCount_pos beta)) ≤
                        1 →
                      AppliedModelingLib.pmfExp other reward ≤ AppliedModelingLib.pmfExp rewardPolicy reward) ∧
                  AppliedModelingLib.finiteKLDivergence welfarePolicy
                        (GolzHaghtalabYang2025Distortion.theorem6Reference
                          (GolzHaghtalabYang2025Distortion.theorem6Epsilon eta)
                          (GolzHaghtalabYang2025Distortion.theorem6Epsilon_pos eta)
                          (GolzHaghtalabYang2025Distortion.theorem6Epsilon_lt_one heta_pos)
                          (GolzHaghtalabYang2025Distortion.theorem6CopyCount_pos beta)) ≤
                      1 ∧
                    (∀
                        (other :
                          PMF
                            (GolzHaghtalabYang2025Distortion.theorem6Alternative
                              (GolzHaghtalabYang2025Distortion.theorem6CopyCount beta))),
                        AppliedModelingLib.finiteKLDivergence other
                              (GolzHaghtalabYang2025Distortion.theorem6Reference
                                (GolzHaghtalabYang2025Distortion.theorem6Epsilon eta)
                                (GolzHaghtalabYang2025Distortion.theorem6Epsilon_pos eta)
                                (GolzHaghtalabYang2025Distortion.theorem6Epsilon_lt_one heta_pos)
                                (GolzHaghtalabYang2025Distortion.theorem6CopyCount_pos beta)) ≤
                            1 →
                          AppliedModelingLib.pmfExp other (GolzHaghtalabYang2025Distortion.theorem6Welfare beta) ≤
                            AppliedModelingLib.pmfExp welfarePolicy
                              (GolzHaghtalabYang2025Distortion.theorem6Welfare beta)) ∧
                      Real.exp beta / (44 * beta) ≤
                        AppliedModelingLib.pmfExp welfarePolicy (GolzHaghtalabYang2025Distortion.theorem6Welfare beta) /
                          AppliedModelingLib.pmfExp rewardPolicy (GolzHaghtalabYang2025Distortion.theorem6Welfare beta)
\end{ReviewVerbatim}

\paragraph{Lean proof endpoint (built separately)}\begin{ReviewVerbatim}
GolzHaghtalabYang2025Distortion.theorem6_rlhf_distortion
\end{ReviewVerbatim}

\paragraph{Recorded source-to-Spec screening}
\textbf{Verdict:} \texttt{matches}
\\
{\raggedright
\textbf{Reason:} Matches the source-\allowbreak{}defined asymptotic reading:\allowbreak{} Theorem 6 asserts a sequence of finite alignment problems with e\textasciicircum{}\{Ω(β)\} distortion, and Appendix F.\allowbreak{}2 itself requires m -\allowbreak{} 2 ≥ 4 exp(β).\allowbreak{} The expanded Spec makes that sequence explicit through m(β), witnesses the asymptotic tail by β ≥ 10, and derives the auxiliary exp(β) ≥ 100 regime internally;\allowbreak{} it does not add a fixed-\allowbreak{}m or extra main-\allowbreak{}text premise.\allowbreak{}
\par}
\reviewmatch{Matches source input}
\noindent\textbf{Reviewer annotation}\par
\reviewerbox
\clearpage
\hypertarget{source-claim-5067ef3ec8c43462}{}
\section*{Review row 10}
\textbf{Source locator:} {\footnotesize\raggedright cited publication, lines 543-\allowbreak{}558\par}
\paragraph{Verbatim source input}
\begin{ReviewVerbatim}
Theorem 7 (NLHF Distortion Upper Bound). For any instance D and any m, data distribution µ,
temperature β of the Bradley-Terry model, and any reference policy πref and KL budget τ , we have
                   1+e−β
dist(NLHF) ≤ β2 · 1−e −β . In the finite-sample regime, we have

                             −β [U+0001 control character]
                                                                          [U+0010 control character] q                           [U+0011 control character]
                                                                                  log(mn)       log(mn)
AvgUtiln (NLHF) ≥ β2 · 1−e
                         1+e −β    · max   ?
                                          π ∈Bτ (πref ) AvgUtil(π ?
                                                                    ) − O  1
                                                                           β  n·min{1, d·µ2 }
                                                                                              + n·βµ 2    .
                                                                                                     min         min


[Next verbatim source excerpt]

      If τ = 0, the only feasible policy is πref , and the claim clearly holds for λ = ∞. We also assume that
πref (a) > 0 for all a ∈ M . Otherwise if πref (a) = 0 for some a ∈ M , then both the regularized and constrained
versions forbid any policy to put nonzero probability on a, which leads to an effectively smaller candidate set.

[Next verbatim source excerpt]

Let A = {1, . . . , m} be a finite set of alternatives. The population of users is described by a
probability distribution D over utility vectors u = (u(1), . . . , u(m)), whose entries 0 ≤ u(x) ≤ 1
indicate a user’s utility for alternative x.3 The objective is to find an alternative x such that its average
utility AvgUtil(x) := Eu∼D [u(x)] across the user population (also known as the utilitarian social
welfare) is as high as possible. We extend this notation to probability distributions π over alternatives
by setting AvgUtil(π) := Ex∼π [AvgUtil(x)].
Both voting rules and alignment methods observe comparisons from n users. We model each
i = 1, . . . , n as a fresh user with independently drawn utility vector ui ∼ D. For exposition, we
assume that each user i provides an equal number d ≥ 1 of pairwise comparisons. For each i,
and for j = 1, . . . , d, we independently draw alternatives xji , yij from a fixed distribution µ over

    3
     Our assumption that utilities be in [0, 1] is weaker than any of the three assumptions — unit-sum, unit-range,
or approval [23] — made in classic distortion to avoid trivial infinite lower bounds.


                                                           3

[form-feed in source extraction]
alternatives, in which the minimum probability mass µmin := minx∈A µ(x) is positive.4 User i
then compares each pair {xji , yij } (for j = 1, . . . , d) through a Bradley-Terry model based on i’s
utilities: i prefers xji over yij (written “xji [U+001F control character]i yij ”) with probability σ β · (ui (xji ) − ui (yij )) , where
                                                                                                        [U+0001 control character]

σ(t) := 1/(1+e−t ) is the logistic sigmoid function and β > 0 is a temperature parameter, and prefers
yij over xji (“yij [U+001F control character]i xji ”) otherwise.5 Whereas this specifies the marginal probability of each pairwise
comparison, we make no assumption about the correlation between i’s choices. For example, i might
derive the pairwise comparisons from a Plackett-Luce ranking, ensuring that      [U+0002 control character] the user’s comparisons  [U+0001 control character][U+0003 control character]
are always consistent.6 We set p(x [U+001F control character] y) for the expected win rate Eu∼D σ β · (u(x) − u(y)) .

[Next verbatim source excerpt]

Alignment Setting. The alignment setting generalizes the social choice setting in two ways: first,
user utilities ui (y | x) may depend on a state x; second, the goal in determining a policy π is not
purely to maximize the reward AvgUtil(π | x), but a trade-off between this reward and the goal of
remaining close to a reference policy πref (· | x) ∈ ∆(A) in terms of the KL divergence DKL (· k πref ).
For theoretical tractability, we focus our analysis on a single state x, which we from here on omit
from the notation. Conceptually, this treatment of alignment on a state-by-state basis corresponds to
an assumption that our policy class is expressive enough so that it can take the optimal distribution of
actions at each state7 and abstracts from the generalization problem of estimating the population’s
preference between a pair of alternatives at the given state x based on preferences in similar states x0 .
Having set aside the dependency between states, we focus on how the regularization with respect
to a reference policy impacts the ability to optimize the average utility of three alignment methods:
RLHF, DPO, and NLHF. In addition to the pairwise comparisons, these methods take in a reference
policy πref ∈ ∆(A) and a KL bound τ ≥ 0, and map these inputs to a policy π in the KL-ball
Bτ (πref ) := {π ∈ ∆(A) | DKL (π k πref ) ≤ τ } around the reference policy. RLHF, DPO, and NLHF
are typically implemented with a KL regularization rather than our KL-constrained formulation,
which we adopt to enable a comparison on equal terms. In Section 5, we show that these perspectives
are equivalent, and that distortion upper bounds carry over to regularized alignment methods.

   4
     In Section 5, we discuss how most of our results extend more general distributions over comparison pairs,
which can, for example. capture k-wise comparisons.
   5
     Should we sample the same alternative x = xji = yij twice for a pair, the user is not asked for a pairwise
comparison. We record this as “x [U+001F control character]i x”, in a slight abuse of notation.
   6
     In particular, if we sample the same unordered pair twice for a user, the answers can be perfectly correlated.
   7
     This assumption is common in the literature; see, for example, Rafailov et al. [43]’s application of first-order
optimality conditions of PPO loss minimization.


                                                           4

[form-feed in source extraction]
The RLHF method first estimates rewards for each alternative, using maximum likelihood estimation
assuming that comparisons were generated by a single Bradley–Terry model:
                  r := argmaxr∈Rm 1≤i≤n,1≤j≤d log σ(r(xji ) − r(yij )) .8
                                     P                                     [U+0001 control character]

Next, RLHF uses PPO [44] to compute the policy πRLHF := argmaxπ∈Bτ (πref ) Ex∼π [r(x)] with
maximum expected reward within the KL-ball. In our setup, DPO is equivalent to RLHF (see
Appendix F.3) and thus has the same distortion.
The alignment method NLHF was inspired in part by a desire to better align with the
preferences of a heterogeneous group [36]. NLHF naturally adapts the definition of maxi-
mal lotteries by constraining both players’ mixed strategies to the KL-ball, i.e., πNLHF :=
argmaxπ1 ∈Bτ (πref ) minπ2 ∈Bτ (πref ) Ex1 ∼π1 ,x2 ∼π2 [Mx1 ,x2 ].
To generalize the definition of distortion to alignment methods, we set the maximum average utility of
any policy in the KL-ball as the benchmark. For fixed m, D, β, d, µ and correlation between pairwise
comparisons, the distortion of alignment method f is
                                                             maxπ∈Bτ (πref ) AvgUtil(π)
                      dist(f ) = supD,πref ,τ lim supn→∞       AvgUtiln (f (·,πref ,τ )) .
\end{ReviewVerbatim}



\paragraph{Expanded PaperInterface specification}
\begin{ReviewVerbatim}
(∀ {User Alternative : Type} [inst : Fintype User] [inst_1 : DecidableEq User] [inst_2 : Fintype Alternative]
    [inst_3 : DecidableEq Alternative] (population : PMF User)
    (utility : AppliedModelingLib.Alignment.Welfare.FiniteUtilityProfile User Alternative),
    AppliedModelingLib.Alignment.Welfare.UnitIntervalUtilityProfile utility →
      ∀ {btScale : ℝ},
        0 < btScale →
          ∀ (reference : PMF Alternative) (klBudget : ℝ) (nlhfPolicy benchmark : PMF Alternative),
            0 ≤ klBudget →
              AppliedModelingLib.PMFFullSupport reference →
                AppliedModelingLib.Alignment.Welfare.IsConstrainedPopulationBradleyTerryMaximin population utility
                    btScale reference klBudget nlhfPolicy →
                  AppliedModelingLib.Alignment.Welfare.InFiniteKLBall reference klBudget benchmark →
                    AppliedModelingLib.Alignment.Welfare.sigmoidChordSlope btScale / (1 / 4) *
                        AppliedModelingLib.Alignment.Welfare.policyAverageUtility population utility benchmark ≤
                      AppliedModelingLib.Alignment.Welfare.policyAverageUtility population utility nlhfPolicy) ∧
  ∀ {User Alternative : Type} [inst : Fintype User] [inst_1 : DecidableEq User] [inst_2 : Fintype Alternative]
    [inst_3 : DecidableEq Alternative] [Nonempty Alternative] {users comparisonsPerUser : ℕ} (population : PMF User)
    (utility : AppliedModelingLib.Alignment.Welfare.FiniteUtilityProfile User Alternative),
    AppliedModelingLib.Alignment.Welfare.UnitIntervalUtilityProfile utility →
      ∀ {btScale : ℝ},
        0 < btScale →
          ∀ (responseLaw : PMF (GolzHaghtalabYang2025Distortion.theorem2UserResponseTable Alternative))
            (preference : AppliedModelingLib.Learning.HumanFeedback.PairwisePreference PUnit.{1} Alternative),
            GolzHaghtalabYang2025Distortion.theorem2UserResponseCalibrated responseLaw preference →
              ∀ (sampling : PMF Alternative),
                0 < users →
                  0 < comparisonsPerUser →
                    ∀ (minimumMass : ℝ),
                      0 < minimumMass →
                        (∀ (alternative : Alternative), minimumMass ≤ (sampling alternative).toReal) →
                          ∀ (reference : PMF Alternative) (klBudget : ℝ) (hbudget : 0 ≤ klBudget),
                            AppliedModelingLib.PMFFullSupport reference →
                              ∀ (delta : ℝ),
                                0 < delta →
                                  delta ≤ 1 →
                                    ↑(Fintype.card Alternative) ^ 2 *
                                          Real.exp (-↑(users * comparisonsPerUser) * minimumMass ^ 2 / 8) ≤
                                        delta →
                                      preference =
                                          AppliedModelingLib.Alignment.Welfare.populationBradleyTerryPreference
                                            population utility btScale →
                                        have failureBudget : ℝ := delta / ↑(Fintype.card Alternative) ^ 2;
                                        ∃ (benchmark : PMF Alternative),
                                          AppliedModelingLib.finiteKLDivergence benchmark reference ≤ klBudget ∧
                                            (∀ (other : PMF Alternative),
                                                AppliedModelingLib.finiteKLDivergence other reference ≤ klBudget →
                                                  AppliedModelingLib.Alignment.Welfare.policyAverageUtility population
                                                      utility other ≤
                                                    AppliedModelingLib.Alignment.Welfare.policyAverageUtility population
                                                      utility benchmark) ∧
                                              (1 - delta) *
                                                  (AppliedModelingLib.Alignment.Welfare.sigmoidChordSlope btScale /
                                                        (1 / 4) *
                                                      AppliedModelingLib.Alignment.Welfare.policyAverageUtility
                                                        population utility benchmark -
                                                    GolzHaghtalabYang2025Distortion.theorem10SourceRate users
                                                        comparisonsPerUser minimumMass failureBudget /
                                                      (btScale * (1 / 4))) ≤
                                                AppliedModelingLib.pmfExp
                                                  (GolzHaghtalabYang2025Distortion.theorem7IidUserBatchLaw responseLaw
                                                    sampling users comparisonsPerUser)
                                                  fun
                                                    (sample :
                                                      Fin users →
                                                        GolzHaghtalabYang2025Distortion.theorem2UserResponseTable
                                                            Alternative ×
                                                          (Fin comparisonsPerUser → Alternative × Alternative)) =>
                                                  AppliedModelingLib.Alignment.Welfare.policyAverageUtility population
                                                    utility
                                                    (GolzHaghtalabYang2025Distortion.theorem7IidUserEmpiricalPolicy
                                                      reference klBudget hbudget sample)
\end{ReviewVerbatim}

\paragraph{Retained governing declarations}
\begin{itemize}\raggedright
\item \hyperlink{library-prerequisite-56c0734f8877ddbe}{Applied\allowbreak{}Modeling\allowbreak{}Lib.\allowbreak{}Alignment.\allowbreak{}Welfare.\allowbreak{}Finite\allowbreak{}Utility\allowbreak{}Profile}
\item \hyperlink{library-prerequisite-1d625de57bf58a6a}{Applied\allowbreak{}Modeling\allowbreak{}Lib.\allowbreak{}Alignment.\allowbreak{}Welfare.\allowbreak{}In\allowbreak{}Finite\allowbreak{}KL\allowbreak{}Ball}
\item \hyperlink{governing-declaration-2b862fbe69d5aa04}{Applied\allowbreak{}Modeling\allowbreak{}Lib.\allowbreak{}Alignment.\allowbreak{}Welfare.\allowbreak{}Is\allowbreak{}Constrained\allowbreak{}Population\allowbreak{}Bradley\allowbreak{}Terry\allowbreak{}Maximin}
\item \hyperlink{library-prerequisite-cbe8afa0d6d1ba78}{Applied\allowbreak{}Modeling\allowbreak{}Lib.\allowbreak{}Alignment.\allowbreak{}Welfare.\allowbreak{}Unit\allowbreak{}Interval\allowbreak{}Utility\allowbreak{}Profile}
\item \hyperlink{library-prerequisite-184c74a4db4311c2}{Applied\allowbreak{}Modeling\allowbreak{}Lib.\allowbreak{}Alignment.\allowbreak{}Welfare.\allowbreak{}policy\allowbreak{}Average\allowbreak{}Utility}
\item \hyperlink{library-prerequisite-cf06d15d7f8b7c31}{Applied\allowbreak{}Modeling\allowbreak{}Lib.\allowbreak{}Alignment.\allowbreak{}Welfare.\allowbreak{}population\allowbreak{}Bradley\allowbreak{}Terry\allowbreak{}Preference}
\item \hyperlink{governing-declaration-a4d623b8c608b566}{Applied\allowbreak{}Modeling\allowbreak{}Lib.\allowbreak{}Alignment.\allowbreak{}Welfare.\allowbreak{}sigmoid\allowbreak{}Chord\allowbreak{}Slope}
\item \hyperlink{library-prerequisite-7e2db20cf087d9ba}{Applied\allowbreak{}Modeling\allowbreak{}Lib.\allowbreak{}Learning.\allowbreak{}Human\allowbreak{}Feedback.\allowbreak{}Pairwise\allowbreak{}Preference}
\item \hyperlink{library-prerequisite-268c959ee7d2d148}{Applied\allowbreak{}Modeling\allowbreak{}Lib.\allowbreak{}PMF\allowbreak{}Full\allowbreak{}Support}
\item \hyperlink{library-prerequisite-d3573881ae45d4c6}{Applied\allowbreak{}Modeling\allowbreak{}Lib.\allowbreak{}finite\allowbreak{}KL\allowbreak{}Divergence}
\item \hyperlink{governing-declaration-8e4b4389a4f6454b}{Applied\allowbreak{}Modeling\allowbreak{}Lib.\allowbreak{}pmf\allowbreak{}Exp}
\item \hyperlink{governing-declaration-aaf6dc2063f9888f}{Golz\allowbreak{}Haghtalab\allowbreak{}Yang2025\allowbreak{}Distortion.\allowbreak{}theorem10\allowbreak{}Source\allowbreak{}Rate}
\item \hyperlink{governing-declaration-080b9388e51a6018}{Golz\allowbreak{}Haghtalab\allowbreak{}Yang2025\allowbreak{}Distortion.\allowbreak{}theorem2\allowbreak{}User\allowbreak{}Response\allowbreak{}Calibrated}
\item \hyperlink{governing-declaration-c5d0273cbd7ddcc9}{Golz\allowbreak{}Haghtalab\allowbreak{}Yang2025\allowbreak{}Distortion.\allowbreak{}theorem2\allowbreak{}User\allowbreak{}Response\allowbreak{}Table}
\item \hyperlink{governing-declaration-485a7e8af83d2b11}{Golz\allowbreak{}Haghtalab\allowbreak{}Yang2025\allowbreak{}Distortion.\allowbreak{}theorem7\allowbreak{}Iid\allowbreak{}User\allowbreak{}Batch\allowbreak{}Law}
\item \hyperlink{governing-declaration-71b9e9aed968af43}{Golz\allowbreak{}Haghtalab\allowbreak{}Yang2025\allowbreak{}Distortion.\allowbreak{}theorem7\allowbreak{}Iid\allowbreak{}User\allowbreak{}Empirical\allowbreak{}Policy}
\end{itemize}
\paragraph{Lean-elaborated claim roles}\begin{ReviewVerbatim}
1. conclusion: (∀ {User Alternative : Type} [inst : Fintype User] [inst_1 : DecidableEq User] [inst_2 : Fintype Alternative]
    [inst_3 : DecidableEq Alternative] (population : PMF User)
    (utility : AppliedModelingLib.Alignment.Welfare.FiniteUtilityProfile User Alternative),
    AppliedModelingLib.Alignment.Welfare.UnitIntervalUtilityProfile utility →
      ∀ {btScale : ℝ},
        0 < btScale →
          ∀ (reference : PMF Alternative) (klBudget : ℝ) (nlhfPolicy benchmark : PMF Alternative),
            0 ≤ klBudget →
              AppliedModelingLib.PMFFullSupport reference →
                AppliedModelingLib.Alignment.Welfare.IsConstrainedPopulationBradleyTerryMaximin population utility
                    btScale reference klBudget nlhfPolicy →
                  AppliedModelingLib.Alignment.Welfare.InFiniteKLBall reference klBudget benchmark →
                    AppliedModelingLib.Alignment.Welfare.sigmoidChordSlope btScale / (1 / 4) *
                        AppliedModelingLib.Alignment.Welfare.policyAverageUtility population utility benchmark ≤
                      AppliedModelingLib.Alignment.Welfare.policyAverageUtility population utility nlhfPolicy) ∧
  ∀ {User Alternative : Type} [inst : Fintype User] [inst_1 : DecidableEq User] [inst_2 : Fintype Alternative]
    [inst_3 : DecidableEq Alternative] [Nonempty Alternative] {users comparisonsPerUser : ℕ} (population : PMF User)
    (utility : AppliedModelingLib.Alignment.Welfare.FiniteUtilityProfile User Alternative),
    AppliedModelingLib.Alignment.Welfare.UnitIntervalUtilityProfile utility →
      ∀ {btScale : ℝ},
        0 < btScale →
          ∀ (responseLaw : PMF (GolzHaghtalabYang2025Distortion.theorem2UserResponseTable Alternative))
            (preference : AppliedModelingLib.Learning.HumanFeedback.PairwisePreference PUnit.{1} Alternative),
            GolzHaghtalabYang2025Distortion.theorem2UserResponseCalibrated responseLaw preference →
              ∀ (sampling : PMF Alternative),
                0 < users →
                  0 < comparisonsPerUser →
                    ∀ (minimumMass : ℝ),
                      0 < minimumMass →
                        (∀ (alternative : Alternative), minimumMass ≤ (sampling alternative).toReal) →
                          ∀ (reference : PMF Alternative) (klBudget : ℝ) (hbudget : 0 ≤ klBudget),
                            AppliedModelingLib.PMFFullSupport reference →
                              ∀ (delta : ℝ),
                                0 < delta →
                                  delta ≤ 1 →
                                    ↑(Fintype.card Alternative) ^ 2 *
                                          Real.exp (-↑(users * comparisonsPerUser) * minimumMass ^ 2 / 8) ≤
                                        delta →
                                      preference =
                                          AppliedModelingLib.Alignment.Welfare.populationBradleyTerryPreference
                                            population utility btScale →
                                        have failureBudget : ℝ := delta / ↑(Fintype.card Alternative) ^ 2;
                                        ∃ (benchmark : PMF Alternative),
                                          AppliedModelingLib.finiteKLDivergence benchmark reference ≤ klBudget ∧
                                            (∀ (other : PMF Alternative),
                                                AppliedModelingLib.finiteKLDivergence other reference ≤ klBudget →
                                                  AppliedModelingLib.Alignment.Welfare.policyAverageUtility population
                                                      utility other ≤
                                                    AppliedModelingLib.Alignment.Welfare.policyAverageUtility population
                                                      utility benchmark) ∧
                                              (1 - delta) *
                                                  (AppliedModelingLib.Alignment.Welfare.sigmoidChordSlope btScale /
                                                        (1 / 4) *
                                                      AppliedModelingLib.Alignment.Welfare.policyAverageUtility
                                                        population utility benchmark -
                                                    GolzHaghtalabYang2025Distortion.theorem10SourceRate users
                                                        comparisonsPerUser minimumMass failureBudget /
                                                      (btScale * (1 / 4))) ≤
                                                AppliedModelingLib.pmfExp
                                                  (GolzHaghtalabYang2025Distortion.theorem7IidUserBatchLaw responseLaw
                                                    sampling users comparisonsPerUser)
                                                  fun
                                                    (sample :
                                                      Fin users →
                                                        GolzHaghtalabYang2025Distortion.theorem2UserResponseTable
                                                            Alternative ×
                                                          (Fin comparisonsPerUser → Alternative × Alternative)) =>
                                                  AppliedModelingLib.Alignment.Welfare.policyAverageUtility population
                                                    utility
                                                    (GolzHaghtalabYang2025Distortion.theorem7IidUserEmpiricalPolicy
                                                      reference klBudget hbudget sample)
\end{ReviewVerbatim}

\paragraph{Lean proof endpoint (built separately)}\begin{ReviewVerbatim}
GolzHaghtalabYang2025Distortion.theorem7_nlhf_distortion
\end{ReviewVerbatim}

\paragraph{Recorded source-to-Spec screening}
\textbf{Verdict:} \texttt{matches}
\\
{\raggedright
\textbf{Reason:} The constrained-\allowbreak{}NLHF population factor and finite-\allowbreak{}sample expected-\allowbreak{}welfare guarantee over the KL ball are retained.\allowbreak{}
\par}
\reviewmatch{Matches source input}
\noindent\textbf{Reviewer annotation}\par
\reviewerbox
\clearpage
\hypertarget{source-claim-067d55f193fc0110}{}
\section*{Review row 11}
\textbf{Source locator:} {\footnotesize\raggedright cited publication, lines 444-\allowbreak{}449\par}
\paragraph{Verbatim source input}
\begin{ReviewVerbatim}
Corollary 4 (of Theorem 7). The Maximal Lotteries voting rule has a distortion of β2 · 1−e
                                                                                         1+e
                                                                                             −β .

The Borda rule, in contrast, does not match this optimal distortion bound, as shown by the following
bound that we prove and state in full detail in Appendix E.2.

[Next verbatim source excerpt]

Let A = {1, . . . , m} be a finite set of alternatives. The population of users is described by a
probability distribution D over utility vectors u = (u(1), . . . , u(m)), whose entries 0 ≤ u(x) ≤ 1
indicate a user’s utility for alternative x.3 The objective is to find an alternative x such that its average
utility AvgUtil(x) := Eu∼D [u(x)] across the user population (also known as the utilitarian social
welfare) is as high as possible. We extend this notation to probability distributions π over alternatives
by setting AvgUtil(π) := Ex∼π [AvgUtil(x)].
Both voting rules and alignment methods observe comparisons from n users. We model each
i = 1, . . . , n as a fresh user with independently drawn utility vector ui ∼ D. For exposition, we
assume that each user i provides an equal number d ≥ 1 of pairwise comparisons. For each i,
and for j = 1, . . . , d, we independently draw alternatives xji , yij from a fixed distribution µ over

    3
     Our assumption that utilities be in [0, 1] is weaker than any of the three assumptions — unit-sum, unit-range,
or approval [23] — made in classic distortion to avoid trivial infinite lower bounds.


                                                           3

[form-feed in source extraction]
alternatives, in which the minimum probability mass µmin := minx∈A µ(x) is positive.4 User i
then compares each pair {xji , yij } (for j = 1, . . . , d) through a Bradley-Terry model based on i’s
utilities: i prefers xji over yij (written “xji [U+001F control character]i yij ”) with probability σ β · (ui (xji ) − ui (yij )) , where
                                                                                                        [U+0001 control character]

σ(t) := 1/(1+e−t ) is the logistic sigmoid function and β > 0 is a temperature parameter, and prefers
yij over xji (“yij [U+001F control character]i xji ”) otherwise.5 Whereas this specifies the marginal probability of each pairwise
comparison, we make no assumption about the correlation between i’s choices. For example, i might
derive the pairwise comparisons from a Plackett-Luce ranking, ensuring that      [U+0002 control character] the user’s comparisons  [U+0001 control character][U+0003 control character]
are always consistent.6 We set p(x [U+001F control character] y) for the expected win rate Eu∼D σ β · (u(x) − u(y)) .

[Next verbatim source excerpt]

Social Choice Setting. A voting rule f observes the sampled pairwise comparisons {xji [U+001F control character]i
yij }i∈[n],j∈[d] and maps them to a probability distribution over alternatives. For some m, D, β, d, µ,
and
  [U+0002 control character] the correlation between comparisons,     the average utility of f for n samples is AvgUtiln (f ) :=
E AvgUtil f ({xji [U+001F control character]i yij }i,j ) , where the expectation is taken over the pairwise comparisons. The
                               [U+0001 control character][U+0003 control character]

distortion of f on D is the competitive ratio between AvgUtiln (f ) and the optimal average utility
maxx∈A AvgUtil(x) in the limit of n → ∞ samples, and the distortion of f the worst-case distortion
over all D:

            dist(f, D) := lim supn→∞ maxAvgUtil
                                        x∈A AvgUtil(x)
                                                (f )   ,               dist(f ) = supD dist(f, D).
                                                       n



For alternatives x, y, let #(x [U+001F control character] y) := {(i, j) | xji = x, yij = y} denote the number of pairwise
comparisons in which x beat y. The (normalized) Borda score [45] of alternative x is
                                                       P
                                                         y∈A #(x[U+001F control character]y)
                                   BC(x) := P         #(x[U+001F control character]y)+
                                                              P             ,
                                                  y∈A           y6=x #(y[U+001F control character]x)


i.e., the fraction of pairwise comparisons involving x in which it wins. The Borda voting rule chooses
the winner uniformly among all alternatives with maximum Borda score. The Maximal Lotteries vot-
ing rule first computes the margin matrix M ∈ Rm×m , where Mx,y = #(x[U+001F control character]y)−#(y[U+001F control character]x)
                                                                          #(x[U+001F control character]y)+#(y[U+001F control character]x) . It then con-
siders a symmetric two-player zero-sum game in which player 1 selects alternative x1 , player 2 selects
alternative x2 , and the payoffs are Mx1 ,x2 for player 1 and Mx2 ,x1 = −Mx1 ,x2 for player 2. The max-
imal lotteries rule returns a distribution π ∈ argmaxπ1 ∈∆(A) minπ2 ∈∆(A) Ex1 ∼π1 ,x2 ∼π2 [Mx1 ,x2 ],
i.e., a mixed strategy in Nash equilibrium. When several such π exist, our results hold for any choice.
\end{ReviewVerbatim}



\paragraph{Expanded PaperInterface specification}
\begin{ReviewVerbatim}
(∀ {User Alternative : Type} [inst : Fintype User] [inst_1 : DecidableEq User] [inst_2 : Fintype Alternative]
    [inst_3 : DecidableEq Alternative] (population : PMF User)
    (utility : AppliedModelingLib.Alignment.Welfare.FiniteUtilityProfile User Alternative),
    AppliedModelingLib.Alignment.Welfare.UnitIntervalUtilityProfile utility →
      ∀ {btScale : ℝ},
        0 < btScale →
          ∀ (maximalLottery benchmark : PMF Alternative),
            (∀ (opponent : PMF Alternative),
                0 ≤
                  AppliedModelingLib.Alignment.Welfare.populationBradleyTerryPolicyMargin population utility btScale
                    maximalLottery opponent) →
              AppliedModelingLib.Alignment.Welfare.sigmoidChordSlope btScale / (1 / 4) *
                  AppliedModelingLib.Alignment.Welfare.policyAverageUtility population utility benchmark ≤
                AppliedModelingLib.Alignment.Welfare.policyAverageUtility population utility maximalLottery) ∧
  ∀ {User Alternative : Type} [inst : Fintype User] [inst_1 : DecidableEq User] [inst_2 : Fintype Alternative]
    [inst_3 : DecidableEq Alternative] [inst_4 : Nonempty Alternative] {users comparisonsPerUser : ℕ}
    (population : PMF User) (utility : AppliedModelingLib.Alignment.Welfare.FiniteUtilityProfile User Alternative),
    AppliedModelingLib.Alignment.Welfare.UnitIntervalUtilityProfile utility →
      ∀ {btScale : ℝ},
        0 < btScale →
          ∀ (responseLaw : PMF (GolzHaghtalabYang2025Distortion.theorem2UserResponseTable Alternative))
            (preference : AppliedModelingLib.Learning.HumanFeedback.PairwisePreference PUnit.{1} Alternative),
            GolzHaghtalabYang2025Distortion.theorem2UserResponseCalibrated responseLaw preference →
              ∀ (sampling : PMF Alternative),
                0 < users →
                  0 < comparisonsPerUser →
                    ∀ (minimumMass : ℝ),
                      0 < minimumMass →
                        (∀ (alternative : Alternative), minimumMass ≤ (sampling alternative).toReal) →
                          ∀ (delta : ℝ),
                            0 < delta →
                              delta ≤ 1 →
                                ↑(Fintype.card Alternative) ^ 2 *
                                      Real.exp (-↑(users * comparisonsPerUser) * minimumMass ^ 2 / 8) ≤
                                    delta →
                                  preference =
                                      AppliedModelingLib.Alignment.Welfare.populationBradleyTerryPreference population
                                        utility btScale →
                                    (∀
                                        (sample :
                                          Fin users →
                                            GolzHaghtalabYang2025Distortion.theorem2UserResponseTable Alternative ×
                                              (Fin comparisonsPerUser → Alternative × Alternative)),
                                        GolzHaghtalabYang2025Distortion.IsEmpiricalMaximalLottery
                                          (fun (first second : Alternative) =>
                                            GolzHaghtalabYang2025Distortion.theorem10IidUserPaperWinRate first second
                                              sample)
                                          (GolzHaghtalabYang2025Distortion.corollary4IidUserEmpiricalMaximalLottery
                                            sample)) ∧
                                      have failureBudget : ℝ := delta / ↑(Fintype.card Alternative) ^ 2;
                                      ∃ (benchmark : PMF Alternative),
                                        (∀ (other : PMF Alternative),
                                            AppliedModelingLib.Alignment.Welfare.policyAverageUtility population utility
                                                other ≤
                                              AppliedModelingLib.Alignment.Welfare.policyAverageUtility population
                                                utility benchmark) ∧
                                          (1 - delta) *
                                              (AppliedModelingLib.Alignment.Welfare.sigmoidChordSlope btScale /
                                                    (1 / 4) *
                                                  AppliedModelingLib.Alignment.Welfare.policyAverageUtility population
                                                    utility benchmark -
                                                GolzHaghtalabYang2025Distortion.theorem10SourceRate users
                                                    comparisonsPerUser minimumMass failureBudget /
                                                  (btScale * (1 / 4))) ≤
                                            AppliedModelingLib.pmfExp
                                              (GolzHaghtalabYang2025Distortion.theorem7IidUserBatchLaw responseLaw
                                                sampling users comparisonsPerUser)
                                              fun
                                                (sample :
                                                  Fin users →
                                                    GolzHaghtalabYang2025Distortion.theorem2UserResponseTable
                                                        Alternative ×
                                                      (Fin comparisonsPerUser → Alternative × Alternative)) =>
                                              AppliedModelingLib.Alignment.Welfare.policyAverageUtility population
                                                utility
                                                (GolzHaghtalabYang2025Distortion.corollary4IidUserEmpiricalMaximalLottery
                                                  sample)
\end{ReviewVerbatim}

\paragraph{Retained governing declarations}
\begin{itemize}\raggedright
\item \hyperlink{library-prerequisite-56c0734f8877ddbe}{Applied\allowbreak{}Modeling\allowbreak{}Lib.\allowbreak{}Alignment.\allowbreak{}Welfare.\allowbreak{}Finite\allowbreak{}Utility\allowbreak{}Profile}
\item \hyperlink{library-prerequisite-cbe8afa0d6d1ba78}{Applied\allowbreak{}Modeling\allowbreak{}Lib.\allowbreak{}Alignment.\allowbreak{}Welfare.\allowbreak{}Unit\allowbreak{}Interval\allowbreak{}Utility\allowbreak{}Profile}
\item \hyperlink{library-prerequisite-184c74a4db4311c2}{Applied\allowbreak{}Modeling\allowbreak{}Lib.\allowbreak{}Alignment.\allowbreak{}Welfare.\allowbreak{}policy\allowbreak{}Average\allowbreak{}Utility}
\item \hyperlink{governing-declaration-5c3611f5a76c4e77}{Applied\allowbreak{}Modeling\allowbreak{}Lib.\allowbreak{}Alignment.\allowbreak{}Welfare.\allowbreak{}population\allowbreak{}Bradley\allowbreak{}Terry\allowbreak{}Policy\allowbreak{}Margin}
\item \hyperlink{library-prerequisite-cf06d15d7f8b7c31}{Applied\allowbreak{}Modeling\allowbreak{}Lib.\allowbreak{}Alignment.\allowbreak{}Welfare.\allowbreak{}population\allowbreak{}Bradley\allowbreak{}Terry\allowbreak{}Preference}
\item \hyperlink{governing-declaration-a4d623b8c608b566}{Applied\allowbreak{}Modeling\allowbreak{}Lib.\allowbreak{}Alignment.\allowbreak{}Welfare.\allowbreak{}sigmoid\allowbreak{}Chord\allowbreak{}Slope}
\item \hyperlink{library-prerequisite-7e2db20cf087d9ba}{Applied\allowbreak{}Modeling\allowbreak{}Lib.\allowbreak{}Learning.\allowbreak{}Human\allowbreak{}Feedback.\allowbreak{}Pairwise\allowbreak{}Preference}
\item \hyperlink{governing-declaration-8e4b4389a4f6454b}{Applied\allowbreak{}Modeling\allowbreak{}Lib.\allowbreak{}pmf\allowbreak{}Exp}
\item \hyperlink{paper-prerequisite-53af1666fc341b0f}{Golz\allowbreak{}Haghtalab\allowbreak{}Yang2025\allowbreak{}Distortion.\allowbreak{}Is\allowbreak{}Empirical\allowbreak{}Maximal\allowbreak{}Lottery}
\item \hyperlink{paper-prerequisite-54cac251e94c8e83}{Golz\allowbreak{}Haghtalab\allowbreak{}Yang2025\allowbreak{}Distortion.\allowbreak{}corollary4\allowbreak{}Iid\allowbreak{}User\allowbreak{}Empirical\allowbreak{}Maximal\allowbreak{}Lottery}
\item \hyperlink{governing-declaration-a1441da0ef6b2fa2}{Golz\allowbreak{}Haghtalab\allowbreak{}Yang2025\allowbreak{}Distortion.\allowbreak{}theorem10\allowbreak{}Iid\allowbreak{}User\allowbreak{}Paper\allowbreak{}Win\allowbreak{}Rate}
\item \hyperlink{governing-declaration-aaf6dc2063f9888f}{Golz\allowbreak{}Haghtalab\allowbreak{}Yang2025\allowbreak{}Distortion.\allowbreak{}theorem10\allowbreak{}Source\allowbreak{}Rate}
\item \hyperlink{governing-declaration-080b9388e51a6018}{Golz\allowbreak{}Haghtalab\allowbreak{}Yang2025\allowbreak{}Distortion.\allowbreak{}theorem2\allowbreak{}User\allowbreak{}Response\allowbreak{}Calibrated}
\item \hyperlink{governing-declaration-c5d0273cbd7ddcc9}{Golz\allowbreak{}Haghtalab\allowbreak{}Yang2025\allowbreak{}Distortion.\allowbreak{}theorem2\allowbreak{}User\allowbreak{}Response\allowbreak{}Table}
\item \hyperlink{governing-declaration-485a7e8af83d2b11}{Golz\allowbreak{}Haghtalab\allowbreak{}Yang2025\allowbreak{}Distortion.\allowbreak{}theorem7\allowbreak{}Iid\allowbreak{}User\allowbreak{}Batch\allowbreak{}Law}
\end{itemize}
\paragraph{Lean-elaborated claim roles}\begin{ReviewVerbatim}
1. conclusion: (∀ {User Alternative : Type} [inst : Fintype User] [inst_1 : DecidableEq User] [inst_2 : Fintype Alternative]
    [inst_3 : DecidableEq Alternative] (population : PMF User)
    (utility : AppliedModelingLib.Alignment.Welfare.FiniteUtilityProfile User Alternative),
    AppliedModelingLib.Alignment.Welfare.UnitIntervalUtilityProfile utility →
      ∀ {btScale : ℝ},
        0 < btScale →
          ∀ (maximalLottery benchmark : PMF Alternative),
            (∀ (opponent : PMF Alternative),
                0 ≤
                  AppliedModelingLib.Alignment.Welfare.populationBradleyTerryPolicyMargin population utility btScale
                    maximalLottery opponent) →
              AppliedModelingLib.Alignment.Welfare.sigmoidChordSlope btScale / (1 / 4) *
                  AppliedModelingLib.Alignment.Welfare.policyAverageUtility population utility benchmark ≤
                AppliedModelingLib.Alignment.Welfare.policyAverageUtility population utility maximalLottery) ∧
  ∀ {User Alternative : Type} [inst : Fintype User] [inst_1 : DecidableEq User] [inst_2 : Fintype Alternative]
    [inst_3 : DecidableEq Alternative] [inst_4 : Nonempty Alternative] {users comparisonsPerUser : ℕ}
    (population : PMF User) (utility : AppliedModelingLib.Alignment.Welfare.FiniteUtilityProfile User Alternative),
    AppliedModelingLib.Alignment.Welfare.UnitIntervalUtilityProfile utility →
      ∀ {btScale : ℝ},
        0 < btScale →
          ∀ (responseLaw : PMF (GolzHaghtalabYang2025Distortion.theorem2UserResponseTable Alternative))
            (preference : AppliedModelingLib.Learning.HumanFeedback.PairwisePreference PUnit.{1} Alternative),
            GolzHaghtalabYang2025Distortion.theorem2UserResponseCalibrated responseLaw preference →
              ∀ (sampling : PMF Alternative),
                0 < users →
                  0 < comparisonsPerUser →
                    ∀ (minimumMass : ℝ),
                      0 < minimumMass →
                        (∀ (alternative : Alternative), minimumMass ≤ (sampling alternative).toReal) →
                          ∀ (delta : ℝ),
                            0 < delta →
                              delta ≤ 1 →
                                ↑(Fintype.card Alternative) ^ 2 *
                                      Real.exp (-↑(users * comparisonsPerUser) * minimumMass ^ 2 / 8) ≤
                                    delta →
                                  preference =
                                      AppliedModelingLib.Alignment.Welfare.populationBradleyTerryPreference population
                                        utility btScale →
                                    (∀
                                        (sample :
                                          Fin users →
                                            GolzHaghtalabYang2025Distortion.theorem2UserResponseTable Alternative ×
                                              (Fin comparisonsPerUser → Alternative × Alternative)),
                                        GolzHaghtalabYang2025Distortion.IsEmpiricalMaximalLottery
                                          (fun (first second : Alternative) =>
                                            GolzHaghtalabYang2025Distortion.theorem10IidUserPaperWinRate first second
                                              sample)
                                          (GolzHaghtalabYang2025Distortion.corollary4IidUserEmpiricalMaximalLottery
                                            sample)) ∧
                                      have failureBudget : ℝ := delta / ↑(Fintype.card Alternative) ^ 2;
                                      ∃ (benchmark : PMF Alternative),
                                        (∀ (other : PMF Alternative),
                                            AppliedModelingLib.Alignment.Welfare.policyAverageUtility population utility
                                                other ≤
                                              AppliedModelingLib.Alignment.Welfare.policyAverageUtility population
                                                utility benchmark) ∧
                                          (1 - delta) *
                                              (AppliedModelingLib.Alignment.Welfare.sigmoidChordSlope btScale /
                                                    (1 / 4) *
                                                  AppliedModelingLib.Alignment.Welfare.policyAverageUtility population
                                                    utility benchmark -
                                                GolzHaghtalabYang2025Distortion.theorem10SourceRate users
                                                    comparisonsPerUser minimumMass failureBudget /
                                                  (btScale * (1 / 4))) ≤
                                            AppliedModelingLib.pmfExp
                                              (GolzHaghtalabYang2025Distortion.theorem7IidUserBatchLaw responseLaw
                                                sampling users comparisonsPerUser)
                                              fun
                                                (sample :
                                                  Fin users →
                                                    GolzHaghtalabYang2025Distortion.theorem2UserResponseTable
                                                        Alternative ×
                                                      (Fin comparisonsPerUser → Alternative × Alternative)) =>
                                              AppliedModelingLib.Alignment.Welfare.policyAverageUtility population
                                                utility
                                                (GolzHaghtalabYang2025Distortion.corollary4IidUserEmpiricalMaximalLottery
                                                  sample)
\end{ReviewVerbatim}

\paragraph{Lean proof endpoint (built separately)}\begin{ReviewVerbatim}
GolzHaghtalabYang2025Distortion.corollary4_maximal_lottery
\end{ReviewVerbatim}

\paragraph{Recorded source-to-Spec screening}
\textbf{Verdict:} \texttt{matches}
\\
{\raggedright
\textbf{Reason:} The Maximal-\allowbreak{}Lotteries factor and finite expected-\allowbreak{}welfare route are retained.\allowbreak{} Positive-\allowbreak{}count margins are the printed ratios;\allowbreak{} the explicitly disclosed neutral zero margin totalizes only unobserved off-\allowbreak{}diagonal pairs and adds no stochastic premise.\allowbreak{}
\par}
\reviewmatch{Matches source input}
\noindent\textbf{Reviewer annotation}\par
\reviewerbox
\clearpage
\hypertarget{source-claim-9824dd64b1ddd3ad}{}
\section*{Review row 12}
\textbf{Source locator:} {\footnotesize\raggedright cited publication, lines 2166-\allowbreak{}2178\par}
\paragraph{Verbatim source input}
\begin{ReviewVerbatim}
Proposition 13 (Equivalence between Constrained and Regularized NLHF.). Let πτ be the output of
the τ -constrained NLHF method defined in Section 2, i.e.,
                                πτ = argmax                 min               E          [Mx1 ,x2 ] ,                         (14)
                                         π1 ∈Bτ (πref ) π2 ∈Bτ (πref ) x1 ∼π1 ,x2 ∼π2

and let π̃λ be the output of the λ-regularized NLHF method [36], i.e.,
      π̃λ = argmax          min             E         [Mx1 ,x2 ] − λ · DKL (π1 k πref ) + λ · DKL (π2 k πref ).               (15)
            π1 ∈∆(M ) π2 ∈∆(M ) x1 ∼π1 ,x2 ∼π2

Then, for each λ ∈ [0, ∞] with solution π̃λ for Eq. (15), we have that π̃λ is also an optimal solution
to the τ -constrained optimization problem in Eq. (14), where τ = DKL (π̃λ k πref ). Conversely, for
each τ ≥ 0 with solution πτ for Eq. (14), there exists λ ∈ [0, ∞] such that πτ is also an optimal
solution to the λ-regularized optimization problem in Eq. (15).

[Next verbatim source excerpt]

Let A = {1, . . . , m} be a finite set of alternatives. The population of users is described by a
probability distribution D over utility vectors u = (u(1), . . . , u(m)), whose entries 0 ≤ u(x) ≤ 1
indicate a user’s utility for alternative x.3 The objective is to find an alternative x such that its average
utility AvgUtil(x) := Eu∼D [u(x)] across the user population (also known as the utilitarian social
welfare) is as high as possible. We extend this notation to probability distributions π over alternatives
by setting AvgUtil(π) := Ex∼π [AvgUtil(x)].
Both voting rules and alignment methods observe comparisons from n users. We model each
i = 1, . . . , n as a fresh user with independently drawn utility vector ui ∼ D. For exposition, we
assume that each user i provides an equal number d ≥ 1 of pairwise comparisons. For each i,
and for j = 1, . . . , d, we independently draw alternatives xji , yij from a fixed distribution µ over

    3
     Our assumption that utilities be in [0, 1] is weaker than any of the three assumptions — unit-sum, unit-range,
or approval [23] — made in classic distortion to avoid trivial infinite lower bounds.


                                                           3

[form-feed in source extraction]
alternatives, in which the minimum probability mass µmin := minx∈A µ(x) is positive.4 User i
then compares each pair {xji , yij } (for j = 1, . . . , d) through a Bradley-Terry model based on i’s
utilities: i prefers xji over yij (written “xji [U+001F control character]i yij ”) with probability σ β · (ui (xji ) − ui (yij )) , where
                                                                                                        [U+0001 control character]

σ(t) := 1/(1+e−t ) is the logistic sigmoid function and β > 0 is a temperature parameter, and prefers
yij over xji (“yij [U+001F control character]i xji ”) otherwise.5 Whereas this specifies the marginal probability of each pairwise
comparison, we make no assumption about the correlation between i’s choices. For example, i might
derive the pairwise comparisons from a Plackett-Luce ranking, ensuring that      [U+0002 control character] the user’s comparisons  [U+0001 control character][U+0003 control character]
are always consistent.6 We set p(x [U+001F control character] y) for the expected win rate Eu∼D σ β · (u(x) − u(y)) .

[Next verbatim source excerpt]

Alignment Setting. The alignment setting generalizes the social choice setting in two ways: first,
user utilities ui (y | x) may depend on a state x; second, the goal in determining a policy π is not
purely to maximize the reward AvgUtil(π | x), but a trade-off between this reward and the goal of
remaining close to a reference policy πref (· | x) ∈ ∆(A) in terms of the KL divergence DKL (· k πref ).
For theoretical tractability, we focus our analysis on a single state x, which we from here on omit
from the notation. Conceptually, this treatment of alignment on a state-by-state basis corresponds to
an assumption that our policy class is expressive enough so that it can take the optimal distribution of
actions at each state7 and abstracts from the generalization problem of estimating the population’s
preference between a pair of alternatives at the given state x based on preferences in similar states x0 .
Having set aside the dependency between states, we focus on how the regularization with respect
to a reference policy impacts the ability to optimize the average utility of three alignment methods:
RLHF, DPO, and NLHF. In addition to the pairwise comparisons, these methods take in a reference
policy πref ∈ ∆(A) and a KL bound τ ≥ 0, and map these inputs to a policy π in the KL-ball
Bτ (πref ) := {π ∈ ∆(A) | DKL (π k πref ) ≤ τ } around the reference policy. RLHF, DPO, and NLHF
are typically implemented with a KL regularization rather than our KL-constrained formulation,
which we adopt to enable a comparison on equal terms. In Section 5, we show that these perspectives
are equivalent, and that distortion upper bounds carry over to regularized alignment methods.

   4
     In Section 5, we discuss how most of our results extend more general distributions over comparison pairs,
which can, for example. capture k-wise comparisons.
   5
     Should we sample the same alternative x = xji = yij twice for a pair, the user is not asked for a pairwise
comparison. We record this as “x [U+001F control character]i x”, in a slight abuse of notation.
   6
     In particular, if we sample the same unordered pair twice for a user, the answers can be perfectly correlated.
   7
     This assumption is common in the literature; see, for example, Rafailov et al. [43]’s application of first-order
optimality conditions of PPO loss minimization.


                                                           4

[form-feed in source extraction]
The RLHF method first estimates rewards for each alternative, using maximum likelihood estimation
assuming that comparisons were generated by a single Bradley–Terry model:
                  r := argmaxr∈Rm 1≤i≤n,1≤j≤d log σ(r(xji ) − r(yij )) .8
                                     P                                     [U+0001 control character]

Next, RLHF uses PPO [44] to compute the policy πRLHF := argmaxπ∈Bτ (πref ) Ex∼π [r(x)] with
maximum expected reward within the KL-ball. In our setup, DPO is equivalent to RLHF (see
Appendix F.3) and thus has the same distortion.
The alignment method NLHF was inspired in part by a desire to better align with the
preferences of a heterogeneous group [36]. NLHF naturally adapts the definition of maxi-
mal lotteries by constraining both players’ mixed strategies to the KL-ball, i.e., πNLHF :=
argmaxπ1 ∈Bτ (πref ) minπ2 ∈Bτ (πref ) Ex1 ∼π1 ,x2 ∼π2 [Mx1 ,x2 ].
To generalize the definition of distortion to alignment methods, we set the maximum average utility of
any policy in the KL-ball as the benchmark. For fixed m, D, β, d, µ and correlation between pairwise
comparisons, the distortion of alignment method f is
                                                             maxπ∈Bτ (πref ) AvgUtil(π)
                      dist(f ) = supD,πref ,τ lim supn→∞       AvgUtiln (f (·,πref ,τ )) .
\end{ReviewVerbatim}



\paragraph{Expanded PaperInterface specification}
\begin{ReviewVerbatim}
(∀ {Alternative : Type} [inst : Fintype Alternative] [inst_1 : DecidableEq Alternative]
    (preference : AppliedModelingLib.Learning.HumanFeedback.PairwisePreference PUnit.{1} Alternative)
    (reference policy : PMF Alternative) (klWeight : WithTop ℝ),
    AppliedModelingLib.PMFFullSupport reference →
      AppliedModelingLib.GameTheory.PreferenceGame.IsContextFreeExtendedRegularizedPreferenceGameEquilibrium preference
          reference policy klWeight →
        AppliedModelingLib.GameTheory.PreferenceGame.IsFiniteKLPreferenceMaximin preference reference
          (AppliedModelingLib.finiteKLDivergence policy reference) policy) ∧
  ∀ {Alternative : Type} [inst : Fintype Alternative] [inst_1 : DecidableEq Alternative]
    (preference : AppliedModelingLib.Learning.HumanFeedback.PairwisePreference PUnit.{1} Alternative)
    (reference policy : PMF Alternative) (klBudget : ℝ),
    0 ≤ klBudget →
      AppliedModelingLib.PMFFullSupport reference →
        AppliedModelingLib.GameTheory.PreferenceGame.IsFiniteKLPreferenceMaximin preference reference klBudget policy →
          ∃ (klWeight : WithTop ℝ),
            AppliedModelingLib.GameTheory.PreferenceGame.IsContextFreeExtendedRegularizedPreferenceGameEquilibrium
              preference reference policy klWeight
\end{ReviewVerbatim}

\paragraph{Retained governing declarations}
\begin{itemize}\raggedright
\item \hyperlink{governing-declaration-578b1c4cb8565a0d}{Applied\allowbreak{}Modeling\allowbreak{}Lib.\allowbreak{}Game\allowbreak{}Theory.\allowbreak{}Preference\allowbreak{}Game.\allowbreak{}Is\allowbreak{}Context\allowbreak{}Free\allowbreak{}Extended\allowbreak{}Regularized\allowbreak{}Preference\allowbreak{}Game\allowbreak{}Equilibrium}
\item \hyperlink{governing-declaration-d8827c2845a356ad}{Applied\allowbreak{}Modeling\allowbreak{}Lib.\allowbreak{}Game\allowbreak{}Theory.\allowbreak{}Preference\allowbreak{}Game.\allowbreak{}Is\allowbreak{}Finite\allowbreak{}KL\allowbreak{}Preference\allowbreak{}Maximin}
\item \hyperlink{library-prerequisite-7e2db20cf087d9ba}{Applied\allowbreak{}Modeling\allowbreak{}Lib.\allowbreak{}Learning.\allowbreak{}Human\allowbreak{}Feedback.\allowbreak{}Pairwise\allowbreak{}Preference}
\item \hyperlink{library-prerequisite-268c959ee7d2d148}{Applied\allowbreak{}Modeling\allowbreak{}Lib.\allowbreak{}PMF\allowbreak{}Full\allowbreak{}Support}
\item \hyperlink{library-prerequisite-d3573881ae45d4c6}{Applied\allowbreak{}Modeling\allowbreak{}Lib.\allowbreak{}finite\allowbreak{}KL\allowbreak{}Divergence}
\end{itemize}
\paragraph{Lean-elaborated claim roles}\begin{ReviewVerbatim}
1. conclusion: (∀ {Alternative : Type} [inst : Fintype Alternative] [inst_1 : DecidableEq Alternative]
    (preference : AppliedModelingLib.Learning.HumanFeedback.PairwisePreference PUnit.{1} Alternative)
    (reference policy : PMF Alternative) (klWeight : WithTop ℝ),
    AppliedModelingLib.PMFFullSupport reference →
      AppliedModelingLib.GameTheory.PreferenceGame.IsContextFreeExtendedRegularizedPreferenceGameEquilibrium preference
          reference policy klWeight →
        AppliedModelingLib.GameTheory.PreferenceGame.IsFiniteKLPreferenceMaximin preference reference
          (AppliedModelingLib.finiteKLDivergence policy reference) policy) ∧
  ∀ {Alternative : Type} [inst : Fintype Alternative] [inst_1 : DecidableEq Alternative]
    (preference : AppliedModelingLib.Learning.HumanFeedback.PairwisePreference PUnit.{1} Alternative)
    (reference policy : PMF Alternative) (klBudget : ℝ),
    0 ≤ klBudget →
      AppliedModelingLib.PMFFullSupport reference →
        AppliedModelingLib.GameTheory.PreferenceGame.IsFiniteKLPreferenceMaximin preference reference klBudget policy →
          ∃ (klWeight : WithTop ℝ),
            AppliedModelingLib.GameTheory.PreferenceGame.IsContextFreeExtendedRegularizedPreferenceGameEquilibrium
              preference reference policy klWeight
\end{ReviewVerbatim}

\paragraph{Lean proof endpoint (built separately)}\begin{ReviewVerbatim}
GolzHaghtalabYang2025Distortion.proposition13_nlhf_equivalence
\end{ReviewVerbatim}

\paragraph{Recorded source-to-Spec screening}
\textbf{Verdict:} \texttt{matches}
\\
{\raggedright
\textbf{Reason:} Both regularized-\allowbreak{}to-\allowbreak{}constrained and constrained-\allowbreak{}to-\allowbreak{}regularized NLHF directions are formalized, including extended regularization weights.\allowbreak{}
\par}
\reviewmatch{Matches source input}
\noindent\textbf{Reviewer annotation}\par
\reviewerbox
\clearpage
\hypertarget{source-claim-fd234c18942cf2cc}{}
\section*{Review row 13}
\textbf{Source locator:} {\footnotesize\raggedright cited publication, lines 607-\allowbreak{}636\par}
\paragraph{Verbatim source input}
\begin{ReviewVerbatim}
Corollary 8. If λ-regularized NLHF (for any λ ≥ 0) returns a policy π̃NLHF , this policy’s average
                               −β
utility is at least a β2 · 1−e
                           1+e−β
                                  fraction of the optimal average utility of any policy π with DKL (π k πref ) ≤
DKL (π̃NLHF k πref ) (minus finite-sample errors, see Theorem 7).
   10
      Why not define distortion by flexibly selecting the benchmark based on the output policy’s KL divergence?
For this definition, an alignment method that always returns the reference policy would spuriously achieve
distortion 1: because its KL divergence is 0, we would benchmark it only against the reference policy, i.e., itself.


                                                               9

[form-feed in source extraction]
Sampling of Comparison Pairs. In our model, we assume that each voter provides d pairwise
comparisons, where both members xji , yij of each comparison pair are sampled i.i.d. from µ. More
generally, we can model {xji , yij } ∼ ν where ν is a distribution over unordered alternative pairs, or
even a distribution over d pairs of alternatives from which {x1i , yi1 , . . . , xdi , yid } are sampled. (To keep
the alignment methods well defined, we assume that each comparison pair has positive probability of
                         [U+0001 control character] of these models can, for example, express k-wise (rather than pairwise)
being sampled.) The latter
comparisons, if d = 2r are all pairs inside a randomly chosen set of r alternatives. Almost all of
our results continue to hold in these general models: the lower bound for all alignment methods that
satisfy the Condorcet loser criterion in the social choice setting (Theorem 3), the exponential lower
bound for RLHF (Theorem 6), and the upper bound for NLHF/Maximal Lotteries (Theorem 7)11 .
Given that our proofs continue to work out, the only “disadvantage” of these stronger models for
sampling comparison pairs is that, without a distribution µ, the Borda voting rule is no longer defined
(and we see no obvious way to generalize the Borda–MLE equivalence [46, 42]). It seems that RLHF
does not only become harder to analyze under these comparison-pair models, but actually performs
worse: we show in Appendix G that RLHF can have a distortion that is not bounded in β in these
extended models, leading to an even clearer separation with NLHF.

[Next verbatim source excerpt]

      If τ = 0, the only feasible policy is πref , and the claim clearly holds for λ = ∞. We also assume that
πref (a) > 0 for all a ∈ M . Otherwise if πref (a) = 0 for some a ∈ M , then both the regularized and constrained
versions forbid any policy to put nonzero probability on a, which leads to an effectively smaller candidate set.

[Next verbatim source excerpt]

Let A = {1, . . . , m} be a finite set of alternatives. The population of users is described by a
probability distribution D over utility vectors u = (u(1), . . . , u(m)), whose entries 0 ≤ u(x) ≤ 1
indicate a user’s utility for alternative x.3 The objective is to find an alternative x such that its average
utility AvgUtil(x) := Eu∼D [u(x)] across the user population (also known as the utilitarian social
welfare) is as high as possible. We extend this notation to probability distributions π over alternatives
by setting AvgUtil(π) := Ex∼π [AvgUtil(x)].
Both voting rules and alignment methods observe comparisons from n users. We model each
i = 1, . . . , n as a fresh user with independently drawn utility vector ui ∼ D. For exposition, we
assume that each user i provides an equal number d ≥ 1 of pairwise comparisons. For each i,
and for j = 1, . . . , d, we independently draw alternatives xji , yij from a fixed distribution µ over

    3
     Our assumption that utilities be in [0, 1] is weaker than any of the three assumptions — unit-sum, unit-range,
or approval [23] — made in classic distortion to avoid trivial infinite lower bounds.


                                                           3

[form-feed in source extraction]
alternatives, in which the minimum probability mass µmin := minx∈A µ(x) is positive.4 User i
then compares each pair {xji , yij } (for j = 1, . . . , d) through a Bradley-Terry model based on i’s
utilities: i prefers xji over yij (written “xji [U+001F control character]i yij ”) with probability σ β · (ui (xji ) − ui (yij )) , where
                                                                                                        [U+0001 control character]

σ(t) := 1/(1+e−t ) is the logistic sigmoid function and β > 0 is a temperature parameter, and prefers
yij over xji (“yij [U+001F control character]i xji ”) otherwise.5 Whereas this specifies the marginal probability of each pairwise
comparison, we make no assumption about the correlation between i’s choices. For example, i might
derive the pairwise comparisons from a Plackett-Luce ranking, ensuring that      [U+0002 control character] the user’s comparisons  [U+0001 control character][U+0003 control character]
are always consistent.6 We set p(x [U+001F control character] y) for the expected win rate Eu∼D σ β · (u(x) − u(y)) .

[Next verbatim source excerpt]

Alignment Setting. The alignment setting generalizes the social choice setting in two ways: first,
user utilities ui (y | x) may depend on a state x; second, the goal in determining a policy π is not
purely to maximize the reward AvgUtil(π | x), but a trade-off between this reward and the goal of
remaining close to a reference policy πref (· | x) ∈ ∆(A) in terms of the KL divergence DKL (· k πref ).
For theoretical tractability, we focus our analysis on a single state x, which we from here on omit
from the notation. Conceptually, this treatment of alignment on a state-by-state basis corresponds to
an assumption that our policy class is expressive enough so that it can take the optimal distribution of
actions at each state7 and abstracts from the generalization problem of estimating the population’s
preference between a pair of alternatives at the given state x based on preferences in similar states x0 .
Having set aside the dependency between states, we focus on how the regularization with respect
to a reference policy impacts the ability to optimize the average utility of three alignment methods:
RLHF, DPO, and NLHF. In addition to the pairwise comparisons, these methods take in a reference
policy πref ∈ ∆(A) and a KL bound τ ≥ 0, and map these inputs to a policy π in the KL-ball
Bτ (πref ) := {π ∈ ∆(A) | DKL (π k πref ) ≤ τ } around the reference policy. RLHF, DPO, and NLHF
are typically implemented with a KL regularization rather than our KL-constrained formulation,
which we adopt to enable a comparison on equal terms. In Section 5, we show that these perspectives
are equivalent, and that distortion upper bounds carry over to regularized alignment methods.

   4
     In Section 5, we discuss how most of our results extend more general distributions over comparison pairs,
which can, for example. capture k-wise comparisons.
   5
     Should we sample the same alternative x = xji = yij twice for a pair, the user is not asked for a pairwise
comparison. We record this as “x [U+001F control character]i x”, in a slight abuse of notation.
   6
     In particular, if we sample the same unordered pair twice for a user, the answers can be perfectly correlated.
   7
     This assumption is common in the literature; see, for example, Rafailov et al. [43]’s application of first-order
optimality conditions of PPO loss minimization.


                                                           4

[form-feed in source extraction]
The RLHF method first estimates rewards for each alternative, using maximum likelihood estimation
assuming that comparisons were generated by a single Bradley–Terry model:
                  r := argmaxr∈Rm 1≤i≤n,1≤j≤d log σ(r(xji ) − r(yij )) .8
                                     P                                     [U+0001 control character]

Next, RLHF uses PPO [44] to compute the policy πRLHF := argmaxπ∈Bτ (πref ) Ex∼π [r(x)] with
maximum expected reward within the KL-ball. In our setup, DPO is equivalent to RLHF (see
Appendix F.3) and thus has the same distortion.
The alignment method NLHF was inspired in part by a desire to better align with the
preferences of a heterogeneous group [36]. NLHF naturally adapts the definition of maxi-
mal lotteries by constraining both players’ mixed strategies to the KL-ball, i.e., πNLHF :=
argmaxπ1 ∈Bτ (πref ) minπ2 ∈Bτ (πref ) Ex1 ∼π1 ,x2 ∼π2 [Mx1 ,x2 ].
To generalize the definition of distortion to alignment methods, we set the maximum average utility of
any policy in the KL-ball as the benchmark. For fixed m, D, β, d, µ and correlation between pairwise
comparisons, the distortion of alignment method f is
                                                             maxπ∈Bτ (πref ) AvgUtil(π)
                      dist(f ) = supD,πref ,τ lim supn→∞       AvgUtiln (f (·,πref ,τ )) .
\end{ReviewVerbatim}



\paragraph{Expanded PaperInterface specification}
\begin{ReviewVerbatim}
∀ {User Alternative : Type} [inst : Fintype User] [inst_1 : DecidableEq User] [inst_2 : Fintype Alternative]
  [inst_3 : DecidableEq Alternative] (population : PMF User)
  (utility : AppliedModelingLib.Alignment.Welfare.FiniteUtilityProfile User Alternative),
  AppliedModelingLib.Alignment.Welfare.UnitIntervalUtilityProfile utility →
    ∀ {btScale : ℝ},
      0 < btScale →
        ∀ (reference policy : PMF Alternative) (klRegularization : ℝ),
          0 ≤ klRegularization →
            AppliedModelingLib.PMFFullSupport reference →
              AppliedModelingLib.GameTheory.PreferenceGame.IsRegularizedPreferenceGameEquilibrium (PMF.pure PUnit.unit)
                  (AppliedModelingLib.Alignment.Welfare.populationBradleyTerryPreference population utility btScale)
                  (AppliedModelingLib.Alignment.Welfare.contextFreeAlternativePolicy reference)
                  (AppliedModelingLib.Alignment.Welfare.contextFreeAlternativePolicy policy) klRegularization →
                ∀ (benchmark : PMF Alternative),
                  AppliedModelingLib.finiteKLDivergence benchmark reference ≤
                      AppliedModelingLib.finiteKLDivergence policy reference →
                    AppliedModelingLib.Alignment.Welfare.sigmoidChordSlope btScale / (1 / 4) *
                        AppliedModelingLib.Alignment.Welfare.policyAverageUtility population utility benchmark ≤
                      AppliedModelingLib.Alignment.Welfare.policyAverageUtility population utility policy
\end{ReviewVerbatim}

\paragraph{Retained governing declarations}
\begin{itemize}\raggedright
\item \hyperlink{library-prerequisite-56c0734f8877ddbe}{Applied\allowbreak{}Modeling\allowbreak{}Lib.\allowbreak{}Alignment.\allowbreak{}Welfare.\allowbreak{}Finite\allowbreak{}Utility\allowbreak{}Profile}
\item \hyperlink{library-prerequisite-cbe8afa0d6d1ba78}{Applied\allowbreak{}Modeling\allowbreak{}Lib.\allowbreak{}Alignment.\allowbreak{}Welfare.\allowbreak{}Unit\allowbreak{}Interval\allowbreak{}Utility\allowbreak{}Profile}
\item \hyperlink{governing-declaration-b7cf9e04e23161e5}{Applied\allowbreak{}Modeling\allowbreak{}Lib.\allowbreak{}Alignment.\allowbreak{}Welfare.\allowbreak{}context\allowbreak{}Free\allowbreak{}Alternative\allowbreak{}Policy}
\item \hyperlink{library-prerequisite-184c74a4db4311c2}{Applied\allowbreak{}Modeling\allowbreak{}Lib.\allowbreak{}Alignment.\allowbreak{}Welfare.\allowbreak{}policy\allowbreak{}Average\allowbreak{}Utility}
\item \hyperlink{library-prerequisite-cf06d15d7f8b7c31}{Applied\allowbreak{}Modeling\allowbreak{}Lib.\allowbreak{}Alignment.\allowbreak{}Welfare.\allowbreak{}population\allowbreak{}Bradley\allowbreak{}Terry\allowbreak{}Preference}
\item \hyperlink{governing-declaration-a4d623b8c608b566}{Applied\allowbreak{}Modeling\allowbreak{}Lib.\allowbreak{}Alignment.\allowbreak{}Welfare.\allowbreak{}sigmoid\allowbreak{}Chord\allowbreak{}Slope}
\item \hyperlink{governing-declaration-49e6e1fe2a2cc849}{Applied\allowbreak{}Modeling\allowbreak{}Lib.\allowbreak{}Game\allowbreak{}Theory.\allowbreak{}Preference\allowbreak{}Game.\allowbreak{}Is\allowbreak{}Regularized\allowbreak{}Preference\allowbreak{}Game\allowbreak{}Equilibrium}
\item \hyperlink{library-prerequisite-268c959ee7d2d148}{Applied\allowbreak{}Modeling\allowbreak{}Lib.\allowbreak{}PMF\allowbreak{}Full\allowbreak{}Support}
\item \hyperlink{library-prerequisite-d3573881ae45d4c6}{Applied\allowbreak{}Modeling\allowbreak{}Lib.\allowbreak{}finite\allowbreak{}KL\allowbreak{}Divergence}
\end{itemize}
\paragraph{Lean-elaborated claim roles}\begin{ReviewVerbatim}
1. parameter: Type
2. parameter: Type
3. parameter: Fintype User
4. parameter: DecidableEq User
5. parameter: Fintype Alternative
6. parameter: DecidableEq Alternative
7. parameter: PMF User
8. parameter: AppliedModelingLib.Alignment.Welfare.FiniteUtilityProfile User Alternative
9. assumption: AppliedModelingLib.Alignment.Welfare.UnitIntervalUtilityProfile utility
10. parameter: ℝ
11. assumption: 0 < btScale
12. parameter: PMF Alternative
13. parameter: PMF Alternative
14. parameter: ℝ
15. assumption: 0 ≤ klRegularization
16. assumption: AppliedModelingLib.PMFFullSupport reference
17. assumption: AppliedModelingLib.GameTheory.PreferenceGame.IsRegularizedPreferenceGameEquilibrium (PMF.pure PUnit.unit)
  (AppliedModelingLib.Alignment.Welfare.populationBradleyTerryPreference population utility btScale)
  (AppliedModelingLib.Alignment.Welfare.contextFreeAlternativePolicy reference)
  (AppliedModelingLib.Alignment.Welfare.contextFreeAlternativePolicy policy) klRegularization
18. parameter: PMF Alternative
19. assumption: AppliedModelingLib.finiteKLDivergence benchmark reference ≤ AppliedModelingLib.finiteKLDivergence policy reference
20. conclusion: AppliedModelingLib.Alignment.Welfare.sigmoidChordSlope btScale / (1 / 4) *
    AppliedModelingLib.Alignment.Welfare.policyAverageUtility population utility benchmark ≤
  AppliedModelingLib.Alignment.Welfare.policyAverageUtility population utility policy
\end{ReviewVerbatim}

\paragraph{Lean proof endpoint (built separately)}\begin{ReviewVerbatim}
GolzHaghtalabYang2025Distortion.corollary8_regularized_nlhf
\end{ReviewVerbatim}

\paragraph{Recorded source-to-Spec screening}
\textbf{Verdict:} \texttt{matches}
\\
{\raggedright
\textbf{Reason:} Regularized NLHF is compared with every policy of no greater KL divergence and receives the Theorem-\allowbreak{}7 welfare factor.\allowbreak{}
\par}
\reviewmatch{Matches source input}
\noindent\textbf{Reviewer annotation}\par
\reviewerbox
\clearpage
\hypertarget{source-claim-ddcbcef3944ac5d9}{}
\section*{Review row 14}
\textbf{Source locator:} {\footnotesize\raggedright cited publication, lines 2263-\allowbreak{}2280\par}
\paragraph{Verbatim source input}
\begin{ReviewVerbatim}
Proposition 14 (Equivalence between Constrained and Regularized RLHF). Let r be the MLE reward
learned from the comparison data, and let πτ be the output of the τ -constrained RLHF method defined
in Section 2, i.e.,
                                         πτ = argmax E [r(x)] ,                                             (19)
                                               π∈Bτ (πref ) x∼π

and let π̃λ be the output of the λ-regularized RLHF method, i.e.,
                              π̃λ = argmax E [r(x)] − λ · DKL (π k πref ).                                  (20)
                                     π∈∆(M ) x∼π

Then, for each λ ∈ [0, ∞] with solution π̃λ for Eq. (20), we have that π̃λ is also an optimal solution
to the τ -constrained optimization problem in Eq. (19), where τ = DKL (π̃λ k πref ). Conversely, for
each τ ≥ 0 with solution πτ for Eq. (19), there exists λ ∈ [0, ∞] such that πτ is also an optimal
solution to the λ-regularized optimization problem in Eq. (20).
   14
      If τ = 0, the only feasible policy is πref , and the claim clearly holds for λ = ∞. We also assume that
πref (a) > 0 for all a ∈ M . Otherwise if πref (a) = 0 for some a ∈ M , then both the regularized and constrained
versions forbid any policy to put nonzero probability on a, which leads to an effectively smaller candidate set.

[Next verbatim source excerpt]

Let A = {1, . . . , m} be a finite set of alternatives. The population of users is described by a
probability distribution D over utility vectors u = (u(1), . . . , u(m)), whose entries 0 ≤ u(x) ≤ 1
indicate a user’s utility for alternative x.3 The objective is to find an alternative x such that its average
utility AvgUtil(x) := Eu∼D [u(x)] across the user population (also known as the utilitarian social
welfare) is as high as possible. We extend this notation to probability distributions π over alternatives
by setting AvgUtil(π) := Ex∼π [AvgUtil(x)].
Both voting rules and alignment methods observe comparisons from n users. We model each
i = 1, . . . , n as a fresh user with independently drawn utility vector ui ∼ D. For exposition, we
assume that each user i provides an equal number d ≥ 1 of pairwise comparisons. For each i,
and for j = 1, . . . , d, we independently draw alternatives xji , yij from a fixed distribution µ over

    3
     Our assumption that utilities be in [0, 1] is weaker than any of the three assumptions — unit-sum, unit-range,
or approval [23] — made in classic distortion to avoid trivial infinite lower bounds.


                                                           3

[form-feed in source extraction]
alternatives, in which the minimum probability mass µmin := minx∈A µ(x) is positive.4 User i
then compares each pair {xji , yij } (for j = 1, . . . , d) through a Bradley-Terry model based on i’s
utilities: i prefers xji over yij (written “xji [U+001F control character]i yij ”) with probability σ β · (ui (xji ) − ui (yij )) , where
                                                                                                        [U+0001 control character]

σ(t) := 1/(1+e−t ) is the logistic sigmoid function and β > 0 is a temperature parameter, and prefers
yij over xji (“yij [U+001F control character]i xji ”) otherwise.5 Whereas this specifies the marginal probability of each pairwise
comparison, we make no assumption about the correlation between i’s choices. For example, i might
derive the pairwise comparisons from a Plackett-Luce ranking, ensuring that      [U+0002 control character] the user’s comparisons  [U+0001 control character][U+0003 control character]
are always consistent.6 We set p(x [U+001F control character] y) for the expected win rate Eu∼D σ β · (u(x) − u(y)) .

[Next verbatim source excerpt]

Alignment Setting. The alignment setting generalizes the social choice setting in two ways: first,
user utilities ui (y | x) may depend on a state x; second, the goal in determining a policy π is not
purely to maximize the reward AvgUtil(π | x), but a trade-off between this reward and the goal of
remaining close to a reference policy πref (· | x) ∈ ∆(A) in terms of the KL divergence DKL (· k πref ).
For theoretical tractability, we focus our analysis on a single state x, which we from here on omit
from the notation. Conceptually, this treatment of alignment on a state-by-state basis corresponds to
an assumption that our policy class is expressive enough so that it can take the optimal distribution of
actions at each state7 and abstracts from the generalization problem of estimating the population’s
preference between a pair of alternatives at the given state x based on preferences in similar states x0 .
Having set aside the dependency between states, we focus on how the regularization with respect
to a reference policy impacts the ability to optimize the average utility of three alignment methods:
RLHF, DPO, and NLHF. In addition to the pairwise comparisons, these methods take in a reference
policy πref ∈ ∆(A) and a KL bound τ ≥ 0, and map these inputs to a policy π in the KL-ball
Bτ (πref ) := {π ∈ ∆(A) | DKL (π k πref ) ≤ τ } around the reference policy. RLHF, DPO, and NLHF
are typically implemented with a KL regularization rather than our KL-constrained formulation,
which we adopt to enable a comparison on equal terms. In Section 5, we show that these perspectives
are equivalent, and that distortion upper bounds carry over to regularized alignment methods.

   4
     In Section 5, we discuss how most of our results extend more general distributions over comparison pairs,
which can, for example. capture k-wise comparisons.
   5
     Should we sample the same alternative x = xji = yij twice for a pair, the user is not asked for a pairwise
comparison. We record this as “x [U+001F control character]i x”, in a slight abuse of notation.
   6
     In particular, if we sample the same unordered pair twice for a user, the answers can be perfectly correlated.
   7
     This assumption is common in the literature; see, for example, Rafailov et al. [43]’s application of first-order
optimality conditions of PPO loss minimization.


                                                           4

[form-feed in source extraction]
The RLHF method first estimates rewards for each alternative, using maximum likelihood estimation
assuming that comparisons were generated by a single Bradley–Terry model:
                  r := argmaxr∈Rm 1≤i≤n,1≤j≤d log σ(r(xji ) − r(yij )) .8
                                     P                                     [U+0001 control character]

Next, RLHF uses PPO [44] to compute the policy πRLHF := argmaxπ∈Bτ (πref ) Ex∼π [r(x)] with
maximum expected reward within the KL-ball. In our setup, DPO is equivalent to RLHF (see
Appendix F.3) and thus has the same distortion.
The alignment method NLHF was inspired in part by a desire to better align with the
preferences of a heterogeneous group [36]. NLHF naturally adapts the definition of maxi-
mal lotteries by constraining both players’ mixed strategies to the KL-ball, i.e., πNLHF :=
argmaxπ1 ∈Bτ (πref ) minπ2 ∈Bτ (πref ) Ex1 ∼π1 ,x2 ∼π2 [Mx1 ,x2 ].
To generalize the definition of distortion to alignment methods, we set the maximum average utility of
any policy in the KL-ball as the benchmark. For fixed m, D, β, d, µ and correlation between pairwise
comparisons, the distortion of alignment method f is
                                                             maxπ∈Bτ (πref ) AvgUtil(π)
                      dist(f ) = supD,πref ,τ lim supn→∞       AvgUtiln (f (·,πref ,τ )) .
\end{ReviewVerbatim}



\paragraph{Expanded PaperInterface specification}
\begin{ReviewVerbatim}
(∀ {Alternative : Type} [inst : Fintype Alternative] [inst_1 : DecidableEq Alternative]
    (reference policy : PMF Alternative) (score : Alternative → ℝ) (klWeight : WithTop ℝ),
    AppliedModelingLib.PMFFullSupport reference →
      AppliedModelingLib.IsFiniteKLExtendedRegularizedUtilityMax reference policy score klWeight →
        AppliedModelingLib.finiteKLDivergence policy reference ≤
            AppliedModelingLib.finiteKLDivergence policy reference ∧
          ∀ (other : PMF Alternative),
            AppliedModelingLib.finiteKLDivergence other reference ≤
                AppliedModelingLib.finiteKLDivergence policy reference →
              AppliedModelingLib.pmfExp other score ≤ AppliedModelingLib.pmfExp policy score) ∧
  ∀ {Alternative : Type} [inst : Fintype Alternative] [inst_1 : DecidableEq Alternative]
    (reference policy : PMF Alternative) (score : Alternative → ℝ) (klBudget : ℝ),
    0 ≤ klBudget →
      AppliedModelingLib.PMFFullSupport reference →
        AppliedModelingLib.finiteKLDivergence policy reference ≤ klBudget →
          (∀ (other : PMF Alternative),
              AppliedModelingLib.finiteKLDivergence other reference ≤ klBudget →
                AppliedModelingLib.pmfExp other score ≤ AppliedModelingLib.pmfExp policy score) →
            ∃ (klWeight : WithTop ℝ),
              AppliedModelingLib.IsFiniteKLExtendedRegularizedUtilityMax reference policy score klWeight
\end{ReviewVerbatim}

\paragraph{Retained governing declarations}
\begin{itemize}\raggedright
\item \hyperlink{governing-declaration-603c21f4ad34025c}{Applied\allowbreak{}Modeling\allowbreak{}Lib.\allowbreak{}Is\allowbreak{}Finite\allowbreak{}KL\allowbreak{}Extended\allowbreak{}Regularized\allowbreak{}Utility\allowbreak{}Max}
\item \hyperlink{library-prerequisite-268c959ee7d2d148}{Applied\allowbreak{}Modeling\allowbreak{}Lib.\allowbreak{}PMF\allowbreak{}Full\allowbreak{}Support}
\item \hyperlink{library-prerequisite-d3573881ae45d4c6}{Applied\allowbreak{}Modeling\allowbreak{}Lib.\allowbreak{}finite\allowbreak{}KL\allowbreak{}Divergence}
\item \hyperlink{governing-declaration-8e4b4389a4f6454b}{Applied\allowbreak{}Modeling\allowbreak{}Lib.\allowbreak{}pmf\allowbreak{}Exp}
\end{itemize}
\paragraph{Lean-elaborated claim roles}\begin{ReviewVerbatim}
1. conclusion: (∀ {Alternative : Type} [inst : Fintype Alternative] [inst_1 : DecidableEq Alternative]
    (reference policy : PMF Alternative) (score : Alternative → ℝ) (klWeight : WithTop ℝ),
    AppliedModelingLib.PMFFullSupport reference →
      AppliedModelingLib.IsFiniteKLExtendedRegularizedUtilityMax reference policy score klWeight →
        AppliedModelingLib.finiteKLDivergence policy reference ≤
            AppliedModelingLib.finiteKLDivergence policy reference ∧
          ∀ (other : PMF Alternative),
            AppliedModelingLib.finiteKLDivergence other reference ≤
                AppliedModelingLib.finiteKLDivergence policy reference →
              AppliedModelingLib.pmfExp other score ≤ AppliedModelingLib.pmfExp policy score) ∧
  ∀ {Alternative : Type} [inst : Fintype Alternative] [inst_1 : DecidableEq Alternative]
    (reference policy : PMF Alternative) (score : Alternative → ℝ) (klBudget : ℝ),
    0 ≤ klBudget →
      AppliedModelingLib.PMFFullSupport reference →
        AppliedModelingLib.finiteKLDivergence policy reference ≤ klBudget →
          (∀ (other : PMF Alternative),
              AppliedModelingLib.finiteKLDivergence other reference ≤ klBudget →
                AppliedModelingLib.pmfExp other score ≤ AppliedModelingLib.pmfExp policy score) →
            ∃ (klWeight : WithTop ℝ),
              AppliedModelingLib.IsFiniteKLExtendedRegularizedUtilityMax reference policy score klWeight
\end{ReviewVerbatim}

\paragraph{Lean proof endpoint (built separately)}\begin{ReviewVerbatim}
GolzHaghtalabYang2025Distortion.proposition14_rlhf_equivalence
\end{ReviewVerbatim}

\paragraph{Recorded source-to-Spec screening}
\textbf{Verdict:} \texttt{matches}
\\
{\raggedright
\textbf{Reason:} This formalizes the analogous two RLHF directions at radius equal to the output policy's KL divergence.\allowbreak{}
\par}
\reviewmatch{Matches source input}
\noindent\textbf{Reviewer annotation}\par
\reviewerbox
\clearpage
\hypertarget{source-claim-f1c80c98e2c74036}{}
\section*{Review row 15}
\textbf{Source locator:} {\footnotesize\raggedright cited publication, lines 637-\allowbreak{}640\par}
\paragraph{Verbatim source input}
\begin{ReviewVerbatim}
Theorem 9 (Unbounded Distortion of RLHF Under Correlated Sampling).             For any β > 0, there
exists a sequence of alignment instances and distributions ν ∈ ∆ A2 over comparison pairs such
                                                                       [U+0001 control character][U+0001 control character]
that RLHF’s distortion is unbounded.

[Next verbatim source excerpt]

      If τ = 0, the only feasible policy is πref , and the claim clearly holds for λ = ∞. We also assume that
πref (a) > 0 for all a ∈ M . Otherwise if πref (a) = 0 for some a ∈ M , then both the regularized and constrained
versions forbid any policy to put nonzero probability on a, which leads to an effectively smaller candidate set.

[Next verbatim source excerpt]


[form-feed in source extraction]
Sampling of Comparison Pairs. In our model, we assume that each voter provides d pairwise
comparisons, where both members xji , yij of each comparison pair are sampled i.i.d. from µ. More
generally, we can model {xji , yij } ∼ ν where ν is a distribution over unordered alternative pairs, or
even a distribution over d pairs of alternatives from which {x1i , yi1 , . . . , xdi , yid } are sampled. (To keep
the alignment methods well defined, we assume that each comparison pair has positive probability of
                         [U+0001 control character] of these models can, for example, express k-wise (rather than pairwise)
being sampled.) The latter

[Next verbatim source excerpt]

that RLHF’s distortion is unbounded.

Proof. We construct our sequence of instances by taking increasingly long prefixes of the sequence
in Lemma 15, i.e., by considering the alternatives a1 , . . . , am for increasing m. Rescaling by some
constants, we can define the MLE rewards in RLHF as
                    X #(x [U+001F control character] y)                                 [U+0001 control character] #(y [U+001F control character] x)                         [U+0001 control character]
    r = argmax                       log σ(r(x) − r(y)) +                      log σ(r(y) − r(x)) .
          r∈Rm                nd                                       nd
                {x,y}∈(A2)


Following Siththaranjan et al. [46], we apply the first-order optimality conditions to obtain that, for
each alternative x,
                  X #(x [U+001F control character] y) X #(x [U+001F control character] y) + #(y [U+001F control character] x)
                                  =                               · σ(r(x) − r(y)).
                         nd                        nd
                      y6=x                    y6=x


                                                              38

[form-feed in source extraction]
By the strong law of large numbers, as the number n of samples goes to infinity (regardless of d), the
sample fraction #(x[U+001F control character]y)
                   nd    converges almost surely to its expected value ν({x, y}) · p(x [U+001F control character] y). Hence,
as n → ∞, the rewards (a random variable depending on the random pairwise comparisons) will
satisfy that
                                                   a.s.
                  X                                     X
                       ν({x, y}) · σ(r(x) − r(y)) −→        ν({x, y}) · p(x [U+001F control character] y).                (22)
                     y6=x                                     y6=x

Consider a distribution ν over pairs of alternatives that assigns each pair of adjacent alternatives
                              1−[U+000F control character]
{at , at+1 } a probability of m−1 of being drawn for comparison, and all other pairs a probability
         [U+000F control character]
of m −(m−1) , where [U+000F control character] > 0 is a small value, dependent on the current m, to be determined in the

[Next verbatim source excerpt]

  Proof. Our population D will be a uniform distribution over three utility vectors, u1 , u2 , and u3 . We
  define the sequence of alternatives and prove the claim by induction over t ≥ 1.
  For t = 1, set u1 (a1 ) = u2 (a1 ) = u3 (a1 ) := 1/3, which clearly satisfies the first claim.
  Now, let t ≥ 2, and suppose that we have defined the utilities for alternatives a1 , . . . , at−1 and
  established the claims for all t0 < t. We define utilities for at and extend the claims to t. Let uA
  denote the utility vector among u1 , u2 , u3 with the highest utility for at−1 , and denote the other two
                                                                                            ∆/2    3[U+0001 control character]
  utility vectors by uB , uC . For convenience, set ∆ := uA (at−1 ) · β and ∆0 := log (e2 (e∆+1)+3) .

  uA (at ) := uA (at−1 )−∆/β = 0,                                              uB (at ) := uB (at−1 )+∆0 /β,                                           uC (at ) := uC (at−1 )+∆0 /β.

[Next verbatim source excerpt]

  It will be useful to derive an alternative expression for ∆0 :
                           [U+0012 control character] ∆/2
                                    + 1)3                    (e∆/2 + 1)3
                                           [U+0013 control character]       [U+0012 control character]                       [U+0013 control character]
                             (e
                ∆0 = log                     = log   e ∆/2
                              2 (e∆ + 3)                   2 e∆/2 (e∆ + 3)
                                     ∆       2 e∆/2 (e∆ + 3)     ∆       (e∆/2 + 1)3 + (e∆/2 − 1)3
                                   =   − log      ∆/2      3
                                                              =    − log
                                     2         (e     + 1)       2              (e∆/2 + 1)3
                                     ∆
                                             [U+0012 control character]     [U+0010 control character] e∆/2 − 1 [U+0011 control character]3 [U+0013 control character]
                                                                     ∆
                                                                        − log 1 + tanh(∆/4)3 .
                                                                                              [U+0001 control character]
                                   =   − log 1 + ∆/2               =
                                     2               e    +1         2

[Next verbatim source excerpt]

Let A = {1, . . . , m} be a finite set of alternatives. The population of users is described by a
probability distribution D over utility vectors u = (u(1), . . . , u(m)), whose entries 0 ≤ u(x) ≤ 1
indicate a user’s utility for alternative x.3 The objective is to find an alternative x such that its average
utility AvgUtil(x) := Eu∼D [u(x)] across the user population (also known as the utilitarian social
welfare) is as high as possible. We extend this notation to probability distributions π over alternatives
by setting AvgUtil(π) := Ex∼π [AvgUtil(x)].
Both voting rules and alignment methods observe comparisons from n users. We model each
i = 1, . . . , n as a fresh user with independently drawn utility vector ui ∼ D. For exposition, we
assume that each user i provides an equal number d ≥ 1 of pairwise comparisons. For each i,
and for j = 1, . . . , d, we independently draw alternatives xji , yij from a fixed distribution µ over

    3
     Our assumption that utilities be in [0, 1] is weaker than any of the three assumptions — unit-sum, unit-range,
or approval [23] — made in classic distortion to avoid trivial infinite lower bounds.


                                                           3

[form-feed in source extraction]
alternatives, in which the minimum probability mass µmin := minx∈A µ(x) is positive.4 User i
then compares each pair {xji , yij } (for j = 1, . . . , d) through a Bradley-Terry model based on i’s
utilities: i prefers xji over yij (written “xji [U+001F control character]i yij ”) with probability σ β · (ui (xji ) − ui (yij )) , where
                                                                                                        [U+0001 control character]

σ(t) := 1/(1+e−t ) is the logistic sigmoid function and β > 0 is a temperature parameter, and prefers
yij over xji (“yij [U+001F control character]i xji ”) otherwise.5 Whereas this specifies the marginal probability of each pairwise
comparison, we make no assumption about the correlation between i’s choices. For example, i might
derive the pairwise comparisons from a Plackett-Luce ranking, ensuring that      [U+0002 control character] the user’s comparisons  [U+0001 control character][U+0003 control character]
are always consistent.6 We set p(x [U+001F control character] y) for the expected win rate Eu∼D σ β · (u(x) − u(y)) .

[Next verbatim source excerpt]

Alignment Setting. The alignment setting generalizes the social choice setting in two ways: first,
user utilities ui (y | x) may depend on a state x; second, the goal in determining a policy π is not
purely to maximize the reward AvgUtil(π | x), but a trade-off between this reward and the goal of
remaining close to a reference policy πref (· | x) ∈ ∆(A) in terms of the KL divergence DKL (· k πref ).
For theoretical tractability, we focus our analysis on a single state x, which we from here on omit
from the notation. Conceptually, this treatment of alignment on a state-by-state basis corresponds to
an assumption that our policy class is expressive enough so that it can take the optimal distribution of
actions at each state7 and abstracts from the generalization problem of estimating the population’s
preference between a pair of alternatives at the given state x based on preferences in similar states x0 .
Having set aside the dependency between states, we focus on how the regularization with respect
to a reference policy impacts the ability to optimize the average utility of three alignment methods:
RLHF, DPO, and NLHF. In addition to the pairwise comparisons, these methods take in a reference
policy πref ∈ ∆(A) and a KL bound τ ≥ 0, and map these inputs to a policy π in the KL-ball
Bτ (πref ) := {π ∈ ∆(A) | DKL (π k πref ) ≤ τ } around the reference policy. RLHF, DPO, and NLHF
are typically implemented with a KL regularization rather than our KL-constrained formulation,
which we adopt to enable a comparison on equal terms. In Section 5, we show that these perspectives
are equivalent, and that distortion upper bounds carry over to regularized alignment methods.

   4
     In Section 5, we discuss how most of our results extend more general distributions over comparison pairs,
which can, for example. capture k-wise comparisons.
   5
     Should we sample the same alternative x = xji = yij twice for a pair, the user is not asked for a pairwise
comparison. We record this as “x [U+001F control character]i x”, in a slight abuse of notation.
   6
     In particular, if we sample the same unordered pair twice for a user, the answers can be perfectly correlated.
   7
     This assumption is common in the literature; see, for example, Rafailov et al. [43]’s application of first-order
optimality conditions of PPO loss minimization.


                                                           4

[form-feed in source extraction]
The RLHF method first estimates rewards for each alternative, using maximum likelihood estimation
assuming that comparisons were generated by a single Bradley–Terry model:
                  r := argmaxr∈Rm 1≤i≤n,1≤j≤d log σ(r(xji ) − r(yij )) .8
                                     P                                     [U+0001 control character]

Next, RLHF uses PPO [44] to compute the policy πRLHF := argmaxπ∈Bτ (πref ) Ex∼π [r(x)] with
maximum expected reward within the KL-ball. In our setup, DPO is equivalent to RLHF (see
Appendix F.3) and thus has the same distortion.
The alignment method NLHF was inspired in part by a desire to better align with the
preferences of a heterogeneous group [36]. NLHF naturally adapts the definition of maxi-
mal lotteries by constraining both players’ mixed strategies to the KL-ball, i.e., πNLHF :=
argmaxπ1 ∈Bτ (πref ) minπ2 ∈Bτ (πref ) Ex1 ∼π1 ,x2 ∼π2 [Mx1 ,x2 ].
To generalize the definition of distortion to alignment methods, we set the maximum average utility of
any policy in the KL-ball as the benchmark. For fixed m, D, β, d, µ and correlation between pairwise
comparisons, the distortion of alignment method f is
                                                             maxπ∈Bτ (πref ) AvgUtil(π)
                      dist(f ) = supD,πref ,τ lim supn→∞       AvgUtiln (f (·,πref ,τ )) .
\end{ReviewVerbatim}



\paragraph{Expanded PaperInterface specification}
\begin{ReviewVerbatim}
∀ {beta : ℝ},
  0 < beta →
    ∀ (bound : ℝ),
      ∃ (m : ℕ) (hm : 3 ≤ m) (epsilon : NNReal) (_ : 0 < epsilon) (_ : epsilon < 1),
        have population : PMF (Fin 3) := GolzHaghtalabYang2025Distortion.lemma15Population;
        have utility : AppliedModelingLib.Alignment.Welfare.FiniteUtilityProfile (Fin 3) (Fin m) :=
          GolzHaghtalabYang2025Distortion.theorem9PrefixProfile beta m;
        have comparisonSampling : AppliedModelingLib.Learning.HumanFeedback.CorrelatedPairSampling (Fin m) :=
          GolzHaghtalabYang2025Distortion.theorem9ComparisonSampling m hm epsilon;
        have reference : PMF (Fin m) :=
          GolzHaghtalabYang2025Distortion.theorem9UniformReference m
            (GolzHaghtalabYang2025Distortion.theorem9Spec._proof_3 m hm);
        have klBudget : ℝ := Real.log ↑m;
        AppliedModelingLib.Alignment.Welfare.UnitIntervalUtilityProfile utility ∧
          AppliedModelingLib.PMFFullSupport reference ∧
            comparisonSampling.IsSymmetric ∧
              comparisonSampling.HasZeroDiagonal ∧
                comparisonSampling.HasFullOffDiagonalSupport ∧
                  0 ≤ klBudget ∧
                    ∀ᵐ (path :
                      ℕ →
                        AppliedModelingLib.Learning.HumanFeedback.BinaryPairwiseReport
                          (Fin
                            m)) ∂AppliedModelingLib.finitePMFIidPathMeasure
                        (GolzHaghtalabYang2025Distortion.theorem9OneComparisonReportLaw beta m hm epsilon),
                      ∀ᶠ (horizon : ℕ) in Filter.atTop,
                        (∃ (reward :
                            AppliedModelingLib.Learning.HumanFeedback.PairwiseCountDataset.ScoreVector (Fin m)),
                            (AppliedModelingLib.Learning.HumanFeedback.PairwiseCountDataset.ofBinaryReports
                                  (AppliedModelingLib.finitePMFIidPathPrefix path horizon)).isPairwiseMLE
                              AppliedModelingLib.Learning.HumanFeedback.PairwiseCountDataset.sigmoidCDFLikePairwiseLink.toFun
                              ⟨0, GolzHaghtalabYang2025Distortion.theorem9Spec._proof_4 m hm⟩ reward) ∧
                          ∀
                            (reward :
                              AppliedModelingLib.Learning.HumanFeedback.PairwiseCountDataset.ScoreVector (Fin m)),
                            (AppliedModelingLib.Learning.HumanFeedback.PairwiseCountDataset.ofBinaryReports
                                    (AppliedModelingLib.finitePMFIidPathPrefix path horizon)).isPairwiseMLE
                                AppliedModelingLib.Learning.HumanFeedback.PairwiseCountDataset.sigmoidCDFLikePairwiseLink.toFun
                                ⟨0, GolzHaghtalabYang2025Distortion.theorem9Spec._proof_4 m hm⟩ reward →
                              ∃ (rewardPolicy : PMF (Fin m)) (welfarePolicy : PMF (Fin m)),
                                AppliedModelingLib.finiteKLDivergence rewardPolicy reference ≤ klBudget ∧
                                  (∀ (other : PMF (Fin m)),
                                      AppliedModelingLib.finiteKLDivergence other reference ≤ klBudget →
                                        AppliedModelingLib.pmfExp other reward ≤
                                          AppliedModelingLib.pmfExp rewardPolicy reward) ∧
                                    AppliedModelingLib.finiteKLDivergence welfarePolicy reference ≤ klBudget ∧
                                      (∀ (other : PMF (Fin m)),
                                          AppliedModelingLib.finiteKLDivergence other reference ≤ klBudget →
                                            AppliedModelingLib.Alignment.Welfare.policyAverageUtility population utility
                                                other ≤
                                              AppliedModelingLib.Alignment.Welfare.policyAverageUtility population
                                                utility welfarePolicy) ∧
                                        0 <
                                            AppliedModelingLib.Alignment.Welfare.policyAverageUtility population utility
                                              rewardPolicy ∧
                                          bound <
                                            AppliedModelingLib.Alignment.Welfare.policyAverageUtility population utility
                                                welfarePolicy /
                                              AppliedModelingLib.Alignment.Welfare.policyAverageUtility population
                                                utility rewardPolicy
\end{ReviewVerbatim}

\paragraph{Retained governing declarations}
\begin{itemize}\raggedright
\item \hyperlink{library-prerequisite-56c0734f8877ddbe}{Applied\allowbreak{}Modeling\allowbreak{}Lib.\allowbreak{}Alignment.\allowbreak{}Welfare.\allowbreak{}Finite\allowbreak{}Utility\allowbreak{}Profile}
\item \hyperlink{library-prerequisite-cbe8afa0d6d1ba78}{Applied\allowbreak{}Modeling\allowbreak{}Lib.\allowbreak{}Alignment.\allowbreak{}Welfare.\allowbreak{}Unit\allowbreak{}Interval\allowbreak{}Utility\allowbreak{}Profile}
\item \hyperlink{library-prerequisite-184c74a4db4311c2}{Applied\allowbreak{}Modeling\allowbreak{}Lib.\allowbreak{}Alignment.\allowbreak{}Welfare.\allowbreak{}policy\allowbreak{}Average\allowbreak{}Utility}
\item \hyperlink{governing-declaration-bc3d6ca8e7e0ac9d}{Applied\allowbreak{}Modeling\allowbreak{}Lib.\allowbreak{}Learning.\allowbreak{}Human\allowbreak{}Feedback.\allowbreak{}Binary\allowbreak{}Pairwise\allowbreak{}Report}
\item \hyperlink{governing-declaration-588529a2add6b6a9}{Applied\allowbreak{}Modeling\allowbreak{}Lib.\allowbreak{}Learning.\allowbreak{}Human\allowbreak{}Feedback.\allowbreak{}Correlated\allowbreak{}Pair\allowbreak{}Sampling}
\item \hyperlink{governing-declaration-4ded7e185e3a56c5}{Applied\allowbreak{}Modeling\allowbreak{}Lib.\allowbreak{}Learning.\allowbreak{}Human\allowbreak{}Feedback.\allowbreak{}Correlated\allowbreak{}Pair\allowbreak{}Sampling.\allowbreak{}Has\allowbreak{}Full\allowbreak{}Off\allowbreak{}Diagonal\allowbreak{}Support}
\item \hyperlink{governing-declaration-28f7d7214bf02489}{Applied\allowbreak{}Modeling\allowbreak{}Lib.\allowbreak{}Learning.\allowbreak{}Human\allowbreak{}Feedback.\allowbreak{}Correlated\allowbreak{}Pair\allowbreak{}Sampling.\allowbreak{}Has\allowbreak{}Zero\allowbreak{}Diagonal}
\item \hyperlink{governing-declaration-11432d593d43fd7d}{Applied\allowbreak{}Modeling\allowbreak{}Lib.\allowbreak{}Learning.\allowbreak{}Human\allowbreak{}Feedback.\allowbreak{}Correlated\allowbreak{}Pair\allowbreak{}Sampling.\allowbreak{}Is\allowbreak{}Symmetric}
\item \hyperlink{governing-declaration-3449f3559185d0fd}{Applied\allowbreak{}Modeling\allowbreak{}Lib.\allowbreak{}Learning.\allowbreak{}Human\allowbreak{}Feedback.\allowbreak{}Pairwise\allowbreak{}Count\allowbreak{}Dataset.\allowbreak{}CDF\allowbreak{}Like\allowbreak{}Pairwise\allowbreak{}Link}
\item \hyperlink{governing-declaration-87eda978eb1b0c9d}{Applied\allowbreak{}Modeling\allowbreak{}Lib.\allowbreak{}Learning.\allowbreak{}Human\allowbreak{}Feedback.\allowbreak{}Pairwise\allowbreak{}Count\allowbreak{}Dataset.\allowbreak{}Score\allowbreak{}Vector}
\item \hyperlink{governing-declaration-cba21d23e0e5aa10}{Applied\allowbreak{}Modeling\allowbreak{}Lib.\allowbreak{}Learning.\allowbreak{}Human\allowbreak{}Feedback.\allowbreak{}Pairwise\allowbreak{}Count\allowbreak{}Dataset.\allowbreak{}is\allowbreak{}Pairwise\allowbreak{}MLE}
\item \hyperlink{governing-declaration-449f87e91bf37ce0}{Applied\allowbreak{}Modeling\allowbreak{}Lib.\allowbreak{}Learning.\allowbreak{}Human\allowbreak{}Feedback.\allowbreak{}Pairwise\allowbreak{}Count\allowbreak{}Dataset.\allowbreak{}of\allowbreak{}Binary\allowbreak{}Reports}
\item \hyperlink{governing-declaration-ed6af38b0210f5ec}{Applied\allowbreak{}Modeling\allowbreak{}Lib.\allowbreak{}Learning.\allowbreak{}Human\allowbreak{}Feedback.\allowbreak{}Pairwise\allowbreak{}Count\allowbreak{}Dataset.\allowbreak{}sigmoid\allowbreak{}CDF\allowbreak{}Like\allowbreak{}Pairwise\allowbreak{}Link}
\item \hyperlink{library-prerequisite-268c959ee7d2d148}{Applied\allowbreak{}Modeling\allowbreak{}Lib.\allowbreak{}PMF\allowbreak{}Full\allowbreak{}Support}
\item \hyperlink{library-prerequisite-d3573881ae45d4c6}{Applied\allowbreak{}Modeling\allowbreak{}Lib.\allowbreak{}finite\allowbreak{}KL\allowbreak{}Divergence}
\item \hyperlink{governing-declaration-779fe1555fbbcea5}{Applied\allowbreak{}Modeling\allowbreak{}Lib.\allowbreak{}finite\allowbreak{}PMF\allowbreak{}Iid\allowbreak{}Path\allowbreak{}Measure}
\item \hyperlink{governing-declaration-89219e02ea5145c7}{Applied\allowbreak{}Modeling\allowbreak{}Lib.\allowbreak{}finite\allowbreak{}PMF\allowbreak{}Iid\allowbreak{}Path\allowbreak{}Prefix}
\item \hyperlink{governing-declaration-8e4b4389a4f6454b}{Applied\allowbreak{}Modeling\allowbreak{}Lib.\allowbreak{}pmf\allowbreak{}Exp}
\item \hyperlink{governing-declaration-877274851b50572c}{Golz\allowbreak{}Haghtalab\allowbreak{}Yang2025\allowbreak{}Distortion.\allowbreak{}lemma15\allowbreak{}Population}
\item \hyperlink{governing-declaration-9f5c608cbac97312}{Golz\allowbreak{}Haghtalab\allowbreak{}Yang2025\allowbreak{}Distortion.\allowbreak{}theorem9\allowbreak{}Comparison\allowbreak{}Sampling}
\item \hyperlink{governing-declaration-e63d2c356e95b6b5}{Golz\allowbreak{}Haghtalab\allowbreak{}Yang2025\allowbreak{}Distortion.\allowbreak{}theorem9\allowbreak{}One\allowbreak{}Comparison\allowbreak{}Report\allowbreak{}Law}
\item \hyperlink{governing-declaration-a2115bb0ddee2a78}{Golz\allowbreak{}Haghtalab\allowbreak{}Yang2025\allowbreak{}Distortion.\allowbreak{}theorem9\allowbreak{}Prefix\allowbreak{}Profile}
\item \hyperlink{governing-declaration-35b3f5ca150bffc2}{Golz\allowbreak{}Haghtalab\allowbreak{}Yang2025\allowbreak{}Distortion.\allowbreak{}theorem9\allowbreak{}Uniform\allowbreak{}Reference}
\end{itemize}
\paragraph{Lean-elaborated claim roles}\begin{ReviewVerbatim}
1. parameter: ℝ
2. assumption: 0 < beta
3. parameter: ℝ
4. conclusion: ∃ (m : ℕ) (hm : 3 ≤ m) (epsilon : NNReal) (_ : 0 < epsilon) (_ : epsilon < 1),
  have population : PMF (Fin 3) := GolzHaghtalabYang2025Distortion.lemma15Population;
  have utility : AppliedModelingLib.Alignment.Welfare.FiniteUtilityProfile (Fin 3) (Fin m) :=
    GolzHaghtalabYang2025Distortion.theorem9PrefixProfile beta m;
  have comparisonSampling : AppliedModelingLib.Learning.HumanFeedback.CorrelatedPairSampling (Fin m) :=
    GolzHaghtalabYang2025Distortion.theorem9ComparisonSampling m hm epsilon;
  have reference : PMF (Fin m) :=
    GolzHaghtalabYang2025Distortion.theorem9UniformReference m
      (GolzHaghtalabYang2025Distortion.theorem9Spec._proof_3 m hm);
  have klBudget : ℝ := Real.log ↑m;
  AppliedModelingLib.Alignment.Welfare.UnitIntervalUtilityProfile utility ∧
    AppliedModelingLib.PMFFullSupport reference ∧
      comparisonSampling.IsSymmetric ∧
        comparisonSampling.HasZeroDiagonal ∧
          comparisonSampling.HasFullOffDiagonalSupport ∧
            0 ≤ klBudget ∧
              ∀ᵐ (path :
                ℕ →
                  AppliedModelingLib.Learning.HumanFeedback.BinaryPairwiseReport
                    (Fin
                      m)) ∂AppliedModelingLib.finitePMFIidPathMeasure
                  (GolzHaghtalabYang2025Distortion.theorem9OneComparisonReportLaw beta m hm epsilon),
                ∀ᶠ (horizon : ℕ) in Filter.atTop,
                  (∃ (reward : AppliedModelingLib.Learning.HumanFeedback.PairwiseCountDataset.ScoreVector (Fin m)),
                      (AppliedModelingLib.Learning.HumanFeedback.PairwiseCountDataset.ofBinaryReports
                            (AppliedModelingLib.finitePMFIidPathPrefix path horizon)).isPairwiseMLE
                        AppliedModelingLib.Learning.HumanFeedback.PairwiseCountDataset.sigmoidCDFLikePairwiseLink.toFun
                        ⟨0, GolzHaghtalabYang2025Distortion.theorem9Spec._proof_4 m hm⟩ reward) ∧
                    ∀ (reward : AppliedModelingLib.Learning.HumanFeedback.PairwiseCountDataset.ScoreVector (Fin m)),
                      (AppliedModelingLib.Learning.HumanFeedback.PairwiseCountDataset.ofBinaryReports
                              (AppliedModelingLib.finitePMFIidPathPrefix path horizon)).isPairwiseMLE
                          AppliedModelingLib.Learning.HumanFeedback.PairwiseCountDataset.sigmoidCDFLikePairwiseLink.toFun
                          ⟨0, GolzHaghtalabYang2025Distortion.theorem9Spec._proof_4 m hm⟩ reward →
                        ∃ (rewardPolicy : PMF (Fin m)) (welfarePolicy : PMF (Fin m)),
                          AppliedModelingLib.finiteKLDivergence rewardPolicy reference ≤ klBudget ∧
                            (∀ (other : PMF (Fin m)),
                                AppliedModelingLib.finiteKLDivergence other reference ≤ klBudget →
                                  AppliedModelingLib.pmfExp other reward ≤
                                    AppliedModelingLib.pmfExp rewardPolicy reward) ∧
                              AppliedModelingLib.finiteKLDivergence welfarePolicy reference ≤ klBudget ∧
                                (∀ (other : PMF (Fin m)),
                                    AppliedModelingLib.finiteKLDivergence other reference ≤ klBudget →
                                      AppliedModelingLib.Alignment.Welfare.policyAverageUtility population utility
                                          other ≤
                                        AppliedModelingLib.Alignment.Welfare.policyAverageUtility population utility
                                          welfarePolicy) ∧
                                  0 <
                                      AppliedModelingLib.Alignment.Welfare.policyAverageUtility population utility
                                        rewardPolicy ∧
                                    bound <
                                      AppliedModelingLib.Alignment.Welfare.policyAverageUtility population utility
                                          welfarePolicy /
                                        AppliedModelingLib.Alignment.Welfare.policyAverageUtility population utility
                                          rewardPolicy
\end{ReviewVerbatim}

\paragraph{Lean proof endpoint (built separately)}\begin{ReviewVerbatim}
GolzHaghtalabYang2025Distortion.theorem9_correlated_sampling_unbounded
\end{ReviewVerbatim}

\paragraph{Recorded source-to-Spec screening}
\textbf{Verdict:} \texttt{matches}
\\
{\raggedright
\textbf{Reason:} This constructs correlated comparison sampling with unbounded RLHF distortion, as in the source claim.\allowbreak{}
\par}
\reviewmatch{Matches source input}
\noindent\textbf{Reviewer annotation}\par
\reviewerbox
\clearpage
\hypertarget{source-claim-cc5e88fff44863d2}{}
\section*{Review row 16}
\textbf{Source locator:} {\footnotesize\raggedright cited publication, lines 2114-\allowbreak{}2149\par}
\paragraph{Verbatim source input}
\begin{ReviewVerbatim}
F.3   Equivalence of DPO and RLHF under Heterogeneous Preferences

In this section, we formalize the observation that DPO and RLHF are equivalent under heterogeneous
preferences. This is consistent with the result by Shirali et al. [45], which shows that DPO also aligns
with the Borda count. We start by recalling the DPO objective [43].12
                                                                                               !
                                      X                         π(xji )             π(yij )
              LDPO (π; πref ) = −               log σ β log                − β log               .
                                  1≤i≤n,1≤j≤d
                                                               πref (xji )         πref (yij )
Now we perform a change of variables to transform π into the following form:
                                                                    π(x)
              π(x) = πref (x) · exp(r̂(x)/β) where r̂(x) := β log            , ∀x ∈ A.
                                                                   πref (x)
Substituting this into the DPO objective, we get:
                                             X
                                                   log σ(r̂(xji ) − r̂(yij )) ,
                                                                             [U+0001 control character]
                       LDPO (π; πref ) = −
                                             1≤i≤n,1≤j≤d
which is exactly the MLE objective for reward learning in RLHF, repeated here for convenience:
                                          X
                                                   log σ(r(xji ) − r(yij )) .
                                                                           [U+0001 control character]
                       LMLE (r) := −
                                           1≤i≤n,1≤j≤d
Therefore, there is a one-to-one correspondence between the MLE reward r? =
argminr∈Rm LMLE (r)13 , and the DPO policy πDPO = argminπ LDPO (π; πref ) as:
                                 πDPO (x) = πref (x) · exp(r? (x)/β).
On the other hand, the RLHF policy with regularization parameter λ (see Appendix F.4 for the
equivalence between regularized and constrained versions of RLHF) is also given by
                                 πRLHF (x) = πref (x) · exp(r? (x)/λ).
Therefore, the DPO policy πDPO and the RLHF policy πRLHF are equivalent when the parameter β
in the DPO objective is equal to λ in the RLHF objective. Notably, both policies are different from
the optimal policy π ? ∝ πref · exp(AvgUtil(x)/λ0 ) (for a potentially different λ0 that make the KL
constraint tight) as r? 6= AvgUtil, which has also been pointed out by Shirali et al. [45].
\end{ReviewVerbatim}



\paragraph{Expanded PaperInterface specification}
\begin{ReviewVerbatim}
∀ {Alternative : Type u_1} [inst : Fintype Alternative] [inst_1 : DecidableEq Alternative] {horizon : ℕ}
  (sample : Fin horizon → AppliedModelingLib.Learning.HumanFeedback.BinaryPairwiseReport Alternative)
  (scoreReference : Alternative) (reference : PMF Alternative)
  (reward : AppliedModelingLib.Learning.HumanFeedback.PairwiseCountDataset.ScoreVector Alternative) {klWeight : ℝ},
  (AppliedModelingLib.Learning.HumanFeedback.PairwiseCountDataset.ofBinaryReports sample).isPairwiseMLE Real.sigmoid
      scoreReference reward →
    0 < klWeight →
      AppliedModelingLib.PMFFullSupport reference →
        GolzHaghtalabYang2025Distortion.sourceDpoPolicy reference reward klWeight =
            GolzHaghtalabYang2025Distortion.sourceRlhfPolicy reference reward klWeight ∧
          GolzHaghtalabYang2025Distortion.sourceDpoBinaryReportLoss sample reference
                (GolzHaghtalabYang2025Distortion.sourceDpoPolicy reference reward klWeight) klWeight =
              GolzHaghtalabYang2025Distortion.sourceRlhfBinaryReportLoss sample reward ∧
            (AppliedModelingLib.Learning.HumanFeedback.RewardEquivalent
                (AppliedModelingLib.Learning.HumanFeedback.dpoImplicitReward
                  (AppliedModelingLib.Alignment.Welfare.contextFreeAlternativePolicy
                    (GolzHaghtalabYang2025Distortion.sourceDpoPolicy reference reward klWeight))
                  (AppliedModelingLib.Alignment.Welfare.contextFreeAlternativePolicy reference) klWeight)
                fun (x : PUnit.{1}) => reward) ∧
              GolzHaghtalabYang2025Distortion.IsSourceDpoLossMinimizer sample reference klWeight
                  (GolzHaghtalabYang2025Distortion.sourceDpoPolicy reference reward klWeight) ∧
                ∀ (policy : PMF Alternative),
                  GolzHaghtalabYang2025Distortion.IsSourceDpoLossMinimizer sample reference klWeight policy →
                    (AppliedModelingLib.Learning.HumanFeedback.PairwiseCountDataset.ofBinaryReports
                            sample).isPairwiseMLE
                        Real.sigmoid scoreReference
                        (GolzHaghtalabYang2025Distortion.sourceDpoNormalizedMLEReward scoreReference reference policy
                          klWeight) ∧
                      GolzHaghtalabYang2025Distortion.sourceDpoPolicy reference
                          (GolzHaghtalabYang2025Distortion.sourceDpoNormalizedMLEReward scoreReference reference policy
                            klWeight)
                          klWeight =
                        policy
\end{ReviewVerbatim}

\paragraph{Retained governing declarations}
\begin{itemize}\raggedright
\item \hyperlink{governing-declaration-b7cf9e04e23161e5}{Applied\allowbreak{}Modeling\allowbreak{}Lib.\allowbreak{}Alignment.\allowbreak{}Welfare.\allowbreak{}context\allowbreak{}Free\allowbreak{}Alternative\allowbreak{}Policy}
\item \hyperlink{governing-declaration-bc3d6ca8e7e0ac9d}{Applied\allowbreak{}Modeling\allowbreak{}Lib.\allowbreak{}Learning.\allowbreak{}Human\allowbreak{}Feedback.\allowbreak{}Binary\allowbreak{}Pairwise\allowbreak{}Report}
\item \hyperlink{governing-declaration-87eda978eb1b0c9d}{Applied\allowbreak{}Modeling\allowbreak{}Lib.\allowbreak{}Learning.\allowbreak{}Human\allowbreak{}Feedback.\allowbreak{}Pairwise\allowbreak{}Count\allowbreak{}Dataset.\allowbreak{}Score\allowbreak{}Vector}
\item \hyperlink{governing-declaration-cba21d23e0e5aa10}{Applied\allowbreak{}Modeling\allowbreak{}Lib.\allowbreak{}Learning.\allowbreak{}Human\allowbreak{}Feedback.\allowbreak{}Pairwise\allowbreak{}Count\allowbreak{}Dataset.\allowbreak{}is\allowbreak{}Pairwise\allowbreak{}MLE}
\item \hyperlink{governing-declaration-449f87e91bf37ce0}{Applied\allowbreak{}Modeling\allowbreak{}Lib.\allowbreak{}Learning.\allowbreak{}Human\allowbreak{}Feedback.\allowbreak{}Pairwise\allowbreak{}Count\allowbreak{}Dataset.\allowbreak{}of\allowbreak{}Binary\allowbreak{}Reports}
\item \hyperlink{governing-declaration-16b50286b4cd31a4}{Applied\allowbreak{}Modeling\allowbreak{}Lib.\allowbreak{}Learning.\allowbreak{}Human\allowbreak{}Feedback.\allowbreak{}Reward\allowbreak{}Equivalent}
\item \hyperlink{governing-declaration-95571fb7643367ad}{Applied\allowbreak{}Modeling\allowbreak{}Lib.\allowbreak{}Learning.\allowbreak{}Human\allowbreak{}Feedback.\allowbreak{}dpo\allowbreak{}Implicit\allowbreak{}Reward}
\item \hyperlink{library-prerequisite-268c959ee7d2d148}{Applied\allowbreak{}Modeling\allowbreak{}Lib.\allowbreak{}PMF\allowbreak{}Full\allowbreak{}Support}
\item \hyperlink{governing-declaration-4cdae0c3705e3d00}{Golz\allowbreak{}Haghtalab\allowbreak{}Yang2025\allowbreak{}Distortion.\allowbreak{}Is\allowbreak{}Source\allowbreak{}Dpo\allowbreak{}Loss\allowbreak{}Minimizer}
\item \hyperlink{governing-declaration-0f2c5b70a888d5b3}{Golz\allowbreak{}Haghtalab\allowbreak{}Yang2025\allowbreak{}Distortion.\allowbreak{}source\allowbreak{}Dpo\allowbreak{}Binary\allowbreak{}Report\allowbreak{}Loss}
\item \hyperlink{governing-declaration-d547690922524bc2}{Golz\allowbreak{}Haghtalab\allowbreak{}Yang2025\allowbreak{}Distortion.\allowbreak{}source\allowbreak{}Dpo\allowbreak{}Normalized\allowbreak{}MLE\allowbreak{}Reward}
\item \hyperlink{governing-declaration-7f6992add20edc0b}{Golz\allowbreak{}Haghtalab\allowbreak{}Yang2025\allowbreak{}Distortion.\allowbreak{}source\allowbreak{}Dpo\allowbreak{}Policy}
\item \hyperlink{governing-declaration-9dc701ec2b5767ff}{Golz\allowbreak{}Haghtalab\allowbreak{}Yang2025\allowbreak{}Distortion.\allowbreak{}source\allowbreak{}Rlhf\allowbreak{}Binary\allowbreak{}Report\allowbreak{}Loss}
\item \hyperlink{governing-declaration-55b65ae526406849}{Golz\allowbreak{}Haghtalab\allowbreak{}Yang2025\allowbreak{}Distortion.\allowbreak{}source\allowbreak{}Rlhf\allowbreak{}Policy}
\end{itemize}
\paragraph{Lean-elaborated claim roles}\begin{ReviewVerbatim}
1. parameter: Type u_1
2. parameter: Fintype Alternative
3. parameter: DecidableEq Alternative
4. parameter: ℕ
5. parameter: Fin horizon → AppliedModelingLib.Learning.HumanFeedback.BinaryPairwiseReport Alternative
6. parameter: Alternative
7. parameter: PMF Alternative
8. parameter: AppliedModelingLib.Learning.HumanFeedback.PairwiseCountDataset.ScoreVector Alternative
9. parameter: ℝ
10. assumption: (AppliedModelingLib.Learning.HumanFeedback.PairwiseCountDataset.ofBinaryReports sample).isPairwiseMLE Real.sigmoid
  scoreReference reward
11. assumption: 0 < klWeight
12. assumption: AppliedModelingLib.PMFFullSupport reference
13. conclusion: GolzHaghtalabYang2025Distortion.sourceDpoPolicy reference reward klWeight =
    GolzHaghtalabYang2025Distortion.sourceRlhfPolicy reference reward klWeight ∧
  GolzHaghtalabYang2025Distortion.sourceDpoBinaryReportLoss sample reference
        (GolzHaghtalabYang2025Distortion.sourceDpoPolicy reference reward klWeight) klWeight =
      GolzHaghtalabYang2025Distortion.sourceRlhfBinaryReportLoss sample reward ∧
    (AppliedModelingLib.Learning.HumanFeedback.RewardEquivalent
        (AppliedModelingLib.Learning.HumanFeedback.dpoImplicitReward
          (AppliedModelingLib.Alignment.Welfare.contextFreeAlternativePolicy
            (GolzHaghtalabYang2025Distortion.sourceDpoPolicy reference reward klWeight))
          (AppliedModelingLib.Alignment.Welfare.contextFreeAlternativePolicy reference) klWeight)
        fun (x : PUnit.{1}) => reward) ∧
      GolzHaghtalabYang2025Distortion.IsSourceDpoLossMinimizer sample reference klWeight
          (GolzHaghtalabYang2025Distortion.sourceDpoPolicy reference reward klWeight) ∧
        ∀ (policy : PMF Alternative),
          GolzHaghtalabYang2025Distortion.IsSourceDpoLossMinimizer sample reference klWeight policy →
            (AppliedModelingLib.Learning.HumanFeedback.PairwiseCountDataset.ofBinaryReports sample).isPairwiseMLE
                Real.sigmoid scoreReference
                (GolzHaghtalabYang2025Distortion.sourceDpoNormalizedMLEReward scoreReference reference policy
                  klWeight) ∧
              GolzHaghtalabYang2025Distortion.sourceDpoPolicy reference
                  (GolzHaghtalabYang2025Distortion.sourceDpoNormalizedMLEReward scoreReference reference policy
                    klWeight)
                  klWeight =
                policy
\end{ReviewVerbatim}

\paragraph{Lean proof endpoint (built separately)}\begin{ReviewVerbatim}
GolzHaghtalabYang2025Distortion.appendixF3_dpo_rlhf_equivalence
\end{ReviewVerbatim}

\paragraph{Recorded source-to-Spec screening}
\textbf{Verdict:} \texttt{matches}
\\
{\raggedright
\textbf{Reason:} The documented corrected finite target normalizes the exponential tilt, equates literal DPO and MLE losses, and preserves policy/\allowbreak{}minimizer correspondence under the stated positive-\allowbreak{}weight, full-\allowbreak{}support domain.\allowbreak{}
\par}
\reviewmatch{Matches source input}
\noindent\textbf{Reviewer annotation}\par
\reviewerbox

\end{document}
